\chapt{अरण्यकाण्डः}

\sect{प्रथमः सर्गः}

\twolineshloka
{प्रविश्य तु महारण्यं दण्डकारण्यमात्मवान्}
{ददर्श रामो दुर्धर्षस्तापसाश्रममण्डलम्} %||3-1-1||

\twolineshloka
{कुशचीरपरिक्षिप्तं ब्राह्म्या लक्ष्म्या समावृतम्}
{यथा प्रदीप्तं दुर्धर्शं गगने सूर्यमण्डलम्} %||3-1-2||

\twolineshloka
{शरण्यं सर्वभूतानां सुसमृष्टाजिरं सदा}
{पूजितं चोपनृत्तं च नित्यमप्सरसां गणैः} %||3-1-3||

\twolineshloka
{विशालैरग्निशरणैः स्रुग्भाण्डैरजिनैः कुशैः}
{समिद्भिस्तोयकलशैः फलमूलैश्च शोभितम्} %||3-1-4||

\twolineshloka
{आरण्यैश्च महावृक्षैः पुण्यैः स्वादुफलैर्वृतम्}
{बलिहोमार्चितं पुण्यं ब्रह्मघोषनिनादितम्} %||3-1-5||

\twolineshloka
{पुष्पैर्वन्यैः परिक्षिप्तं पद्मिन्या च सपद्मया}
{फलमूलाशनैर्दान्तैश्चीरकृष्णाजिनाम्बरैः} %||3-1-6||

\twolineshloka
{सूर्यवैश्वानराभैश्च पुराणैर्मुनिभिर्वृतम्}
{पुण्यैश नियताहारैः शोभितं परमर्षिभिः} %||3-1-7||

\twolineshloka
{तद्ब्रह्मभवनप्रख्यं ब्रह्मघोषनिनादितम्}
{ब्रह्मविद्भिर्महाभागैर्ब्राह्मणैरुपशोभितम्} %||3-1-8||

\twolineshloka
{तद्दृष्ट्वा राघवः श्रीमांस्तापसाश्रममण्डलम्}
{अभ्यगच्छन्महातेजा विज्यं कृत्वा महद्धनुः} %||3-1-9||

\twolineshloka
{दिव्यज्ञानोपपन्नास्ते रामं दृष्ट्वा महर्षयः}
{अभ्यगच्छंस्तदा प्रीता वैदेहीं च यशस्विनीम्} %||3-1-10||

\twolineshloka
{ते तं सोममिवोद्यन्तं दृष्ट्वा वै धर्मचारिणः}
{मङ्गलानि प्रयुञ्जानाः प्रत्यगृह्णन्दृढव्रताः} %||3-1-11||

\twolineshloka
{रूपसंहननं लक्ष्मीं सौकुमार्यं सुवेषताम्}
{ददृशुर्विस्मिताकारा रामस्य वनवासिनः} %||3-1-12||

\twolineshloka
{वैदेहीं लक्ष्मणं रामं नेत्रैरनिमिषैरिव}
{आश्चर्यभूतान्ददृशुः सर्वे ते वनचारिणः} %||3-1-13||

\twolineshloka
{अत्रैनं हि महाभागाः सर्वभूतहिते रताः}
{अतिथिं पर्णशालायां राघवं संन्यवेशयन्} %||3-1-14||

\twolineshloka
{ततो रामस्य सत्कृत्य विधिना पावकोपमाः}
{आजह्रुस्ते महाभागाः सलिलं धर्मचारिणः} %||3-1-15||

\twolineshloka
{मूलं पुष्पं फलं वन्यमाश्रमं च महात्मनः}
{निवेदयीत्वा धर्मज्ञास्ततः प्राञ्जलयोऽब्रुवन्} %||3-1-16||

\twolineshloka
{धर्मपालो जनस्यास्य शरण्यश्च महायशाः}
{पूजनीयश्च मान्यश्च राजा दण्डधरो गुरुः} %||3-1-17||

\twolineshloka
{इन्द्रस्यैव चतुर्भागः प्रजा रक्षति राघव}
{राजा तस्माद्वनान्भोगान्भुङ्क्ते लोकनमस्कृतः} %||3-1-18||

\twolineshloka
{ते वयं भवता रक्ष्या भवद्विषयवासिनः}
{नगरस्थो वनस्थो वा त्वं नो राजा जनेश्वरः} %||3-1-19||

\twolineshloka
{न्यस्तदण्डा वयं राजञ्जितक्रोधा जितेन्द्रियाः}
{रक्षितव्यास्त्वया शश्वद्गर्भभूतास्तपोधनाः} %||3-1-20||

\twolineshloka
{एवमुक्त्वा फलैर्मूलैः पुष्पैर्वन्यैश्च राघवम्}
{अन्यैश्च विविधाहारैः सलक्ष्मणमपूजयन्} %||3-1-21||

\twolineshloka
{तथान्ये तापसाः सिद्धा रामं वैश्वानरोपमाः}
{न्यायवृत्ता यथान्यायं तर्पयामासुरीश्वरम्} %||3-2-22||


{॥इत्यार्षे श्रीमद्रामायणे वाल्मीकीये आदिकाव्ये अरण्यकाण्डे प्रथमः सर्गः॥}

\sect{द्वितीयः सर्गः}

\twolineshloka
{कृतातिथ्योऽथ रामस्तु सूर्यस्योदयनं प्रति}
{आमन्त्र्य स मुनीन्सर्वान्वनमेवान्वगाहत} %||3-2-1||

\twolineshloka
{नानामृगगणाकीर्णं शार्दूलवृकसेवितम्}
{ध्वस्तवृक्षलतागुल्मं दुर्दर्श सलिलाशयम्} %||3-2-2||

\twolineshloka
{निष्कूजनानाशकुनि झिल्लिका गणनादितम्}
{लक्ष्मणानुगतो रामो वनमध्यं ददर्श ह} %||3-2-3||

\twolineshloka
{वनमध्ये तु काकुत्स्थस्तस्मिन्घोरमृगायुते}
{ददर्श गिरिशृङ्गाभं पुरुषादं महास्वनम्} %||3-2-4||

\twolineshloka
{गभीराक्षं महावक्त्रं विकटं विषमोदरम्}
{बीभत्सं विषमं दीर्घं विकृतं घोरदर्शनम्} %||3-2-5||

\twolineshloka
{वसानं चर्मवैयाघ्रं वसार्द्रं रुधिरोक्षितम्}
{त्रासनं सर्वभूतानां व्यादितास्यमिवान्तकम्} %||3-2-6||

\twolineshloka
{त्रीन्सिंहांश्चतुरो व्याघ्रान्द्वौ वृकौ पृषतान्दश}
{सविषाणं वसादिग्धं गजस्य च शिरो महत्} %||3-2-7||

\twolineshloka
{अवसज्यायसे शूले विनदन्तं महास्वनम्}
{स रामं लक्ष्मणं चैव सीतां दृष्ट्वा च मैथिलीम्} %||3-2-8||

\twolineshloka
{अभ्यधावत्सुसङ्क्रुद्धः प्रजाः काल इवान्तकः}
{स कृत्वा भैरवं नादं चालयन्निव मेदिनीम्} %||3-2-9||

\twolineshloka
{अङ्गेनादाय वैदेहीमपक्रम्य ततोऽब्रवीत्}
{युवां जटाचीरधरौ सभार्यौ क्षीणजीवितौ} %||3-2-10||

\twolineshloka
{प्रविष्टौ दण्डकारण्यं शरचापासिधारिणौ}
{कथं तापसयोर्वां च वासः प्रमदया सह} %||3-2-11||

\twolineshloka
{अधर्मचारिणौ पापौ कौ युवां मुनिदूषकौ}
{अहं वनमिदं दुर्गं विराघो नाम राक्षसः} %||3-2-12||

\threelineshloka
{चरामि सायुधो नित्यमृषिमांसानि भक्षयन्}
{इयं नारी वरारोहा मम भर्या भविष्यति}
{युवयोः पापयोश्चाहं पास्यामि रुधिरं मृधे} %||3-2-13||

\threelineshloka
{तस्यैवं ब्रुवतो धृष्टं विराधस्य दुरात्मनः}
{श्रुत्वा सगर्वितं वाक्यं सम्भ्रान्ता जनकात्मजा}
{सीता प्रावेपतोद्वेगात्प्रवाते कदली यथा} %||3-2-14||

\twolineshloka
{तां दृष्ट्वा राघवः सीतां विराधाङ्कगतां शुभाम्}
{अब्रवील्लक्ष्मणं वाक्यं मुखेन परिशुष्यता} %||3-2-15||

\threelineshloka
{पश्य सौम्य नरेन्द्रस्य जनकस्यात्मसम्भवाम्}
{मम भार्यां शुभाचारां विराधाङ्के प्रवेशिताम्}
{अत्यन्त सुखसंवृद्धां राजपुत्रीं यशस्विनीम्} %||3-2-16||

\twolineshloka
{यदभिप्रेतमस्मासु प्रियं वर वृतं च यत्}
{कैकेय्यास्तु सुसंवृत्तं क्षिप्रमद्यैव लक्ष्मण} %||3-2-17||

\threelineshloka
{या न तुष्यति राज्येन पुत्रार्थे दीर्घदर्शिनी}
{ययाहं सर्वभूतानां हितः प्रस्थापितो वनम्}
{अद्येदानीं सकामा सा या माता मम मध्यमा} %||3-2-18||

\twolineshloka
{परस्पर्शात्तु वैदेह्या न दुःखतरमस्ति मे}
{पितुर्विनाशात्सौमित्रे स्वराज्यहरणात्तथा} %||3-2-19||

\twolineshloka
{इति ब्रुवति काकुत्स्थे बाष्पशोकपरिप्लुते}
{अब्रवील्लक्ष्मणः क्रुद्धो रुद्धो नाग इव श्वसन्} %||3-2-20||

\twolineshloka
{अनाथ इव भूतानां नाथस्त्वं वासवोपमः}
{मया प्रेष्येण काकुत्स्थ किमर्थं परितप्स्यसे} %||3-2-21||

\twolineshloka
{शरेण निहतस्याद्य मया क्रुद्धेन रक्षसः}
{विराधस्य गतासोर्हि मही पास्यति शोणितम्} %||3-2-22||

\twolineshloka
{राज्यकामे मम क्रोधो भरते यो बभूव ह}
{तं विराधे विमोक्ष्यामि वज्री वज्रमिवाचले} %||3-2-23||

\fourlineindentedshloka
{मम भुजबलवेगवेगितः}
{ पततु शरोऽस्य महान्महोरसि}
{व्यपनयतु तनोश्च जीवितम्}
{ पततु ततश्च महीं विघूर्णितः} %||3-3-24||


{॥इत्यार्षे श्रीमद्रामायणे वाल्मीकीये आदिकाव्ये अरण्यकाण्डे द्वितीयः सर्गः॥}

\sect{तृतीयः सर्गः}

\twolineshloka
{अथोवाच पुनर्वाक्यं विराधः पूरयन्वनम्}
{आत्मानं पृच्छते ब्रूतं कौ युवां क्व गमिष्यथः} %||3-3-1||

\twolineshloka
{तमुवाच ततो रामो राक्षसं ज्वलिताननम्}
{पृच्छन्तं सुमहातेजा इक्ष्वाकुकुलमात्मनः} %||3-3-2||

\twolineshloka
{क्षत्रियो वृत्तसम्पन्नौ विद्धि नौ वनगोचरौ}
{त्वां तु वेदितुमिच्छावः कस्त्वं चरसि दण्डकान्} %||3-3-3||

\twolineshloka
{तमुवाच विराधस्तु रामं सत्यपराक्रमम्}
{हन्त वक्ष्यामि ते राजन्निबोध मम राघव} %||3-3-4||

\twolineshloka
{पुत्रः किल जयस्याहं माता मम शतह्रदा}
{विराध इति मामाहुः पृथिव्यां सर्वराक्षसाः} %||3-3-5||

\twolineshloka
{तपसा चापि मे प्राप्ता ब्रह्मणो हि प्रसादजा}
{शस्त्रेणावध्यता लोकेऽच्छेद्याभेद्यत्वमेव च} %||3-3-6||

\twolineshloka
{उत्सृज्य प्रमदामेनामनपेक्षौ यथागतम्}
{त्वरमाणौ पालयेथां न वां जीवितमाददे} %||3-3-7||

\twolineshloka
{तं रामः प्रत्युवाचेदं कोपसंरक्तलोचनः}
{राक्षसं विकृताकारं विराधं पापचेतसम्} %||3-3-8||

\twolineshloka
{क्षुद्र धिक्त्वां तु हीनार्थं मृत्युमन्वेषसे ध्रुवम्}
{रणे सम्प्राप्स्यसे तिष्ठ न मे जीवन्गमिष्यसि} %||3-3-9||

\twolineshloka
{ततः सज्यं धनुः कृत्वा रामः सुनिशिताञ्शरान्}
{सुशीघ्रमभिसन्धाय राक्षसं निजघान ह} %||3-3-10||

\twolineshloka
{धनुषा ज्यागुणवता सप्तबाणान्मुमोच ह}
{रुक्मपुङ्खान्महावेगान्सुपर्णानिलतुल्यगान्} %||3-3-11||

\twolineshloka
{ते शरीरं विराधस्य भित्त्वा बर्हिणवाससः}
{निपेतुः शोणितादिग्धा धरण्यां पावकोपमाः} %||3-3-12||

\twolineshloka
{स विनद्य महानादं शूलं शक्रध्वजोपमम्}
{प्रगृह्याशोभत तदा व्यात्तानन इवान्तकः} %||3-3-13||

\twolineshloka
{तच्छूलं वज्रसङ्काशं गगने ज्वलनोपमम्}
{द्वाभ्यां शराभ्यां चिच्छेद रामः शस्त्रभृतां वरः} %||3-3-14||

\twolineshloka
{तस्य रौद्रस्य सौमित्रिर्बाहुं सव्यं बभञ्ज ह}
{रामस्तु दक्षिणं बाहुं तरसा तस्य रक्षसः} %||3-3-15||

\threelineshloka
{स भग्नबाहुः संविग्नो निपपाताशु राक्षसः}
{धरण्यां मेघसङ्काशो वज्रभिन्न इवाचलः}
{इदं प्रोवाच काकुत्स्थं विराधः पुरुषर्षभम्} %||3-3-16||

\twolineshloka
{कौसल्या सुप्रजास्तात रामस्त्वं विदितो मया}
{वैदेही च महाभागा लक्ष्मणश्च महायशाः} %||3-3-17||

\twolineshloka
{अभिशापादहं घोरां प्रविष्टो राक्षसीं तनुम्}
{तुम्बुरुर्नाम गन्धर्वः शप्तो वैश्वरणेन हि} %||3-3-18||

\twolineshloka
{प्रसाद्यमानश्च मया सोऽब्रवीन्मां महायशाः}
{यदा दाशरथी रामस्त्वां वधिष्यति संयुगे} %||3-3-19||

\twolineshloka
{तदा प्रकृतिमापन्नो भवान्स्वर्गं गमिष्यति}
{इति वैश्रवणो राजा रम्भासक्तमुवाच ह} %||3-3-20||

\threelineshloka
{अनुपस्थीयमानो मां सङ्क्रुद्धो व्यजहार ह}
{तव प्रसादान्मुक्तोऽहमभिशापात्सुदारुणात्}
{भवनं स्वं गमिष्यामि स्वस्ति वोऽस्तु परन्तप} %||3-3-21||

\twolineshloka
{इतो वसति धर्मात्मा शरभङ्गः प्रतापवान्}
{अध्यर्धयोजने तात महर्षिः सूर्यसंनिभः} %||3-3-22||

\twolineshloka
{तं क्षिप्रमभिगच्छ त्वं स ते श्रेयो विधास्यति}
{अवटे चापि मां राम निक्षिप्य कुशली व्रज} %||3-3-23||

\twolineshloka
{रक्षसां गतसत्त्वानामेष धर्मः सनातनः}
{अवटे ये निधीयन्ते तेषां लोकाः सनातनाः} %||3-3-24||

\twolineshloka
{एवमुक्त्वा तु काकुत्स्थं विराधः शरपीडितः}
{बभूव स्वर्गसम्प्राप्तो न्यस्तदेहो महाबलः} %||3-3-25||

\twolineshloka
{तं मुक्तकण्ठमुत्क्षिप्य शङ्कुकर्णं महास्वनम्}
{विराधं प्राक्षिपच्छ्वभ्रे नदन्तं भैरवस्वनम्} %||3-3-26||

\fourlineindentedshloka
{ततस्तु तौ काञ्चनचित्रकार्मुकौ}
{ निहत्य रक्षः परिगृह्य मैथिलीम्}
{विजह्रतुस्तौ मुदितौ महावने}
{ दिवि स्थितौ चन्द्रदिवाकराविव} %||3-4-27||


{॥इत्यार्षे श्रीमद्रामायणे वाल्मीकीये आदिकाव्ये अरण्यकाण्डे तृतीयः सर्गः॥}

\sect{चतुर्थः सर्गः}

\threelineshloka
{हत्वा तु तं भीमबलं विराधं राक्षसं वने}
{ततः सीतां परिष्वज्य समाश्वास्य च वीर्यवान्}
{अब्रवील्लक्ष्मणां रामो भ्रातरं दीप्ततेजसम्} %||3-4-1||

\twolineshloka
{कष्टं वनमिदं दुर्गं न च स्मो वनगोचराः}
{अभिगच्छामहे शीघ्रं शरभङ्गं तपोधनम्} %||3-4-2||

{आश्रमं शरभङ्गस्य राघवोऽभिजगाम ह॥\devanumber{3}॥} %||3-4-3|| (Check)

\addtocounter{shlokacount}{1}

\twolineshloka
{तस्य देवप्रभावस्य तपसा भावितात्मनः}
{समीपे शरभङ्गस्य ददर्श महदद्भुतम्} %||3-4-4||

\twolineshloka
{विभ्राजमानं वपुषा सूर्यवैश्वानरोपमम्}
{असंस्पृशन्तं वसुधां ददर्श विबुधेश्वरम्} %||3-4-5||

\twolineshloka
{सुप्रभाभरणं देवं विरजोऽम्बरधारिणम्}
{तद्विधैरेव बहुभिः पूज्यमानं महात्मभिः} %||3-4-6||

\twolineshloka
{हरिभिर्वाजिभिर्युक्तमन्तरिक्षगतं रथम्}
{ददर्शादूरतस्तस्य तरुणादित्यसंनिभम्} %||3-4-7||

\twolineshloka
{पाण्डुराभ्रघनप्रख्यं चन्द्रमण्डलसंनिभम्}
{अपश्यद्विमलं छत्रं चित्रमाल्योपशोभितम्} %||3-4-8||

\twolineshloka
{चामरव्यजने चाग्र्ये रुक्मदण्डे महाधने}
{गृहीते वननारीभ्यां धूयमाने च मूर्धनि} %||3-4-9||

\twolineshloka
{गन्धर्वामरसिद्धाश्च बहवः परमर्षयः}
{अन्तरिक्षगतं देवं वाग्भिरग्र्याभिरीडिरे} %||3-4-10||

\threelineshloka
{दृष्ट्वा शतक्रतुं तत्र रामो लक्ष्मणमब्रवीत्}
{ये हयाः पुरुहूतस्य पुरा शक्रस्य नः श्रुताः}
{अन्तरिक्षगता दिव्यास्त इमे हरयो ध्रुवम्} %||3-4-11||

\twolineshloka
{इमे च पुरुषव्याघ्र ये तिष्ठन्त्यभितो रथम्}
{शतं शतं कुण्डलिनो युवानः खड्गपाणयः} %||3-4-12||

\twolineshloka
{उरोदेशेषु सर्वेषां हारा ज्वलनसंनिभाः}
{रूपं बिभ्रति सौमित्रे पञ्चविंशतिवार्षिकम्} %||3-4-13||

\twolineshloka
{एतद्धि किल देवानां वयो भवति नित्यदा}
{यथेमे पुरुषव्याघ्रा दृश्यन्ते प्रियदर्शनाः} %||3-4-14||

\twolineshloka
{इहैव सह वैदेह्या मुहूर्तं तिष्ठ लक्ष्मण}
{यावज्जनाम्यहं व्यक्तं क एष द्युतिमान्रथे} %||3-4-15||

\twolineshloka
{तमेवमुक्त्वा सौमित्रिमिहैव स्थीयतामिति}
{अभिचक्राम काकुत्स्थः शरभङ्गाश्रमं प्रति} %||3-4-16||

\twolineshloka
{ततः समभिगच्छन्तं प्रेक्ष्य रामं शचीपतिः}
{शरभङ्गमनुज्ञाप्य विबुधानिदमब्रवीत्} %||3-4-17||

\twolineshloka
{इहोपयात्यसौ रामो यावन्मां नाभिभाषते}
{निष्ठां नयत तावत्तु ततो मां द्रष्टुमर्हति} %||3-4-18||

\twolineshloka
{जितवन्तं कृतार्थं च द्रष्टाहमचिरादिमम्}
{कर्म ह्यनेन कर्तव्यं महदन्यैः सुदुष्करम्} %||3-4-19||

\twolineshloka
{इति वज्री तमामन्त्र्य मानयित्वा च तापसम्}
{रथेन हरियुक्तेन ययौ दिवमरिन्दमः} %||3-4-20||

\twolineshloka
{प्रयाते तु सहस्राक्षे राघवः सपरिच्छदः}
{अग्निहोत्रमुपासीनं शरभङ्गमुपागमत्} %||3-4-21||

\twolineshloka
{तस्य पादौ च सङ्गृह्य रामः सीता च लक्ष्मणः}
{निषेदुस्तदनुज्ञाता लब्धवासा निमन्त्रिताः} %||3-4-22||

\twolineshloka
{ततः शक्रोपयानं तु पर्यपृच्छत्स राघवः}
{शरभङ्गश्च तत्सर्वं राघवाय न्यवेदयत्} %||3-4-23||

\twolineshloka
{मामेष वरदो राम ब्रह्मलोकं निनीषति}
{जितमुग्रेण तपसा दुष्प्रापमकृतात्मभिः} %||3-4-24||

\twolineshloka
{अहं ज्ञात्वा नरव्याघ्र वर्तमानमदूरतः}
{ब्रह्मलोकं न गच्छामि त्वामदृष्ट्वा प्रियातिथिम्} %||3-4-25||

\threelineshloka
{समागम्य गमिष्यामि त्रिदिवं देवसेवितम्}
{अक्षया नरशार्दूल जिता लोका मया शुभाः}
{ब्राह्म्याश्च नाकपृष्ठ्याश्च प्रतिगृह्णीष्व मामकान्} %||3-4-26||

\twolineshloka
{एवमुक्तो नरव्याघ्रः सर्वशास्त्रविशारदः}
{ऋषिणा शरभङ्गेन राघवो वाक्यमब्रवीत्} %||3-4-27||

\twolineshloka
{अहमेवाहरिष्यामि सर्वाँल्लोकान्महामुने}
{आवासं त्वहमिच्छामि प्रदिष्टमिह कानने} %||3-4-28||

\twolineshloka
{राघवेणैवमुक्तस्तु शक्रतुल्यबलेन वै}
{शरभङ्गो महाप्राज्ञः पुनरेवाब्रवीद्वचः} %||3-4-29||

\twolineshloka
{सुतीक्ष्णमभिगच्छ त्वं शुचौ देशे तपस्विनम्}
{रमणीये वनोद्देशे स ते वासं विधास्यति} %||3-4-30||

\twolineshloka
{एष पन्था नरव्याघ्र मुहूर्तं पश्य तात माम्}
{यावज्जहामि गात्राणि जीर्णं त्वचमिवोरगः} %||3-4-31||

\twolineshloka
{ततोऽग्निं स समाधाय हुत्वा चाज्येन मन्त्रवित्}
{शरभङ्गो महातेजाः प्रविवेश हुताशनम्} %||3-4-32||

\twolineshloka
{तस्य रोमाणि केशांश्च ददाहाग्निर्महात्मनः}
{जीर्णं त्वचं तथास्थीनि यच्च मांसं च शोणितम्} %||3-4-33||

\twolineshloka
{स च पावकसङ्काशः कुमारः समपद्यत}
{उत्थायाग्निचयात्तस्माच्छरभङ्गो व्यरोचत} %||3-4-34||

\twolineshloka
{स लोकानाहिताग्नीनामृषीणां च महात्मनाम्}
{देवानां च व्यतिक्रम्य ब्रह्मलोकं व्यरोहत} %||3-4-35||

\fourlineindentedshloka
{स पुण्यकर्मा भुवने द्विजर्षभः}
{ पितामहं सानुचरं ददर्श ह}
{पितामहश्चापि समीक्ष्य तं द्विजम्}
{ ननन्द सुस्वागतमित्युवाच ह} %||3-5-36||


{॥इत्यार्षे श्रीमद्रामायणे वाल्मीकीये आदिकाव्ये अरण्यकाण्डे चतुर्थः सर्गः॥}

\sect{पञ्चमः सर्गः}

\twolineshloka
{शरभङ्गे दिवं प्राप्ते मुनिसङ्घाः समागताः}
{अभ्यगच्छन्त काकुत्स्थं रामं ज्वलिततेजसम्} %||3-5-1||

\twolineshloka
{वैखानसा वालखिल्याः सम्प्रक्षाला मरीचिपाः}
{अश्मकुट्टाश्च बहवः पत्राहाराश्च तापसाः} %||3-5-2||

\twolineshloka
{दन्तोलूखलिनश्चैव तथैवोन्मज्जकाः परे}
{मुनयः सलिलाहारा वायुभक्षास्तथापरे} %||3-5-3||

\twolineshloka
{आकाशनिलयाश्चैव तथा स्थण्डिलशायिनः}
{तथोर्ध्ववासिनो दान्तास्तथार्द्रपटवाससः} %||3-5-4||

\threelineshloka
{सजपाश्च तपोनित्यास्तथा पञ्चतपोऽन्विताः}
{सर्वे ब्राह्म्या श्रिया जुष्टा दृढयोगसमाहिताः}
{शरभङ्गाश्रमे राममभिजग्मुश्च तापसाः} %||3-5-5||

\twolineshloka
{अभिगम्य च धर्मज्ञा रामं धर्मभृतां वरम्}
{ऊचुः परमधर्मज्ञमृषिसङ्घाः समाहिताः} %||3-5-6||

\twolineshloka
{त्वमिक्ष्वाकुकुलस्यास्य पृथिव्याश्च महारथः}
{प्रधानश्चासि नाथश्च देवानां मघवानिव} %||3-5-7||

\twolineshloka
{विश्रुतस्त्रिषु लोकेषु यशसा विक्रमेण च}
{पितृव्रतत्वं सत्यं च त्वयि धर्मश्च पुष्कलः} %||3-5-8||

\twolineshloka
{त्वामासाद्य महात्मानं धर्मज्ञं धर्मवत्सलम्}
{अर्थित्वान्नाथ वक्ष्यामस्तच्च नः क्षन्तुमर्हसि} %||3-5-9||

\twolineshloka
{अधार्मस्तु महांस्तात भवेत्तस्य महीपतेः}
{यो हरेद्बलिषड्भागं न च रक्षति पुत्रवत्} %||3-5-10||

\twolineshloka
{युञ्जानः स्वानिव प्राणान्प्राणैरिष्टान्सुतानिव}
{नित्ययुक्तः सदा रक्षन्सर्वान्विषयवासिनः} %||3-5-11||

\twolineshloka
{प्राप्नोति शाश्वतीं राम कीर्तिं स बहुवार्षिकीम्}
{ब्रह्मणः स्थानमासाद्य तत्र चापि महीयते} %||3-5-12||

\twolineshloka
{यत्करोति परं धर्मं मुनिर्मूलफलाशनः}
{तत्र राज्ञश्चतुर्भागः प्रजा धर्मेण रक्षतः} %||3-5-13||

\twolineshloka
{सोऽयं ब्राह्मणभूयिष्ठो वानप्रस्थगणो महान्}
{त्वन्नाथोऽनाथवद्राम राक्षसैर्वध्यते भृशम्} %||3-5-14||

\twolineshloka
{एहि पश्य शरीराणि मुनीनां भावितात्मनाम्}
{हतानां राक्षसैर्घोरैर्बहूनां बहुधा वने} %||3-5-15||

\twolineshloka
{पम्पानदीनिवासानामनुमन्दाकिनीमपि}
{चित्रकूटालयानां च क्रियते कदनं महत्} %||3-5-16||

\twolineshloka
{एवं वयं न मृष्यामो विप्रकारं तपस्विनम्}
{क्रियमाणं वने घोरं रक्षोभिर्भीमकर्मभिः} %||3-5-17||

\twolineshloka
{ततस्त्वां शरणार्थं च शरण्यं समुपस्थिताः}
{परिपालय नो राम वध्यमानान्निशाचरैः} %||3-5-18||

\threelineshloka
{एतच्छ्रुत्वा तु काकुत्स्थस्तापसानां तपस्विनाम्}
{इदं प्रोवाच धर्मात्मा सर्वानेव तपस्विनः}
{नैवमर्हथ मां वक्तुमाज्ञाप्योऽहं तपस्विनम्} %||3-5-19||

\threelineshloka
{भवतामर्थसिद्ध्यर्थमागतोऽहं यदृच्छया}
{तस्य मेऽयं वने वासो भविष्यति महाफलः}
{तपस्विनां रणे शत्रून्हन्तुमिच्छामि राक्षसान्} %||3-5-20||

\fourlineindentedshloka
{दत्त्वा वरं चापि तपोधनानाम्}
{ धर्मे धृतात्मा सहलक्ष्मणेन}
{तपोधनैश्चापि सहार्य वृत्तः}
{ सुतीष्क्णमेवाभिजगाम वीरः} %||3-6-21||


{॥इत्यार्षे श्रीमद्रामायणे वाल्मीकीये आदिकाव्ये अरण्यकाण्डे पञ्चमः सर्गः॥}

\sect{षष्ठः सर्गः}

\twolineshloka
{रामस्तु सहितो भ्रात्रा सीतया च परन्तपः}
{सुतीक्ष्णस्याश्रमपदं जगाम सह तैर्द्विजैः} %||3-6-1||

\twolineshloka
{स गत्वा दूरमध्वानं नदीस्तीर्त्व बहूदकाः}
{ददर्श विपुलं शैलं महामेघमिवोन्नतम्} %||3-6-2||

\twolineshloka
{ततस्तदिक्ष्वाकुवरौ सततं विविधैर्द्रुमैः}
{काननं तौ विविशतुः सीतया सह राघवौ} %||3-6-3||

\twolineshloka
{प्रविष्टस्तु वनं घोरं बहुपुष्पफलद्रुमम्}
{ददर्शाश्रममेकान्ते चीरमालापरिष्कृतम्} %||3-6-4||

\twolineshloka
{तत्र तापसमासीनं मलपङ्कजटाधरम्}
{रामः सुतीक्ष्णं विधिवत्तपोवृद्धमभाषत} %||3-6-5||

\twolineshloka
{रामोऽहमस्मि भगवन्भवन्तं द्रष्टुमागतः}
{तन्माभिवद धर्मज्ञ महर्षे सत्यविक्रम} %||3-6-6||

\twolineshloka
{स निरीक्ष्य ततो वीरं रामं धर्मभृतां वरम्}
{समाश्लिष्य च बाहुभ्यामिदं वचनमब्रवीत्} %||3-6-7||

\twolineshloka
{स्वागतं खलु ते वीर राम धर्मभृतां वर}
{आश्रमोऽयं त्वयाक्रान्तः सनाथ इव साम्प्रतम्} %||3-6-8||

\twolineshloka
{प्रतीक्षमाणस्त्वामेव नारोहेऽहं महायशः}
{देवलोकमितो वीर देहं त्यक्त्वा महीतले} %||3-6-9||

\threelineshloka
{चित्रकूटमुपादाय राज्यभ्रष्टोऽसि मे श्रुतः}
{इहोपयातः काकुत्स्थो देवराजः शतक्रतुः}
{सर्वाँल्लोकाञ्जितानाह मम पुण्येन कर्मणा} %||3-6-10||

\twolineshloka
{तेषु देवर्षिजुष्टेषु जितेषु तपसा मया}
{मत्प्रसादात्सभार्यस्त्वं विहरस्व सलक्ष्मणः} %||3-6-11||

\twolineshloka
{तमुग्रतपसं दीप्तं महर्षिं सत्यवादिनम्}
{प्रत्युवाचात्मवान्रामो ब्रह्माणमिव वासवः} %||3-6-12||

\twolineshloka
{अहमेवाहरिष्यामि स्वयं लोकान्महामुने}
{आवासं त्वहमिच्छामि प्रदिष्टमिह कानने} %||3-6-13||

\twolineshloka
{भवान्सर्वत्र कुशलः सर्वभूतहिते रतः}
{आख्यातः शरभङ्गेन गौतमेन महात्मना} %||3-6-14||

\twolineshloka
{एवमुक्तस्तु रामेण महर्षिर्लोकविश्रुतः}
{अब्रवीन्मधुरं वाक्यं हर्षेण महताप्लुतः} %||3-6-15||

\twolineshloka
{अयमेवाश्रमो राम गुणवान्रम्यतामिह}
{ऋषिसङ्घानुचरितः सदा मूलफलैर्युतः} %||3-6-16||

\twolineshloka
{इममाश्रममागम्य मृगसङ्घा महायशाः}
{अटित्वा प्रतिगच्छन्ति लोभयित्वाकुतोभयाः} %||3-6-17||

\twolineshloka
{तच्छ्रुत्वा वचनं तस्य महर्षेर्लक्ष्मणाग्रजः}
{उवाच वचनं धीरो विकृष्य सशरं धनुः} %||3-6-18||

\twolineshloka
{तानहं सुमहाभाग मृगसङ्घान्समागतान्}
{हन्यां निशितधारेण शरेणाशनिवर्चसा} %||3-6-19||

\twolineshloka
{भवांस्तत्राभिषज्येत किं स्यात्कृच्छ्रतरं ततः}
{एतस्मिन्नाश्रमे वासं चिरं तु न समर्थये} %||3-6-20||

\twolineshloka
{तमेवमुक्त्वा वरदं रामः सन्ध्यामुपागमत्}
{अन्वास्य पश्चिमां सन्ध्यां तत्र वासमकल्पयत्} %||3-6-21||

\fourlineindentedshloka
{ततः शुभं तापसभोज्यमन्नम्}
{ स्वयं सुतीक्ष्णः पुरुषर्षभाभ्याम्}
{ताभ्यां सुसत्कृत्य ददौ महात्मा}
{ सन्ध्यानिवृत्तौ रजनीं समीक्ष्य} %||3-7-22||


{॥इत्यार्षे श्रीमद्रामायणे वाल्मीकीये आदिकाव्ये अरण्यकाण्डे षष्ठः सर्गः॥}

\sect{सप्तमः सर्गः}

\twolineshloka
{रामस्तु सहसौमित्रिः सुतीक्ष्णेनाभिपूजितः}
{परिणम्य निशां तत्र प्रभाते प्रत्यबुध्यत} %||3-7-1||

\twolineshloka
{उत्थाय तु यथाकालं राघवः सह सीतया}
{उपास्पृशत्सुशीतेन जलेनोत्पलगन्धिना} %||3-7-2||

\twolineshloka
{अथ तेऽग्निं सुरांश्चैव वैदेही रामलक्ष्मणौ}
{काल्यं विधिवदभ्यर्च्य तपस्विशरणे वने} %||3-7-3||

\twolineshloka
{उदयन्न्तं दिनकरं दृष्ट्वा विगतकल्मषाः}
{सुतीक्ष्णमभिगम्येदं श्लक्ष्णं वचनमब्रुवन्} %||3-7-4||

\twolineshloka
{सुखोषिताः स्म भगवंस्त्वया पूज्येन पूजिताः}
{आपृच्छामः प्रयास्यामो मुनयस्त्वरयन्ति नः} %||3-7-5||

\twolineshloka
{त्वरामहे वयं द्रष्टुं कृत्स्नमाश्रममण्डलम्}
{ऋषीणां पुण्यशीलानां दण्डकारण्यवासिनाम्} %||3-7-6||

\twolineshloka
{अभ्यनुज्ञातुमिच्छामः सहैभिर्मुनिपुङ्गवैः}
{धर्मनित्यैस्तपोदान्तैर्विशिखैरिव पावकैः} %||3-7-7||

\twolineshloka
{अविषह्यातपो यावत्सूर्यो नातिविराजिते}
{अमार्गेणागतां लक्ष्मीं प्राप्येवान्वयवर्जितः} %||3-7-8||

\twolineshloka
{तावदिच्छामहे गन्तुमित्युक्त्वा चरणौ मुनेः}
{ववन्दे सहसौमित्रिः सीतया सह राघवः} %||3-7-9||

\twolineshloka
{तौ संस्पृशन्तौ चरणावुत्थाप्य मुनिपुङ्गवः}
{गाढमालिङ्ग्य सस्नेहमिदं वचनमब्रवीत्} %||3-7-10||

\twolineshloka
{अरिष्टं गच्छ पन्थानं राम सौमित्रिणा सह}
{सीतया चानया सार्धं छाययेवानुवृत्तया} %||3-7-11||

\twolineshloka
{पश्याश्रमपदं रम्यं दण्डकारण्यवासिनाम्}
{एषां तपस्विनां वीर तपसा भावितात्मनाम्} %||3-7-12||

\twolineshloka
{सुप्राज्यफलमूलानि पुष्पितानि वनानि च}
{प्रशान्तमृगयूथानि शान्तपक्षिगणानि च} %||3-7-13||

\twolineshloka
{फुल्लपङ्कजषडानि प्रसन्नसलिलानि च}
{कारण्डवविकीर्णानि तटाकानि सरांसि च} %||3-7-14||

\twolineshloka
{द्रक्ष्यसे दृष्टिरम्याणि गिरिप्रस्रवणानि च}
{रमणीयान्यरण्यानि मयूराभिरुतानि च} %||3-7-15||

\twolineshloka
{गम्यतां वत्स सौमित्रे भवानपि च गच्छतु}
{आगन्तव्यं च ते दृष्ट्वा पुनरेवाश्रमं मम} %||3-7-16||

\twolineshloka
{एवमुक्तस्तथेत्युक्त्वा काकुत्स्थः सहलक्ष्मणः}
{प्रदक्षिणं मुनिं कृता प्रस्थातुमुपचक्रमे} %||3-7-17||

\twolineshloka
{ततः शुभतरे तूणी धनुषी चायतेक्षणा}
{ददौ सीता तयोर्भ्रात्रोः खड्गौ च विमलौ ततः} %||3-7-18||

\twolineshloka
{आबध्य च शुभे तूणी चापे चादाय सस्वने}
{निष्क्रान्तावाश्रमाद्गन्तुमुभौ तौ रामलक्ष्मणौ} %||3-8-19||


{॥इत्यार्षे श्रीमद्रामायणे वाल्मीकीये आदिकाव्ये अरण्यकाण्डे सप्तमः सर्गः॥}

\sect{अष्टमः सर्गः}

\twolineshloka
{सुतीक्ष्णेनाभ्यनुज्ञातं प्रस्थितं रघुनन्दनम्}
{वैदेही स्निग्धया वाचा भर्तारमिदमब्रवीत्} %||3-8-1||

\twolineshloka
{अयं धर्मः सुसूक्ष्मेण विधिना प्राप्यते महान्}
{निवृत्तेन च शक्योऽयं व्यसनात्कामजादिह} %||3-8-2||

\threelineshloka
{त्रीण्येव व्यसनान्यत्र कामजानि भवन्त्युत}
{मिथ्या वाक्यं परमकं तस्माद्गुरुतरावुभौ}
{परदाराभिगमनं विना वैरं च रौद्रता} %||3-8-3||

\twolineshloka
{मिथ्यावाक्यं न ते भूतं न भविष्यति राघव}
{कुतोऽभिलषणं स्त्रीणां परेषां धर्मनाशनम्} %||3-8-4||

\twolineshloka
{तच्च सर्वं महाबाहो शक्यं वोढुं जितेन्द्रियैः}
{तव वश्येन्द्रियत्वं च जानामि शुभदर्शन} %||3-8-5||

\twolineshloka
{तृतीयं यदिदं रौद्रं परप्राणाभिहिंसनम्}
{निर्वैरं क्रियते मोहात्तच्च ते समुपस्थितम्} %||3-8-6||

\twolineshloka
{प्रतिज्ञातस्त्वया वीर दण्डकारण्यवासिनाम्}
{ऋषीणां रक्षणार्थाय वधः संयति रक्षसाम्} %||3-8-7||

\twolineshloka
{एतन्निमित्तं च वनं दण्डका इति विश्रुतम्}
{प्रस्थितस्त्वं सह भ्रात्रा धृतबाणशरासनः} %||3-8-8||

\twolineshloka
{ततस्त्वां प्रस्थितं दृष्ट्वा मम चिन्ताकुलं मनः}
{त्वद्वृत्तं चिन्तयन्त्या वै भवेन्निःश्रेयसं हितम्} %||3-8-9||

\twolineshloka
{न हि मे रोचते वीर गमनं दण्डकान्प्रति}
{कारणं तत्र वक्ष्यामि वदन्त्याः श्रूयतां मम} %||3-8-10||

\twolineshloka
{त्वं हि बाणधनुष्पाणिर्भ्रात्रा सह वनं गतः}
{दृष्ट्वा वनचरान्सर्वान्कच्चित्कुर्याः शरव्ययम्} %||3-8-11||

\twolineshloka
{क्षत्रियाणामिह धनुर्हुताशस्येन्धनानि च}
{समीपतः स्थितं तेजोबलमुच्छ्रयते भृशम्} %||3-8-12||

\twolineshloka
{पुरा किल महाबाहो तपस्वी सत्यवाक्शुचिः}
{कस्मिंश्चिदभवत्पुण्ये वने रतमृगद्विजे} %||3-8-13||

\twolineshloka
{तस्यैव तपसो विघ्नं कर्तुमिन्द्रः शचीपतिः}
{खड्गपाणिरथागच्छदाश्रमं भट रूपधृक्} %||3-8-14||

\twolineshloka
{तस्मिंस्तदाश्रमपदे निहितः खड्ग उत्तमः}
{स न्यासविधिना दत्तः पुण्ये तपसि तिष्ठतः} %||3-8-15||

\twolineshloka
{स तच्छस्त्रमनुप्राप्य न्यासरक्षणतत्परः}
{वने तु विचरत्येव रक्षन्प्रत्ययमात्मनः} %||3-8-16||

\twolineshloka
{यत्र गच्छत्युपादातुं मूलानि च फलानि च}
{न विना याति तं खड्गं न्यासरक्षणतत्परः} %||3-8-17||

\twolineshloka
{नित्यं शस्त्रं परिवहन्क्रमेण स तपोधनः}
{चकार रौद्रीं स्वां बुद्धिं त्यक्त्वा तपसि निश्चयम्} %||3-8-18||

\twolineshloka
{ततः स रौद्राभिरतः प्रमत्तोऽधर्मकर्षितः}
{तस्य शस्त्रस्य संवासाज्जगाम नरकं मुनिः} %||3-8-19||

\twolineshloka
{स्नेहाच्च बहुमानाच्च स्मारये त्वां न शिक्षये}
{न कथञ्चन सा कार्या हृहीतधनुषा त्वया} %||3-8-20||

\twolineshloka
{बुद्धिर्वैरं विना हन्तुं राक्षसान्दण्डकाश्रितान्}
{अपराधं विना हन्तुं लोकान्वीर न कामये} %||3-8-21||

\twolineshloka
{क्षत्रियाणां तु वीराणां वनेषु नियतात्मनाम्}
{धनुषा कार्यमेतावदार्तानामभिरक्षणम्} %||3-8-22||

\twolineshloka
{क्व च शस्त्रं क्व च वनं क्व च क्षात्रं तपः क्व च}
{व्याविद्धमिदमस्माभिर्देशधर्मस्तु पूज्यताम्} %||3-8-23||

\twolineshloka
{तदार्यकलुषा बुद्धिर्जायते शस्त्रसेवनात्}
{पुनर्गत्वा त्वयोध्यायां क्षत्रधर्मं चरिष्यसि} %||3-8-24||

\twolineshloka
{अक्षया तु भवेत्प्रीतिः श्वश्रू श्वशुरयोर्मम}
{यदि राज्यं हि संन्यस्य भवेस्त्वं निरतो मुनिः} %||3-8-25||

\twolineshloka
{धर्मादर्थः प्रभवति धर्मात्प्रभवते सुखम्}
{धर्मेण लभते सर्वं धर्मसारमिदं जगत्} %||3-8-26||

\twolineshloka
{आत्मानं नियमैस्तैस्तैः कर्षयित्वा प्रयत्नतः}
{प्राप्यते निपुणैर्धर्मो न सुखाल्लभ्यते सुखम्} %||3-8-27||

\twolineshloka
{नित्यं शुचिमतिः सौम्य चर धर्मं तपोवने}
{सर्वं हि विदितं तुभ्यं त्रैलोक्यमपि तत्त्वतः} %||3-8-28||

\fourlineindentedshloka
{स्त्रीचापलादेतदुदाहृतं मे}
{ धर्मं च वक्तुं तव कः समर्थः}
{विचार्य बुद्ध्या तु सहानुजेन}
{ यद्रोचते तत्कुरु माचिरेण} %||3-9-29||


{॥इत्यार्षे श्रीमद्रामायणे वाल्मीकीये आदिकाव्ये अरण्यकाण्डे अष्टमः सर्गः॥}

\sect{नवमः सर्गः}

\twolineshloka
{वाक्यमेतत्तु वैदेह्या व्याहृतं भर्तृभक्तया}
{श्रुत्वा धर्मे स्थितो रामः प्रत्युवाचाथ मैथिलीम्} %||3-9-1||

\twolineshloka
{हितमुक्तं त्वया देवि स्निग्धया सदृशं वचः}
{कुलं व्यपदिशन्त्या च धर्मज्ञे जनकात्मजे} %||3-9-2||

\twolineshloka
{किं तु वक्ष्याम्यहं देवि त्वयैवोक्तमिदं वचः}
{क्षत्रियैर्धार्यते चापो नार्तशब्दो भवेदिति} %||3-9-3||

\twolineshloka
{ते चार्ता दण्डकारण्ये मुनयः संशितव्रताः}
{मां सीते स्वयमागम्य शरण्याः शरणं गताः} %||3-9-4||

\twolineshloka
{वसन्तो धर्मनिरता वने मूलफलाशनाः}
{न लभन्ते सुखं भीता राक्षसैः क्रूरकर्मभिः} %||3-9-5||

\twolineshloka
{काले काले च निरता नियमैर्विविधैर्वने}
{भक्ष्यन्ते राक्षसैर्भीमैर्नरमांसोपजीविभिः} %||3-9-6||

\twolineshloka
{ते भक्ष्यमाणा मुनयो दण्डकारण्यवासिनः}
{अस्मानभ्यवपद्येति मामूचुर्द्विजसत्तमाः} %||3-9-7||

\twolineshloka
{मया तु वचनं श्रुत्वा तेषामेवं मुखाच्च्युतम्}
{कृत्वा चरणशुश्रूषां वाक्यमेतदुदाहृतम्} %||3-9-8||

\threelineshloka
{प्रसीदन्तु भवन्तो मे ह्रीरेषा हि ममातुला}
{यदीदृशैरहं विप्रैरुपस्थेयैरुपस्थितः}
{किं करोमीति च मया व्याहृतं द्विजसंनिधौ} %||3-9-9||

\threelineshloka
{सर्वैरेव समागम्य वागियं समुदाहृता}
{राक्षसैर्दण्डकारण्ये बहुभिः कामरूपिभिः}
{अर्दिताः स्म भृशं राम भवान्नस्त्रातुमर्हति} %||3-9-10||

\twolineshloka
{होमकाले तु सम्प्राप्ते पर्वकालेषु चानघ}
{धर्षयन्ति स्म दुर्धर्षा राक्षसाः पिशिताशनाः} %||3-9-11||

\twolineshloka
{राक्षसैर्धर्षितानां च तापसानां तपस्विनाम्}
{गतिं मृगयमाणानां भवान्नः परमा गतिः} %||3-9-12||

\twolineshloka
{कामं तपः प्रभावेन शक्ता हन्तुं निशाचरान्}
{चिरार्जितं तु नेच्छामस्तपः खण्डयितुं वयम्} %||3-9-13||

\twolineshloka
{बहुविघ्नं तपोनित्यं दुश्चरं चैव राघव}
{तेन शापं न मुञ्चामो भक्ष्यमाणाश्च राक्षसैः} %||3-9-14||

\twolineshloka
{तदर्द्यमानान्रक्षोभिर्दण्डकारण्यवासिभिः}
{रक्षनस्त्वं सह भ्रात्रा त्वन्नाथा हि वयं वने} %||3-9-15||

\twolineshloka
{मया चैतद्वचः श्रुत्वा कार्त्स्न्येन परिपालनम्}
{ऋषीणां दण्डकारण्ये संश्रुतं जनकात्मजे} %||3-9-16||

\twolineshloka
{संश्रुत्य च न शक्ष्यामि जीवमानः प्रतिश्रवम्}
{मुनीनामन्यथा कर्तुं सत्यमिष्टं हि मे सदा} %||3-9-17||

\twolineshloka
{अप्यहं जीवितं जह्यां त्वां वा सीते सलक्ष्मणाम्}
{न तु प्रतिज्ञां संश्रुत्य ब्राह्मणेभ्यो विशेषतः} %||3-9-18||

\twolineshloka
{तदवश्यं मया कार्यमृषीणां परिपालनम्}
{अनुक्तेनापि वैदेहि प्रतिज्ञाय तु किं पुनः} %||3-9-19||

\threelineshloka
{मम स्नेहाच्च सौहार्दादिदमुक्तं त्वया वचः}
{परितुष्टोऽस्म्यहं सीते न ह्यनिष्टोऽनुशिष्यते}
{सदृशं चानुरूपं च कुलस्य तव शोभने} %||3-9-20||

\fourlineindentedshloka
{इत्येवमुक्त्वा वचनं महात्मा}
{ सीतां प्रियां मैथिल राजपुत्रीम्}
{रामो धनुष्मान्सहलक्ष्मणेन}
{ जगाम रम्याणि तपोवनानि} %||3-10-21||


{॥इत्यार्षे श्रीमद्रामायणे वाल्मीकीये आदिकाव्ये अरण्यकाण्डे नवमः सर्गः॥}

\sect{दशमः सर्गः}

\twolineshloka
{अग्रतः प्रययौ रामः सीता मध्ये सुमध्यमा}
{पृष्ठतस्तु धनुष्पाणिर्लक्ष्मणोऽनुजगाम ह} %||3-10-1||

\twolineshloka
{तौ पश्यमानौ विविधाञ्शैलप्रस्थान्वनानि च}
{नदीश्च विविधा रम्या जग्मतुः सह सीतया} %||3-10-2||

\twolineshloka
{सारसांश्चक्रवाकांश्च नदीपुलिनचारिणः}
{सरांसि च सपद्मानि युतानि जलजैः खगैः} %||3-10-3||

\twolineshloka
{यूथबद्धांश्च पृषतान्मदोन्मत्तान्विषाणिनः}
{महिषांश्च वराहांश्च गजांश्च द्रुमवैरिणः} %||3-10-4||

\twolineshloka
{ते गत्वा दूरमध्वानं लम्बमाने दिवाकरे}
{ददृशुः सहिता रम्यं तटाकं योजनायतम्} %||3-10-5||

\twolineshloka
{पद्मपुष्करसम्बाधं गजयूथैरलङ्कृतम्}
{सारसैर्हंसकादम्बैः सङ्कुलं जलचारिभिः} %||3-10-6||

\twolineshloka
{प्रसन्नसलिले रम्यतस्मिन्सरसि शुश्रुवे}
{गीतवादित्रनिर्घोषो न तु कश्चन दृश्यते} %||3-10-7||

\twolineshloka
{ततः कौतूहलाद्रामो लक्ष्मणश्च महारथः}
{मुनिं धर्मभृतं नाम प्रष्टुं समुपचक्रमे} %||3-10-8||

\twolineshloka
{इदमत्यद्भुतं श्रुत्वा सर्वेषां नो महामुने}
{कौतूहलं महज्जातं किमिदं साधु कथ्यताम्} %||3-10-9||

\twolineshloka
{तेनैवमुक्तो धर्मात्मा राघवेण मुनिस्तदा}
{प्रभावं सरसः कृत्स्नमाख्यातुमुपचक्रमे} %||3-10-10||

\twolineshloka
{इदं पञ्चाप्सरो नाम तटाकं सार्वकालिकम्}
{निर्मितं तपसा राम मुनिना माण्डकर्णिना} %||3-10-11||

\twolineshloka
{स हि तेपे तपस्तीव्रं माण्डकर्णिर्महामुनिः}
{दशवर्षसहस्राणि वायुभक्षो जलाश्रय} %||3-10-12||

\threelineshloka
{ततः प्रव्यथिताः सर्वे देवाः साग्निपुरोगमाः}
{अब्रुवन्वचनं सर्वे परस्पर समागताः}
{अस्मकं कस्यचित्स्थानमेष प्रार्थयते मुनिः} %||3-10-13||

\twolineshloka
{ततः कर्तुं तपोविघ्नं सर्वैर्देवैर्नियोजिताः}
{प्रधानाप्सरसः पञ्चविद्युच्चलितवर्चसः} %||3-10-14||

\twolineshloka
{अप्सरोभिस्ततस्ताभिर्मुनिर्दृष्टपरावरः}
{नीतो मदनवश्यत्वं सुराणां कार्यसिद्धये} %||3-10-15||

\twolineshloka
{ताश्चैवाप्सरसः पञ्चमुनेः पत्नीत्वमागताः}
{तटाके निर्मितं तासामस्मिन्नन्तर्हितं गृहम्} %||3-10-16||

\twolineshloka
{तत्रैवाप्सरसः पञ्चनिवसन्त्यो यथासुखम्}
{रमयन्ति तपोयोगान्मुनिं यौवनमास्थितम्} %||3-10-17||

\twolineshloka
{तासां सङ्क्रीडमानानामेष वादित्रनिःस्वनः}
{श्रूयते भूषणोन्मिश्रो गीतशब्दो मनोहरः} %||3-10-18||

\twolineshloka
{आश्चर्यमिति तस्यैतद्वचनं भावितात्मनः}
{राघवः प्रतिजग्राह सह भ्रात्रा महायशाः} %||3-10-19||

\twolineshloka
{एवं कथयमानस्य ददर्शाश्रममण्डलम्}
{कुशचीरपरिक्षिप्तं नानावृक्षसमावृतम्} %||3-10-20||

\twolineshloka
{प्रविश्य सह वैदेह्या लक्ष्मणेन च राघवः}
{तदा तस्मिन्स काकुत्स्थः श्रीमत्याश्रममण्डले} %||3-10-21||

\twolineshloka
{उषित्वा सुसुखं तत्र पूर्ज्यमानो महर्षिभिः}
{जगाम चाश्रमांस्तेषां पर्यायेण तपस्विनाम्} %||3-10-22||

\twolineshloka
{येषामुषितवान्पूर्वं सकाशे स महास्त्रवित्}
{क्वचित्परिदशान्मासानेकं संवत्सरं क्वचित्} %||3-10-23||

\twolineshloka
{क्वचिच्च चतुरो मासान्पञ्चषट्चापरान्क्वचित्}
{अपरत्राधिकान्मासानध्यर्धमधिकं क्वचित्} %||3-10-24||

\threelineshloka
{त्रीन्मासानष्टमासांश्च राघवो न्यवसत्सुखम्}
{तथा संवसतस्तस्य मुनीनामाश्रमेषु वै}
{रमतश्चानुकुल्येन ययुः संवत्सरा दश} %||3-10-25||

\twolineshloka
{परिसृत्य च धर्मज्ञो राघवः सह सीतया}
{सुतीक्ष्णस्याश्रमं श्रीमान्पुनरेवाजगाम ह} %||3-10-26||

\twolineshloka
{स तमाश्रममागम्य मुनिभिः प्रतिपूजितः}
{तत्रापि न्यवसद्रामः कञ्चित्कालमरिन्दमः} %||3-10-27||

\twolineshloka
{अथाश्रमस्थो विनयात्कदाचित्तं महामुनिम्}
{उपासीनः स काकुत्स्थः सुतीक्ष्णमिदमब्रवीत्} %||3-10-28||

\twolineshloka
{अस्मिन्नरण्ये भगवन्नगस्त्यो मुनिसत्तमः}
{वसतीति मया नित्यं कथाः कथयतां श्रुतम्} %||3-10-29||

\twolineshloka
{न तु जानामि तं देशं वनस्यास्य महत्तया}
{कुत्राश्रमपदं पुण्यं महर्षेस्तस्य धीमतः} %||3-10-30||

\twolineshloka
{प्रसादात्तत्र भवतः सानुजः सह सीतया}
{अगस्त्यमभिगच्छेयमभिवादयितुं मुनिम्} %||3-10-31||

\twolineshloka
{मनोरथो महानेष हृदि सम्परिवर्तते}
{यदहं तं मुनिवरं शुश्रूषेयमपि स्वयम्} %||3-10-32||

\twolineshloka
{इति रामस्य स मुनिः श्रुत्वा धर्मात्मनो वचः}
{सुतीक्ष्णः प्रत्युवाचेदं प्रीतो दशरथात्मजम्} %||3-10-33||

\twolineshloka
{अहमप्येतदेव त्वां वक्तुकामः सलक्ष्मणम्}
{अगस्त्यमभिगच्छेति सीतया सह राघव} %||3-10-34||

\twolineshloka
{दिष्ट्या त्विदानीमर्थेऽस्मिन्स्वयमेव ब्रवीषि माम्}
{अहमाख्यामि ते वत्स यत्रागस्त्यो महामुनिः} %||3-10-35||

\twolineshloka
{योजनान्याश्रमात्तात याहि चत्वारि वै ततः}
{दक्षिणेन महाञ्श्रीमानगस्त्यभ्रातुराश्रमः} %||3-10-36||

\twolineshloka
{स्थलप्राये वनोद्देशे पिप्पलीवनशोभिते}
{बहुपुष्पफले रम्ये नानाशकुनिनादिते} %||3-10-37||

\twolineshloka
{पद्मिन्यो विविधास्तत्र प्रसन्नसलिलाः शिवाः}
{हंसकारण्डवाकीर्णाश्चक्रवाकोपशोभिताः} %||3-10-38||

\twolineshloka
{तत्रैकां रजनीमुष्य प्रभाते राम गम्यताम्}
{दक्षिणां दिशमास्थाय वनखण्डस्य पार्श्वतः} %||3-10-39||

\threelineshloka
{तत्रागस्त्याश्रमपदं गत्वा योजनमन्तरम्}
{रमणीये वनोद्देशे बहुपादप संवृते}
{रंस्यते तत्र वैदेही लक्ष्मणश्च त्वया सह} %||3-10-40||

\threelineshloka
{स हि रम्यो वनोद्देशो बहुपादपसङ्कुलः}
{यदि बुद्धिः कृता द्रष्टुमगस्त्यं तं महामुनिम्}
{अद्यैव गमने बुद्धिं रोचयस्व महायशः} %||3-10-41||

\twolineshloka
{इति रामो मुनेः श्रुत्वा सह भ्रात्राभिवाद्य च}
{प्रतस्थेऽगस्त्यमुद्दिश्य सानुजः सह सीतया} %||3-10-42||

\twolineshloka
{पश्यन्वनानि चित्राणि पर्वपांश्चाभ्रसंनिभान्}
{सरांसि सरितश्चैव पथि मार्गवशानुगाः} %||3-10-43||

\twolineshloka
{सुतीक्ष्णेनोपदिष्टेन गत्वा तेन पथा सुखम्}
{इदं परमसंहृष्टो वाक्यं लक्ष्मणमब्रवीत्} %||3-10-44||

\twolineshloka
{एतदेवाश्रमपदं नूनं तस्य महात्मनः}
{अगस्त्यस्य मुनेर्भ्रातुर्दृश्यते पुण्यकर्मणः} %||3-10-45||

\twolineshloka
{यथा हीमे वनस्यास्य ज्ञाताः पथि सहस्रशः}
{संनताः फलभरेण पुष्पभारेण च द्रुमाः} %||3-10-46||

\twolineshloka
{पिप्पलीनां च पक्वानां वनादस्मादुपागतः}
{गन्धोऽयं पवनोत्क्षिप्तः सहसा कटुकोदयः} %||3-10-47||

\twolineshloka
{तत्र तत्र च दृश्यन्ते सङ्क्षिप्ताः काष्ठसञ्चयाः}
{लूनाश्च पथि दृश्यन्ते दर्भा वैदूर्यवर्चसः} %||3-10-48||

\twolineshloka
{एतच्च वनमध्यस्थं कृष्णाभ्रशिखरोपमम्}
{पावकस्याश्रमस्थस्य धूमाग्रं सम्प्रदृश्यते} %||3-10-49||

\twolineshloka
{विविक्तेषु च तीर्थेषु कृतस्नाना द्विजातयः}
{पुष्पोपहारं कुर्वन्ति कुसुमैः स्वयमार्जितैः} %||3-10-50||

\twolineshloka
{तत्सुतीक्ष्णस्य वचनं यथा सौम्य मया श्रुतम्}
{अगस्त्यस्याश्रमो भ्रातुर्नूनमेष भविष्यति} %||3-10-51||

\twolineshloka
{निगृह्य तरसा मृत्युं लोकानां हितकाम्यया}
{यस्य भ्रात्रा कृतेयं दिक्शरण्या पुण्यकर्मणा} %||3-10-52||

\twolineshloka
{इहैकदा किल क्रूरो वातापिरपि चेल्वलः}
{भ्रातरौ सहितावास्तां ब्राह्मणघ्नौ महासुरौ} %||3-10-53||

\twolineshloka
{धारयन्ब्राह्मणं रूपमिल्वलः संस्कृतं वदन्}
{आमन्त्रयति विप्रान्स श्राद्धमुद्दिश्य निर्घृणः} %||3-10-54||

\twolineshloka
{भ्रातरं संस्कृतं भ्राता ततस्तं मेषरूपिणम्}
{तान्द्विजान्भोजयामास श्राद्धदृष्टेन कर्मणा} %||3-10-55||

\twolineshloka
{ततो भुक्तवतां तेषां विप्राणामिल्वलोऽब्रवीत्}
{वातापे निष्क्रमस्वेति स्वरेण महता वदन्} %||3-10-56||

\twolineshloka
{ततो भ्रातुर्वचः श्रुत्वा वातापिर्मेषवन्नदन्}
{भित्त्वा भित्वा शरीराणि ब्राह्मणानां विनिष्पतत्} %||3-10-57||

\twolineshloka
{ब्राह्मणानां सहस्राणि तैरेवं कामरूपिभिः}
{विनाशितानि संहत्य नित्यशः पिशिताशनैः} %||3-10-58||

\twolineshloka
{अगस्त्येन तदा देवैः प्रार्थितेन महर्षिणा}
{अनुभूय किल श्राद्धे भक्षितः स महासुरः} %||3-10-59||

\twolineshloka
{ततः सम्पन्नमित्युक्त्वा दत्त्वा हस्तावसेचनम्}
{भ्रातरं निष्क्रमस्वेति इल्वलः सोऽभ्यभाषत} %||3-10-60||

\twolineshloka
{तं तथा भाषमाणं तु भ्रातरं विप्रघातिनम्}
{अब्रवीत्प्रहसन्धीमानगस्त्यो मुनिसत्तमः} %||3-10-61||

\twolineshloka
{कुतो निष्क्रमितुं शक्तिर्मया जीर्णस्य रक्षसः}
{भ्रातुस्ते मेष रूपस्य गतस्य यमसादनम्} %||3-10-62||

\twolineshloka
{अथ तस्य वचः श्रुत्वा भ्रातुर्निधनसंश्रितम्}
{प्रधर्षयितुमारेभे मुनिं क्रोधान्निशाचरः} %||3-10-63||

\twolineshloka
{सोऽभ्यद्रवद्द्विजेन्द्रं तं मुनिना दीप्ततेजसा}
{चक्षुषानलकल्पेन निर्दग्धो निधनं गतः} %||3-10-64||

\twolineshloka
{तस्यायमाश्रमो भ्रातुस्तटाकवनशोभितः}
{विप्रानुकम्पया येन कर्मेदं दुष्करं कृतम्} %||3-10-65||

\twolineshloka
{एवं कथयमानस्य तस्य सौमित्रिणा सह}
{रामस्यास्तं गतः सूर्यः सन्ध्याकालोऽभ्यवर्तत} %||3-10-66||

\twolineshloka
{उपास्य पश्चिमां सन्ध्यां सह भ्रात्रा यथाविधि}
{प्रविवेशाश्रमपदं तमृषिं चाभ्यवादयन्} %||3-10-67||

\twolineshloka
{सम्यक्प्रतिगृहीतस्तु मुनिना तेन राघवः}
{न्यवसत्तां निशामेकां प्राश्य मूलफलानि च} %||3-10-68||

\twolineshloka
{तस्यां रात्र्यां व्यतीतायां विमले सूर्यमण्डले}
{भ्रातरं तमगस्त्यस्य आमन्त्रयत राघवः} %||3-10-69||

\twolineshloka
{अभिवादये त्वा भगवन्सुखमध्युषितो निशाम्}
{आमन्त्रये त्वां गच्छामि गुरुं ते द्रष्टुमग्रजम्} %||3-10-70||

\twolineshloka
{गम्यतामिति तेनोक्तो जगाम रघुनन्दनः}
{यथोद्दिष्टेन मार्गेण वनं तच्चावलोकयन्} %||3-10-71||

\twolineshloka
{नीवारान्पनसांस्तालांस्तिमिशान्वञ्जुलान्धवान्}
{चिरिबिल्वान्मधूकांश्च बिल्वानपि च तिन्दुकान्} %||3-10-72||

\twolineshloka
{पुष्पितान्पुष्पिताग्राभिर्लताभिरनुवेष्टितान्}
{ददर्श रामः शतशस्तत्र कान्तारपादपान्} %||3-10-73||

\twolineshloka
{हस्तिहस्तैर्विमृदितान्वानरैरुपशोभितान्}
{मत्तैः शकुनिसङ्घैश्च शतशः प्रतिनादितान्} %||3-10-74||

\twolineshloka
{ततोऽब्रवीत्समीपस्थं रामो राजीवलोचनः}
{पृष्ठतोऽनुगतं वीरं लक्ष्मणं लक्ष्मिवर्धनम्} %||3-10-75||

\twolineshloka
{स्निग्धपत्रा यथा वृक्षा यथा क्षान्ता मृगद्विजाः}
{आश्रमो नातिदूरस्थो महर्षेर्भावितात्मनः} %||3-10-76||

\twolineshloka
{अगस्त्य इति विख्यातो लोके स्वेनैव कर्मणा}
{आश्रमो दृश्यते तस्य परिश्रान्त श्रमापहः} %||3-10-77||

\twolineshloka
{प्राज्यधूमाकुलवनश्चीरमालापरिष्कृतः}
{प्रशान्तमृगयूथश्च नानाशकुनिनादितः} %||3-10-78||

\twolineshloka
{निगृह्य तरसा मृत्युं लोकानां हितकाम्यया}
{दक्षिणा दिक्कृता येन शरण्या पुण्यकर्मणा} %||3-10-79||

\twolineshloka
{तस्येदमाश्रमपदं प्रभावाद्यस्य राक्षसैः}
{दिगियं दक्षिणा त्रासाद्दृश्यते नोपभुज्यते} %||3-10-80||

\twolineshloka
{यदा प्रभृति चाक्रान्ता दिगियं पुण्यकर्मणा}
{तदा प्रभृति निर्वैराः प्रशान्ता रजनीचराः} %||3-10-81||

\twolineshloka
{नाम्ना चेयं भगवतो दक्षिणा दिक्प्रदक्षिणा}
{प्रथिता त्रिषु लोकेषु दुर्धर्षा क्रूरकर्मभिः} %||3-10-82||

\twolineshloka
{मार्गं निरोद्धुं सततं भास्करस्याचलोत्तमः}
{सन्देशं पालयंस्तस्य विन्ध्यशौलो न वर्धते} %||3-10-83||

\twolineshloka
{अयं दीर्घायुषस्तस्य लोके विश्रुतकर्मणः}
{अगस्त्यस्याश्रमः श्रीमान्विनीतमृगसेवितः} %||3-10-84||

\twolineshloka
{एष लोकार्चितः साधुर्हिते नित्यं रतः सताम्}
{अस्मानधिगतानेष श्रेयसा योजयिष्यति} %||3-10-85||

\twolineshloka
{आराधयिष्याम्यत्राहमगस्त्यं तं महामुनिम्}
{शेषं च वनवासस्य सौम्य वत्स्याम्यहं प्रभो} %||3-10-86||

\twolineshloka
{अत्र देवाः सगन्धर्वाः सिद्धाश्च परमर्षयः}
{अगस्त्यं नियताहारं सततं पर्युपासते} %||3-10-87||

\twolineshloka
{नात्र जीवेन्मृषावादी क्रूरो वा यदि वा शठः}
{नृशंसः कामवृत्तो वा मुनिरेष तथाविधः} %||3-10-88||

\twolineshloka
{अत्र देवाश्च यक्षाश्च नागाश्च पतगैः सह}
{वसन्ति नियताहारा धर्ममाराधयिष्णवः} %||3-10-89||

\twolineshloka
{अत्र सिद्धा महात्मानो विमानैः सूर्यसंनिभैः}
{त्यक्त्वा देहान्नवैर्देहैः स्वर्याताः परमर्षयः} %||3-10-90||

\twolineshloka
{यक्षत्वममरत्वं च राज्यानि विविधानि च}
{अत्र देवाः प्रयच्छन्ति भूतैराराधिताः शुभैः} %||3-10-91||

\twolineshloka
{आगताः स्माश्रमपदं सौमित्रे प्रविशाग्रतः}
{निवेदयेह मां प्राप्तमृषये सह सीतया} %||3-11-92||


{॥इत्यार्षे श्रीमद्रामायणे वाल्मीकीये आदिकाव्ये अरण्यकाण्डे दशमः सर्गः॥}

\sect{एकादशः सर्गः}

\twolineshloka
{स प्रविश्याश्रमपदं लक्ष्मणो राघवानुजः}
{अगस्त्यशिष्यमासाद्य वाक्यमेतदुवाच ह} %||3-11-1||

\twolineshloka
{राजा दशरथो नाम ज्येष्ठस्तस्य सुतो बली}
{रामः प्राप्तो मुनिं द्रष्टुं भार्यया सह सीतया} %||3-11-2||

\twolineshloka
{लक्ष्मणो नाम तस्याहं भ्राता त्ववरजो हितः}
{अनुकूलश्च भक्तश्च यदि ते श्रोत्रमागतः} %||3-11-3||

\twolineshloka
{ते वयं वनमत्युग्रं प्रविष्टाः पितृशासनात्}
{द्रष्टुमिच्छामहे सर्वे भगवन्तं निवेद्यताम्} %||3-11-4||

\twolineshloka
{तस्य तद्वचनं श्रुत्वा लक्ष्मणस्य तपोधनः}
{तथेत्युक्त्वाग्निशरणं प्रविवेश निवेदितुम्} %||3-11-5||

\twolineshloka
{स प्रविश्य मुनिश्रेष्ठं तपसा दुष्प्रधर्षणम्}
{कृताञ्जलिरुवाचेदं रामागमनमञ्जसा} %||3-11-6||

\twolineshloka
{पुत्रौ दशरथस्येमौ रामो लक्ष्मण एव च}
{प्रविष्टावाश्रमपदं सीतया सह भार्यया} %||3-11-7||

\twolineshloka
{द्रष्टुं भवन्तमायातौ शुश्रूषार्थमरिन्दमौ}
{यदत्रानन्तरं तत्त्वमाज्ञापयितुमर्हसि} %||3-11-8||

\twolineshloka
{ततः शिष्यादुपश्रुत्य प्राप्तं रामं सलक्ष्मणम्}
{वैदेहीं च महाभागामिदं वचनमब्रवीत्} %||3-11-9||

\twolineshloka
{दिष्ट्या रामश्चिरस्याद्य द्रष्टुं मां समुपागतः}
{मनसा काङ्क्षितं ह्यस्य मयाप्यागमनं प्रति} %||3-11-10||

\twolineshloka
{गम्यतां सत्कृतो रामः सभार्यः सहलक्ष्मणः}
{प्रवेश्यतां समीपं मे किं चासौ न प्रवेशितः} %||3-11-11||

\twolineshloka
{एवमुक्तस्तु मुनिना धर्मज्ञेन महात्मना}
{अभिवाद्याब्रवीच्छिष्यस्तथेति नियताञ्जलिः} %||3-11-12||

\twolineshloka
{ततो निष्क्रम्य सम्भ्रान्तः शिष्यो लक्ष्मणमब्रवीत्}
{क्वासौ रामो मुनिं द्रष्टुमेतु प्रविशतु स्वयम्} %||3-11-13||

\twolineshloka
{ततो गत्वाश्रमपदं शिष्येण सह लक्ष्मणः}
{दर्शयामास काकुत्स्थं सीतां च जनकात्मजाम्} %||3-11-14||

\twolineshloka
{तं शिष्यः प्रश्रितं वाक्यमगस्त्यवचनं ब्रुवन्}
{प्रावेशयद्यथान्यायं सत्कारार्थं सुसत्कृतम्} %||3-11-15||

\twolineshloka
{प्रविवेश ततो रामः सीतया सहलक्ष्मणः}
{प्रशान्तहरिणाकीर्णमाश्रमं ह्यवलोकयन्} %||3-11-16||

\twolineshloka
{स तत्र ब्रह्मणः स्थानमग्नेः स्थानं तथैव च}
{विष्णोः स्थानं महेन्द्रस्य स्थानं चैव विवस्वतः} %||3-11-17||

\twolineshloka
{सोमस्थानं भगस्थानं स्थानं कौबेरमेव च}
{धातुर्विधातुः स्थानं च वायोः स्थानं तथैव च} %||3-11-18||

\threelineshloka
{ततः शिष्यैः परिवृतो मुनिरप्यभिनिष्पतत्}
{तं ददर्शाग्रतो रामो मुनीनां दीप्ततेजसम्}
{अब्रवीद्वचनं वीरो लक्ष्मणं लक्ष्मिवर्धनम्} %||3-11-19||

\twolineshloka
{एष लक्ष्मण निष्क्रामत्यगस्त्यो भगवानृषिः}
{औदार्येणावगच्छामि निधानं तपसामिमम्} %||3-11-20||

\twolineshloka
{एवमुक्त्वा महाबाहुरगस्त्यं सूर्यवर्चसम्}
{जग्राह परमप्रीतस्तस्य पादौ परन्तपः} %||3-11-21||

\twolineshloka
{अभिवाद्य तु धर्मात्मा तस्थौ रामः कृताञ्जलिः}
{सीतया सह वैदेह्या तदा राम सलक्ष्मणः} %||3-11-22||

\twolineshloka
{प्रतिगृह्य च काकुत्स्थमर्चयित्वासनोदकैः}
{कुशलप्रश्नमुक्त्वा च आस्यतामिति सोऽब्रवीत्} %||3-11-23||

\twolineshloka
{अग्निं हुत्वा प्रदायार्घ्यमतिथिं प्रतिपूज्य च}
{वानप्रस्थेन धर्मेण स तेषां भोजनं ददौ} %||3-11-24||

\twolineshloka
{प्रथमं चोपविश्याथ धर्मज्ञो मुनिपुङ्गवः}
{उवाच राममासीनं प्राञ्जलिं धर्मकोविदम्} %||3-11-25||

\twolineshloka
{अन्यथा खलु काकुत्स्थ तपस्वी समुदाचरन्}
{दुःसाक्षीव परे लोके स्वानि मांसानि भक्षयेत्} %||3-11-26||

\twolineshloka
{राजा सर्वस्य लोकस्य धर्मचारी महारथः}
{पूजनीयश्च मान्यश्च भवान्प्राप्तः प्रियातिथिः} %||3-11-27||

\twolineshloka
{एवमुक्त्वा फलैर्मूलैः पुष्पैश्चान्यैश्च राघवम्}
{पूजयित्वा यथाकामं पुनरेव ततोऽब्रवीत्} %||3-11-28||

\twolineshloka
{इदं दिव्यं महच्चापं हेमवज्रविभूषितम्}
{वैष्णवं पुरुषव्याघ्र निर्मितं विश्वकर्मणा} %||3-11-29||

\twolineshloka
{अमोघः सूर्यसङ्काशो ब्रह्मदत्तः शरोत्तमः}
{दत्तो मम महेन्द्रेण तूणी चाक्षयसायकौ} %||3-11-30||

\twolineshloka
{सम्पूर्णौ निशितैर्बाणैर्ज्वलद्भिरिव पावकैः}
{महाराजत कोशोऽयमसिर्हेमविभूषितः} %||3-11-31||

\twolineshloka
{अनेन धनुषा राम हत्वा सङ्ख्ये महासुरान्}
{आजहार श्रियं दीप्तां पुरा विष्णुर्दिवौकसाम्} %||3-11-32||

\twolineshloka
{तद्धनुस्तौ च तूणीरौ शरं खड्गं च मानद}
{जयाय प्रतिगृह्णीष्व वज्रं वज्रधरो यथा} %||3-11-33||

\twolineshloka
{एवमुक्त्वा महातेजाः समस्तं तद्वरायुधम्}
{दत्त्वा रामाय भगवानगस्त्यः पुनरब्रवीत्} %||3-12-34||


{॥इत्यार्षे श्रीमद्रामायणे वाल्मीकीये आदिकाव्ये अरण्यकाण्डे एकादशः सर्गः॥}

\sect{द्वादशः सर्गः}

\twolineshloka
{राम प्रीतोऽस्मि भद्रं ते परितुष्टोऽस्मि लक्ष्मण}
{अभिवादयितुं यन्मां प्राप्तौ स्थः सह सीतया} %||3-12-1||

\twolineshloka
{अध्वश्रमेण वां खेदो बाधते प्रचुरश्रमः}
{व्यक्तमुत्कण्ठते चापि मैथिली जनकात्मजा} %||3-12-2||

\twolineshloka
{एषा हि सुकुमारी च दुःखैश्च न विमानिता}
{प्राज्यदोषं वनं प्रप्ता भर्तृस्नेहप्रचोदिता} %||3-12-3||

\twolineshloka
{यथैषा रमते राम इह सीता तथा कुरु}
{दुष्करं कृतवत्येषा वने त्वामनुगच्छती} %||3-12-4||

\twolineshloka
{एषा हि प्रकृतिः स्त्रीणामासृष्टे रघुनन्दन}
{समस्थमनुरज्यन्ते विषमस्थं त्यजन्ति च} %||3-12-5||

\twolineshloka
{शतह्रदानां लोलत्वं शस्त्राणां तीक्ष्णतां तथा}
{गरुडानिलयोः शैघ्र्यमनुगच्छन्ति योषितः} %||3-12-6||

\twolineshloka
{इयं तु भवतो भार्या दोषैरेतैर्विवर्जिताः}
{श्लाघ्या च व्यपदेश्या च यथा देवी ह्यरुन्धती} %||3-12-7||

\twolineshloka
{अलङ्कृतोऽयं देशश्च यत्र सौमित्रिणा सह}
{वैदेह्या चानया राम वत्स्यसि त्वमरिन्दम} %||3-12-8||

\twolineshloka
{एवमुक्तस्तु मुनिना राघवः संयताञ्जलिः}
{उवाच प्रश्रितं वाक्यमृषिं दीप्तमिवानलम्} %||3-12-9||

\twolineshloka
{धन्योऽस्म्यनुगृहीतोऽस्मि यस्य मे मुनिपुङ्गवः}
{गुणैः सभ्रातृभार्यस्य वरदः परितुष्यति} %||3-12-10||

\twolineshloka
{किं तु व्यादिश मे देशं सोदकं बहुकाननम्}
{यत्राश्रमपदं कृत्वा वसेयं निरतः सुखम्} %||3-12-11||

\twolineshloka
{ततोऽब्रवीन्मुनि श्रेष्ठः श्रुत्वा रामस्य भाषितम्}
{ध्यात्वा मुहूर्तं धर्मात्मा धीरो धीरतरं वचः} %||3-12-12||

\twolineshloka
{इतो द्वियोजने तात बहुमूलफलोदकः}
{देशो बहुमृगः श्रीमान्पञ्चवट्यभिविश्रुतः} %||3-12-13||

\twolineshloka
{तत्र गत्वाश्रमपदं कृत्वा सौमित्रिणा सह}
{रमस्व त्वं पितुर्वाक्यं यथोक्तमनुपालयन्} %||3-12-14||

\twolineshloka
{विदितो ह्येष वृत्तान्तो मम सर्वस्तवानघ}
{तपसश्च प्रभावेन स्नेहाद्दशरथस्य च} %||3-12-15||

\twolineshloka
{हृदयस्थश्च ते छन्दो विज्ञातस्तपसा मया}
{इह वासं प्रतिज्ञाय मया सह तपोवने} %||3-12-16||

\twolineshloka
{अतश्च त्वामहं ब्रूमि गच्छ पञ्चवटीमिति}
{स हि रम्यो वनोद्देशो मैथिली तत्र रंस्यते} %||3-12-17||

\twolineshloka
{स देशः श्लाघनीयश्च नातिदूरे च राघव}
{गोदावर्याः समीपे च मैथिली तत्र रंस्यते} %||3-12-18||

\twolineshloka
{प्राज्यमूलफलैश्चैव नानाद्विज गणैर्युतः}
{विविक्तश्च महाबाहो पुण्यो रम्यस्तथैव च} %||3-12-19||

\twolineshloka
{भवानपि सदारश्च शक्तश्च परिरक्षणे}
{अपि चात्र वसन्रामस्तापसान्पालयिष्यसि} %||3-12-20||

\twolineshloka
{एतदालक्ष्यते वीर मधुकानां महद्वनम्}
{उत्तरेणास्य गन्तव्यं न्यग्रोधमभिगच्छता} %||3-12-21||

\twolineshloka
{ततः स्थलमुपारुह्य पर्वतस्याविदूरतः}
{ख्यातः पञ्चवटीत्येव नित्यपुष्पितकाननः} %||3-12-22||

\twolineshloka
{अगस्त्येनैवमुक्तस्तु रामः सौमित्रिणा सह}
{सात्कृत्यामन्त्रयामास तमृषिं सत्यवादिनम्} %||3-12-23||

\twolineshloka
{तौ तु तेनाभ्यनुज्ञातौ कृतपादाभिवन्दनौ}
{तदाश्रमात्पञ्चवटीं जग्मतुः सह सीतया} %||3-12-24||

\fourlineindentedshloka
{गृहीतचापौ तु नराधिपात्मजौ}
{ विषक्ततूणी समरेष्वकातरौ}
{यथोपदिष्टेन पथा महर्षिणा}
{ प्रजग्मतुः पञ्चवटीं समाहितौ} %||3-13-25||


{॥इत्यार्षे श्रीमद्रामायणे वाल्मीकीये आदिकाव्ये अरण्यकाण्डे द्वादशः सर्गः॥}

\sect{त्रयोदशः सर्गः}

\twolineshloka
{अथ पञ्चवटीं गच्छन्नन्तरा रघुनन्दनः}
{आससाद महाकायं गृध्रं भीमपराक्रमम्} %||3-13-1||

\twolineshloka
{तं दृष्ट्वा तौ महाभागौ वनस्थं रामलक्ष्मणौ}
{मेनाते राक्षसं पक्षिं ब्रुवाणौ को भवानिति} %||3-13-2||

\twolineshloka
{स तौ मधुरया वाचा सौम्यया प्रीणयन्निव}
{उवाच वत्स मां विद्धि वयस्यं पितुरात्मनः} %||3-13-3||

\twolineshloka
{स तं पितृसखं बुद्ध्वा पूजयामास राघवः}
{स तस्य कुलमव्यग्रमथ पप्रच्छ नाम च} %||3-13-4||

\twolineshloka
{रामस्य वचनं श्रुत्वा कुलमात्मानमेव च}
{आचचक्षे द्विजस्तस्मै सर्वभूतसमुद्भवम्} %||3-13-5||

\twolineshloka
{पूर्वकाले महाबाहो ये प्रजापतयोऽभवन्}
{तान्मे निगदतः सर्वानादितः शृणु राघव} %||3-13-6||

\twolineshloka
{कर्दमः प्रथमस्तेषां विकृतस्तदनन्तरम्}
{शेषश्च संश्रयश्चैव बहुपुत्रश्च वीर्यवान्} %||3-13-7||

\twolineshloka
{स्थाणुर्मरीचिरत्रिश्च क्रतुश्चैव महाबलः}
{पुलस्त्यश्चाङ्गिराश्चैव प्रचेताः पुलहस्तथा} %||3-13-8||

\twolineshloka
{दक्षो विवस्वानपरोऽरिष्टनेमिश्च राघव}
{कश्यपश्च महातेजास्तेषामासीच्च पश्चिमः} %||3-13-9||

\twolineshloka
{प्रजापतेस्तु दक्षस्य बभूवुरिति नः श्रुतम्}
{षष्टिर्दुहितरो राम यशस्विन्यो महायशः} %||3-13-10||

\twolineshloka
{कश्यपः प्रतिजग्राह तासामष्टौ सुमध्यमाः}
{अदितिं च दितिं चैव दनूमपि च कालकाम्} %||3-13-11||

\twolineshloka
{ताम्रां क्रोधवशां चैव मनुं चाप्यनलामपि}
{तास्तु कन्यास्ततः प्रीतः कश्यपः पुनरब्रवीत्} %||3-13-12||

\twolineshloka
{पुत्रांस्त्रैलोक्यभर्तॄन्वै जनयिष्यथ मत्समान्}
{अदितिस्तन्मना राम दितिश्च दनुरेव च} %||3-13-13||

\twolineshloka
{कालका च महाबाहो शेषास्त्वमनसोऽभवन्}
{अदित्यां जज्ञिरे देवास्त्रयस्त्रिंशदरिन्दम} %||3-13-14||

\twolineshloka
{आदित्या वसवो रुद्रा अश्विनौ च परन्तप}
{दितिस्त्वजनयत्पुत्रान्दैत्यांस्तात यशस्विनः} %||3-13-15||

\twolineshloka
{तेषामियं वसुमती पुरासीत्सवनार्णवा}
{दनुस्त्वजनयत्पुत्रमश्वग्रीवमरिन्दम} %||3-13-16||

\twolineshloka
{नरकं कालकं चैव कालकापि व्यजायत}
{क्रौञ्चीं भासीं तथा श्येनीं धृतराष्ट्रीं तथा शुकीम्} %||3-13-17||

\twolineshloka
{ताम्रापि सुषुवे कन्याः पञ्चैता लोकविश्रुताः}
{उलूकाञ्जनयत्क्रौञ्ची भासी भासान्व्यजायत} %||3-13-18||

\twolineshloka
{श्येनी श्येनांश्च गृध्रांश्च व्यजायत सुतेजसः}
{धृतराष्ट्री तु हंसांश्च कलहंसांश्च सर्वशः} %||3-13-19||

\twolineshloka
{चक्रवाकांश्च भद्रं ते विजज्ञे सापि भामिनी}
{शुकी नतां विजज्ञे तु नताया विनता सुता} %||3-13-20||

\twolineshloka
{दशक्रोधवशा राम विजज्ञेऽप्यात्मसम्भवाः}
{मृगीं च मृगमन्दां च हरीं भद्रमदामपि} %||3-13-21||

\twolineshloka
{मातङ्गीमथ शार्दूलीं श्वेतां च सुरभीं तथा}
{सर्वलक्षणसम्पन्नां सुरसां कद्रुकामपि} %||3-13-22||

\twolineshloka
{अपत्यं तु मृगाः सर्वे मृग्या नरवरोत्तम}
{ऋष्काश्च मृगमन्दायाः सृमराश्चमरास्तथा} %||3-13-23||

\twolineshloka
{ततस्त्विरावतीं नाम जज्ञे भद्रमदा सुताम्}
{तस्यास्त्वैरावतः पुत्रो लोकनाथो महागजः} %||3-13-24||

\twolineshloka
{हर्याश्च हरयोऽपत्यं वानराश्च तपस्विनः}
{गोलाङ्गूलांश्च शार्दूली व्याघ्रांश्चाजनयत्सुतान्} %||3-13-25||

\twolineshloka
{मातङ्ग्यास्त्वथ मातङ्गा अपत्यं मनुजर्षभ}
{दिशागजं तु श्वेताक्षं श्वेता व्यजनयत्सुतम्} %||3-13-26||

\twolineshloka
{ततो दुहितरौ राम सुरभिर्देव्यजायत}
{रोहिणीं नाम भद्रं ते गन्धर्वीं च यशस्विनीम्} %||3-13-27||

\twolineshloka
{रोहिण्यजनयद्गा वै गन्धर्वी वाजिनः सुतान्}
{सुरसाजनयन्नागान्राम कद्रूश्च पन्नगान्} %||3-13-28||

\twolineshloka
{मनुर्मनुष्याञ्जनयत्कश्यपस्य महात्मनः}
{ब्राह्मणान्क्षत्रियान्वैश्याञ्शूद्रांश्च मनुजर्षभ} %||3-13-29||

\twolineshloka
{मुखतो ब्राह्मणा जाता उरसः क्षत्रियास्तथा}
{ऊरुभ्यां जज्ञिरे वैश्याः पद्भ्यां शूद्रा इति श्रुतिः} %||3-13-30||

\twolineshloka
{सर्वान्पुण्यफलान्वृक्षाननलापि व्यजायत}
{विनता च शुकी पौत्री कद्रूश्च सुरसा स्वसा} %||3-13-31||

\twolineshloka
{कद्रूर्नागसहस्क्रं तु विजज्ञे धरणीधरम्}
{द्वौ पुत्रौ विनतायास्तु गरुडोऽरुण एव च} %||3-13-32||

\twolineshloka
{तस्माज्जातोऽहमरुणात्सम्पातिश्च ममाग्रजः}
{जटायुरिति मां विद्धि श्येनीपुत्रमरिन्दम} %||3-13-33||

\twolineshloka
{सोऽहं वाससहायस्ते भविष्यामि यदीच्छसि}
{सीतां च तात रक्षिष्ये त्वयि याते सलक्ष्मणे} %||3-13-34||

\fourlineindentedshloka
{जटायुषं तु प्रतिपूज्य राघवो}
{ मुदा परिष्वज्य च संनतोऽभवत्}
{पितुर्हि शुश्राव सखित्वमात्मवा}
{ञ्जटायुषा सङ्कथितं पुनः पुनः} %||3-13-35||

\fourlineindentedshloka
{स तत्र सीतां परिदाय मैथिलीम्}
{ सहैव तेनातिबलेन पक्षिणा}
{जगाम तां पञ्चवटीं सलक्ष्मणो}
{ रिपून्दिधक्षञ्शलभानिवानलः} %||3-14-36||


{॥इत्यार्षे श्रीमद्रामायणे वाल्मीकीये आदिकाव्ये अरण्यकाण्डे त्रयोदशः सर्गः॥}

\sect{चतुर्दशः सर्गः}

\twolineshloka
{ततः पञ्चवटीं गत्वा नानाव्यालमृगायुताम्}
{उवाच भ्रातरं रामो लक्ष्मणं दीप्ततेजसम्} %||3-14-1||

\twolineshloka
{आगताः स्म यथोद्दिष्टममुं देशं महर्षिणा}
{अयं पञ्चवटी देशः सौम्य पुष्पितकाननः} %||3-14-2||

\twolineshloka
{सर्वतश्चार्यतां दृष्टिः कानने निपुणो ह्यसि}
{आश्रमः कतरस्मिन्नो देशे भवति सम्मतः} %||3-14-3||

\twolineshloka
{रमते यत्र वैदेही त्वमहं चैव लक्ष्मण}
{तादृशो दृश्यतां देशः संनिकृष्टजलाशयः} %||3-14-4||

\twolineshloka
{वनरामण्यकं यत्र जलरामण्यकं तथा}
{संनिकृष्टं च यत्र स्यात्समित्पुष्पकुशोदकम्} %||3-14-5||

\twolineshloka
{एवमुक्तस्तु रामेण लक्मणः संयताञ्जलिः}
{सीता समक्षं काकुत्स्थमिदं वचनमब्रवीत्} %||3-14-6||

\twolineshloka
{परवानस्मि काकुत्स्थ त्वयि वर्षशतं स्थिते}
{स्वयं तु रुचिरे देशे क्रियतामिति मां वद} %||3-14-7||

\twolineshloka
{सुप्रीतस्तेन वाक्येन लक्ष्मणस्य महाद्युतिः}
{विमृशन्रोचयामास देशं सर्वगुणान्वितम्} %||3-14-8||

\twolineshloka
{स तं रुचिरमाक्रम्य देशमाश्रमकर्मणि}
{हस्ते गृहीत्वा हस्तेन रामः सौमित्रिमब्रवीत्} %||3-14-9||

\twolineshloka
{अयं देशः समः श्रीमान्पुष्पितैर्तरुभिर्वृतः}
{इहाश्रमपदं सौम्य यथावत्कर्तुमर्हसि} %||3-14-10||

\twolineshloka
{इयमादित्यसङ्काशैः पद्मैः सुरभिगन्धिभिः}
{अदूरे दृश्यते रम्या पद्मिनी पद्मशोभिता} %||3-14-11||

\twolineshloka
{यथाख्यातमगस्त्येन मुनिना भावितात्मना}
{इयं गोदावरी रम्या पुष्पितैस्तरुभिर्वृता} %||3-14-12||

\twolineshloka
{हंसकारण्डवाकीर्णा चक्रवाकोपशोभिता}
{नातिदूरे न चासन्ने मृगयूथनिपीडिता} %||3-14-13||

\twolineshloka
{मयूरनादिता रम्याः प्रांशवो बहुकन्दराः}
{दृश्यन्ते गिरयः सौम्य फुल्लैस्तरुभिरावृताः} %||3-14-14||

\twolineshloka
{सौवर्णे राजतैस्ताम्रैर्देशे देशे च धातुभिः}
{गवाक्षिता इवाभान्ति गजाः परमभक्तिभिः} %||3-14-15||

\twolineshloka
{सालैस्तालैस्तमालैश्च खर्जूरैः पनसाम्रकैः}
{नीवारैस्तिमिशैश्चैव पुंनागैश्चोपशोभिताः} %||3-14-16||

\twolineshloka
{चूतैरशोकैस्तिलकैश्चम्पकैः केतकैरपि}
{पुष्पगुल्मलतोपेतैस्तैस्तैस्तरुभिरावृताः} %||3-14-17||

\twolineshloka
{चन्दनैः स्यन्दनैर्नीपैः पनसैर्लकुचैरपि}
{धवाश्वकर्णखदिरैः शमीकिंशुकपाटलैः} %||3-14-18||

\twolineshloka
{इदं पुण्यमिदं मेध्यमिदं बहुमृगद्विजम्}
{इह वत्स्याम सौमित्रे सार्धमेतेन पक्षिणा} %||3-14-19||

\twolineshloka
{एवमुक्तस्तु रामेण लक्ष्मणः परवीरहा}
{अचिरेणाश्रमं भ्रातुश्चकार सुमहाबलः} %||3-14-20||

\twolineshloka
{पर्णशालां सुविपुलां तत्र सङ्घातमृत्तिकाम्}
{सुस्तम्भां मस्करैर्दीर्घैः कृतवंशां सुशोभनाम्} %||3-14-21||

\twolineshloka
{स गत्वा लक्ष्मणः श्रीमान्नदीं गोदावरीं तदा}
{स्नात्वा पद्मानि चादाय सफलः पुनरागतः} %||3-14-22||

\twolineshloka
{ततः पुष्पबलिं कृत्वा शान्तिं च स यथाविधि}
{दर्शयामास रामाय तदाश्रमपदं कृतम्} %||3-14-23||

\twolineshloka
{स तं दृष्ट्वा कृतं सौम्यमाश्रमं सह सीतया}
{राघवः पर्णशालायां हर्षमाहारयत्परम्} %||3-14-24||

\twolineshloka
{सुसंहृष्टः परिष्वज्य बाहुभ्यां लक्ष्मणं तदा}
{अतिस्निग्धं च गाढं च वचनं चेदमब्रवीत्} %||3-14-25||

\twolineshloka
{प्रीतोऽस्मि ते महत्कर्म त्वया कृतमिदं प्रभो}
{प्रदेयो यन्निमित्तं ते परिष्वङ्गो मया कृतः} %||3-14-26||

\twolineshloka
{भावज्ञेन कृतज्ञेन धर्मज्ञेन च लक्ष्मण}
{त्वया पुत्रेण धर्मात्मा न संवृत्तः पिता मम} %||3-14-27||

\twolineshloka
{एवं लक्ष्मणमुक्त्वा तु राघवो लक्ष्मिवर्धनः}
{तस्मिन्देशे बहुफले न्यवसत्स सुखं वशी} %||3-14-28||

\twolineshloka
{कञ्चित्कालं स धर्मात्मा सीतया लक्ष्मणेन च}
{अन्वास्यमानो न्यवसत्स्वर्गलोके यथामरः} %||3-15-29||


{॥इत्यार्षे श्रीमद्रामायणे वाल्मीकीये आदिकाव्ये अरण्यकाण्डे चतुर्दशः सर्गः॥}

\sect{पञ्चदशः सर्गः}

\twolineshloka
{वसतस्तस्य तु सुखं राघवस्य महात्मनः}
{शरद्व्यपाये हेमन्त ऋतुरिष्टः प्रवर्तते} %||3-15-1||

\twolineshloka
{स कदाचित्प्रभातायां शर्वर्यां रघुनन्दनः}
{प्रययावभिषेकार्थं रम्यां गोदावरीं नदीम्} %||3-15-2||

\twolineshloka
{प्रह्वः कलशहस्तस्तं सीतया सह वीर्यवान्}
{पृष्ठतोऽनुव्रजन्भ्राता सौमित्रिरिदमब्रवीत्} %||3-15-3||

\twolineshloka
{अयं स कालः सम्प्राप्तः प्रियो यस्ते प्रियंवद}
{अलङ्कृत इवाभाति येन संवत्सरः शुभः} %||3-15-4||

\twolineshloka
{नीहारपरुषो लोकः पृथिवी सस्यमालिनी}
{जलान्यनुपभोग्यानि सुभगो हव्यवाहनः} %||3-15-5||

\twolineshloka
{नवाग्रयणपूजाभिरभ्यर्च्य पितृदेवताः}
{कृताग्रयणकाः काले सन्तो विगतकल्मषाः} %||3-15-6||

\twolineshloka
{प्राज्यकामा जनपदाः सम्पन्नतरगोरसाः}
{विचरन्ति महीपाला यात्रार्थं विजिगीषवः} %||3-15-7||

\twolineshloka
{सेवमाने दृढं सूर्ये दिशमन्तकसेविताम्}
{विहीनतिलकेव स्त्री नोत्तरा दिक्प्रकाशते} %||3-15-8||

\twolineshloka
{प्रकृत्या हिमकोशाढ्यो दूरसूर्यश्च साम्प्रतम्}
{यथार्थनामा सुव्यक्तं हिमवान्हिमवान्गिरिः} %||3-15-9||

\twolineshloka
{अत्यन्तसुखसञ्चारा मध्याह्ने स्पर्शतः सुखाः}
{दिवसाः सुभगादित्याश्छायासलिलदुर्भगाः} %||3-15-10||

\twolineshloka
{मृदुसूर्याः सनीहाराः पटुशीताः समारुताः}
{शून्यारण्या हिमध्वस्ता दिवसा भान्ति साम्प्रतम्} %||3-15-11||

\twolineshloka
{निवृत्ताकाशशयनाः पुष्यनीता हिमारुणाः}
{शीता वृद्धतरायामास्त्रियामा यान्ति साम्प्रतम्} %||3-15-12||

\twolineshloka
{रविसङ्क्रान्तसौभाग्यस्तुषारारुणमण्डलः}
{निःश्वासान्ध इवादर्शश्चन्द्रमा न प्रकाशते} %||3-15-13||

\twolineshloka
{ज्योत्स्ना तुषारमलिना पौर्णमास्यां न राजते}
{सीतेव चातपश्यामा लक्ष्यते न तु शोभते} %||3-15-14||

\twolineshloka
{प्रकृत्या शीतलस्पर्शो हिमविद्धश्च साम्प्रतम्}
{प्रवाति पश्चिमो वायुः काले द्विगुणशीतलः} %||3-15-15||

\twolineshloka
{बाष्पच्छन्नान्यरण्यानि यवगोधूमवन्ति च}
{शोभन्तेऽभ्युदिते सूर्ये नदद्भिः क्रौञ्चसारसैः} %||3-15-16||

\twolineshloka
{खर्जूरपुष्पाकृतिभिः शिरोभिः पूर्णतण्डुलैः}
{शोभन्ते किं चिदालम्बाः शालयः कनकप्रभाः} %||3-15-17||

\twolineshloka
{मयूखैरुपसर्पद्भिर्हिमनीहारसंवृतैः}
{दूरमभ्युदितः सूर्यः शशाङ्क इव लक्ष्यते} %||3-15-18||

\twolineshloka
{अग्राह्यवीर्यः पूर्वाह्णे मध्याह्ने स्पर्शतः सुखः}
{संरक्तः किञ्चिदापाण्डुरातपः शोभते क्षितौ} %||3-15-19||

\twolineshloka
{अवश्यायनिपातेन किञ्चित्प्रक्लिन्नशाद्वला}
{वनानां शोभते भूमिर्निविष्टतरुणातपा} %||3-15-20||

\twolineshloka
{अवश्यायतमोनद्धा नीहारतमसावृताः}
{प्रसुप्ता इव लक्ष्यन्ते विपुष्पा वनराजयः} %||3-15-21||

\twolineshloka
{बाष्पसञ्छन्नसलिला रुतविज्ञेयसारसाः}
{हिमार्द्रवालुकैस्तीरैः सरितो भान्ति साम्प्रतम्} %||3-15-22||

\twolineshloka
{तुषारपतनाच्चैव मृदुत्वाद्भास्करस्य च}
{शैत्यादगाग्रस्थमपि प्रायेण रसवज्जलम्} %||3-15-23||

\twolineshloka
{जराजर्जरितैः पर्णैः शीर्णकेसरकर्णिकैः}
{नालशेषा हिमध्वस्ता न भान्ति कमलाकराः} %||3-15-24||

\twolineshloka
{अस्मिंस्तु पुरुषव्याघ्र काले दुःखसमन्वितः}
{तपश्चरति धर्मात्मा त्वद्भक्त्या भरतः पुरे} %||3-15-25||

\twolineshloka
{त्यक्त्वा राज्यं च मानं च भोगांश्च विविधान्बहून्}
{तपस्वी नियताहारः शेते शीते महीतले} %||3-15-26||

\twolineshloka
{सोऽपि वेलामिमां नूनमभिषेकार्थमुद्यतः}
{वृतः प्रकृतिभिर्नित्यं प्रयाति सरयूं नदीम्} %||3-15-27||

\twolineshloka
{अत्यन्तसुखसंवृद्धः सुकुमारो हिमार्दितः}
{कथं त्वपररात्रेषु सरयूमवगाहते} %||3-15-28||

\twolineshloka
{पद्मपत्रेक्षणः श्यामः श्रीमान्निरुदरो महान्}
{धर्मज्ञः सत्यवादी च ह्रीनिषेधो जितेन्द्रियः} %||3-15-29||

\twolineshloka
{प्रियाभिभाषी मधुरो दीर्घबाहुररिन्दमः}
{सन्त्यज्य विविधान्सौख्यानार्यं सर्वात्मनाश्रितः} %||3-15-30||

\twolineshloka
{जितः स्वर्गस्तव भ्रात्रा भरतेन महात्मना}
{वनस्थमपि तापस्ये यस्त्वामनुविधीयते} %||3-15-31||

\twolineshloka
{न पित्र्यमनुवर्न्तन्ते मातृकं द्विपदा इति}
{ख्यातो लोकप्रवादोऽयं भरतेनान्यथाकृतः} %||3-15-32||

\twolineshloka
{भर्ता दशरथो यस्याः साधुश्च भरतः सुतः}
{कथं नु साम्बा कैकेयी तादृशी क्रूरदर्शिनी} %||3-15-33||

\twolineshloka
{इत्येवं लक्ष्मणे वाक्यं स्नेहाद्ब्रुवति धार्मिके}
{परिवादं जनन्यास्तमसहन्राघवोऽब्रवीत्} %||3-15-34||

\twolineshloka
{न तेऽम्बा मध्यमा तात गर्हितव्या कथञ्चन}
{तामेवेक्ष्वाकुनाथस्य भरतस्य कथां कुरु} %||3-15-35||

\twolineshloka
{निश्चितापि हि मे बुद्धिर्वनवासे दृढव्रता}
{भरतस्नेहसन्तप्ता बालिशीक्रियते पुनः} %||3-15-36||

\twolineshloka
{इत्येवं विलपंस्तत्र प्राप्य गोदावरीं नदीम्}
{चक्रेऽभिषेकं काकुत्स्थः सानुजः सह सीतया} %||3-15-37||

\twolineshloka
{तर्पयित्वाथ सलिलैस्ते पितॄन्दैवतानि च}
{स्तुवन्ति स्मोदितं सूर्यं देवताश्च समाहिताः} %||3-15-38||

\fourlineindentedshloka
{कृताभिषेकः स रराज रामः}
{ सीताद्वितीयः सह लक्ष्मणेन}
{कृताभिषेकस्त्वगराजपुत्र्या}
{ रुद्रः सनन्दिर्भगवानिवेशः} %||3-16-39||


{॥इत्यार्षे श्रीमद्रामायणे वाल्मीकीये आदिकाव्ये अरण्यकाण्डे पञ्चदशः सर्गः॥}

\sect{षोडशः सर्गः}

\twolineshloka
{कृताभिषेको रामस्तु सीता सौमित्रिरेव च}
{तस्माद्गोदावरीतीरात्ततो जग्मुः स्वमाश्रमम्} %||3-16-1||

\twolineshloka
{आश्रमं तमुपागम्य राघवः सहलक्ष्मणः}
{कृत्वा पौर्वाह्णिकं कर्म पर्णशालामुपागमत्} %||3-16-2||

\threelineshloka
{स रामः पर्णशालायामासीनः सह सीतया}
{विरराज महाबाहुश्चित्रया चन्द्रमा इव}
{लक्ष्मणेन सह भ्रात्रा चकार विविधाः कथाः} %||3-16-3||

\twolineshloka
{तदासीनस्य रामस्य कथासंसक्तचेतसः}
{तं देशं राक्षसी काचिदाजगाम यदृच्छया} %||3-16-4||

\twolineshloka
{सा तु शूर्पणखा नाम दशग्रीवस्य रक्षसः}
{भगिनी राममासाद्य ददर्श त्रिदशोपमम्} %||3-16-5||

\twolineshloka
{सिंहोरस्कं महाबाहुं पद्मपत्रनिभेक्षणम्}
{सुकुमारं महासत्त्वं पार्थिवव्यञ्जनान्वितम्} %||3-16-6||

\twolineshloka
{राममिन्दीवरश्यामं कन्दर्पसदृशप्रभम्}
{बभूवेन्द्रोपमं दृष्ट्वा राक्षसी काममोहिता} %||3-16-7||

\twolineshloka
{सुमुखं दुर्मुखी रामं वृत्तमध्यं महोदरी}
{विशालाक्षं विरूपाक्षी सुकेशं ताम्रमूर्धजा} %||3-16-8||

\twolineshloka
{प्रियरूपं विरूपा सा सुस्वरं भैरवस्वना}
{तरुणं दारुणा वृद्धा दक्षिणं वामभाषिणी} %||3-16-9||

\twolineshloka
{न्यायवृत्तं सुदुर्वृत्ता प्रियमप्रियदर्शना}
{शरीरजसमाविष्टा राक्षसी राममब्रवीत्} %||3-16-10||

\twolineshloka
{जटी तापसरूपेण सभार्यः शरचापधृक्}
{आगतस्त्वमिमं देशं कथं राक्षससेवितम्} %||3-16-11||

\twolineshloka
{एवमुक्तस्तु राक्षस्या शूर्पणख्या परन्तपः}
{ऋजुबुद्धितया सर्वमाख्यातुमुपचक्रमे} %||3-16-12||

\twolineshloka
{आसीद्दशरथो नाम राजा त्रिदशविक्रमः}
{तस्याहमग्रजः पुत्रो रामो नाम जनैः श्रुतः} %||3-16-13||

\twolineshloka
{भ्रातायं लक्ष्मणो नाम यवीयान्मामनुव्रतः}
{इयं भार्या च वैदेही मम सीतेति विश्रुता} %||3-16-14||

\twolineshloka
{नियोगात्तु नरेन्द्रस्य पितुर्मातुश्च यन्त्रितः}
{धर्मार्थं धर्मकाङ्क्षी च वनं वस्तुमिहागतः} %||3-16-15||

\twolineshloka
{त्वां तु वेदितुमिच्छामि कथ्यतां कासि कस्य वा}
{इह वा किंनिमित्तं त्वमागता ब्रूहि तत्त्वतः} %||3-16-16||

\twolineshloka
{साब्रवीद्वचनं श्रुत्वा राक्षसी मदनार्दिता}
{श्रूयतां राम वक्ष्यामि तत्त्वार्थं वचनं मम} %||3-16-17||

\twolineshloka
{अहं शूर्पणखा नाम राक्षसी कामरूपिणी}
{अरण्यं विचरामीदमेका सर्वभयङ्करा} %||3-16-18||

\twolineshloka
{रावणो नाम मे भ्राता राक्षसो राक्षसेश्वरः}
{प्रवृद्धनिद्रश्च सदा कुम्भकर्णो महाबलः} %||3-16-19||

\twolineshloka
{विभीषणस्तु धर्मात्मा न तु राक्षसचेष्टितः}
{प्रख्यातवीर्यौ च रणे भ्रातरौ खरदूषणौ} %||3-16-20||

\threelineshloka
{तानहं समतिक्रान्ता राम त्वापूर्वदर्शनात्}
{समुपेतास्मि भावेन भर्तारं पुरुषोत्तमम्}
{चिराय भव भर्ता मे सीतया किं करिष्यसि} %||3-16-21||

\twolineshloka
{विकृता च विरूपा च न सेयं सदृशी तव}
{अहमेवानुरूपा ते भार्यारूपेण पश्य माम्} %||3-16-22||

\twolineshloka
{इमां विरूपामसतीं करालां निर्णतोदरीम्}
{अनेन सह ते भ्रात्रा भक्षयिष्यामि मानुषीम्} %||3-16-23||

\twolineshloka
{ततः पर्वतशृङ्गाणि वनानि विविधानि च}
{पश्यन्सह मया कान्त दण्डकान्विचरिष्यसि} %||3-16-24||

\twolineshloka
{इत्येवमुक्तः काकुत्स्थः प्रहस्य मदिरेक्षणाम्}
{इदं वचनमारेभे वक्तुं वाक्यविशारदः} %||3-17-25||


{॥इत्यार्षे श्रीमद्रामायणे वाल्मीकीये आदिकाव्ये अरण्यकाण्डे षोडशः सर्गः॥}

\sect{सप्तदशः सर्गः}

\twolineshloka
{तां तु शूर्पणखां रामः कामपाशावपाशिताम्}
{स्वेच्छया श्लक्ष्णया वाचा स्मितपूर्वमथाब्रवीत्} %||3-17-1||

\twolineshloka
{कृतदारोऽस्मि भवति भार्येयं दयिता मम}
{त्वद्विधानां तु नारीणां सुदुःखा ससपत्नता} %||3-17-2||

\twolineshloka
{अनुजस्त्वेष मे भ्राता शीलवान्प्रियदर्शनः}
{श्रीमानकृतदारश्च लक्ष्मणो नाम वीर्यवान्} %||3-17-3||

\twolineshloka
{अपूर्वी भार्यया चार्थी तरुणः प्रियदर्शनः}
{अनुरूपश्च ते भर्ता रूपस्यास्य भविष्यति} %||3-17-4||

\twolineshloka
{एनं भज विशालाक्षि भर्तारं भ्रातरं मम}
{असपत्ना वरारोहे मेरुमर्कप्रभा यथा} %||3-17-5||

\twolineshloka
{इति रामेण सा प्रोक्ता राक्षसी काममोहिता}
{विसृज्य रामं सहसा ततो लक्ष्मणमब्रवीत्} %||3-17-6||

\twolineshloka
{अस्य रूपस्य ते युक्ता भार्याहं वरवर्णिनी}
{मया सह सुखं सर्वान्दण्डकान्विचरिष्यसि} %||3-17-7||

\twolineshloka
{एवमुक्तस्तु सौमित्री राक्षस्या वाक्यकोविदः}
{ततः शूर्पणखीं स्मित्वा लक्ष्मणो युक्तमब्रवीत्} %||3-17-8||

\twolineshloka
{कथं दासस्य मे दासी भार्या भवितुमिच्छसि}
{सोऽहमार्येण परवान्भ्रात्रा कमलवर्णिनी} %||3-17-9||

\twolineshloka
{समृद्धार्थस्य सिद्धार्था मुदितामलवर्णिनी}
{आर्यस्य त्वं विशालाक्षि भार्या भव यवीयसी} %||3-17-10||

\twolineshloka
{एतां विरूपामसतीं करालां निर्णतोदरीम्}
{भार्यां वृद्धां परित्यज्य त्वामेवैष भजिष्यति} %||3-17-11||

\twolineshloka
{को हि रूपमिदं श्रेष्ठं सन्त्यज्य वरवर्णिनि}
{मानुषेषु वरारोहे कुर्याद्भावं विचक्षणः} %||3-17-12||

\twolineshloka
{इति सा लक्ष्मणेनोक्ता कराला निर्णतोदरी}
{मन्यते तद्वचः सत्यं परिहासाविचक्षणा} %||3-17-13||

\twolineshloka
{सा रामं पर्णशालायामुपविष्टं परन्तपम्}
{सीतया सह दुर्धर्षमब्रवीत्काममोहिता} %||3-17-14||

\twolineshloka
{इमां विरूपामसतीं करालां निर्णतोदरीम्}
{वृद्धां भार्यामवष्टभ्य न मां त्वं बहु मन्यसे} %||3-17-15||

\twolineshloka
{अद्येमां भक्षयिष्यामि पश्यतस्तव मानुषीम्}
{त्वया सह चरिष्यामि निःसपत्ना यथासुखम्} %||3-17-16||

\twolineshloka
{इत्युक्त्वा मृगशावाक्षीमलातसदृशेक्षणा}
{अभ्यधावत्सुसङ्क्रुद्धा महोल्का रोहिणीमिव} %||3-17-17||

\twolineshloka
{तां मृत्युपाशप्रतिमामापतन्तीं महाबलः}
{निगृह्य रामः कुपितस्ततो लक्ष्मणमब्रवीत्} %||3-17-18||

\twolineshloka
{क्रूरैरनार्यैः सौमित्रे परिहासः कथञ्चन}
{न कार्यः पश्य वैदेहीं कथञ्चित्सौम्य जीवतीम्} %||3-17-19||

\twolineshloka
{इमां विरूपामसतीमतिमत्तां महोदरीम्}
{राक्षसीं पुरुषव्याघ्र विरूपयितुमर्हसि} %||3-17-20||

\twolineshloka
{इत्युक्तो लक्ष्मणस्तस्याः क्रुद्धो रामस्य पश्यतः}
{उद्धृत्य खड्गं चिच्छेद कर्णनासं महाबलः} %||3-17-21||

\twolineshloka
{निकृत्तकर्णनासा तु विस्वरं सा विनद्य च}
{यथागतं प्रदुद्राव घोरा शूर्पणखा वनम्} %||3-17-22||

\twolineshloka
{सा विरूपा महाघोरा राक्षसी शोणितोक्षिता}
{ननाद विविधान्नादान्यथा प्रावृषि तोयदः} %||3-17-23||

\twolineshloka
{सा विक्षरन्ती रुधिरं बहुधा घोरदर्शना}
{प्रगृह्य बाहू गर्जन्ती प्रविवेश महावनम्} %||3-17-24||

\fourlineindentedshloka
{ततस्तु सा राक्षससङ्घसंवृतम्}
{ खरं जनस्थानगतं विरूपिता}
{उपेत्य तं भ्रातरमुग्रतेजसम्}
{ पपात भूमौ गगनाद्यथाशनिः} %||3-17-25||

\fourlineindentedshloka
{ततः सभार्यं भयमोहमूर्छिता}
{ सलक्ष्मणं राघवमागतं वनम्}
{विरूपणं चात्मनि शोणितोक्षिता}
{ शशंस सर्वं भगिनी खरस्य सा} %||3-18-26||


{॥इत्यार्षे श्रीमद्रामायणे वाल्मीकीये आदिकाव्ये अरण्यकाण्डे सप्तदशः सर्गः॥}

\sect{अष्टादशः सर्गः}

\twolineshloka
{तां तथा पतितां दृष्ट्वा विरूपां शोणितोक्षिताम्}
{भगिनीं क्रोधसन्तप्तः खरः पप्रच्छ राक्षसः} %||3-18-1||

\twolineshloka
{बलविक्रमसम्पन्ना कामगा कामरूपिणी}
{इमामवस्थां नीता त्वं केनान्तकसमा गता} %||3-18-2||

\twolineshloka
{देवगन्धर्वभूतानामृषीणां च महात्मनाम्}
{कोऽयमेवं महावीर्यस्त्वां विरूपां चकार ह} %||3-18-3||

\twolineshloka
{न हि पश्याम्यहं लोके यः कुर्यान्मम विप्रियम्}
{अन्तरेन सहस्राक्षं महेन्द्रं पाकशासनम्} %||3-18-4||

\twolineshloka
{अद्याहं मार्गणैः प्राणानादास्ये जीवितान्तकैः}
{सलिले क्षीरमासक्तं निष्पिबन्निव सारसः} %||3-18-5||

\twolineshloka
{निहतस्य मया सङ्ख्ये शरसङ्कृत्तमर्मणः}
{सफेनं रुधिरं रक्तं मेदिनी कस्य पास्यति} %||3-18-6||

\twolineshloka
{कस्य पत्ररथाः कायान्मांसमुत्कृत्य सङ्गताः}
{प्रहृष्टा भक्षयिष्यन्ति निहतस्य मया रणे} %||3-18-7||

\twolineshloka
{तं न देवा न गन्धर्वा न पिशाचा न राक्षसाः}
{मयापकृष्टं कृपणं शक्तास्त्रातुं महाहवे} %||3-18-8||

\twolineshloka
{उपलभ्य शनैः संज्ञां तं मे शंसितुमर्हसि}
{येन त्वं दुर्विनीतेन वने विक्रम्य निर्जिता} %||3-18-9||

\twolineshloka
{इति भ्रातुर्वचः श्रुत्वा क्रुद्धस्य च विशेषतः}
{ततः शूर्पणखा वाक्यं सबाष्पमिदमब्रवीत्} %||3-18-10||

\twolineshloka
{तरुणौ रूपसम्पन्नौ सुकूमारौ महाबलौ}
{पुण्डरीकविशालाक्षौ चीरकृष्णाजिनाम्बरौ} %||3-18-11||

\twolineshloka
{गन्धर्वराजप्रतिमौ पार्थिवव्यञ्जनान्वितौ}
{देवौ वा मानुषौ वा तौ न तर्कयितुमुत्सहे} %||3-18-12||

\twolineshloka
{तरुणी रूपसम्पन्ना सर्वाभरणभूषिता}
{दृष्टा तत्र मया नारी तयोर्मध्ये सुमध्यमा} %||3-18-13||

\twolineshloka
{ताभ्यामुभाभ्यां सम्भूय प्रमदामधिकृत्य ताम्}
{इमामवस्थां नीताहं यथानाथासती तथा} %||3-18-14||

\twolineshloka
{तस्याश्चानृजुवृत्तायास्तयोश्च हतयोरहम्}
{सफेनं पातुमिच्छामि रुधिरं रणमूर्धनि} %||3-18-15||

\twolineshloka
{एष मे प्रथमः कामः कृतस्तात त्वया भवेत्}
{तस्यास्तयोश्च रुधिरं पिबेयमहमाहवे} %||3-18-16||

\twolineshloka
{इति तस्यां ब्रुवाणायां चतुर्दश महाबलान्}
{व्यादिदेश खरः क्रुद्धो राक्षसानन्तकोपमान्} %||3-18-17||

\twolineshloka
{मानुषौ शस्त्रसम्पन्नौ चीरकृष्णाजिनाम्बरौ}
{प्रविष्टौ दण्डकारण्यं घोरं प्रमदया सह} %||3-18-18||

\twolineshloka
{तौ हत्वा तां च दुर्वृत्तामुपावर्तितुमर्हथ}
{इयं च रुधिरं तेषां भगिनी मम पास्यति} %||3-18-19||

\twolineshloka
{मनोरथोऽयमिष्टोऽस्या भगिन्या मम राक्षसाः}
{शीघ्रं सम्पद्यतां गत्वा तौ प्रमथ्य स्वतेजसा} %||3-18-20||

\twolineshloka
{इति प्रतिसमादिष्टा राक्षसास्ते चतुर्दश}
{तत्र जग्मुस्तया सार्धं घना वातेरिता यथा} %||3-19-21||


{॥इत्यार्षे श्रीमद्रामायणे वाल्मीकीये आदिकाव्ये अरण्यकाण्डे अष्टादशः सर्गः॥}

\sect{एकोनविंशः सर्गः}

\twolineshloka
{ततः शूर्पणखा घोरा राघवाश्रममागता}
{रक्षसामाचचक्षे तौ भ्रातरौ सह सीतया} %||3-19-1||

\twolineshloka
{ते रामं पर्णशालायामुपविष्टं महाबलम्}
{ददृशुः सीतया सार्धं वैदेह्या लक्ष्मणेन च} %||3-19-2||

\twolineshloka
{तान्दृष्ट्वा राघवः श्रीमानागतां तां च राक्षसीम्}
{अब्रवीद्भ्रातरं रामो लक्ष्मणं दीप्ततेजसम्} %||3-19-3||

\twolineshloka
{मुहूर्तं भव सौमित्रे सीतायाः प्रत्यनन्तरः}
{इमानस्या वधिष्यामि पदवीमागतानिह} %||3-19-4||

\twolineshloka
{वाक्यमेतत्ततः श्रुत्वा रामस्य विदितात्मनः}
{तथेति लक्ष्मणो वाक्यं रामस्य प्रत्यपूजयत्} %||3-19-5||

\twolineshloka
{राघवोऽपि महच्चापं चामीकरविभूषितम्}
{चकार सज्यं धर्मात्मा तानि रक्षांसि चाब्रवीत्} %||3-19-6||

\twolineshloka
{पुत्रौ दशरथस्यावां भ्रातरौ रामलक्ष्मणौ}
{प्रविष्टौ सीतया सार्धं दुश्चरं दण्डकावनम्} %||3-19-7||

\twolineshloka
{फलमूलाशनौ दान्तौ तापसौ धर्मचारिणौ}
{वसन्तौ दण्डकारण्ये किमर्थमुपहिंसथ} %||3-19-8||

\twolineshloka
{युष्मान्पापात्मकान्हन्तुं विप्रकारान्महावने}
{ऋषीणां तु नियोगेन प्राप्तोऽहं सशरासनः} %||3-19-9||

\twolineshloka
{तिष्ठतैवात्र सन्तुष्टा नोपसर्पितुमर्हथ}
{यदि प्राणैरिहार्थो वो निवर्तध्वं निशाचराः} %||3-19-10||

\twolineshloka
{तस्य तद्वचनं श्रुत्वा राक्षसास्ते चतुर्दश}
{ऊचुर्वाचं सुसङ्क्रुद्धा ब्रह्मघ्नः शूलपाणयः} %||3-19-11||

\twolineshloka
{संरक्तनयना घोरा रामं रक्तान्तलोचनम्}
{परुषा मधुराभाषं हृष्टादृष्टपराक्रमम्} %||3-19-12||

\twolineshloka
{क्रोधमुत्पाद्य नो भर्तुः खरस्य सुमहात्मनः}
{त्वमेव हास्यसे प्राणानद्यास्माभिर्हतो युधि} %||3-19-13||

\twolineshloka
{का हि ते शक्तिरेकस्य बहूनां रणमूर्धनि}
{अस्माकमग्रतः स्थातुं किं पुनर्योद्धुमाहवे} %||3-19-14||

\twolineshloka
{एभिर्बाहुप्रयुक्तैर्नः परिघैः शूलपट्टिशैः}
{प्राणांस्त्यक्ष्यसि वीर्यं च धनुश्च करपीडितम्} %||3-19-15||

\threelineshloka
{इत्येवमुक्त्वा संरब्धा राक्षसास्ते चतुर्दश}
{उद्यतायुधनिस्त्रिंशा राममेवाभिदुद्रुवुः}
{चिक्षिपुस्तानि शूलानि राघवं प्रति दुर्जयम्} %||3-19-16||

\twolineshloka
{तानि शूलानि काकुत्स्थः समस्तानि चतुर्दश}
{तावद्भिरेव चिच्छेद शरैः काञ्चनभूषणैः} %||3-19-17||

\twolineshloka
{ततः पश्चान्महातेजा नाराचान्सूर्यसंनिभान्}
{जग्राह परमक्रुद्धश्चतुर्दश शिलाशितान्} %||3-19-18||

\twolineshloka
{गृहीत्वा धनुरायम्य लक्ष्यानुद्दिश्य राक्षसान्}
{मुमोच राघवो बाणान्वज्रानिव शतक्रतुः} %||3-19-19||

\twolineshloka
{रुक्मपुङ्खाश्च विशिखाः प्रदीप्ता हेमभूषणाः}
{अन्तरिक्षे महोल्कानां बभूवुस्तुल्यवर्चसः} %||3-19-20||

\twolineshloka
{ते भित्त्वा रक्षसां वेगाद्वक्षांसि रुधिराप्लुताः}
{विनिष्पेतुस्तदा भूमौ न्यमज्जन्ताशनिस्वनाः} %||3-19-21||

\twolineshloka
{ते भिन्नहृदया भूमौ छिन्नमूला इव द्रुमाः}
{निपेतुः शोणितार्द्राङ्गा विकृता विगतासवः} %||3-19-22||

\twolineshloka
{तान्भूमौ पतितान्दृष्ट्वा राक्षसी क्रोधमूर्छिता}
{परित्रस्ता पुनस्तत्र व्यसृजद्भैरवं रवम्} %||3-19-23||

\threelineshloka
{सा नदन्ती महानादं जवाच्छूर्पणखा पुनः}
{उपगम्य खरं सा तु किञ्चित्संशुष्क शोणिता}
{पपात पुनरेवार्ता सनिर्यासेव वल्लरी} %||3-19-24||

\fourlineindentedshloka
{निपातितान्प्रेक्ष्य रणे तु राक्षसा}
{न्प्रधाविता शूर्पणखा पुनस्ततः}
{वधं च तेषां निखिलेन रक्षसाम्}
{ शशंस सर्वं भगिनी खरस्य सा} %||3-20-25||


{॥इत्यार्षे श्रीमद्रामायणे वाल्मीकीये आदिकाव्ये अरण्यकाण्डे एकोनविंशः सर्गः॥}

\sect{विंशः सर्गः}

\twolineshloka
{स पुनः पतितां दृष्ट्वा क्रोधाच्छूर्पणखां खरः}
{उवाच व्यक्तता वाचा तामनर्थार्थमागताम्} %||3-20-1||

\twolineshloka
{मया त्विदानीं शूरास्ते राक्षसा रुधिराशनाः}
{त्वत्प्रियार्थं विनिर्दिष्टाः किमर्थं रुद्यते पुनः} %||3-20-2||

\twolineshloka
{भक्ताश्चैवानुरक्ताश्च हिताश्च मम नित्यशः}
{घ्नन्तोऽपि न निहन्तव्या न न कुर्युर्वचो मम} %||3-20-3||

\twolineshloka
{किमेतच्छ्रोतुमिच्छामि कारणं यत्कृते पुनः}
{हा नाथेति विनर्दन्ती सर्पवद्वेष्टसे क्षितौ} %||3-20-4||

\twolineshloka
{अनाथवद्विलपसि किं नु नाथे मयि स्थिते}
{उत्तिष्ठोत्तिष्ठ मा भैषीर्वैक्लव्यं त्यज्यतामिह} %||3-20-5||

\twolineshloka
{इत्येवमुक्ता दुर्धर्षा खरेण परिसान्त्विता}
{विमृज्य नयने सास्रे खरं भ्रातरमब्रवीत्} %||3-20-6||

\twolineshloka
{प्रेषिताश्च त्वया शूरा राक्षसास्ते चतुर्दश}
{निहन्तुं राघवं घोरा मत्प्रियार्थं सलक्ष्मणम्} %||3-20-7||

\twolineshloka
{ते तु रामेण सामर्षाः शूलपट्टिशपाणयः}
{समरे निहताः सर्वे सायकैर्मर्मभेदिभिः} %||3-20-8||

\twolineshloka
{तान्भूमौ पतितान्दृष्ट्वा क्षणेनैव महाबलान्}
{रामस्य च महत्कर्म महांस्त्रासोऽभवन्मम} %||3-20-9||

\twolineshloka
{सास्मि भीता समुद्विग्ना विषण्णा च निशाचर}
{शरणं त्वां पुनः प्राप्ता सर्वतो भयदर्शिनी} %||3-20-10||

\twolineshloka
{विषादनक्राध्युषिते परित्रासोर्मिमालिनि}
{किं मां न त्रायसे मग्नां विपुले शोकसागरे} %||3-20-11||

\twolineshloka
{एते च निहता भूमौ रामेण निशितैः शरैः}
{ये च मे पदवीं प्राप्ता राक्षसाः पिशिताशनाः} %||3-20-12||

\threelineshloka
{मयि ते यद्यनुक्रोशो यदि रक्षःसु तेषु च}
{रामेण यदि शक्तिस्ते तेजो वास्ति निशाचर}
{दण्डकारण्यनिलयं जहि राक्षसकण्टकम्} %||3-20-13||

\twolineshloka
{यदि रामं ममामित्रमद्य त्वं न वधिष्यसि}
{तव चैवाग्रतः प्राणांस्त्यक्ष्यामि निरपत्रपा} %||3-20-14||

\twolineshloka
{बुद्ध्याहमनुपश्यामि न त्वं रामस्य संयुगे}
{स्थातुं प्रतिमुखे शक्तः सचापस्य महारणे} %||3-20-15||

\twolineshloka
{शूरमानी न शूरस्त्वं मिथ्यारोपितविक्रमः}
{मानुषौ यन्न शक्नोषि हन्तुं तौ रामलक्ष्मणौ} %||3-20-16||

\twolineshloka
{अपयाहि जनस्थानात्त्वरितः सहबान्धवः}
{निःसत्त्वस्याल्पवीर्यस्य वासस्ते कीदृशस्त्विह} %||3-20-17||

\threelineshloka
{रामतेजोऽभिभूतो हि त्वं क्षिप्रं विनशिष्यसि}
{स हि तेजःसमायुक्तो रामो दशरथात्मजः}
{भ्राता चास्य महावीर्यो येन चास्मि विरूपिता} %||3-21-18||


{॥इत्यार्षे श्रीमद्रामायणे वाल्मीकीये आदिकाव्ये अरण्यकाण्डे विंशः सर्गः॥}

\sect{एकविंशः सर्गः}

\twolineshloka
{एवमाधर्षितः शूरः शूर्पणख्या खरस्तदा}
{उवाच रक्षसां मध्ये खरः खरतरं वचः} %||3-21-1||

\twolineshloka
{तवापमानप्रभवः क्रोधोऽयमतुलो मम}
{न शक्यते धारयितुं लवणाम्भ इवोत्थितम्} %||3-21-2||

\twolineshloka
{न रामं गणये वीर्यान्मानुषं क्षीणजीवितम्}
{आत्मा दुश्चरितैः प्राणान्हतो योऽद्य विमोक्ष्यति} %||3-21-3||

\twolineshloka
{बाष्पः संह्रियतामेष सम्भ्रमश्च विमुच्यताम्}
{अहं रामः सह भ्रात्रा नयामि यमसादनम्} %||3-21-4||

\twolineshloka
{परश्वधहतस्याद्य मन्दप्राणस्य भूतले}
{रामस्य रुधिरं रक्तमुष्णं पास्यसि राक्षसि} %||3-21-5||

\twolineshloka
{सा प्रहृष्ट्वा वचः श्रुत्वा खरस्य वदनाच्च्युतम्}
{प्रशशंस पुनर्मौर्ख्याद्भ्रातरं रक्षसां वरम्} %||3-21-6||

\twolineshloka
{तया परुषितः पूर्वं पुनरेव प्रशंसितः}
{अब्रवीद्दूषणं नाम खरः सेनापतिं तदा} %||3-21-7||

\twolineshloka
{चतुर्दश सहस्राणि मम चित्तानुवर्तिनाम्}
{रक्षसीं भीमवेगानां समरेष्वनिवर्तिनाम्} %||3-21-8||

\twolineshloka
{नीलजीमूतवर्णानां घोराणां क्रूरकर्मणाम्}
{लोकसिंहाविहाराणां बलिनामुग्रतेजसाम्} %||3-21-9||

\twolineshloka
{तेषां शार्दूलदर्पाणां महास्यानां महौजसाम्}
{सर्वोद्योगमुदीर्णानां रक्षसां सौम्य कारय} %||3-21-10||

\twolineshloka
{उपस्थापय मे क्षिप्रं रथं सौम्य धनूंषि च}
{शरांश्च चित्रान्खड्गांश्च शक्तीश्च विविधाः शिताः} %||3-21-11||

\twolineshloka
{अग्रे निर्यातुमिच्छामि पौलस्त्यानां महात्मनाम्}
{वधार्थं दुर्विनीतस्य रामस्य रणकोविदः} %||3-21-12||

\twolineshloka
{इति तस्य ब्रुवाणस्य सूर्यवर्णं महारथम्}
{सदश्वैः शबलैर्युक्तमाचचक्षेऽथ दूषणः} %||3-21-13||

\twolineshloka
{तं मेरुशिखराकारं तप्तकाञ्चनभूषणम्}
{हेमचक्रमसम्बाधं वैदूर्यमय कूबरम्} %||3-21-14||

\twolineshloka
{मत्स्यैः पुष्पैर्द्रुमैः शैलैश्चन्द्रसूर्यैश्च काञ्चनैः}
{माङ्गल्यैः पक्षिसङ्घैश्च ताराभिश्च समावृतम्} %||3-21-15||

\twolineshloka
{ध्वजनिस्त्रिंशसम्पन्नं किङ्किणीकविभूषितम्}
{सदश्वयुक्तं सोऽमर्षादारुरोह रथं खरः} %||3-21-16||

\twolineshloka
{निशाम्य तं रथगतं राक्षसा भीमविक्रमाः}
{तस्थुः सम्परिवार्यैनं दूषणं च महाबलम्} %||3-21-17||

\twolineshloka
{खरस्तु तान्महेष्वासान्घोरचर्मायुधध्वजान्}
{निर्यातेत्यब्रवीद्दृष्ट्वा रथस्थः सर्वराक्षसान्} %||3-21-18||

\twolineshloka
{ततस्तद्राक्षसं सैन्यं घोरचर्मायुधध्वजम्}
{निर्जगाम जनस्थानान्महानादं महाजवम्} %||3-21-19||

\twolineshloka
{मुद्गरैः पट्टिशैः शूलैः सुतीक्ष्णैश्च परश्वधैः}
{खड्गैश्चक्रैश्च हस्तस्थैर्भ्राजमानैश्च तोमरैः} %||3-21-20||

\twolineshloka
{शक्तिभिः पतिघैर्घोरैरतिमात्रैश्च कार्मुकैः}
{गदासिमुसलैर्वज्रैर्गृहीतैर्भीमदर्शनैः} %||3-21-21||

\twolineshloka
{राक्षसानां सुघोराणां सहस्राणि चतुर्दश}
{निर्यातानि जनस्थानात्खरचित्तानुवर्तिनाम्} %||3-21-22||

\twolineshloka
{तांस्त्वभिद्रवतो दृष्ट्वा राक्षसान्भीमविक्रमान्}
{खरस्यापि रथः किञ्चिज्जगाम तदनन्तरम्} %||3-21-23||

\twolineshloka
{ततस्ताञ्शबलानश्वांस्तप्तकाञ्चनभूषितान्}
{खरस्य मतमाज्ञाय सारथिः समचोदयत्} %||3-21-24||

\twolineshloka
{स चोदितो रथः शीघ्रं खरस्य रिपुघातिनः}
{शब्देनापूरयामास दिशश्च प्रतिशस्तथा} %||3-21-25||

\fourlineindentedshloka
{प्रवृद्धमन्युस्तु खरः खरस्वनो}
{ रिपोर्वधार्थं त्वरितो यथान्तकः}
{अचूचुदत्सारथिमुन्नदन्पुनर्-}
{ महाबलो मेघ इवाश्मवर्षवान्} %||3-22-26||


{॥इत्यार्षे श्रीमद्रामायणे वाल्मीकीये आदिकाव्ये अरण्यकाण्डे एकविंशः सर्गः॥}

\sect{द्वाविंशः सर्गः}

\twolineshloka
{तत्प्रयातं बलं घोरमशिवं शोणितोदकम्}
{अभ्यवर्षन्महामेघस्तुमुलो गर्दभारुणः} %||3-22-1||

\twolineshloka
{निपेतुस्तुरगास्तस्य रथयुक्ता महाजवाः}
{समे पुष्पचिते देशे राजमार्गे यदृच्छया} %||3-22-2||

\twolineshloka
{श्यामं रुधिरपर्यन्तं बभूव परिवेषणम्}
{अलातचक्रप्रतिमं प्रतिगृह्य दिवाकरम्} %||3-22-3||

\twolineshloka
{ततो ध्वजमुपागम्य हेमदण्डं समुच्छ्रितम्}
{समाक्रम्य महाकायस्तस्थौ गृध्रः सुदारुणः} %||3-22-4||

\twolineshloka
{जनस्थानसमीपे च समाक्रम्य खरस्वनाः}
{विस्वरान्विविधांश्चक्रुर्मांसादा मृगपक्षिणः} %||3-22-5||

\twolineshloka
{व्याजह्रुश्च पदीप्तायां दिशि वै भैरवस्वनम्}
{अशिवा यातु दाहानां शिवा घोरा महास्वनाः} %||3-22-6||

\twolineshloka
{प्रभिन्नगिरिसङ्काशास्तोयशोषितधारिणः}
{आकाशं तदनाकाशं चक्रुर्भीमा बलाहकाः} %||3-22-7||

\twolineshloka
{बभूव तिमिरं घोरमुद्धतं रोमहर्षणम्}
{दिशो वा विदिशो वापि सुव्यक्तं न चकाशिरे} %||3-22-8||

\twolineshloka
{क्षतजार्द्रसवर्णाभा सन्ध्याकालं विना बभौ}
{खरस्याभिमुखं नेदुस्तदा घोरा मृगाः खगाः} %||3-22-9||

\twolineshloka
{नित्याशिवकरा युद्धे शिवा घोरनिदर्शनाः}
{नेदुर्बलस्याभिमुखं ज्वालोद्गारिभिराननैः} %||3-22-10||

\twolineshloka
{कबन्धः परिघाभासो दृश्यते भास्करान्तिके}
{जग्राह सूर्यं स्वर्भानुरपर्वणि महाग्रहः} %||3-22-11||

\twolineshloka
{प्रवाति मारुतः शीघ्रं निष्प्रभोऽभूद्दिवाकरः}
{उत्पेतुश्च विना रात्रिं ताराः खद्योतसप्रभाः} %||3-22-12||

\twolineshloka
{संलीनमीनविहगा नलिन्यः पुष्पपङ्कजाः}
{तस्मिन्क्षणे बभूवुश्च विना पुष्पफलैर्द्रुमाः} %||3-22-13||

\twolineshloka
{उद्धूतश्च विना वातं रेणुर्जलधरारुणः}
{वीचीकूचीति वाश्यन्तो बभूवुस्तत्र सारिकाः} %||3-22-14||

\twolineshloka
{उल्काश्चापि सनिर्घोषा निपेतुर्घोरदर्शनाः}
{प्रचचाल मही चापि सशैलवनकानना} %||3-22-15||

\twolineshloka
{खरस्य च रथस्थस्य नर्दमानस्य धीमतः}
{प्राकम्पत भुजः सव्यः खरश्चास्यावसज्जत} %||3-22-16||

\twolineshloka
{सास्रा सम्पद्यते दृष्टिः पश्यमानस्य सर्वतः}
{ललाटे च रुजा जाता न च मोहान्न्यवर्तत} %||3-22-17||

\twolineshloka
{तान्समीक्ष्य महोत्पातानुत्थितान्रोमहर्षणान्}
{अब्रवीद्राक्षसान्सर्वान्प्रहसन्स खरस्तदा} %||3-22-18||

\twolineshloka
{महोत्पातानिमान्सर्वानुत्थितान्घोरदर्शनान्}
{न चिन्तयाम्यहं वीर्याद्बलवान्दुर्बलानिव} %||3-22-19||

\twolineshloka
{तारा अपि शरैस्तीक्ष्णैः पातयेयं नभस्तलात्}
{मृत्युं मरणधर्मेण सङ्क्रुद्धो योजयाम्यहम्} %||3-22-20||

\twolineshloka
{राघवं तं बलोत्सिक्तं भ्रातरं चापि लक्ष्मणम्}
{अहत्वा सायकैस्तीक्ष्णैर्नोपावर्तितुमुत्सहे} %||3-22-21||

\twolineshloka
{सकामा भगिनी मेऽस्तु पीत्वा तु रुधिरं तयोः}
{यन्निमित्तं तु रामस्य लक्ष्मणस्य विपर्ययः} %||3-22-22||

\twolineshloka
{न क्वचित्प्राप्तपूर्वो मे संयुगेषु पराजयः}
{युष्माकमेतत्प्रत्यक्षं नानृतं कथयाम्यहम्} %||3-22-23||

\twolineshloka
{देवराजमपि क्रुद्धो मत्तैरावतयायिनम्}
{वज्रहस्तं रणे हन्यां किं पुनस्तौ च मानुषौ} %||3-22-24||

\twolineshloka
{सा तस्य गर्जितं श्रुत्वा राक्षसस्य महाचमूः}
{प्रहर्षमतुलं लेभे मृत्युपाशावपाशिता} %||3-22-25||

\twolineshloka
{समेयुश्च महात्मानो युद्धदर्शनकाङ्क्षिणः}
{ऋषयो देवगन्धर्वाः सिद्धाश्च सह चारणैः} %||3-22-26||

\twolineshloka
{समेत्य चोरुः सहितास्तेऽन्यायं पुण्यकर्मणः}
{स्वस्ति गोब्राह्मणेभ्योऽस्तु लोकानां ये च सम्मताः} %||3-22-27||

\twolineshloka
{जयतां राघवो युद्धे पौलस्त्यान्रजनीचरान्}
{चक्रा हस्तो यथा युद्धे सर्वानसुरपुङ्गवान्} %||3-22-28||

\twolineshloka
{एतच्चान्यच्च बहुशो ब्रुवाणाः परमर्षयः}
{ददृशुर्वाहिनीं तेषां राक्षसानां गतायुषाम्} %||3-22-29||

\twolineshloka
{रथेन तु खरो वेगात्सैन्यस्याग्राद्विनिःसृतः}
{तं दृष्ट्वा राक्षसं भूयो राक्षसाश्च विनिःसृताः} %||3-22-30||

\twolineshloka
{श्येन गामी पृथुग्रीवो यज्ञशत्रुर्विहङ्गमः}
{दुर्जयः करवीराक्षः परुषः कालकार्मुकः} %||3-22-31||

\twolineshloka
{मेघमाली महामाली सर्पास्यो रुधिराशनः}
{द्वादशैते महावीर्याः प्रतस्थुरभितः खरम्} %||3-22-32||

\twolineshloka
{महाकपालः स्थूलाक्षः प्रमाथी त्रिशिरास्तथा}
{चत्वार एते सेनाग्र्या दूषणं पृष्ठतोऽन्वयुः} %||3-22-33||

\fourlineindentedshloka
{सा भीमवेगा समराभिकामा}
{ सुदारुणा राक्षसवीर सेना}
{तौ राजपुत्रौ सहसाभ्युपेता}
{ मालाग्रहाणामिव चन्द्रसूर्यौ} %||3-23-34||


{॥इत्यार्षे श्रीमद्रामायणे वाल्मीकीये आदिकाव्ये अरण्यकाण्डे द्वाविंशः सर्गः॥}

\sect{त्रयोविंशः सर्गः}

\twolineshloka
{आश्रमं प्रति याते तु खरे खरपराक्रमे}
{तानेवौत्पातिकान्रामः सह भ्रात्रा ददर्श ह} %||3-23-1||

\twolineshloka
{तानुत्पातान्महाघोरानुत्थितान्रोमहर्षणान्}
{प्रजानामहितान्दृष्ट्वा वाक्यं लक्ष्मणमब्रवीत्} %||3-23-2||

\twolineshloka
{इमान्पश्य महाबाहो सर्वभूतापहारिणः}
{समुत्थितान्महोत्पातान्संहर्तुं सर्वराक्षसान्} %||3-23-3||

\twolineshloka
{अमी रुधिरधारास्तु विसृजन्तः खरस्वनान्}
{व्योम्नि मेघा विवर्तन्ते परुषा गर्दभारुणाः} %||3-23-4||

\twolineshloka
{सधूमाश्च शराः सर्वे मम युद्धाभिनन्दिनः}
{रुक्मपृष्ठानि चापानि विवेष्टन्ते च लक्ष्मण} %||3-23-5||

\twolineshloka
{यादृशा इह कूजन्ति पक्षिणो वनचारिणः}
{अग्रतो नो भयं प्राप्तं संशयो जीवितस्य च} %||3-23-6||

\twolineshloka
{सम्प्रहारस्तु सुमहान्भविष्यति न संशयः}
{अयमाख्याति मे बाहुः स्फुरमाणो मुहुर्मुहुः} %||3-23-7||

\twolineshloka
{संनिकर्षे तु नः शूर जयं शत्रोः पराजयम्}
{सुप्रभं च प्रसन्नं च तव वक्त्रं हि लक्ष्यते} %||3-23-8||

\twolineshloka
{उद्यतानां हि युद्धार्थं येषां भवति लक्ष्मणः}
{निष्प्रभं वदनं तेषां भवत्यायुः परिक्षयः} %||3-23-9||

\twolineshloka
{अनागतविधानं तु कर्तव्यं शुभमिच्छता}
{आपदं शङ्कमानेन पुरुषेण विपश्चिता} %||3-23-10||

\twolineshloka
{तस्माद्गृहीत्वा वैदेहीं शरपाणिर्धनुर्धरः}
{गुहामाश्रयशैलस्य दुर्गां पादपसङ्कुलाम्} %||3-23-11||

\twolineshloka
{प्रतिकूलितुमिच्छामि न हि वाक्यमिदं त्वया}
{शापितो मम पादाभ्यां गम्यतां वत्स माचिरम्} %||3-23-12||

\twolineshloka
{एवमुक्तस्तु रामेण लक्ष्मणः सह सीतया}
{शरानादाय चापं च गुहां दुर्गां समाश्रयत्} %||3-23-13||

\twolineshloka
{तस्मिन्प्रविष्टे तु गुहां लक्ष्मणे सह सीतया}
{हन्त निर्युक्तमित्युक्त्वा रामः कवचमाविशत्} %||3-23-14||

\twolineshloka
{सा तेनाग्निनिकाशेन कवचेन विभूषितः}
{बभूव रामस्तिमिरे विधूमोऽग्निरिवोत्थितः} %||3-23-15||

\twolineshloka
{स चापमुद्यम्य महच्छरानादाय वीर्यवान्}
{बभूवावस्थितस्तत्र ज्यास्वनैः पूरयन्दिशः} %||3-23-16||

\twolineshloka
{ततो देवाः सगन्धर्वाः सिद्धाश्च सह चारणैः}
{ऊचुः परमसन्त्रस्ता गुह्यकाश्च परस्परम्} %||3-23-17||

\twolineshloka
{चतुर्दश सहस्राणि रक्षसां भीमकर्मणाम्}
{एकश्च रामो धर्मात्मा कथं युद्धं भविष्यति} %||3-23-18||

\twolineshloka
{ततो गम्भीरनिर्ह्रादं घोरवर्मायुधध्वजम्}
{अनीकं यातुधानानां समन्तात्प्रत्यदृश्यत} %||3-23-19||

\twolineshloka
{सिंहनादं विसृजतामन्योन्यमभिगर्जताम्}
{चापानि विस्फरयतां जृम्भतां चाप्यभीक्ष्णशः} %||3-23-20||

\twolineshloka
{विप्रघुष्टस्वनानां च दुन्दुभींश्चापि निघ्नताम्}
{तेषां सुतुमुलः शब्दः पूरयामास तद्वनम्} %||3-23-21||

\twolineshloka
{तेन शब्देन वित्रस्ताः श्वापदा वनचारिणः}
{दुद्रुवुर्यत्र निःशब्दं पृष्ठतो नावलोकयन्} %||3-23-22||

\twolineshloka
{तत्त्वनीकं महावेगं रामं समुपसर्पत}
{घृतनानाप्रहरणं गम्भीरं सागरोपमम्} %||3-23-23||

\twolineshloka
{रामोऽपि चारयंश्चक्षुः सर्वतो रणपण्डितः}
{ददर्श खरसैन्यं तद्युद्धाभिमुखमुद्यतम्} %||3-23-24||

\twolineshloka
{वितत्य च धनुर्भीमं तूण्याश्चोद्धृत्य सायकान्}
{क्रोधमाहारयत्तीव्रं वधार्थं सर्वरक्षसाम्} %||3-23-25||

\twolineshloka
{दुष्प्रेक्ष्यः सोऽभवत्क्रुद्धो युगान्ताग्निरिव ज्वलन्}
{तं दृष्ट्वा तेजसाविष्टं प्राव्यथन्वनदेवताः} %||3-23-26||

\twolineshloka
{तस्य क्रुद्धस्य रूपं तु रामस्य ददृशे तदा}
{दक्षस्येव क्रतुं हन्तुमुद्यतस्य पिनाकिनः} %||3-24-27||


{॥इत्यार्षे श्रीमद्रामायणे वाल्मीकीये आदिकाव्ये अरण्यकाण्डे त्रयोविंशः सर्गः॥}

\sect{चतुर्विंशः सर्गः}

\twolineshloka
{अवष्टब्धधनुं रामं क्रुद्धं च रिपुघातिनम्}
{ददर्शाश्रममागम्य खरः सह पुरःसरैः} %||3-24-1||

\twolineshloka
{तं दृष्ट्वा सगुणं चापमुद्यम्य खरनिःस्वनम्}
{रामस्याभिमुखं सूतं चोद्यतामित्यचोदयत्} %||3-24-2||

\twolineshloka
{स खरस्याज्ञया सूतस्तुरगान्समचोदयत्}
{यत्र रामो महाबाहुरेको धुन्वन्धनुः स्थितः} %||3-24-3||

\twolineshloka
{तं तु निष्पतितं दृष्ट्वा सर्वे ते रजनीचराः}
{नर्दमाना महानादं सचिवाः पर्यवारयन्} %||3-24-4||

\twolineshloka
{स तेषां यातुधानानां मध्ये रतो गतः खरः}
{बभूव मध्ये ताराणां लोहिताङ्ग इवोदितः} %||3-24-5||

\twolineshloka
{ततस्तं भीमधन्वानं क्रुद्धाः सर्वे निशाचराः}
{रामं नानाविधैः शस्त्रैरभ्यवर्षन्त दुर्जयम्} %||3-24-6||

\twolineshloka
{मुद्गरैरायसैः शूलैः प्रासैः खड्गैः परश्वधैः}
{राक्षसाः समरे रामं निजघ्नू रोषतत्पराः} %||3-24-7||

\twolineshloka
{ते बलाहकसङ्काशा महानादा महाबलाः}
{अभ्यधावन्त काकुत्स्थं रामं युद्धे जिघांसवः} %||3-24-8||

\twolineshloka
{ते रामे शरवर्षाणि व्यसृजन्रक्षसां गुणाः}
{शैलेन्द्रमिव धाराभिर्वर्षमाणा महाधनाः} %||3-24-9||

\twolineshloka
{स तैः परिवृतो घोरै राघवो रक्षसां गणैः}
{तिथिष्विव महादेवो वृतः पारिषदां गणैः} %||3-24-10||

\twolineshloka
{तानि मुक्तानि शस्त्राणि यातुधानैः स राघवः}
{प्रतिजग्राह विशिखैर्नद्योघानिव सागरः} %||3-24-11||

\twolineshloka
{स तैः प्रहरणैर्घोरैर्भिन्नगात्रो न विव्यथे}
{रामः प्रदीप्तैर्बहुभिर्वज्रैरिव महाचलः} %||3-24-12||

\twolineshloka
{स विद्धः क्षतजादिग्धः सर्वगात्रेषु राघवः}
{बभूव रामः सन्ध्याभ्रैर्दिवाकर इवावृतः} %||3-24-13||

\twolineshloka
{विषेदुर्देवगन्धर्वाः सिद्धाश्च परमर्षयः}
{एकं सहस्त्रैर्बहुभिस्तदा दृष्ट्वा समावृतम्} %||3-24-14||

\twolineshloka
{ततो रामः सुसङ्क्रुद्धो मण्डलीकृतकार्मुकः}
{ससर्ज निशितान्बाणाञ्शतशोऽथ सहस्रशः} %||3-24-15||

\twolineshloka
{दुरावारान्दुर्विषहान्कालपाशोपमान्रणे}
{मुमोच लीलया रामः कङ्कपत्रानजिह्मगान्} %||3-24-16||

\twolineshloka
{ते शराः शत्रुसैन्येषु मुक्ता रामेण लीलया}
{आददू रक्षसां प्राणान्पाशाः कालकृता इव} %||3-24-17||

\twolineshloka
{भित्त्वा राक्षसदेहांस्तांस्ते शरा रुधिराप्लुताः}
{अन्तरिक्षगता रेजुर्दीप्ताग्निसमतेजसः} %||3-24-18||

\twolineshloka
{असङ्ख्येयास्तु रामस्य सायकाश्चापमण्डलात्}
{विनिष्पेतुरतीवोग्रा रक्षः प्राणापहारिणः} %||3-24-19||

\twolineshloka
{तैर्धनूंषि ध्वजाग्राणि वर्माणि च शिरांसि च}
{बहून्सहस्ताभरणानूरून्करिकरोपमान्} %||3-24-20||

\twolineshloka
{ततो नालीकनाराचैस्तीक्ष्णाग्रैश्च विकर्णिभिः}
{भीममार्तस्वरं चक्रुर्भिद्यमाना निशाचराः} %||3-24-21||

\twolineshloka
{तत्सैन्यं निशितैर्बाणैरर्दितं मर्मभेदिभिः}
{रामेण न सुखं लेभे शुष्कं वनमिवाग्निना} %||3-24-22||

\twolineshloka
{केचिद्भीमबलाः शूराः शूलान्खड्गान्परश्वधान्}
{चिक्षिपुः परमक्रुद्धा रामाय रजनीचराः} %||3-24-23||

\twolineshloka
{तानि बाणैर्महाबाहुः शस्त्राण्यावार्य राघवः}
{जहार समरे प्राणांश्चिच्छेद च शिरोधरान्} %||3-24-24||

\twolineshloka
{अवशिष्टाश्च ये तत्र विषण्णाश्च निशाचराः}
{खरमेवाभ्यधावन्त शरणार्थं शरार्दिताः} %||3-24-25||

\twolineshloka
{तान्सर्वान्पुनरादाय समाश्वास्य च दूषणः}
{अभ्यधावत काकुत्स्थं क्रुद्धो रुद्रमिवान्तकः} %||3-24-26||

\twolineshloka
{निवृत्तास्तु पुनः सर्वे दूषणाश्रयनिर्भयाः}
{राममेवाभ्यधावन्त सालतालशिलायुधाः} %||3-24-27||

\twolineshloka
{तद्बभूवाद्भुतं युद्धं तुमुलं रोमहर्षणम्}
{रामस्यास्य महाघोरं पुनस्तेषां च रक्षसाम्} %||3-25-28||


{॥इत्यार्षे श्रीमद्रामायणे वाल्मीकीये आदिकाव्ये अरण्यकाण्डे चतुर्विंशः सर्गः॥}

\sect{पञ्चविंशः सर्गः}

\twolineshloka
{तद्द्रुमाणां शिलानां च वर्षं प्राणहरं महत्}
{प्रतिजग्राह धर्मात्मा राघवस्तीक्ष्णसायकैः} %||3-25-1||

\twolineshloka
{प्रतिगृह्य च तद्वरं निमीलित इवर्षभः}
{रामः क्रोधं परं भेजे वधार्थं सर्वरक्षसाम्} %||3-25-2||

\twolineshloka
{ततः क्रोधसमाविष्टः प्रदीप्त इव तेजसा}
{शरैरभ्यकिरत्सैन्यं सर्वतः सहदूषणम्} %||3-25-3||

\twolineshloka
{ततः सेनापतिः क्रुद्धो दूषणः शत्रुदूषणः}
{जग्राह गिरिशृङ्गाभं परिघं रोमहर्षणम्} %||3-25-4||

\twolineshloka
{वेष्टितं काञ्चनैः पट्टैर्देवसैन्याभिमर्दनम्}
{आयसैः शङ्कुभिस्तीक्ष्णैः कीर्णं परवसोक्षिताम्} %||3-25-5||

\threelineshloka
{वज्राशनिसमस्पर्शं परगोपुरदारणम्}
{तं महोरगसङ्काशं प्रगृह्य परिघं रणे}
{दूषणोऽभ्यपतद्रामं क्रूरकर्मा निशाचरः} %||3-25-6||

\twolineshloka
{तस्याभिपतमानस्य दूषणस्य स राघवः}
{द्वाभ्यां शराभ्यां चिच्छेद सहस्ताभरणौ भुजौ} %||3-25-7||

\twolineshloka
{भ्रष्टस्तस्य महाकायः पपात रणमूर्धनि}
{परिघश्छिन्नहस्तस्य शक्रध्वज इवाग्रतः} %||3-25-8||

\twolineshloka
{स कराभ्यां विकीर्णाभ्यां पपात भुवि दूषणः}
{विषाणाभ्यां विशीर्णाभ्यां मनस्वीव महागजः} %||3-25-9||

\twolineshloka
{दृष्ट्वा तं पतितं भूमौ दूषणं निहतं रणे}
{साधु साध्विति काकुत्स्थं सर्वभूतान्यपूजयन्} %||3-25-10||

\threelineshloka
{एतस्मिन्नन्तरे क्रुद्धास्त्रयः सेनाग्रयायिनः}
{संहत्याभ्यद्रवन्रामं मृत्युपाशावपाशिताः}
{महाकपालः स्थूलाक्षः प्रमाथी च महाबलः} %||3-25-11||

\twolineshloka
{महाकपालो विपुलं शूलमुद्यम्य राक्षसः}
{स्थूलाक्षः पट्टिशं गृह्य प्रमाथी च परश्वधम्} %||3-25-12||

\twolineshloka
{दृष्ट्वैवापततस्तांस्तु राघवः सायकैः शितैः}
{तीक्ष्णाग्रैः प्रतिजग्राह सम्प्राप्तानतिथीनिव} %||3-25-13||

\twolineshloka
{महाकपालस्य शिरश्चिच्छेद रघुनङ्गनः}
{असङ्ख्येयैस्तु बाणौघैः प्रममाथ प्रमाथिनम्} %||3-25-14||

\twolineshloka
{स्थूलाक्षस्याक्षिणी तीक्ष्णैः पूरयामास सायकैः}
{स पपात हतो भूमौ विटपीव महाद्रुमः} %||3-25-15||

\twolineshloka
{ततः पावकसङ्काशैर्हेमवज्रविभूषितैः}
{जघनशेषं तेजस्वी तस्य सैन्यस्य सायकैः} %||3-25-16||

\twolineshloka
{ते रुक्मपुङ्खा विशिखाः सधूमा इव पावकाः}
{निजघ्नुस्तानि रक्षांसि वज्रा इव महाद्रुमान्} %||3-25-17||

\twolineshloka
{रक्षसां तु शतं रामः शतेनैकेन कर्णिना}
{सहस्रं च सहस्रेण जघान रणमूर्धनि} %||3-25-18||

\twolineshloka
{तैर्भिन्नवर्माभरणाश्छिन्नभिन्नशरासनाः}
{निपेतुः शोणितादिग्धा धरण्यां रजनीचराः} %||3-25-19||

\twolineshloka
{तैर्मुक्तकेशैः समरे पतितैः शोणितोक्षितैः}
{आस्तीर्णा वसुधा कृत्स्ना महावेदिः कुशैरिव} %||3-25-20||

\twolineshloka
{क्षणेन तु महाघोरं वनं निहतराक्षसम्}
{बभूव निरय प्रख्यं मांसशोणितकर्दमम्} %||3-25-21||

\twolineshloka
{चतुर्दश सहस्राणि रक्षसां भीमकर्मणाम्}
{हतान्येकेन रामेण मानुषेण पदातिना} %||3-25-22||

\twolineshloka
{तस्य सैन्यस्य सर्वस्य खरः शेषो महारथः}
{राक्षसस्त्रिशिराश्चैव रामश्च रिपुसूदनः} %||3-25-23||

\fourlineindentedshloka
{ततस्तु तद्भीमबलं महाहवे}
{ समीक्ष्य रामेण हतं बलीयसा}
{रथेन रामं महता खरस्ततः}
{ समाससादेन्द्र इवोद्यताशनिः} %||3-26-24||


{॥इत्यार्षे श्रीमद्रामायणे वाल्मीकीये आदिकाव्ये अरण्यकाण्डे पञ्चविंशः सर्गः॥}

\sect{षड्विंशः सर्गः}

\twolineshloka
{खरं तु रामाभिमुखं प्रयान्तं वाहिनीपतिः}
{राक्षसस्त्रिशिरा नाम संनिपत्येदमब्रवीत्} %||3-26-1||

\twolineshloka
{मां नियोजय विक्रान्त संनिवर्तस्व साहसात्}
{पश्य रामं महाबाहुं संयुगे विनिपातितम्} %||3-26-2||

\twolineshloka
{प्रतिजानामि ते सत्यमायुधं चाहमालभे}
{यथा रामं वधिष्यामि वधार्हं सर्वरक्षसाम्} %||3-26-3||

\twolineshloka
{अहं वास्य रणे मृत्युरेष वा समरे मम}
{विनिवर्त्य रणोत्साहं मुहूर्तं प्राश्निको भव} %||3-26-4||

\twolineshloka
{प्रहृष्टो वा हते रामे जनस्थानं प्रयास्यसि}
{मयि वा निहते रामं संयुगायोपयास्यसि} %||3-26-5||

\twolineshloka
{खरस्त्रिशिरसा तेन मृत्युलोभात्प्रसादितः}
{गच्छ युध्येत्यनुज्ञातो राघवाभिमुखो ययौ} %||3-26-6||

\twolineshloka
{त्रिशिराश्च रथेनैव वाजियुक्तेन भास्वता}
{अभ्यद्रवद्रणे रामं त्रिशृङ्ग इव पर्वतः} %||3-26-7||

\twolineshloka
{शरधारा समूहान्स महामेघ इवोत्सृजन्}
{व्यसृजत्सदृशं नादं जलार्द्रस्येव दुन्दुभेः} %||3-26-8||

\twolineshloka
{आगच्छन्तं त्रिशिरसं राक्षसं प्रेक्ष्य राघवः}
{धनुषा प्रतिजग्राह विधुन्वन्सायकाञ्शितान्} %||3-26-9||

\twolineshloka
{स सम्प्रहारस्तुमुलो राम त्रिशिरसोर्महान्}
{बभूवातीव बलिनोः सिंहकुञ्जरयोरिव} %||3-26-10||

\twolineshloka
{ततस्त्रिशिरसा बाणैर्ललाटे ताडितस्त्रिभिः}
{अमर्षी कुपितो रामः संरब्धमिदमब्रवीत्} %||3-26-11||

\threelineshloka
{अहो विक्रमशूरस्य राक्षसस्येदृशं बलम्}
{पुष्पैरिव शरैर्यस्य ललाटेऽस्मि परिक्षतः}
{ममापि प्रतिगृह्णीष्व शरांश्चापगुणच्युतान्} %||3-26-12||

\twolineshloka
{एवमुक्त्वा तु संरब्धः शरानाशीविषोपमान्}
{त्रिशिरो वक्षसि क्रुद्धो निजघान चतुर्दश} %||3-26-13||

\twolineshloka
{चतुर्भिस्तुरगानस्य शरैः संनतपर्वभिः}
{न्यपातयत तेजस्वी चतुरस्तस्य वाजिनः} %||3-26-14||

\twolineshloka
{अष्टभिः सायकैः सूतं रथोपस्थे न्यपातयत्}
{रामश्चिच्छेद बाणेन ध्वजं चास्य समुच्छ्रितम्} %||3-26-15||

\twolineshloka
{ततो हतरथात्तस्मादुत्पतन्तं निशाचरम्}
{बिभेद रामस्तं बाणैर्हृदये सोऽभवज्जडः} %||3-26-16||

\twolineshloka
{सायकैश्चाप्रमेयात्मा सामर्षस्तस्य रक्षसः}
{शिरांस्यपातयत्त्रीणि वेगवद्भिस्त्रिभिः शतैः} %||3-26-17||

\twolineshloka
{स भूमौ शोणितोद्गारी रामबाणाभिपीडितः}
{न्यपतत्पतितैः पूर्वं स्वशिरोभिर्निशाचरः} %||3-26-18||

\twolineshloka
{हतशेषास्ततो भग्ना राक्षसाः खरसंश्रयाः}
{द्रवन्ति स्म न तिष्ठन्ति व्याघ्रत्रस्ता मृगा इव} %||3-26-19||

\twolineshloka
{तान्खरो द्रवतो दृष्ट्वा निवर्त्य रुषितः स्वयम्}
{राममेवाभिदुद्राव राहुश्चन्द्रमसं यथा} %||3-27-20||


{॥इत्यार्षे श्रीमद्रामायणे वाल्मीकीये आदिकाव्ये अरण्यकाण्डे षड्विंशः सर्गः॥}

\sect{सप्तविंशः सर्गः}

\twolineshloka
{निहतं दूषणं दृष्ट्वा रणे त्रिशिरसा सह}
{खरस्याप्यभवत्त्रासो दृष्ट्वा रामस्य विक्रमम्} %||3-27-1||

\twolineshloka
{स दृष्ट्वा राक्षसं सैन्यमविषह्यं महाबलम्}
{हतमेकेन रामेण दूषणस्त्रिशिरा अपि} %||3-27-2||

\twolineshloka
{तद्बलं हतभूयिष्ठं विमनाः प्रेक्ष्य राक्षसः}
{आससाद खरो रामं नमुचिर्वासवं यथा} %||3-27-3||

\twolineshloka
{विकृष्य बलवच्चापं नाराचान्रक्तभोजनान्}
{खरश्चिक्षेप रामाय क्रुद्धानाशीविषानिव} %||3-27-4||

\twolineshloka
{ज्यां विधुन्वन्सुबहुशः शिक्षयास्त्राणि दर्शयन्}
{चचार समरे मार्गाञ्शरै रथगतः खरः} %||3-27-5||

\twolineshloka
{स सर्वाश्च दिशो बाणैः प्रदिशश्च महारथः}
{पूरयामास तं दृष्ट्वा रामोऽपि सुमहद्धनुः} %||3-27-6||

\twolineshloka
{स सायकैर्दुर्विषहैः सस्फुलिङ्गैरिवाग्निभिः}
{नभश्चकाराविवरं पर्जन्य इव वृष्टिभिः} %||3-27-7||

\twolineshloka
{तद्बभूव शितैर्बाणैः खररामविसर्जितैः}
{पर्याकाशमनाकाशं सर्वतः शरसङ्कुलम्} %||3-27-8||

\twolineshloka
{शरजालावृतः सूर्यो न तदा स्म प्रकाशते}
{अन्योन्यवधसंरम्भादुभयोः सम्प्रयुध्यतोः} %||3-27-9||

\twolineshloka
{ततो नालीकनाराचैस्तीक्ष्णाग्रैश्च विकर्णिभिः}
{आजघान रणे रामं तोत्रैरिव महाद्विपम्} %||3-27-10||

\twolineshloka
{तं रथस्थं धनुष्पाणिं राक्षसं पर्यवस्थितम्}
{ददृशुः सर्वभूतानि पाशहस्तमिवान्तकम्} %||3-27-11||

\twolineshloka
{तं सिंहमिव विक्रान्तं सिंहविक्रान्तगामिनम्}
{दृष्ट्वा नोद्विजते रामः सिंहः क्षुद्रमृगं यथा} %||3-27-12||

\twolineshloka
{ततः सूर्यनिकाशेन रथेन महता खरः}
{आससाद रणे रामं पतङ्ग इव पावकम्} %||3-27-13||

\twolineshloka
{ततोऽस्य सशरं चापं मुष्टिदेशे महात्मनः}
{खरश्चिच्छेद रामस्य दर्शयन्पाणिलाघवम्} %||3-27-14||

\twolineshloka
{स पुनस्त्वपरान्सप्त शरानादाय वर्मणि}
{निजघान रणे क्रुद्धः शक्राशनिसमप्रभान्} %||3-27-15||

\twolineshloka
{ततस्तत्प्रहतं बाणैः खरमुक्तैः सुपर्वभिः}
{पपात कवचं भूमौ रामस्यादित्यवर्चसः} %||3-27-16||

\twolineshloka
{स शरैरर्पितः क्रुद्धः सर्वगात्रेषु राघवः}
{रराज समरे रामो विधूमोऽग्निरिव ज्वलन्} %||3-27-17||

\twolineshloka
{ततो गम्भीरनिर्ह्रादं रामः शत्रुनिबर्हणः}
{चकारान्ताय स रिपोः सज्यमन्यन्महद्धनुः} %||3-27-18||

\twolineshloka
{सुमहद्वैष्णवं यत्तदतिसृष्टं महर्षिणा}
{वरं तद्धनुरुद्यम्य खरं समभिधावत} %||3-27-19||

\twolineshloka
{ततः कनकपुङ्खैस्तु शरैः संनतपर्वभिः}
{चिच्छेद रामः सङ्क्रुद्धः खरस्य समरे ध्वजम्} %||3-27-20||

\twolineshloka
{स दर्शनीयो बहुधा विच्छिन्नः काञ्चनो ध्वजः}
{जगाम धरणीं सूर्यो देवतानामिवाज्ञया} %||3-27-21||

\twolineshloka
{तं चतुर्भिः खरः क्रुद्धो रामं गात्रेषु मार्गणैः}
{विव्याध हृदि मर्मज्ञो मातङ्गमिव तोमरैः} %||3-27-22||

\twolineshloka
{स रामो बहुभिर्बाणैः खरकार्मुकनिःसृतैः}
{विद्धो रुधिरसिक्ताङ्गो बभूव रुषितो भृशम्} %||3-27-23||

\twolineshloka
{स धनुर्धन्विनां श्रेष्ठः प्रगृह्य परमाहवे}
{मुमोच परमेष्वासः षट्शरानभिलक्षितान्} %||3-27-24||

\twolineshloka
{शिरस्येकेन बाणेन द्वाभ्यां बाह्वोरथार्पयत्}
{त्रिभिश्चन्द्रार्धवक्त्रैश्च वक्षस्यभिजघान ह} %||3-27-25||

\twolineshloka
{ततः पश्चान्महातेजा नाराचान्भास्करोपमान्}
{जिघांसू राक्षसं क्रुद्धस्त्रयोदश शिलाशितान्} %||3-27-26||

\twolineshloka
{ततोऽस्य युगमेकेन चतुर्भिश्चतुरो हयान्}
{षष्ठेन च शिरः सङ्ख्ये चिच्छेद खरसारथेः} %||3-27-27||

\fourlineshloka
{त्रिभिस्त्रिवेणुं बलवान्द्वाभ्यामक्षं महाबलः}
{द्वादशेन तु बाणेन खरस्य सशरं धनुः}
{छित्त्वा वज्रनिकाशेन राघवः प्रहसन्निव}
{त्रयोदशेनेन्द्रसमो बिभेद समरे खरम्} %||3-27-28||

\twolineshloka
{प्रभग्नधन्वा विरथो हताश्वो हतसारथिः}
{गदापाणिरवप्लुत्य तस्थौ भूमौ खरस्तदा} %||3-27-29||

\fourlineindentedshloka
{तत्कर्म रामस्य महारथस्य}
{ समेत्य देवाश्च महर्षयश्च}
{अपूजयन्प्राञ्जलयः प्रहृष्टाः}
{ तदा विमानाग्रगताः समेताः} %||3-28-30||


{॥इत्यार्षे श्रीमद्रामायणे वाल्मीकीये आदिकाव्ये अरण्यकाण्डे सप्तविंशः सर्गः॥}

\sect{अष्टाविंशः सर्गः}

\twolineshloka
{खरं तु विरथं रामो गदापाणिमवस्थितम्}
{मृदुपूर्वं महातेजाः परुषं वाक्यमब्रवीत्} %||3-28-1||

\twolineshloka
{गजाश्वरथसम्बाधे बले महति तिष्ठता}
{कृतं सुदारुणं कर्म सर्वलोकजुगुप्सितम्} %||3-28-2||

\twolineshloka
{उद्वेजनीयो भूतानां नृशंसः पापकर्मकृत्}
{त्रयाणामपि लोकानामीश्वरोऽपि न तिष्ठति} %||3-28-3||

\twolineshloka
{कर्म लोकविरुद्धं तु कुर्वाणं क्षणदाचर}
{तीक्ष्णं सर्वजनो हन्ति सर्पं दुष्टमिवागतम्} %||3-28-4||

\twolineshloka
{लोभात्पापानि कुर्वाणः कामाद्वा यो न बुध्यते}
{भ्रष्टः पश्यति तस्यान्तं ब्राह्मणी करकादिव} %||3-28-5||

\twolineshloka
{वसतो दण्डकारण्ये तापसान्धर्मचारिणः}
{किं नु हत्वा महाभागान्फलं प्राप्स्यसि राक्षस} %||3-28-6||

\twolineshloka
{न चिरं पापकर्माणः क्रूरा लोकजुगुप्सिताः}
{ऐश्वर्यं प्राप्य तिष्ठन्ति शीर्णमूला इव द्रुमाः} %||3-28-7||

\twolineshloka
{अवश्यं लभते कर्ता फलं पापस्य कर्मणः}
{घोरं पर्यागते काले द्रुमः पुष्पमिवार्तवम्} %||3-28-8||

\twolineshloka
{नचिरात्प्राप्यते लोके पापानां कर्मणां फलम्}
{सविषाणामिवान्नानां भुक्तानां क्षणदाचर} %||3-28-9||

\twolineshloka
{पापमाच्चरतां घोरं लोकस्याप्रियमिच्छताम्}
{अहमासादितो राजा प्राणान्हन्तुं निशाचर} %||3-28-10||

\twolineshloka
{अद्य हि त्वां मया मुक्ताः शराः काञ्चनभूषणाः}
{विदार्य निपतिष्यन्ति वल्मीकमिव पन्नगाः} %||3-28-11||

\twolineshloka
{ये त्वया दण्डकारण्ये भक्षिता धर्मचारिणः}
{तानद्य निहतः सङ्ख्ये ससैन्योऽनुगमिष्यसि} %||3-28-12||

\twolineshloka
{अद्य त्वां निहतं बाणैः पश्यन्तु परमर्षयः}
{निरयस्थं विमानस्था ये त्वया हिंसिताः पुरा} %||3-28-13||

\twolineshloka
{प्रहर त्वं यथाकामं कुरु यत्नं कुलाधम}
{अद्य ते पातयिष्यामि शिरस्तालफलं यथा} %||3-28-14||

\twolineshloka
{एवमुक्तस्तु रामेण क्रुद्धः संरक्तलोचनः}
{प्रत्युवाच ततो रामं प्रहसन्क्रोधमूर्छितः} %||3-28-15||

\twolineshloka
{प्राकृतान्राक्षसान्हत्वा युद्धे दशरथात्मज}
{आत्मना कथमात्मानमप्रशस्यं प्रशंससि} %||3-28-16||

\twolineshloka
{विक्रान्ता बलवन्तो वा ये भवन्ति नरर्षभाः}
{कथयन्ति न ते किञ्चित्तेजसा स्वेन गर्विताः} %||3-28-17||

\twolineshloka
{प्राकृतास्त्वकृतात्मानो लोके क्षत्रियपांसनाः}
{निरर्थकं विकत्थन्ते यथा राम विकत्थसे} %||3-28-18||

\twolineshloka
{कुलं व्यपदिशन्वीरः समरे कोऽभिधास्यति}
{मृत्युकाले हि सम्प्राप्ते स्वयमप्रस्तवे स्तवम्} %||3-28-19||

\twolineshloka
{सर्वथा तु लघुत्वं ते कत्थनेन विदर्शितम्}
{सुवर्णप्रतिरूपेण तप्तेनेव कुशाग्निना} %||3-28-20||

\twolineshloka
{न तु मामिह तिष्ठन्तं पश्यसि त्वं गदाधरम्}
{धराधरमिवाकम्प्यं पर्वतं धातुभिश्चितम्} %||3-28-21||

\twolineshloka
{पर्याप्तोऽहं गदापाणिर्हन्तुं प्राणान्रणे तव}
{त्रयाणामपि लोकानां पाशहस्त इवान्तकः} %||3-28-22||

\twolineshloka
{कामं बह्वपि वक्तव्यं त्वयि वक्ष्यामि न त्वहम्}
{अस्तं गच्छेद्धि सविता युद्धविघ्रस्ततो भवेत्} %||3-28-23||

\twolineshloka
{चतुर्दश सहस्राणि राक्षसानां हतानि ते}
{त्वद्विनाशात्करोम्यद्य तेषामश्रुप्रमार्जनम्} %||3-28-24||

\twolineshloka
{इत्युक्त्वा परमक्रुद्धस्तां गदां परमाङ्गदाम्}
{खरश्चिक्षेप रामाय प्रदीप्तामशनिं यथा} %||3-28-25||

\twolineshloka
{खरबाहुप्रमुक्ता सा प्रदीप्ता महती गदा}
{भस्मवृक्षांश्च गुल्मांश्च कृत्वागात्तत्समीपतः} %||3-28-26||

\twolineshloka
{तामापतन्तीं ज्वलितां मृत्युपाशोपमां गदा}
{अन्तरिक्षगतां रामश्चिच्छेद बहुधा शरैः} %||3-28-27||

\twolineshloka
{सा विशीर्णा शरैर्भिन्ना पपात धरणीतले}
{गदामन्त्रौषधिबलैर्व्यालीव विनिपातिता} %||3-29-28||


{॥इत्यार्षे श्रीमद्रामायणे वाल्मीकीये आदिकाव्ये अरण्यकाण्डे अष्टाविंशः सर्गः॥}

\sect{एकोनत्रिंशः सर्गः}

\twolineshloka
{भित्त्वा तु तां गदां बाणै राघवो धर्मवत्सलः}
{स्मयमानः खरं वाक्यं संरब्धमिदमब्रवीत्} %||3-29-1||

\twolineshloka
{एतत्ते बलसर्वस्वं दर्शितं राक्षसाधम}
{शक्तिहीनतरो मत्तो वृथा त्वमुपगर्जितम्} %||3-29-2||

\twolineshloka
{एषा बाणविनिर्भिन्ना गदा भूमितलं गता}
{अभिधानप्रगल्भस्य तव प्रत्ययघातिनी} %||3-29-3||

\twolineshloka
{यत्त्वयोक्तं विनष्टानामिदमश्रुप्रमार्जनम्}
{राक्षसानां करोमीति मिथ्या तदपि ते वचः} %||3-29-4||

\twolineshloka
{नीचस्य क्षुद्रशीलस्य मिथ्यावृत्तस्य रक्षसः}
{प्राणानपहरिष्यामि गरुत्मानमृतं यथा} %||3-29-5||

\twolineshloka
{अद्य ते भिन्नकण्ठस्य फेनबुद्बुदभूषितम्}
{विदारितस्य मद्बाणैर्मही पास्यति शोणितम्} %||3-29-6||

\twolineshloka
{पांसुरूषितसर्वाङ्गः स्रस्तन्यस्तभुजद्वयः}
{स्वप्स्यसे गां समाश्लिष्य दुर्लभां प्रमदामिव} %||3-29-7||

\twolineshloka
{प्रवृद्धनिद्रे शयिते त्वयि राक्षसपांसने}
{भविष्यन्त्यशरण्यानां शरण्या दण्डका इमे} %||3-29-8||

\twolineshloka
{जनस्थाने हतस्थाने तव राक्षसमच्छरैः}
{निर्भया विचरिष्यन्ति सर्वतो मुनयो वने} %||3-29-9||

\twolineshloka
{अद्य विप्रसरिष्यन्ति राक्षस्यो हतबान्धवाः}
{बाष्पार्द्रवदना दीना भयादन्यभयावहाः} %||3-29-10||

\twolineshloka
{अद्य शोकरसज्ञास्ता भविष्यन्ति निशाचर}
{अनुरूपकुलाः पत्न्यो यासां त्वं पतिरीदृशः} %||3-29-11||

\twolineshloka
{नृशंसशील क्षुद्रात्मन्नित्यं ब्राह्मणकण्टक}
{त्वत्कृते शङ्कितैरग्नौ मुनिभिः पात्यते हविः} %||3-29-12||

\twolineshloka
{तमेवमभिसंरब्धं ब्रुवाणं राघवं रणे}
{खरो निर्भर्त्सयामास रोषात्खरतर स्वनः} %||3-29-13||

\twolineshloka
{दृढं खल्ववलिप्तोऽसि भयेष्वपि च निर्भयः}
{वाच्यावाच्यं ततो हि त्वं मृत्युवश्यो न बुध्यसे} %||3-29-14||

\twolineshloka
{कालपाशपरिक्षिप्ता भवन्ति पुरुषा हि ये}
{कार्याकार्यं न जानन्ति ते निरस्तषडिन्द्रियाः} %||3-29-15||

\twolineshloka
{एवमुक्त्वा ततो रामं संरुध्य भृकुटिं ततः}
{स ददर्श महासालमविदूरे निशाचरः} %||3-29-16||

\twolineshloka
{रणे प्रहरणस्यार्थे सर्वतो ह्यवलोकयन्}
{स तमुत्पाटयामास सन्दृश्य दशनच्छदम्} %||3-29-17||

\twolineshloka
{तं समुत्क्षिप्य बाहुभ्यां विनर्दित्वा महाबलः}
{राममुद्दिश्य चिक्षेप हतस्त्वमिति चाब्रवीत्} %||3-29-18||

\twolineshloka
{तमापतन्तं बाणौघैश्छित्त्वा रामः प्रतापवान्}
{रोषमाहारयत्तीव्रं निहन्तुं समरे खरम्} %||3-29-19||

\twolineshloka
{जातस्वेदस्ततो रामो रोषाद्रक्तान्तलोचनः}
{निर्बिभेद सहस्रेण बाणानां समरे खरम्} %||3-29-20||

\twolineshloka
{तस्य बाणान्तराद्रक्तं बहु सुस्राव फेनिलम्}
{गिरेः प्रस्रवणस्येव तोयधारापरिस्रवः} %||3-29-21||

\twolineshloka
{विह्वलः स कृतो बाणैः खरो रामेण संयुगे}
{मत्तो रुधिरगन्धेन तमेवाभ्यद्रवद्द्रुतम्} %||3-29-22||

\twolineshloka
{तमापतन्तं संरब्धं कृतास्त्रो रुधिराप्लुतम्}
{अपसर्पत्प्रतिपदं किञ्चित्त्वरितविक्रमः} %||3-29-23||

\twolineshloka
{ततः पावकसङ्काशं बधाय समरे शरम्}
{खरस्य रामो जग्राह ब्रह्मदण्डमिवापरम्} %||3-29-24||

\twolineshloka
{स तद्दत्तं मघवता सुरराजेन धीमता}
{सन्दधे च स धर्मात्मा मुमोच च खरं प्रति} %||3-29-25||

\twolineshloka
{स विमुक्तो महाबाणो निर्घातसमनिःस्वनः}
{रामेण धनुरुद्यम्य खरस्योरसि चापतत्} %||3-29-26||

\twolineshloka
{स पपात खरो भूमौ दह्यमानः शराग्निना}
{रुद्रेणैव विनिर्दग्धः श्वेतारण्ये यथान्धकः} %||3-29-27||

\twolineshloka
{स वृत्र इव वज्रेण फेनेन नमुचिर्यथा}
{बलो वेन्द्राशनिहतो निपपात हतः खरः} %||3-29-28||

\twolineshloka
{ततो राजर्षयः सर्वे सङ्गताः परमर्षयः}
{सभाज्य मुदिता राममिदं वचनमब्रुवन्} %||3-29-29||

\twolineshloka
{एतदर्थं महातेजा महेन्द्रः पाकशासनः}
{शरभङ्गाश्रमं पुण्यमाजगाम पुरन्दरः} %||3-29-30||

\twolineshloka
{आनीतस्त्वमिमं देशमुपायेन महर्षिभिः}
{एषां वधार्थं क्रूराणां रक्षसां पापकर्मणाम्} %||3-29-31||

\twolineshloka
{तदिदं नः कृतं कार्यं त्वया दशरथात्मज}
{सुखं धर्मं चरिष्यन्ति दण्डकेषु महर्षयः} %||3-29-32||

\twolineshloka
{एतस्मिन्नन्तरे वीरो लक्ष्मणः सह सीतया}
{गिरिदुर्गाद्विनिष्क्रम्य संविवेशाश्रमं सुखी} %||3-29-33||

\twolineshloka
{ततो रामस्तु विजयी पूज्यमानो महर्षिभिः}
{प्रविवेशाश्रमं वीरो लक्ष्मणेनाभिवादितः} %||3-29-34||

\twolineshloka
{तं दृष्ट्वा शत्रुहन्तारं महर्षीणां सुखावहम्}
{बभूव हृष्टा वैदेही भर्तारं परिषस्वजे} %||3-30-35||


{॥इत्यार्षे श्रीमद्रामायणे वाल्मीकीये आदिकाव्ये अरण्यकाण्डे एकोनत्रिंशः सर्गः॥}

\sect{त्रिंशः सर्गः}

\twolineshloka
{ततः शूर्पणखा दृष्ट्वा सहस्राणि चतुर्दश}
{हतान्येकेन रामेण रक्षसां भीमकर्मणाम्} %||3-30-1||

\twolineshloka
{दूषणं च खरं चैव हतं त्रिशिरसं रणे}
{दृष्ट्वा पुनर्महानादं ननाद जलदोपमा} %||3-30-2||

\twolineshloka
{सा दृष्ट्वा कर्म रामस्य कृतमन्यैः सुदुष्करम्}
{जगाम परमौद्विग्ना लङ्कां रावणपालिताम्} %||3-30-3||

\twolineshloka
{स ददर्श विमानाग्रे रावणं दीप्ततेजसम्}
{उपोपविष्टं सचिवैर्मरुद्भिरिव वासवम्} %||3-30-4||

\twolineshloka
{आसीनं सूर्यसङ्काशे काञ्चने परमासने}
{रुक्मवेदिगतं प्राज्यं ज्वलन्तमिव पावकम्} %||3-30-5||

\twolineshloka
{देवगन्धर्वभूतानामृषीणां च महात्मनाम्}
{अजेयं समरे शूरं व्यात्ताननमिवान्तकम्} %||3-30-6||

\twolineshloka
{देवासुरविमर्देषु वज्राशनिकृतव्रणम्}
{ऐरावतविषाणाग्रैरुत्कृष्टकिणवक्षसम्} %||3-30-7||

\twolineshloka
{विंशद्भुजं दशग्रीवं दर्शनीयपरिच्छदम्}
{विशालवक्षसं वीरं राजलक्ष्मणलक्षितम्} %||3-30-8||

\twolineshloka
{स्निग्धवैदूर्यसङ्काशं तप्तकाञ्चनकुण्डलम्}
{सुभुजं शुक्लदशनं महास्यं पर्वतोपमम्} %||3-30-9||

\twolineshloka
{विष्णुचक्रनिपातैश्च शतशो देवसंयुगे}
{आहताङ्गं समस्तैश्च देवप्रहरणैस्तथा} %||3-30-10||

\twolineshloka
{अक्षोभ्याणां समुद्राणां क्षोभणं क्षिप्रकारिणम्}
{क्षेप्तारं पर्वताग्राणां सुराणां च प्रमर्दनम्} %||3-30-11||

\twolineshloka
{उच्छेत्तारं च धर्माणां परदाराभिमर्शनम्}
{सर्वदिव्यास्त्रयोक्तारं यज्ञविघ्नकरं सदा} %||3-30-12||

\twolineshloka
{पुरीं भोगवतीं गत्वा पराजित्य च वासुकिम्}
{तक्षकस्य प्रियां भार्यां पराजित्य जहार यः} %||3-30-13||

\twolineshloka
{कैलासं पर्वतं गत्वा विजित्य नरवाहनम्}
{विमानं पुष्पकं तस्य कामगं वै जहार यः} %||3-30-14||

\twolineshloka
{वनं चैत्ररथं दिव्यं नलिनीं नन्दनं वनम्}
{विनाशयति यः क्रोधाद्देवोद्यानानि वीर्यवान्} %||3-30-15||

\twolineshloka
{चन्द्रसूर्यौ महाभागावुत्तिष्ठन्तौ परन्तपौ}
{निवारयति बाहुभ्यां यः शैलशिखरोपमः} %||3-30-16||

\twolineshloka
{दशवर्षसहस्राणि तपस्तप्त्वा महावने}
{पुरा स्वयम्भुवे धीरः शिरांस्युपजहार यः} %||3-30-17||

\twolineshloka
{देवदानवगन्धर्वपिशाचपतगोरगैः}
{अभयं यस्य सङ्ग्रामे मृत्युतो मानुषादृते} %||3-30-18||

\twolineshloka
{मन्त्ररभितुष्टं पुण्यमध्वरेषु द्विजातिभिः}
{हविर्धानेषु यः सोममुपहन्ति महाबलः} %||3-30-19||

\threelineshloka
{आप्तयज्ञहरं क्रूरं ब्रह्मघ्नं दुष्टचारिणम्}
{कर्कशं निरनुक्रोशं प्रजानामहिते रतम्}
{रावणं सर्वभूतानां सर्वलोकभयावहम्} %||3-30-20||

\threelineshloka
{राक्षसी भ्रातरं क्रूरं सा ददर्श महाबलम्}
{तं दिव्यवस्त्राभरणं दिव्यमाल्योपशोभितम्}
{राक्षसेन्द्रं महाभागं पौलस्त्य कुलनन्दनम्} %||3-30-21||

\fourlineindentedshloka
{तमब्रवीद्दीप्तविशाललोचनम्}
{ प्रदर्शयित्वा भयमोहमूर्छिता}
{सुदारुणं वाक्यमभीतचारिणी}
{ महात्मना शूर्पणखा विरूपिता} %||3-31-22||


{॥इत्यार्षे श्रीमद्रामायणे वाल्मीकीये आदिकाव्ये अरण्यकाण्डे त्रिंशः सर्गः॥}

\sect{एकत्रिंशः सर्गः}

\twolineshloka
{ततः शूर्पणखा दीना रावणं लोकरावणम्}
{अमात्यमध्ये सङ्क्रुद्धा परुषं वाक्यमब्रवीत्} %||3-31-1||

\twolineshloka
{प्रमत्तः कामभोगेषु स्वैरवृत्तो निरङ्कुशः}
{समुत्पन्नं भयं घोरं बोद्धव्यं नावबुध्यसे} %||3-31-2||

\twolineshloka
{सक्तं ग्राम्येषु भोगेषु कामवृत्तं महीपतिम्}
{लुब्धं न बहु मन्यन्ते श्मशानाग्निमिव प्रजाः} %||3-31-3||

\twolineshloka
{स्वयं कार्याणि यः काले नानुतिष्ठति पार्थिवः}
{स तु वै सह राज्येन तैश्च कार्यैर्विनश्यति} %||3-31-4||

\twolineshloka
{अयुक्तचारं दुर्दर्शमस्वाधीनं नराधिपम्}
{वर्जयन्ति नरा दूरान्नदीपङ्कमिव द्विपाः} %||3-31-5||

\twolineshloka
{ये न रक्षन्ति विषयमस्वाधीना नराधिपः}
{ते न वृद्ध्या प्रकाशन्ते गिरयः सागरे यथा} %||3-31-6||

\twolineshloka
{आत्मवद्भिर्विगृह्य त्वं देवगन्धर्वदानवैः}
{अयुक्तचारश्चपलः कथं राजा भविष्यसि} %||3-31-7||

\twolineshloka
{येषां चारश्च कोशश्च नयश्च जयतां वर}
{अस्वाधीना नरेन्द्राणां प्राकृतैस्ते जनैः समाः} %||3-31-8||

\twolineshloka
{यस्मात्पश्यन्ति दूरस्थान्सर्वानर्थान्नराधिपाः}
{चारेण तस्मादुच्यन्ते राजानो दीर्घचक्षुषः} %||3-31-9||

\twolineshloka
{अयुक्तचारं मन्ये त्वां प्राकृतैः सचिवैर्वृतम्}
{स्वजनं च जनस्थानं हतं यो नावबुध्यसे} %||3-31-10||

\twolineshloka
{चतुर्दश सहस्राणि रक्षसां भीमकर्मणाम्}
{हतान्येकेन रामेण खरश्च सहदूषणः} %||3-31-11||

\twolineshloka
{ऋषीणामभयं दत्तं कृतक्षेमाश्च दण्डकाः}
{धर्षितं च जनस्थानं रामेणाक्लिष्टकर्मणा} %||3-31-12||

\twolineshloka
{त्वं तु लुब्धः प्रमत्तश्च पराधीनश्च रावण}
{विषये स्वे समुत्पन्नं भयं यो नावबुध्यसे} %||3-31-13||

\twolineshloka
{तीक्ष्णमल्पप्रदातारं प्रमत्तं गर्वितं शठम्}
{व्यसने सर्वभूतानि नाभिधावन्ति पार्थिवम्} %||3-31-14||

\twolineshloka
{अतिमानिनमग्राह्यमात्मसम्भावितं नरम्}
{क्रोधनं व्यसने हन्ति स्वजनोऽपि नराधिपम्} %||3-31-15||

\twolineshloka
{नानुतिष्ठति कार्याणि भयेषु न बिभेति च}
{क्षिप्रं राज्याच्च्युतो दीनस्तृणैस्तुल्यो भविष्यति} %||3-31-16||

\twolineshloka
{शुष्ककाष्ठैर्भवेत्कार्यं लोष्टैरपि च पांसुभिः}
{न तु स्थानात्परिभ्रष्टैः कार्यं स्याद्वसुधाधिपैः} %||3-31-17||

\twolineshloka
{उपभुक्तं यथा वासः स्रजो वा मृदिता यथा}
{एवं राज्यात्परिभ्रष्टः समर्थोऽपि निरर्थकः} %||3-31-18||

\twolineshloka
{अप्रमत्तश्च यो राजा सर्वज्ञो विजितेन्द्रियः}
{कृतज्ञो धर्मशीलश्च स राजा तिष्ठते चिरम्} %||3-31-19||

\twolineshloka
{नयनाभ्यां प्रसुप्तोऽपि जागर्ति नयचक्षुषा}
{व्यक्तक्रोधप्रसादश्च स राजा पूज्यते जनैः} %||3-31-20||

\twolineshloka
{त्वं तु रावणदुर्बुद्धिर्गुणैरेतैर्विवर्जितः}
{यस्य तेऽविदितश्चारै रक्षसां सुमहान्वधः} %||3-31-21||

\fourlineindentedshloka
{परावमन्ता विषयेषु सङ्गतो}
{ नदेश कालप्रविभाग तत्त्ववित्}
{अयुक्तबुद्धिर्गुणदोषनिश्चये}
{ विपन्नराज्यो न चिराद्विपत्स्यते} %||3-31-22||

\fourlineindentedshloka
{इति स्वदोषान्परिकीर्तितांस्तया}
{ समीक्ष्य बुद्ध्या क्षणदाचरेश्वरः}
{धनेन दर्पेण बलेन चान्वितो}
{ विचिन्तयामास चिरं स रावणः} %||3-32-23||


{॥इत्यार्षे श्रीमद्रामायणे वाल्मीकीये आदिकाव्ये अरण्यकाण्डे एकत्रिंशः सर्गः॥}

\sect{द्वात्रिंशः सर्गः}

\twolineshloka
{ततः शूर्पणखां क्रुद्धां ब्रुवतीं परुषं वचः}
{अमात्यमध्ये सङ्क्रुद्धः परिपप्रच्छ रावणः} %||3-32-1||

\twolineshloka
{कश्च रामः कथं वीर्यः किं रूपः किं पराक्रमः}
{किमर्थं दण्डकारण्यं प्रविष्टश्च सुदुश्चरम्} %||3-32-2||

\twolineshloka
{आयुधं किं च रामस्य निहता येन राक्षसाः}
{खरश्च निहतं सङ्ख्ये दूषणस्त्रिशिरास्तथा} %||3-32-3||

\twolineshloka
{इत्युक्ता राक्षसेन्द्रेण राक्षसी क्रोधमूर्छिता}
{ततो रामं यथान्यायमाख्यातुमुपचक्रमे} %||3-32-4||

\twolineshloka
{दीर्घबाहुर्विशालाक्षश्चीरकृष्णाजिनाम्बरः}
{कन्दर्पसमरूपश्च रामो दशरथात्मजः} %||3-32-5||

\twolineshloka
{शक्रचापनिभं चापं विकृष्य कनकाङ्गदम्}
{दीप्तान्क्षिपति नाराचान्सर्पानिव महाविषान्} %||3-32-6||

\twolineshloka
{नाददानं शरान्घोरान्न मुञ्चन्तं महाबलम्}
{न कार्मुकं विकर्षन्तं रामं पश्यामि संयुगे} %||3-32-7||

\twolineshloka
{हन्यमानं तु तत्सैन्यं पश्यामि शरवृष्टिभिः}
{इन्द्रेणैवोत्तमं सस्यमाहतं त्वश्मवृष्टिभिः} %||3-32-8||

\twolineshloka
{रक्षसां भीमवीर्याणां सहस्राणि चतुर्दश}
{निहतानि शरैस्तीक्ष्णैस्तेनैकेन पदातिना} %||3-32-9||

\twolineshloka
{अर्धाधिकमुहूर्तेन खरश्च सहदूषणः}
{ऋषीणामभयं दत्तं कृतक्षेमाश्च दण्डकाः} %||3-32-10||

\twolineshloka
{एका कथञ्चिन्मुक्ताहं परिभूय महात्मना}
{स्त्रीवधं शङ्कमानेन रामेण विदितात्मना} %||3-32-11||

\twolineshloka
{भ्राता चास्य महातेजा गुणतस्तुल्यविक्रमः}
{अनुरक्तश्च भक्तश्च लक्ष्मणो नाम वीर्यवान्} %||3-32-12||

\twolineshloka
{अमर्षी दुर्जयो जेता विक्रान्तो बुद्धिमान्बली}
{रामस्य दक्षिणे बाहुर्नित्यं प्राणो बहिष्चरः} %||3-32-13||

\twolineshloka
{रामस्य तु विशालाक्षी धर्मपत्नी यशस्विनी}
{सीता नाम वरारोहा वैदेही तनुमध्यमा} %||3-32-14||

\twolineshloka
{नैव देवी न गन्धर्वा न यक्षी न च किंनरी}
{तथारूपा मया नारी दृष्टपूर्वा महीतले} %||3-32-15||

\twolineshloka
{यस्य सीता भवेद्भार्या यं च हृष्टा परिष्वजेत्}
{अतिजीवेत्स सर्वेषु लोकेष्वपि पुरन्दरात्} %||3-32-16||

\twolineshloka
{सा सुशीला वपुःश्लाघ्या रूपेणाप्रतिमा भुवि}
{तवानुरूपा भार्या सा त्वं च तस्यास्तथा पतिः} %||3-32-17||

\twolineshloka
{तां तु विस्तीर्णजघनां पीनोत्तुङ्गपयोधराम्}
{भार्यार्थे तु तवानेतुमुद्यताहं वराननाम्} %||3-32-18||

\twolineshloka
{तां तु दृष्ट्वाद्य वैदेहीं पूर्णचन्द्रनिभाननाम्}
{मन्मथस्य शराणां च त्वं विधेयो भविष्यसि} %||3-32-19||

\twolineshloka
{यदि तस्यामभिप्रायो भार्यार्थे तव जायते}
{शीघ्रमुद्ध्रियतां पादो जयार्थमिह दक्षिणः} %||3-32-20||

\twolineshloka
{कुरु प्रियं तथा तेषां रक्षसां राक्षसेश्वर}
{वधात्तस्य नृशंसस्य रामस्याश्रमवासिनः} %||3-32-21||

\twolineshloka
{तं शरैर्निशितैर्हत्वा लक्ष्मणं च महारथम्}
{हतनाथां सुखं सीतां यथावदुपभोक्ष्यसे} %||3-32-22||

\twolineshloka
{रोचते यदि ते वाक्यं ममैतद्राक्षसेश्वर}
{क्रियतां निर्विशङ्केन वचनं मम राघव} %||3-32-23||

\fourlineindentedshloka
{निशम्य रामेण शरैरजिह्मगैर्-}
{ हताञ्जनस्थानगतान्निशाचरान्}
{खरं च बुद्ध्वा निहतं च दूषणम्}
{ त्वमद्य कृत्यं प्रतिपत्तुमर्हसि} %||3-33-24||


{॥इत्यार्षे श्रीमद्रामायणे वाल्मीकीये आदिकाव्ये अरण्यकाण्डे द्वात्रिंशः सर्गः॥}

\sect{त्रयस्त्रिंशः सर्गः}

\twolineshloka
{ततः शूर्पणखा वाक्यं तच्छ्रुत्वा रोमहर्षणम्}
{सचिवानभ्यनुज्ञाय कार्यं बुद्ध्वा जगाम ह} %||3-33-1||

\twolineshloka
{तत्कार्यमनुगम्याथ यथावदुपलभ्य च}
{दोषाणां च गुणानां च सम्प्रधार्य बलाबलम्} %||3-33-2||

\twolineshloka
{इति कर्तव्यमित्येव कृत्वा निश्चयमात्मनः}
{स्थिरबुद्धिस्ततो रम्यां यानशालां जगाम ह} %||3-33-3||

\twolineshloka
{यानशालां ततो गत्वा प्रच्छन्नं राक्षसाधिपः}
{सूतं सञ्चोदयामास रथः संयुज्यतामिति} %||3-33-4||

\twolineshloka
{एवमुक्तः क्षणेनैव सारथिर्लघुविक्रमः}
{रथं संयोजयामास तस्याभिमतमुत्तमम्} %||3-33-5||

\twolineshloka
{काञ्चनं रथमास्थाय कामगं रत्नभूषितम्}
{पिशाचवदनैर्युक्तं खरैः कनकभूषणैः} %||3-33-6||

\twolineshloka
{मेघप्रतिमनादेन स तेन धनदानुजः}
{राक्षसाधिपतिः श्रीमान्ययौ नदनदीपतिम्} %||3-33-7||

\twolineshloka
{स श्वेतबालव्यसनः श्वेतच्छत्रो दशाननः}
{स्निग्धवैदूर्यसङ्काशस्तप्तकाञ्चनभूषणः} %||3-33-8||

\twolineshloka
{दशास्यो विंशतिभुजो दर्शनीय परिच्छदः}
{त्रिदशारिर्मुनीन्द्रघ्नो दशशीर्ष इवाद्रिराट्} %||3-33-9||

\twolineshloka
{कामगं रथमास्थाय शुशुभे राक्षसाधिपः}
{विद्युन्मण्डलवान्मेघः सबलाक इवाम्बरे} %||3-33-10||

\twolineshloka
{सशैलं सागरानूपं वीर्यवानवलोकयन्}
{नानापुष्पफलैर्वृक्षैरनुकीर्णं सहस्रशः} %||3-33-11||

\twolineshloka
{शीतमङ्गलतोयाभिः पद्मिनीभिः समन्ततः}
{विशालैराश्रमपदैर्वेदिमद्भिः समावृतम्} %||3-33-12||

\twolineshloka
{कदल्याढकिसम्बाधं नालिकेरोपशोभितम्}
{सालैस्तालैस्तमालैश्च तरुभिश्च सुपुष्पितैः} %||3-33-13||

\twolineshloka
{अत्यन्तनियताहारैः शोभितं परमर्षिभिः}
{नागैः सुपर्णैर्गन्धर्वैः किंनरैश्च सहस्रशः} %||3-33-14||

\twolineshloka
{जितकामैश्च सिद्धैश्च चामणैश्चोपशोभितम्}
{आजैर्वैखानसैर्माषैर्वालखिल्यैर्मरीचिपैः} %||3-33-15||

\twolineshloka
{दिव्याभरणमाल्याभिर्दिव्यरूपाभिरावृतम्}
{क्रीडा रतिविधिज्ञाभिरप्सरोभिः सहस्रशः} %||3-33-16||

\twolineshloka
{सेवितं देवपत्नीभिः श्रीमतीभिः श्रिया वृतम्}
{देवदानवसङ्घैश्च चरितं त्वमृताशिभिः} %||3-33-17||

\twolineshloka
{हंसक्रौञ्चप्लवाकीर्णं सारसैः सम्प्रणादितम्}
{वैदूर्यप्रस्तरं रम्यं स्निग्धं सागरतेजसा} %||3-33-18||

\twolineshloka
{पाण्डुराणि विशालानि दिव्यमाल्ययुतानि च}
{तूर्यगीताभिजुष्टानि विमानानि समन्ततः} %||3-33-19||

\twolineshloka
{तपसा जितलोकानां कामगान्यभिसम्पतन्}
{गन्धर्वाप्सरसश्चैव ददर्श धनदानुजः} %||3-33-20||

\twolineshloka
{निर्यासरसमूलानां चन्दनानां सहस्रशः}
{वनानि पश्यन्सौम्यानि घ्राणतृप्तिकराणि च} %||3-33-21||

\twolineshloka
{अगरूणां च मुख्यानां वनान्युपवनानि च}
{तक्कोलानां च जात्यानां फलानां च सुगन्धिनाम्} %||3-33-22||

\twolineshloka
{पुष्पाणि च तमालस्य गुल्मानि मरिचस्य च}
{मुक्तानां च समूहानि शुष्यमाणानि तीरतः} %||3-33-23||

\twolineshloka
{शङ्खानां प्रस्तरं चैव प्रवालनिचयं तथा}
{काञ्चनानि च शैलानि राजतानि च सर्वशः} %||3-33-24||

\twolineshloka
{प्रस्रवाणि मनोज्ञानि प्रसन्नानि ह्रदानि च}
{धनधान्योपपन्नानि स्त्रीरत्नैरावृतानि च} %||3-33-25||

\twolineshloka
{हस्त्यश्वरथगाढानि नगराण्यवलोकयन्}
{तं समं सर्वतः स्निग्धं मृदुसंस्पर्शमारुतम्} %||3-33-26||

\twolineshloka
{अनूपं सिन्धुराजस्य ददर्श त्रिदिवोपमम्}
{तत्रापश्यत्स मेघाभं न्यग्रोधमृषिभिर्वृतम्} %||3-33-27||

\threelineshloka
{समन्ताद्यस्य ताः शाखाः शतयोजनमायताः}
{यस्य हस्तिनमादाय महाकायं च कच्चपम्}
{भक्षार्थं गरुडः शाखामाजगाम महाबलः} %||3-33-28||

\twolineshloka
{तस्य तां सहसा शाखां भारेण पतगोत्तमः}
{सुपर्णः पर्णबहुलां बभञ्जाथ महाबलः} %||3-33-29||

\twolineshloka
{तत्र वैखानसा माषा वालखिल्या मरीचिपाः}
{अजा बभूवुर्धूम्राश्च सङ्गताः परमर्षयः} %||3-33-30||

\twolineshloka
{तेषां दयार्थं गरुडस्तां शाखां शतयोजनाम्}
{जगामादाय वेगेन तौ चोभौ गजकच्छपौ} %||3-33-31||

\threelineshloka
{एकपादेन धर्मात्मा भक्षयित्वा तदामिषम्}
{निषादविषयं हत्वा शाखया पतगोत्तमः}
{प्रहर्षमतुलं लेभे मोक्षयित्वा महामुनीन्} %||3-33-32||

\twolineshloka
{स तेनैव प्रहर्षेण द्विगुणीकृतविक्रमः}
{अमृतानयनार्थं वै चकार मतिमान्मतिम्} %||3-33-33||

\twolineshloka
{अयोजालानि निर्मथ्य भित्त्वा रत्नगृहं वरम्}
{महेन्द्रभवनाद्गुप्तमाजहारामृतं ततः} %||3-33-34||

\twolineshloka
{तं महर्षिगणैर्जुष्टं सुपर्णकृतलक्षणम्}
{नाम्ना सुभद्रं न्यग्रोधं ददर्श धनदानुजः} %||3-33-35||

\twolineshloka
{तं तु गत्वा परं पारं समुद्रस्य नदीपतेः}
{ददर्शाश्रममेकान्ते पुण्ये रम्ये वनान्तरे} %||3-33-36||

\twolineshloka
{तत्र कृष्णाजिनधरं जटावल्कलधारिणम्}
{ददर्श नियताहारं मारीचं नाम राक्षसम्} %||3-33-37||

\twolineshloka
{स रावणः समागम्य विधिवत्तेन रक्षसा}
{ततः पश्चादिदं वाक्यमब्रवीद्वाक्यकोविदः} %||3-34-38||


{॥इत्यार्षे श्रीमद्रामायणे वाल्मीकीये आदिकाव्ये अरण्यकाण्डे त्रयस्त्रिंशः सर्गः॥}

\sect{चतुस्त्रिंशः सर्गः}

\twolineshloka
{मारीच श्रूयतां तात वचनं मम भाषतः}
{आर्तोऽस्मि मम चार्तस्य भवान्हि परमा गतिः} %||3-34-1||

\twolineshloka
{जानीषे त्वं जनस्थानं भ्राता यत्र खरो मम}
{दूषणश्च महाबाहुः स्वसा शूर्पणखा च मे} %||3-34-2||

\twolineshloka
{त्रिशिराश्च महातेजा राक्षसः पिशिताशनः}
{अन्ये च बहवः शूरा लब्धलक्षा निशाचराः} %||3-34-3||

\twolineshloka
{वसन्ति मन्नियोगेन अधिवासं च राक्षसः}
{बाधमाना महारण्ये मुनीन्ये धर्मचारिणः} %||3-34-4||

\twolineshloka
{चतुर्दश सहस्राणि रक्षसां भीमकर्मणाम्}
{शूराणां लब्धलक्षाणां खरचित्तानुवर्तिनाम्} %||3-34-5||

\twolineshloka
{ते त्विदानीं जनस्थाने वसमाना महाबलाः}
{सङ्गताः परमायत्ता रामेण सह संयुगे} %||3-34-6||

\twolineshloka
{तेन संजातरोषेण रामेण रणमूर्धनि}
{अनुक्त्वा परुषं किञ्चिच्छरैर्व्यापारितं धनुः} %||3-34-7||

\twolineshloka
{चतुर्दश सहस्राणि रक्षसां भीमकर्मणाम्}
{निहतानि शरैस्तीक्ष्णैर्मानुषेण पदातिना} %||3-34-8||

\twolineshloka
{खरश्च निहतः सङ्ख्ये दूषणश्च निपातितः}
{हत्वा त्रिशिरसं चापि निर्भया दण्डकाः कृताः} %||3-34-9||

\twolineshloka
{पित्रा निरस्तः क्रुद्धेन सभार्यः क्षीणजीवितः}
{स हन्ता तस्य सैन्यस्य रामः क्षत्रियपांसनः} %||3-34-10||

\twolineshloka
{अशीलः कर्कशस्तीक्ष्णो मूर्खो लुब्धोऽजितेन्द्रियः}
{त्यक्तधर्मस्त्वधर्मात्मा भूतानामहिते रतः} %||3-34-11||

\twolineshloka
{येन वैरं विनारण्ये सत्त्वमाश्रित्य केवलम्}
{कर्णनासापहारेण भगिनी मे विरूपिता} %||3-34-12||

\twolineshloka
{तस्य भार्यां जनस्थानात्सीतां सुरसुतोपमाम्}
{आनयिष्यामि विक्रम्य सहायस्तत्र मे भव} %||3-34-13||

\twolineshloka
{त्वया ह्यहं सहायेन पार्श्वस्थेन महाबल}
{भ्रातृभिश्च सुरान्युद्धे समग्रान्नाभिचिन्तये} %||3-34-14||

\twolineshloka
{तत्सहायो भव त्वं मे समर्थो ह्यसि राक्षस}
{वीर्ये युद्धे च दर्पे च न ह्यस्ति सदृशस्तव} %||3-34-15||

\twolineshloka
{एतदर्थमहं प्राप्तस्त्वत्समीपं निशाचर}
{शृणु तत्कर्म साहाय्ये यत्कार्यं वचनान्मम} %||3-34-16||

\twolineshloka
{सौवर्णस्त्वं मृगो भूत्वा चित्रो रजतबिन्दुभिः}
{आश्रमे तस्य रामस्य सीतायाः प्रमुखे चर} %||3-34-17||

\twolineshloka
{त्वां तु निःसंशयं सीता दृष्ट्वा तु मृगरूपिणम्}
{गृह्यतामिति भर्तारं लक्ष्मणं चाभिधास्यति} %||3-34-18||

\twolineshloka
{ततस्तयोरपाये तु शून्ये सीतां यथासुखम्}
{निराबाधो हरिष्यामि राहुश्चन्द्रप्रभामिव} %||3-34-19||

\twolineshloka
{ततः पश्चात्सुखं रामे भार्याहरणकर्शिते}
{विस्रब्धं प्रहरिष्यामि कृतार्थेनान्तरात्मना} %||3-34-20||

\twolineshloka
{तस्य रामकथां श्रुत्वा मारीचस्य महात्मनः}
{शुष्कं समभवद्वक्त्रं परित्रस्तो बभूव च} %||3-34-21||

\fourlineindentedshloka
{स रावणं त्रस्तविषण्णचेता}
{ महावने रामपराक्रमज्ञः}
{कृताञ्जलिस्तत्त्वमुवाच वाक्यम्}
{ हितं च तस्मै हितमात्मनश्च} %||3-35-22||


{॥इत्यार्षे श्रीमद्रामायणे वाल्मीकीये आदिकाव्ये अरण्यकाण्डे चतुस्त्रिंशः सर्गः॥}

\sect{पञ्चत्रिंशः सर्गः}

\twolineshloka
{तच्छ्रुत्वा राक्षसेन्द्रस्य वाक्यं वाक्यविशारदः}
{प्रत्युवाच महाप्राज्ञो मारीचो राक्षसेश्वरम्} %||3-35-1||

\twolineshloka
{सुलभाः पुरुषा राजन्सततं प्रियवादिनः}
{अप्रियस्य च पथ्यस्य वक्ता श्रोता च दुर्लभः} %||3-35-2||

\twolineshloka
{न नूनं बुध्यसे रामं महावीर्यं गुणोन्नतम्}
{अयुक्तचारश्चपलो महेन्द्रवरुणोपमम्} %||3-35-3||

\twolineshloka
{अपि स्वस्ति भवेत्तात सर्वेषां भुवि रक्षसाम्}
{अपि रामो न सङ्क्रुद्धः कुर्याल्लोकमराक्षसम्} %||3-35-4||

\twolineshloka
{अपि ते जीवितान्ताय नोत्पन्ना जनकात्मजा}
{अपि सीता निमित्तं च न भवेद्व्यसनं महत्} %||3-35-5||

\twolineshloka
{अपि त्वामीश्वरं प्राप्य कामवृत्तं निरङ्कुशम्}
{न विनश्येत्पुरी लङ्का त्वया सह सराक्षसा} %||3-35-6||

\twolineshloka
{त्वद्विधः कामवृत्तो हि दुःशीलः पापमन्त्रितः}
{आत्मानं स्वजनं राष्ट्रं स राजा हन्ति दुर्मतिः} %||3-35-7||

\twolineshloka
{न च पित्रा परित्यक्तो नामर्यादः कथञ्चन}
{न लुब्धो न च दुःशीलो न च क्षत्रियपांसनः} %||3-35-8||

\twolineshloka
{न च धर्मगुणैर्हीनैः कौसल्यानन्दवर्धनः}
{न च तीक्ष्णो हि भूतानां सर्वेषां च हिते रतः} %||3-35-9||

\twolineshloka
{वञ्चितं पितरं दृष्ट्वा कैकेय्या सत्यवादिनम्}
{करिष्यामीति धर्मात्मा ततः प्रव्रजितो वनम्} %||3-35-10||

\twolineshloka
{कैकेय्याः प्रियकामार्थं पितुर्दशरथस्य च}
{हित्वा राज्यं च भोगांश्च प्रविष्टो दण्डकावनम्} %||3-35-11||

\twolineshloka
{न रामः कर्कशस्तात नाविद्वान्नाजितेन्द्रियः}
{अनृतं न श्रुतं चैव नैव त्वं वक्तुमर्हसि} %||3-35-12||

\twolineshloka
{रामो विग्रहवान्धर्मः साधुः सत्यपराक्रमः}
{राजा सर्वस्य लोकस्य देवानामिव वासवः} %||3-35-13||

\twolineshloka
{कथं त्वं तस्य वैदेहीं रक्षितां स्वेन तेजसा}
{इच्छसि प्रसभं हर्तुं प्रभामिव विवस्वतः} %||3-35-14||

\twolineshloka
{शरार्चिषमनाधृष्यं चापखड्गेन्धनं रणे}
{रामाग्निं सहसा दीप्तं न प्रवेष्टुं त्वमर्हसि} %||3-35-15||

\twolineshloka
{धनुर्व्यादितदीप्तास्यं शरार्चिषममर्षणम्}
{चापबाणधरं वीरं शत्रुसेनापहारिणम्} %||3-35-16||

\twolineshloka
{राज्यं सुखं च सन्त्यज्य जीवितं चेष्टमात्मनः}
{नात्यासादयितुं तात रामान्तकमिहार्हसि} %||3-35-17||

\twolineshloka
{अप्रमेयं हि तत्तेजो यस्य सा जनकात्मजा}
{न त्वं समर्थस्तां हर्तुं रामचापाश्रयां वने} %||3-35-18||

\twolineshloka
{प्राणेभ्योऽपि प्रियतरा भार्या नित्यमनुव्रता}
{दीप्तस्येव हुताशस्य शिखा सीता सुमध्यमा} %||3-35-19||

\threelineshloka
{किमुद्यमं व्यर्थमिमं कृत्वा ते राक्षसाधिप}
{दृष्टश्चेत्त्वं रणे तेन तदन्तं तव जीवितम्}
{जीवितं च सुखं चैव राज्यं चैव सुदुर्लभम्} %||3-35-20||

\twolineshloka
{स सर्वैः सचिवैः सार्धं विभीषणपुरस्कृतैः}
{मन्त्रयित्वा तु धर्मिष्ठैः कृत्वा निश्चयमात्मनः} %||3-35-21||

\threelineshloka
{दोषाणां च गुणानां च सम्प्रधार्य बलाबलम्}
{आत्मनश्च बलं ज्ञात्वा राघवस्य च तत्त्वतः}
{हितं हि तव निश्चित्य क्षमं त्वं कर्तुमर्हसि} %||3-35-22||

\fourlineindentedshloka
{अहं तु मन्ये तव न क्षमं रणे}
{ समागमं कोसलराजसूनुना}
{इदं हि भूयः शृणु वाक्यमुत्तमम्}
{ क्षमं च युक्तं च निशाचराधिप} %||3-36-23||


{॥इत्यार्षे श्रीमद्रामायणे वाल्मीकीये आदिकाव्ये अरण्यकाण्डे पञ्चत्रिंशः सर्गः॥}

\sect{षट्त्रिंशः सर्गः}

\twolineshloka
{कदाचिदप्यहं वीर्यात्पर्यटन्पृथिवीमिमाम्}
{बलं नागसहस्रस्य धारयन्पर्वतोपमः} %||3-36-1||

\threelineshloka
{नीलजीमूतसङ्काशस्तप्तकाञ्चनकुण्डलः}
{भयं लोकस्य जनयन्किरीटी परिघायुधः}
{व्यचरं दण्डकारण्यमृषिमांसानि भक्षयन्} %||3-36-2||

\twolineshloka
{विश्वामित्रोऽथ धर्मात्मा मद्वित्रस्तो महामुनिः}
{स्वयं गत्वा दशरथं नरेन्द्रमिदमब्रवीत्} %||3-36-3||

\twolineshloka
{अयं रक्षतु मां रामः पर्वकाले समाहितः}
{मारीचान्मे भयं घोरं समुत्पन्नं नरेश्वर} %||3-36-4||

\twolineshloka
{इत्येवमुक्तो धर्मात्मा राजा दशरथस्तदा}
{प्रत्युवाच महाभागं विश्वामित्रं महामुनिम्} %||3-36-5||

\threelineshloka
{ऊन षोडश वर्षोऽयमकृतास्त्रश्च राघवः}
{कामं तु मम यत्सैन्यं मया सह गमिष्यति}
{बधिष्यामि मुनिश्रेष्ठ शत्रुं तव यथेप्सितम्} %||3-36-6||

\twolineshloka
{इत्येवमुक्तः स मुनी राजानं पुनरब्रवीत्}
{रामान्नान्यद्बलं लोके पर्याप्तं तस्य रक्षसः} %||3-36-7||

\twolineshloka
{बालोऽप्येष महातेजाः समर्थस्तस्य निग्रहे}
{गमिष्ये राममादाय स्वस्ति तेऽस्तु परन्तपः} %||3-36-8||

\twolineshloka
{इत्येवमुक्त्वा स मुनिस्तमादाय नृपात्मजम्}
{जगाम परमप्रीतो विश्वामित्रः स्वमाश्रमम्} %||3-36-9||

\twolineshloka
{तं तदा दण्डकारण्ये यज्ञमुद्दिश्य दीक्षितम्}
{बभूवावस्थितो रामश्चित्रं विस्फारयन्धनुः} %||3-36-10||

\twolineshloka
{अजातव्यञ्जनः श्रीमान्बालः श्यामः शुभेक्षणः}
{एकवस्त्रधरो धन्वी शिखी कनकमालया} %||3-36-11||

\twolineshloka
{शोभयन्दण्डकारण्यं दीप्तेन स्वेन तेजसा}
{अदृश्यत तदा रामो बालचन्द्र इवोदितः} %||3-36-12||

\twolineshloka
{ततोऽहं मेघसङ्काशस्तप्तकाञ्चनकुण्डलः}
{बली दत्तवरो दर्पादाजगाम तदाश्रमम्} %||3-36-13||

\twolineshloka
{तेन दृष्टः प्रविष्टोऽहं सहसैवोद्यतायुधः}
{मां तु दृष्ट्वा धनुः सज्यमसम्भ्रान्तश्चकार ह} %||3-36-14||

\twolineshloka
{अवजानन्नहं मोहाद्बालोऽयमिति राघवम्}
{विश्वामित्रस्य तां वेदिमध्यधावं कृतत्वरः} %||3-36-15||

\twolineshloka
{तेन मुक्तस्ततो बाणः शितः शत्रुनिबर्हणः}
{तेनाहं ताडितः क्षिप्तः समुद्रे शतयोजने} %||3-36-16||

\threelineshloka
{रामस्य शरवेगेन निरस्तो भ्रान्तचेतनः}
{पातितोऽहं तदा तेन गम्भीरे सागराम्भसि}
{प्राप्य संज्ञां चिरात्तात लङ्कां प्रति गतः पुरीम्} %||3-36-17||

\twolineshloka
{एवमस्मि तदा मुक्तः सहायास्ते निपातिताः}
{अकृतास्त्रेण रामेण बालेनाक्लिष्टकर्मणा} %||3-36-18||

\twolineshloka
{तन्मया वार्यमाणस्त्वं यदि रामेण विग्रहम्}
{करिष्यस्यापदं घोरां क्षिप्रं प्राप्य नशिष्यसि} %||3-36-19||

\twolineshloka
{क्रीडा रतिविधिज्ञानां समाजोत्सवशालिनाम्}
{रक्षसां चैव सन्तापमनर्थं चाहरिष्यसि} %||3-36-20||

\twolineshloka
{हर्म्यप्रासादसम्बाधां नानारत्नविभूषिताम्}
{द्रक्ष्यसि त्वं पुरीं लङ्कां विनष्टां मैथिलीकृते} %||3-36-21||

\twolineshloka
{अकुर्वन्तोऽपि पापानि शुचयः पापसंश्रयात्}
{परपापैर्विनश्यन्ति मत्स्या नागह्रदे यथा} %||3-36-22||

\twolineshloka
{दिव्यचन्दनदिग्धाङ्गान्दिव्याभरणभूषितान्}
{द्रक्ष्यस्यभिहतान्भूमौ तव दोषात्तु राक्षसान्} %||3-36-23||

\twolineshloka
{हृतदारान्सदारांश्च दशविद्रवतो दिशः}
{हतशेषानशरणान्द्रक्ष्यसि त्वं निशाचरान्} %||3-36-24||

\twolineshloka
{शरजालपरिक्षिप्तामग्निज्वालासमावृताम्}
{प्रदग्धभवनां लङ्कां द्रक्ष्यसि त्वमसंशयम्} %||3-36-25||

\twolineshloka
{प्रमदानां सहस्राणि तव राजन्परिग्रहः}
{भव स्वदारनिरतः स्वकुलं रक्षराक्षस} %||3-36-26||

\twolineshloka
{मानं वृद्धिं च राज्यं च जीवितं चेष्टमात्मनः}
{यदीच्छसि चिरं भोक्तुं मा कृथा राम विप्रियम्} %||3-36-27||

\fourlineindentedshloka
{निवार्यमाणः सुहृदा मया भृशम्}
{ प्रसह्य सीतां यदि धर्षयिष्यसि}
{गमिष्यसि क्षीणबलः सबान्धवो}
{ यमक्षयं रामशरात्तजीवितः} %||3-37-28||


{॥इत्यार्षे श्रीमद्रामायणे वाल्मीकीये आदिकाव्ये अरण्यकाण्डे षट्त्रिंशः सर्गः॥}

\sect{सप्तत्रिंशः सर्गः}

\twolineshloka
{एवमस्मि तदा मुक्तः कथञ्चित्तेन संयुगे}
{इदानीमपि यद्वृत्तं तच्छृणुष्व यदुत्तरम्} %||3-37-1||

\twolineshloka
{राक्षसाभ्यामहं द्वाभ्यामनिर्विण्णस्तथा कृतः}
{सहितो मृगरूपाभ्यां प्रविष्टो दण्डकावनम्} %||3-37-2||

\twolineshloka
{दीप्तजिह्वो महाकायस्तीक्ष्णशृण्गो महाबलः}
{व्यचरन्दण्डकारण्यं मांसभक्षो महामृगः} %||3-37-3||

\twolineshloka
{अग्निहोत्रेषु तीर्थेषु चैत्यवृक्षेषु रावण}
{अत्यन्तघोरो व्यचरंस्तापसांस्तान्प्रधर्षयन्} %||3-37-4||

\twolineshloka
{निहत्य दण्डकारण्ये तापसान्धर्मचारिणः}
{रुधिराणि पिबंस्तेषां तथा मांसानि भक्षयन्} %||3-37-5||

\twolineshloka
{ऋषिमांसाशनः क्रूरस्त्रासयन्वनगोचरान्}
{तदा रुधिरमत्तोऽहं व्यचरं दण्डकावनम्} %||3-37-6||

\twolineshloka
{तदाहं दण्डकारण्ये विचरन्धर्मदूषकः}
{आसादयं तदा रामं तापसं धर्ममाश्रितम्} %||3-37-7||

\twolineshloka
{वैदेहीं च महाभागां लक्ष्मणं च महारथम्}
{तापसं नियताहारं सर्वभूतहिते रतम्} %||3-37-8||

\twolineshloka
{सोऽहं वनगतं रामं परिभूय महाबलम्}
{तापसोऽयमिति ज्ञात्वा पूर्ववैरमनुस्मरन्} %||3-37-9||

\twolineshloka
{अभ्यधावं सुसङ्क्रुद्धस्तीक्ष्णशृङ्गो मृगाकृतिः}
{जिघांसुरकृतप्रज्ञस्तं प्रहारमनुस्मरन्} %||3-37-10||

\twolineshloka
{तेन मुक्तास्त्रयो बाणाः शिताः शत्रुनिबर्हणाः}
{विकृष्य बलवच्चापं सुपर्णानिलतुल्यगाः} %||3-37-11||

\twolineshloka
{ते बाणा वज्रसङ्काशाः सुघोरा रक्तभोजनाः}
{आजग्मुः सहिताः सर्वे त्रयः संनतपर्वणः} %||3-37-12||

\twolineshloka
{पराक्रमज्ञो रामस्य शठो दृष्टभयः पुरा}
{समुत्क्रान्तस्ततो मुक्तस्तावुभौ राक्षसौ हतौ} %||3-37-13||

\twolineshloka
{शरेण मुक्तो रामस्य कथञ्चित्प्राप्य जीवितम्}
{इह प्रव्राजितो युक्तस्तापसोऽहं समाहितः} %||3-37-14||

\twolineshloka
{वृक्षे वृक्षे हि पश्यामि चीरकृष्णाजिनाम्बरम्}
{गृहीतधनुषं रामं पाशहस्तमिवान्तकम्} %||3-37-15||

\twolineshloka
{अपि रामसहस्राणि भीतः पश्यामि रावण}
{रामभूतमिदं सर्वमरण्यं प्रतिभाति मे} %||3-37-16||

\twolineshloka
{राममेव हि पश्यामि रहिते राक्षसेश्वर}
{दृष्ट्वा स्वप्नगतं राममुद्भ्रमामि विचेतनः} %||3-37-17||

\twolineshloka
{रकारादीनि नामानि रामत्रस्तस्य रावण}
{रत्नानि च रथाश्चैव त्रासं संजनयन्ति मे} %||3-37-18||

\threelineshloka
{अहं तस्य प्रभावज्ञो न युद्धं तेन ते क्षमम्}
{रणे रामेण युध्यस्व क्षमां वा कुरु राक्षस}
{न ते रामकथा कार्या यदि मां द्रष्टुमिच्छसि} %||3-37-19||

\fourlineindentedshloka
{इदं वचो बन्धुहितार्थिना मया}
{ यथोच्यमानं यदि नाभिपत्स्यसे}
{सबान्धवस्त्यक्ष्यसि जीवितं रणे}
{ हतोऽद्य रामेण शरैरजिह्मगैः} %||3-38-20||


{॥इत्यार्षे श्रीमद्रामायणे वाल्मीकीये आदिकाव्ये अरण्यकाण्डे सप्तत्रिंशः सर्गः॥}

\sect{अष्टात्रिंशः सर्गः}

\twolineshloka
{मारीचेन तु तद्वाक्यं क्षमं युक्तं च रावणः}
{उक्तो न प्रतिजग्राह मर्तुकाम इवौषधम्} %||3-38-1||

\twolineshloka
{तं पथ्यहितवक्तारं मारीचं राक्षसाधिपः}
{अब्रवीत्परुषं वाक्यमयुक्तं कालचोदितः} %||3-38-2||

\twolineshloka
{यत्किलैतदयुक्तार्थं मारीच मयि कथ्यते}
{वाक्यं निष्फलमत्यर्थं बीजमुप्तमिवोषरे} %||3-38-3||

\twolineshloka
{त्वद्वाक्यैर्न तु मां शक्यं भेत्तुं रामस्य संयुगे}
{पापशीलस्य मूर्खस्य मानुषस्य विशेषतः} %||3-38-4||

\twolineshloka
{यस्त्यक्त्वा सुहृदो राज्यं मातरं पितरं तथा}
{स्त्रीवाक्यं प्राकृतं श्रुत्वा वनमेकपदे गतः} %||3-38-5||

\twolineshloka
{अवश्यं तु मया तस्य संयुगे खरघातिनः}
{प्राणैः प्रियतरा सीता हर्तव्या तव संनिधौ} %||3-38-6||

\twolineshloka
{एवं मे निश्चिता बुद्धिर्हृदि मारीच वर्तते}
{न व्यावर्तयितुं शक्या सेन्द्रैरपि सुरासुरैः} %||3-38-7||

\twolineshloka
{दोषं गुणं वा सम्पृष्टस्त्वमेवं वक्तुमर्हसि}
{अपायं वाप्युपायं वा कार्यस्यास्य विनिश्चये} %||3-38-8||

\twolineshloka
{सम्पृष्टेन तु वक्तव्यं सचिवेन विपश्चिता}
{उद्यताञ्जलिना राज्ञो य इच्छेद्भूतिमात्मनः} %||3-38-9||

\twolineshloka
{वाक्यमप्रतिकूलं तु मृदुपूर्वं शुभं हितम्}
{उपचारेण युक्तं च वक्तव्यो वसुधाधिपः} %||3-38-10||

\twolineshloka
{सावमर्दं तु यद्वाक्यं मारीच हितमुच्यते}
{नाभिनन्दति तद्राजा मानार्हो मानवर्जितम्} %||3-38-11||

\threelineshloka
{पञ्चरूपाणि राजानो धारयन्त्यमितौजसः}
{अग्नेरिन्द्रस्य सोमस्य यमस्य वरुणस्य च}
{औष्ण्यं तथा विक्रमं च सौम्यं दण्डं प्रसन्नताम्} %||3-38-12||

\twolineshloka
{तस्मात्सर्वास्ववस्थासु मान्याः पूज्याश्च पार्थिवाः}
{त्वं तु धर्ममविज्ञाय केवलं मोहमास्थितः} %||3-38-13||

\threelineshloka
{अभ्यागतं मां दौरात्म्यात्परुषं वदसीदृशम्}
{गुणदोषौ न पृच्छामि क्षमं चात्मनि राक्षस}
{अस्मिंस्तु स भवान्कृत्ये साहाय्यं कर्तुमर्हति} %||3-38-14||

\twolineshloka
{सौवर्णस्त्वं मृगो भूत्वा चित्रो रजतबिन्दुभिः}
{प्रलोभयित्वा वैदेहीं यथेष्टं गन्तुमर्हसि} %||3-38-15||

\twolineshloka
{त्वां तु मायामृगं दृष्ट्वा काञ्चनं जातविस्मया}
{आनयैनमिति क्षिप्रं रामं वक्ष्यति मैथिली} %||3-38-16||

\twolineshloka
{अपक्रान्ते च काकुत्स्थे लक्ष्मणे च यथासुखम्}
{आनयिष्यामि वैदेहीं सहस्राक्षः शचीमिव} %||3-38-17||

\twolineshloka
{एवं कृत्वा त्विदं कार्यं यथेष्टं गच्छ राक्षस}
{राज्यस्यार्धं प्रदास्यामि मारीच तव सुव्रत} %||3-38-18||

\threelineshloka
{गच्छ सौम्य शिवं मार्गं कार्यस्यास्य विवृद्धये}
{प्राप्य सीतामयुद्धेन वञ्चयित्वा तु राघवम्}
{लङ्कां प्रति गमिष्यामि कृतकार्यः सह त्वया} %||3-38-19||

\twolineshloka
{एतत्कार्यमवश्यं मे बलादपि करिष्यसि}
{राज्ञो हि प्रतिकूलस्थो न जातु सुखमेधते} %||3-38-20||

\fourlineindentedshloka
{आसाद्य तं जीवितसंशयस्ते}
{ मृत्युर्ध्रुवो ह्यद्य मया विरुध्य}
{एतद्यथावत्परिगृह्य बुद्ध्या}
{ यदत्र पथ्यं कुरु तत्तथा त्वम्} %||3-39-21||


{॥इत्यार्षे श्रीमद्रामायणे वाल्मीकीये आदिकाव्ये अरण्यकाण्डे अष्टात्रिंशः सर्गः॥}

\sect{एकोनचत्वारिंशः सर्गः}

\twolineshloka
{आज्ञप्तो राजवद्वाक्यं प्रतिकूलं निशाचरः}
{अब्रवीत्परुषं वाक्यं मारीचो राक्षसाधिपम्} %||3-39-1||

\twolineshloka
{केनायमुपदिष्टस्ते विनाशः पापकर्मणा}
{सपुत्रस्य सराष्ट्रस्य सामात्यस्य निशाचर} %||3-39-2||

\twolineshloka
{कस्त्वया सुखिना राजन्नाभिनन्दति पापकृत्}
{केनेदमुपदिष्टं ते मृत्युद्वारमुपायतः} %||3-39-3||

\twolineshloka
{शत्रवस्तव सुव्यक्तं हीनवीर्या निशाचर}
{इच्छन्ति त्वां विनश्यन्तमुपरुद्धं बलीयसा} %||3-39-4||

\twolineshloka
{केनेदमुपदिष्टं ते क्षुद्रेणाहितवादिना}
{यस्त्वामिच्छति नश्यन्तं स्वकृतेन निशाचर} %||3-39-5||

\twolineshloka
{वध्याः खलु न हन्यन्ते सचिवास्तव रावण}
{ये त्वामुत्पथमारूढं न निगृह्णन्ति सर्वशः} %||3-39-6||

\twolineshloka
{अमात्यैः कामवृत्तो हि राजा कापथमाश्रितः}
{निग्राह्यः सर्वथा सद्भिर्न निग्राह्यो निगृह्यसे} %||3-39-7||

\twolineshloka
{धर्ममर्थं च कामं च यशश्च जयतां वर}
{स्वामिप्रसादात्सचिवाः प्राप्नुवन्ति निशाचर} %||3-39-8||

\twolineshloka
{विपर्यये तु तत्सर्वं व्यर्थं भवति रावण}
{व्यसनं स्वामिवैगुण्यात्प्राप्नुवन्तीतरे जनाः} %||3-39-9||

\twolineshloka
{राजमूलो हि धर्मश्च जयश्च जयतां वर}
{तस्मात्सर्वास्ववस्थासु रक्षितव्यो नराधिपः} %||3-39-10||

\twolineshloka
{राज्यं पालयितुं शक्यं न तीक्ष्णेन निशाचर}
{न चापि प्रतिकूलेन नाविनीतेन राक्षस} %||3-39-11||

\twolineshloka
{ये तीक्ष्णमन्त्राः सचिवा भज्यन्ते सह तेन वै}
{विषमेषु रथाः शीघ्रं मन्दसारथयो यथा} %||3-39-12||

\twolineshloka
{बहवः साधवो लोके युक्तधर्ममनुष्ठिताः}
{परेषामपराधेन विनष्टाः सपरिच्छदाः} %||3-39-13||

\twolineshloka
{स्वामिना प्रतिकूलेन प्रजास्तीक्ष्णेन रावण}
{रक्ष्यमाणा न वर्धन्ते मेषा गोमायुना यथा} %||3-39-14||

\twolineshloka
{अवश्यं विनशिष्यन्ति सर्वे रावण राक्षसाः}
{येषां त्वं कर्कशो राजा दुर्बुद्धिरजितेन्द्रियः} %||3-39-15||

\twolineshloka
{तदिदं काकतालीयं घोरमासादितं त्वया}
{अत्र किं शोभनं यत्त्वं ससैन्यो विनशिष्यसि} %||3-39-16||

\twolineshloka
{मां निहत्य तु रामोऽसौ नचिरात्त्वां वधिष्यति}
{अनेन कृतकृत्योऽस्मि म्रिये यदरिणा हतः} %||3-39-17||

\twolineshloka
{दर्शनादेव रामस्य हतं मामुपधारय}
{आत्मानं च हतं विद्धि हृत्वा सीतां सबान्धवम्} %||3-39-18||

\twolineshloka
{आनयिष्यसि चेत्सीतामाश्रमात्सहितो मया}
{नैव त्वमसि नैवाहं नैव लङ्का न राक्षसाः} %||3-39-19||

\fourlineindentedshloka
{निवार्यमाणस्तु मया हितैषिणा}
{ न मृष्यसे वाक्यमिदं निशाचर}
{परेतकल्पा हि गतायुषो नरा}
{ हितं न गृह्णन्ति सुहृद्भिरीरितम्} %||3-40-20||


{॥इत्यार्षे श्रीमद्रामायणे वाल्मीकीये आदिकाव्ये अरण्यकाण्डे एकोनचत्वारिंशः सर्गः॥}

\sect{चत्वारिंशः सर्गः}

\twolineshloka
{एवमुक्त्वा तु परुषं मारीचो रावणं ततः}
{गच्छावेत्यब्रवीद्दीनो भयाद्रात्रिञ्चरप्रभोः} %||3-40-1||

\twolineshloka
{दृष्टश्चाहं पुनस्तेन शरचापासिधारिणा}
{मद्वधोद्यतशस्त्रेण विनष्टं जीवितं च मे} %||3-40-2||

\twolineshloka
{किं तु कर्तुं मया शक्यमेवं त्वयि दुरात्मनि}
{एष गच्छाम्यहं तात स्वस्ति तेऽस्तु निशाचर} %||3-40-3||

\twolineshloka
{प्रहृष्टस्त्वभवत्तेन वचनेन स राक्षसः}
{परिष्वज्य सुसंश्लिष्टमिदं वचनमब्रवीत्} %||3-40-4||

\twolineshloka
{एतच्छौण्डीर्ययुक्तं ते मच्छन्दादिव भाषितम्}
{इदानीमसि मारीचः पूर्वमन्यो निशाचरः} %||3-40-5||

\twolineshloka
{आरुह्यतामयं शीघ्रं खगो रत्नविभूषितः}
{मया सह रथो युक्तः पिशाचवदनैः खरैः} %||3-40-6||

\twolineshloka
{ततो रावणमारीचौ विमानमिव तं रथम्}
{आरुह्य ययतुः शीघ्रं तस्मादाश्रममण्डलात्} %||3-40-7||

\twolineshloka
{तथैव तत्र पश्यन्तौ पत्तनानि वनानि च}
{गिरींश्च सरितः सर्वा राष्ट्राणि नगराणि च} %||3-40-8||

\twolineshloka
{समेत्य दण्डकारण्यं राघवस्याश्रमं ततः}
{ददर्श सहमरीचो रावणो राक्षसाधिपः} %||3-40-9||

\twolineshloka
{अवतीर्य रथात्तस्मात्ततः काञ्चनभूषणात्}
{हस्ते गृहीत्वा मारीचं रावणो वाक्यमब्रवीत्} %||3-40-10||

\twolineshloka
{एतद्रामाश्रमपदं दृश्यते कदलीवृतम्}
{क्रियतां तत्सखे शीघ्रं यदर्थं वयमागताः} %||3-40-11||

\twolineshloka
{स रावणवचः श्रुत्वा मारीचो राक्षसस्तदा}
{मृगो भूत्वाश्रमद्वारि रामस्य विचचार ह} %||3-40-12||

\twolineshloka
{मणिप्रवरशृङ्गाग्रः सितासितमुखाकृतिः}
{रक्तपद्मोत्पलमुख इन्द्रनीलोत्पलश्रवाः} %||3-40-13||

\twolineshloka
{किञ्चिदभ्युन्नत ग्रीव इन्द्रनीलनिभोदरः}
{मधूकनिभपार्श्वश्च कञ्जकिञ्जल्कसंनिभः} %||3-40-14||

\twolineshloka
{वैदूर्यसङ्काशखुरस्तनुजङ्घः सुसंहतः}
{इन्द्रायुधसवर्णेन पुच्छेनोर्ध्वं विराजितः} %||3-40-15||

\twolineshloka
{मनोहरस्निग्धवर्णो रत्नैर्नानाविधैर्वृतः}
{क्षणेन राक्षसो जातो मृगः परमशोभनः} %||3-40-16||

\twolineshloka
{वनं प्रज्वलयन्रम्यं रामाश्रमपदं च तत्}
{मनोहरं दर्शनीयं रूपं कृत्वा स राक्षसः} %||3-40-17||

\twolineshloka
{प्रलोभनार्थं वैदेह्या नानाधातुविचित्रितम्}
{विचरन्गच्छते सम्यक्शाद्वलानि समन्ततः} %||3-40-18||

\twolineshloka
{रूप्यबिन्दुशतैश्चित्रो भूत्वा च प्रियदर्शनः}
{विटपीनां किसलयान्भङ्क्त्वादन्विचचार ह} %||3-40-19||

\twolineshloka
{कदलीगृहकं गत्वा कर्णिकारानितस्ततः}
{समाश्रयन्मन्दगतिः सीतासन्दर्शनं तदा} %||3-40-20||

\twolineshloka
{राजीवचित्रपृष्ठः स विरराज महामृगः}
{रामाश्रमपदाभ्याशे विचचार यथासुखम्} %||3-40-21||

\twolineshloka
{पुनर्गत्वा निवृत्तश्च विचचार मृगोत्तमः}
{गत्वा मुहूर्तं त्वरया पुनः प्रतिनिवर्तते} %||3-40-22||

\twolineshloka
{विक्रीडंश्च पुनर्भूमौ पुनरेव निषीदति}
{आश्रमद्वारमागम्य मृगयूथानि गच्छति} %||3-40-23||

\twolineshloka
{मृगयूथैरनुगतः पुनरेव निवर्तते}
{सीतादर्शनमाकाङ्क्षन्राक्षसो मृगतां गतः} %||3-40-24||

\twolineshloka
{परिभ्रमति चित्राणि मण्डलानि विनिष्पतन्}
{समुद्वीक्ष्य च सर्वे तं मृगा येऽन्ये वनेचराः} %||3-40-25||

\twolineshloka
{उपगम्य समाघ्राय विद्रवन्ति दिशो दश}
{राक्षसः सोऽपि तान्वन्यान्मृगान्मृगवधे रतः} %||3-40-26||

\twolineshloka
{प्रच्छादनार्थं भावस्य न भक्षयति संस्पृशन्}
{तस्मिन्नेव ततः काले वैदेही शुभलोचना} %||3-40-27||

\twolineshloka
{कुसुमापचये व्यग्रा पादपानत्यवर्तत}
{कर्णिकारानशोकांश्च चूटांश्च मदिरेक्षणा} %||3-40-28||

\threelineshloka
{कुसुमान्यपचिन्वन्ती चचार रुचिरानना}
{अनर्हारण्यवासस्य सा तं रत्नमयं मृगम्}
{मुक्तामणिविचित्राङ्गं ददर्श परमाङ्गना} %||3-40-29||

\twolineshloka
{तं वै रुचिरदण्तौष्ठं रूप्यधातुतनूरुहम्}
{विस्मयोत्फुल्लनयना सस्नेहं समुदैक्षत} %||3-40-30||

\twolineshloka
{स च तां रामदयितां पश्यन्मायामयो मृगः}
{विचचार ततस्तत्र दीपयन्निव तद्वनम्} %||3-40-31||

\twolineshloka
{अदृष्टपूर्वं दृष्ट्वा तं नानारत्नमयं मृगम्}
{विस्मयं परमं सीता जगाम जनकात्मजा} %||3-41-32||


{॥इत्यार्षे श्रीमद्रामायणे वाल्मीकीये आदिकाव्ये अरण्यकाण्डे चत्वारिंशः सर्गः॥}

\sect{एकचत्वारिंशः सर्गः}

\twolineshloka
{सा तं सम्प्रेक्ष्य सुश्रोणी कुसुमानि विचिन्वती}
{हेमराजतवर्णाभ्यां पार्श्वाभ्यामुपशोभितम्} %||3-41-1||

\twolineshloka
{प्रहृष्टा चानवद्याङ्गी मृष्टहाटकवर्णिनी}
{भर्तारमपि चाक्रन्दल्लक्ष्मणं चैव सायुधम्} %||3-41-2||

\twolineshloka
{तयाहूतौ नरव्याघ्रौ वैदेह्या रामलक्ष्मणौ}
{वीक्षमाणौ तु तं देशं तदा ददृशतुर्मृगम्} %||3-41-3||

\twolineshloka
{शङ्कमानस्तु तं दृष्ट्वा लक्ष्मणो राममब्रवीत्}
{तमेवैनमहं मन्ये मारीचं राक्षसं मृगम्} %||3-41-4||

\twolineshloka
{चरन्तो मृगयां हृष्टाः पापेनोपाधिना वने}
{अनेन निहता राम राजानः कामरूपिणा} %||3-41-5||

\twolineshloka
{अस्य मायाविदो मायामृगरूपमिदं कृतम्}
{भानुमत्पुरुषव्याघ्र गन्धर्वपुरसंनिभम्} %||3-41-6||

\twolineshloka
{मृगो ह्येवंविधो रत्नविचित्रो नास्ति राघव}
{जगत्यां जगतीनाथ मायैषा हि न संशयः} %||3-41-7||

\twolineshloka
{एवं ब्रुवाणं काकुत्स्थं प्रतिवार्य शुचिस्मिता}
{उवाच सीता संहृष्टा छद्मना हृतचेतना} %||3-41-8||

\twolineshloka
{आर्यपुत्राभिरामोऽसौ मृगो हरति मे मनः}
{आनयैनं महाबाहो क्रीडार्थं नो भविष्यति} %||3-41-9||

\twolineshloka
{इहाश्रमपदेऽस्माकं बहवः पुण्यदर्शनाः}
{मृगाश्चरन्ति सहिताश्चमराः सृमरास्तथा} %||3-41-10||

\twolineshloka
{ऋक्षाः पृषतसङ्घाश्च वानराः किंनरास्तथा}
{विचरन्ति महाबाहो रूपश्रेष्ठा महाबलाः} %||3-41-11||

\twolineshloka
{न चास्य सदृशो राजन्दृष्टपूर्वो मृगः पुरा}
{तेजसा क्षमया दीप्त्या यथायं मृगसत्तमः} %||3-41-12||

\twolineshloka
{नानावर्णविचित्राङ्गो रत्नबिन्दुसमाचितः}
{द्योतयन्वनमव्यग्रं शोभते शशिसंनिभः} %||3-41-13||

\twolineshloka
{अहो रूपमहो लक्ष्मीः स्वरसम्पच्च शोभना}
{मृगोऽद्भुतो विचित्रोऽसौ हृदयं हरतीव मे} %||3-41-14||

\twolineshloka
{यदि ग्रहणमभ्येति जीवन्नेव मृगस्तव}
{आश्चर्यभूतं भवति विस्मयं जनयिष्यति} %||3-41-15||

\twolineshloka
{समाप्तवनवासानां राज्यस्थानां च नः पुनः}
{अन्तःपुरविभूषार्थो मृग एष भविष्यति} %||3-41-16||

\twolineshloka
{भरतस्यार्यपुत्रस्य श्वश्रूणां मम च प्रभो}
{मृगरूपमिदं दिव्यं विस्मयं जनयिष्यति} %||3-41-17||

\twolineshloka
{जीवन्न यदि तेऽभ्येति ग्रहणं मृगसत्तमः}
{अजिनं नरशार्दूल रुचिरं मे भविष्यति} %||3-41-18||

\twolineshloka
{निहतस्यास्य सत्त्वस्य जाम्बूनदमयत्वचि}
{शष्पबृस्यां विनीतायामिच्छाम्यहमुपासितुम्} %||3-41-19||

\twolineshloka
{कामवृत्तमिदं रौद्रं स्त्रीणामसदृशं मतम्}
{वपुषा त्वस्य सत्त्वस्य विस्मयो जनितो मम} %||3-41-20||

\threelineshloka
{तेन काञ्चनरोम्णा तु मणिप्रवरशृङ्गिणा}
{तरुणादित्यवर्णेन नक्षत्रपथवर्चसा}
{बभूव राघवस्यापि मनो विस्मयमागतम्} %||3-41-21||

\twolineshloka
{एवं सीतावचः श्रुत्वा दृष्ट्वा च मृगमद्भुतम्}
{उवाच राघवो हृष्टो भ्रातरं लक्ष्मणं वचः} %||3-41-22||

\twolineshloka
{पश्य लक्ष्मण वैदेह्याः स्पृहां मृगगतामिमाम्}
{रूपश्रेष्ठतया ह्येष मृगोऽद्य न भविष्यति} %||3-41-23||

\twolineshloka
{न वने नन्दनोद्देशे न चैत्ररथसंश्रये}
{कुतः पृथिव्यां सौमित्रे योऽस्य कश्चित्समो मृगः} %||3-41-24||

\twolineshloka
{प्रतिलोमानुलोमाश्च रुचिरा रोमराजयः}
{शोभन्ते मृगमाश्रित्य चित्राः कनकबिन्दुभिः} %||3-41-25||

\twolineshloka
{पश्यास्य जृम्भमाणस्य दीप्तामग्निशिखोपमाम्}
{जिह्वां मुखान्निःसरन्तीं मेघादिव शतह्रदाम्} %||3-41-26||

\twolineshloka
{मसारगल्वर्कमुखः शङ्खमुक्तानिभोदरः}
{कस्य नामानिरूप्योऽसौ न मनो लोभयेन्मृगः} %||3-41-27||

\twolineshloka
{कस्य रूपमिदं दृष्ट्वा जाम्बूनदमयप्रभम्}
{नानारत्नमयं दिव्यं न मनो विस्मयं व्रजेत्} %||3-41-28||

\twolineshloka
{मांसहेतोरपि मृगान्विहारार्थं च धन्विनः}
{घ्नन्ति लक्ष्मण राजानो मृगयायां महावने} %||3-41-29||

\twolineshloka
{धनानि व्यवसायेन विचीयन्ते महावने}
{धातवो विविधाश्चापि मणिरत्नसुवर्णिनः} %||3-41-30||

\twolineshloka
{तत्सारमखिलं नॄणां धनं निचयवर्धनम्}
{मनसा चिन्तितं सर्वं यथा शुक्रस्य लक्ष्मण} %||3-41-31||

\twolineshloka
{अर्थी येनार्थकृत्येन संव्रजत्यविचारयन्}
{तमर्थमर्थशास्त्रज्ञः प्राहुरर्थ्याश्च लक्ष्मण} %||3-41-32||

\twolineshloka
{एतस्य मृगरत्नस्य परार्ध्ये काञ्चनत्वचि}
{उपवेक्ष्यति वैदेही मया सह सुमध्यमा} %||3-41-33||

\twolineshloka
{न कादली न प्रियकी न प्रवेणी न चाविकी}
{भवेदेतस्य सदृशी स्पर्शनेनेति मे मतिः} %||3-41-34||

\twolineshloka
{एष चैव मृगः श्रीमान्यश्च दिव्यो नभश्चरः}
{उभावेतौ मृगौ दिव्यौ तारामृगमहीमृगौ} %||3-41-35||

\twolineshloka
{यदि वायं तथा यन्मां भवेद्वदसि लक्ष्मण}
{मायैषा राक्षसस्येति कर्तव्योऽस्य वधो मया} %||3-41-36||

\twolineshloka
{एतेन हि नृशंसेन मारीचेनाकृतात्मना}
{वने विचरता पूर्वं हिंसिता मुनिपुङ्गवाः} %||3-41-37||

\twolineshloka
{उत्थाय बहवो येन मृगयायां जनाधिपाः}
{निहताः परमेष्वासास्तस्माद्वध्यस्त्वयं मृगः} %||3-41-38||

\twolineshloka
{पुरस्तादिह वातापिः परिभूय तपस्विनः}
{उदरस्थो द्विजान्हन्ति स्वगर्भोऽश्वतरीमिव} %||3-41-39||

\twolineshloka
{स कदाचिच्चिराल्लोके आससाद महामुनिम्}
{अगस्त्यं तेजसा युक्तं भक्ष्यस्तस्य बभूव ह} %||3-41-40||

\twolineshloka
{समुत्थाने च तद्रूपं कर्तुकामं समीक्ष्य तम्}
{उत्स्मयित्वा तु भगवान्वातापिमिदमब्रवीत्} %||3-41-41||

\twolineshloka
{त्वयाविगण्य वातापे परिभूताश्च तेजसा}
{जीवलोके द्विजश्रेष्ठास्तस्मादसि जरां गतः} %||3-41-42||

\twolineshloka
{एवं तन्न भवेद्रक्षो वातापिरिव लक्ष्मण}
{मद्विधं योऽतिमन्येत धर्मनित्यं जितेन्द्रियम्} %||3-41-43||

\twolineshloka
{भवेद्धतोऽयं वातापिरगस्त्येनेव मा गतिः}
{इह त्वं भव संनद्धो यन्त्रितो रक्ष मैथिलीम्} %||3-41-44||

\twolineshloka
{अस्यामायत्तमस्माकं यत्कृत्यं रघुनन्दन}
{अहमेनं वधिष्यामि ग्रहीष्याम्यथ वा मृगम्} %||3-41-45||

\twolineshloka
{यावद्गच्छामि सौमित्रे मृगमानयितुं द्रुतम्}
{पश्य लक्ष्मण वैदेहीं मृगत्वचि गतस्पृहाम्} %||3-41-46||

\twolineshloka
{त्वचा प्रधानया ह्येष मृगोऽद्य न भविष्यति}
{अप्रमत्तेन ते भाव्यमाश्रमस्थेन सीतया} %||3-41-47||

\twolineshloka
{यावत्पृषतमेकेन सायकेन निहन्म्यहम्}
{हत्वैतच्चर्म आदाय शीघ्रमेष्यामि लक्ष्मण} %||3-41-48||

\fourlineindentedshloka
{प्रदक्षिणेनातिबलेन पक्षिणा}
{ जटायुषा बुद्धिमता च लक्ष्मण}
{भवाप्रमत्तः प्रतिगृह्य मैथिलीम्}
{ प्रतिक्षणं सर्वत एव शङ्कितः} %||3-42-49||


{॥इत्यार्षे श्रीमद्रामायणे वाल्मीकीये आदिकाव्ये अरण्यकाण्डे एकचत्वारिंशः सर्गः॥}

\sect{द्विचत्वारिंशः सर्गः}

\twolineshloka
{तथा तु तं समादिश्य भ्रातरं रघुनन्दनः}
{बबन्धासिं महातेजा जाम्बूनदमयत्सरुम्} %||3-42-1||

\twolineshloka
{ततस्त्रिविणतं चापमादायात्मविभूषणम्}
{आबध्य च कलापौ द्वौ जगामोदग्रविक्रमः} %||3-42-2||

\twolineshloka
{तं वञ्चयानो राजेन्द्रमापतन्तं निरीक्ष्य वै}
{बभूवान्तर्हितस्त्रासात्पुनः सन्दर्शनेऽभवत्} %||3-42-3||

\twolineshloka
{बद्धासिर्धनुरादाय प्रदुद्राव यतो मृगः}
{तं स पश्यति रूपेण द्योतमानमिवाग्रतः} %||3-42-4||

\twolineshloka
{अवेक्ष्यावेक्ष्य धावन्तं धनुष्पाणिर्महावने}
{अतिवृत्तमिषोः पाताल्लोभयानं कदाचन} %||3-42-5||

\twolineshloka
{शङ्कितं तु समुद्भ्रान्तमुत्पतन्तमिवाम्बरे}
{दश्यमानमदृश्यं च नवोद्देशेषु केषुचित्} %||3-42-6||

\twolineshloka
{छिन्नाभ्रैरिव संवीतं शारदं चन्द्रमण्डलम्}
{मुहूर्तादेव ददृशे मुहुर्दूरात्प्रकाशते} %||3-42-7||

\twolineshloka
{दर्शनादर्शनेनैव सोऽपाकर्षत राघवम्}
{आसीत्क्रुद्धस्तु काकुत्स्थो विवशस्तेन मोहितः} %||3-42-8||

\twolineshloka
{अथावतस्थे सुश्रान्तश्छायामाश्रित्य शाद्वले}
{मृगैः परिवृतो वन्यैरदूरात्प्रत्यदृश्यत} %||3-42-9||

\twolineshloka
{दृष्ट्वा रामो महातेजास्तं हन्तुं कृतनिश्चयः}
{सन्धाय सुदृढे चापे विकृष्य बलवद्बली} %||3-42-10||

\twolineshloka
{तमेव मृगमुद्दिश्य ज्वलन्तमिव पन्नगम्}
{मुमोच ज्वलितं दीप्तमस्त्रब्रह्मविनिर्मितम्} %||3-42-11||

\twolineshloka
{स भृशं मृगरूपस्य विनिर्भिद्य शरोत्तमः}
{मारीचस्यैव हृदयं विभेदाशनिसंनिभः} %||3-42-12||

\threelineshloka
{तालमात्रमथोत्पत्य न्यपतत्स शरातुरः}
{व्यनदद्भैरवं नादं धरण्यामल्पजीवितः}
{म्रियमाणस्तु मारीचो जहौ तां कृत्रिमां तनुम्} %||3-42-13||

\twolineshloka
{सम्प्राप्तकालमाज्ञाय चकार च ततः स्वरम्}
{सदृशं राघवस्यैव हा सीते लक्ष्मणेति च} %||3-42-14||

\threelineshloka
{तेन मर्मणि निर्विद्धः शरेणानुपमेन हि}
{मृगरूपं तु तत्त्यक्त्वा राक्षसं रूपमात्मनः}
{चक्रे स सुमहाकायो मारीचो जीवितं त्यजन्} %||3-42-15||

\twolineshloka
{ततो विचित्रकेयूरः सर्वाभरणभूषितः}
{हेममाली महादंष्ट्रो राक्षसोऽभूच्छराहतः} %||3-42-16||

\twolineshloka
{तं दृष्ट्वा पतितं भूमौ राक्षसं घोरदर्शनम्}
{जगाम मनसा सीतां लक्ष्मणस्य वचः स्मरन्} %||3-42-17||

\twolineshloka
{हा सीते लक्ष्मणेत्येवमाक्रुश्य तु महास्वरम्}
{ममार राक्षसः सोऽयं श्रुत्वा सीता कथं भवेत्} %||3-42-18||

\twolineshloka
{लक्ष्मणश्च महाबाहुः कामवस्थां गमिष्यति}
{इति सञ्चिन्त्य धर्मात्मा रामो हृष्टतनूरुहः} %||3-42-19||

\twolineshloka
{तत्र रामं भयं तीव्रमाविवेश विषादजम्}
{राक्षसं मृगरूपं तं हत्वा श्रुत्वा च तत्स्वरम्} %||3-42-20||

\twolineshloka
{निहत्य पृषतं चान्यं मांसमादाय राघवः}
{त्वरमाणो जनस्थानं ससाराभिमुखस्तदा} %||3-43-21||


{॥इत्यार्षे श्रीमद्रामायणे वाल्मीकीये आदिकाव्ये अरण्यकाण्डे द्विचत्वारिंशः सर्गः॥}

\sect{त्रिचत्वारिंशः सर्गः}

\twolineshloka
{आर्तस्वरं तु तं भर्तुर्विज्ञाय सदृशं वने}
{उवाच लक्ष्मणं सीता गच्छ जानीहि राघवम्} %||3-43-1||

\twolineshloka
{न हि मे जीवितं स्थाने हृदयं वावतिष्ठते}
{क्रोशतः परमार्तस्य श्रुतः शब्दो मया भृशम्} %||3-43-2||

\twolineshloka
{आक्रन्दमानं तु वने भ्रातरं त्रातुमर्हसि}
{तं क्षिप्रमभिधाव त्वं भ्रातरं शरणैषिणम्} %||3-43-3||

\twolineshloka
{रक्षसां वशमापन्नं सिंहानामिव गोवृषम्}
{न जगाम तथोक्तस्तु भ्रातुराज्ञाय शासनम्} %||3-43-4||

\twolineshloka
{तमुवाच ततस्तत्र कुपिता जनकात्मजा}
{सौमित्रे मित्ररूपेण भ्रातुस्त्वमसि शत्रुवत्} %||3-43-5||

\twolineshloka
{यस्त्वमस्यामवस्थायां भ्रातरं नाभिपद्यसे}
{इच्छसि त्वं विनश्यन्तं रामं लक्ष्मण मत्कृते} %||3-43-6||

\twolineshloka
{व्यसनं ते प्रियं मन्ये स्नेहो भ्रातरि नास्ति ते}
{तेन तिष्ठसि विस्रब्धस्तमपश्यन्महाद्युतिम्} %||3-43-7||

\twolineshloka
{किं हि संशयमापन्ने तस्मिन्निह मया भवेत्}
{कर्तव्यमिह तिष्ठन्त्या यत्प्रधानस्त्वमागतः} %||3-43-8||

\twolineshloka
{इति ब्रुवाणं वैदेहीं बाष्पशोकपरिप्लुताम्}
{अब्रवील्लक्ष्मणस्त्रस्तां सीतां मृगवधूमिव} %||3-43-9||

\twolineshloka
{देवि देवमनुष्येषु गन्धर्वेषु पतत्रिषु}
{राक्षसेषु पिशाचेषु किंनरेषु मृगेषु च} %||3-43-10||

\twolineshloka
{दानवेषु च घोरेषु न स विद्येत शोभने}
{यो रामं प्रतियुध्येत समरे वासवोपमम्} %||3-43-11||

\twolineshloka
{अवध्यः समरे रामो नैवं त्वं वक्तुमर्हसि}
{न त्वामस्मिन्वने हातुमुत्सहे राघवं विना} %||3-43-12||

\twolineshloka
{अनिवार्यं बलं तस्य बलैर्बलवतामपि}
{त्रिभिर्लोकैः समुद्युक्तैः सेश्वरैः सामरैरपि} %||3-43-13||

\twolineshloka
{हृदयं निर्वृतं तेऽस्तु सन्तापस्त्यज्यतामयम्}
{आगमिष्यति ते भर्ता शीघ्रं हत्वा मृगोत्तमम्} %||3-43-14||

\twolineshloka
{न स तस्य स्वरो व्यक्तं न कश्चिदपि दैवतः}
{गन्धर्वनगरप्रख्या माया सा तस्य रक्षसः} %||3-43-15||

\twolineshloka
{न्यासभूतासि वैदेहि न्यस्ता मयि महात्मना}
{रामेण त्वं वरारोहे न त्वां त्यक्तुमिहोत्सहे} %||3-43-16||

\twolineshloka
{कृतवैराश्च कल्याणि वयमेतैर्निशाचरैः}
{खरस्य निधने देवि जनस्थानवधं प्रति} %||3-43-17||

\twolineshloka
{राक्षसा विधिना वाचो विसृजन्ति महावने}
{हिंसाविहारा वैदेहि न चिन्तयितुमर्हसि} %||3-43-18||

\twolineshloka
{लक्ष्मणेनैवमुक्ता तु क्रुद्धा संरक्तलोचना}
{अब्रवीत्परुषं वाक्यं लक्ष्मणं सत्यवादिनम्} %||3-43-19||

\twolineshloka
{अनार्य करुणारम्भ नृशंस कुलपांसन}
{अहं तव प्रियं मन्ये तेनैतानि प्रभाषसे} %||3-43-20||

\twolineshloka
{नैतच्चित्रं सपत्नेषु पापं लक्ष्मण यद्भवेत्}
{त्वद्विधेषु नृशंसेषु नित्यं प्रच्छन्नचारिषु} %||3-43-21||

\twolineshloka
{सुदुष्टस्त्वं वने राममेकमेकोऽनुगच्छसि}
{मम हेतोः प्रतिच्छन्नः प्रयुक्तो भरतेन वा} %||3-43-22||

\twolineshloka
{कथमिन्दीवरश्यामं रामं पद्मनिभेक्षणम्}
{उपसंश्रित्य भर्तारं कामयेयं पृथग्जनम्} %||3-43-23||

\twolineshloka
{समक्षं तव सौमित्रे प्राणांस्त्यक्ष्ये न संशयः}
{रामं विना क्षणमपि न हि जीवामि भूतले} %||3-43-24||

\twolineshloka
{इत्युक्तः परुषं वाक्यं सीतया सोमहर्षणम्}
{अब्रवील्लक्ष्मणः सीतां प्राञ्जलिर्विजितेन्द्रियः} %||3-43-25||

\twolineshloka
{उत्तरं नोत्सहे वक्तुं दैवतं भवती मम}
{वाक्यमप्रतिरूपं तु न चित्रं स्त्रीषु मैथिलि} %||3-43-26||

\twolineshloka
{स्वभावस्त्वेष नारीणामेषु लोकेषु दृश्यते}
{विमुक्तधर्माश्चपलास्तीक्ष्णा भेदकराः स्त्रियः} %||3-43-27||

\twolineshloka
{उपशृण्वन्तु मे सर्वे साक्षिभूता वनेचराः}
{न्यायवादी यथा वाक्यमुक्तोऽहं परुषं त्वया} %||3-43-28||

\twolineshloka
{धिक्त्वामद्य प्रणश्य त्वं यन्मामेवं विशङ्कसे}
{स्त्रीत्वाद्दुष्टस्वभावेन गुरुवाक्ये व्यवस्थितम्} %||3-43-29||

\twolineshloka
{गमिष्ये यत्र काकुत्स्थः स्वस्ति तेऽस्तु वरानने}
{रक्षन्तु त्वां विशालाक्षि समग्रा वनदेवताः} %||3-43-30||

\twolineshloka
{निमित्तानि हि घोराणि यानि प्रादुर्भवन्ति मे}
{अपि त्वां सह रामेण पश्येयं पुनरागतः} %||3-43-31||

\twolineshloka
{लक्ष्मणेनैवमुक्ता तु रुदती जनकात्मजा}
{प्रत्युवाच ततो वाक्यं तीव्रं बाष्पपरिप्लुता} %||3-43-32||

\twolineshloka
{गोदावरीं प्रवेक्ष्यामि विना रामेण लक्ष्मण}
{आबन्धिष्येऽथवा त्यक्ष्ये विषमे देहमात्मनः} %||3-43-33||

\twolineshloka
{पिबामि वा विषं तीक्ष्णं प्रवेक्ष्यामि हुताशनम्}
{न त्वहं राघवादन्यं पदापि पुरुषं स्पृशे} %||3-43-34||

\twolineshloka
{इति लक्ष्मणमाक्रुश्य सीता दुःखसमन्विता}
{पाणिभ्यां रुदती दुःखादुदरं प्रजघान ह} %||3-43-35||

\fourlineindentedshloka
{तामार्तरूपां विमना रुदन्तीम्}
{ सौमित्रिरालोक्य विशालनेत्राम्}
{आश्वासयामास न चैव भर्तुः}
{ तं भ्रातरं किञ्चिदुवाच सीता} %||3-43-36||

\fourlineindentedshloka
{ततस्तु सीतामभिवाद्य लक्ष्मणः}
{ कृताञ्जलिः किञ्चिदभिप्रणम्य}
{अवेक्षमाणो बहुशश्च मैथिलीम्}
{ जगाम रामस्य समीपमात्मवान्} %||3-44-37||


{॥इत्यार्षे श्रीमद्रामायणे वाल्मीकीये आदिकाव्ये अरण्यकाण्डे त्रिचत्वारिंशः सर्गः॥}

\sect{चतुश्चत्वारिंशः सर्गः}

\twolineshloka
{तया परुषमुक्तस्तु कुपितो राघवानुजः}
{स विकाङ्क्षन्भृशं रामं प्रतस्थे नचिरादिव} %||3-44-1||

\twolineshloka
{तदासाद्य दशग्रीवः क्षिप्रमन्तरमास्थितः}
{अभिचक्राम वैदेहीं परिव्राजकरूपधृक्} %||3-44-2||

\threelineshloka
{श्लक्ष्णकाषायसंवीतः शिखी छत्री उपानही}
{वामे चांसेऽवसज्याथ शुभे यष्टिकमण्डलू}
{परिव्राजकरूपेण वैदेहीं समुपागमत्} %||3-44-3||

\twolineshloka
{तामाससादातिबलो भ्रातृभ्यां रहितां वने}
{रहितां सूर्यचन्द्राभ्यां सन्ध्यामिव महत्तमः} %||3-44-4||

\twolineshloka
{तामपश्यत्ततो बालां राजपुत्रीं यशस्विनीम्}
{रोहिणीं शशिना हीनां ग्रहवद्भृशदारुणः} %||3-44-5||

\twolineshloka
{तमुग्रं पापकर्माणं जनस्थानरुहा द्रुमाः}
{समीक्ष्य न प्रकम्पन्ते न प्रवाति च मारुतः} %||3-44-6||

\twolineshloka
{शीघ्रस्रोताश्च तं दृष्ट्वा वीक्षन्तं रक्तलोचनम्}
{स्तिमितं गन्तुमारेभे भयाद्गोदावरी नदी} %||3-44-7||

\twolineshloka
{रामस्य त्वन्तरं प्रेप्सुर्दशग्रीवस्तदन्तरे}
{उपतस्थे च वैदेहीं भिक्षुरूपेण रावणः} %||3-44-8||

\twolineshloka
{अभव्यो भव्यरूपेण भर्तारमनुशोचतीम्}
{अभ्यवर्तत वैदेहीं चित्रामिव शनैश्चरः} %||3-44-9||

\twolineshloka
{स पापो भव्यरूपेण तृणैः कूप इवावृतः}
{अतिष्ठत्प्रेक्ष्य वैदेहीं रामपत्नीं यशस्विनीम्} %||3-44-10||

\twolineshloka
{शुभां रुचिरदन्तौष्ठीं पूर्णचन्द्रनिभाननाम्}
{आसीनां पर्णशालायां बाष्पशोकाभिपीडिताम्} %||3-44-11||

\twolineshloka
{स तां पद्मपलाशाक्षीं पीतकौशेयवासिनीम्}
{अभ्यगच्छत वैदेहीं दुष्टचेता निशाचरः} %||3-44-12||

\twolineshloka
{स मन्मथशराविष्टो ब्रह्मघोषमुदीरयन्}
{अब्रवीत्प्रश्रितं वाक्यं रहिते राक्षसाधिपः} %||3-44-13||

\twolineshloka
{तामुत्तमां त्रिलोकानां पद्महीनामिव श्रियम्}
{विभ्राजमानां वपुषा रावणः प्रशशंस ह} %||3-44-14||

\twolineshloka
{का त्वं काञ्चनवर्णाभे पीतकौशेयवासिनि}
{कमलानां शुभां मालां पद्मिनीव च बिभ्रती} %||3-44-15||

\twolineshloka
{ह्रीः श्रीः कीर्तिः शुभा लक्ष्मीरप्सरा वा शुभानने}
{भूतिर्वा त्वं वरारोहे रतिर्वा स्वैरचारिणी} %||3-44-16||

\twolineshloka
{समाः शिखरिणः स्निग्धाः पाण्डुरा दशनास्तव}
{विशाले विमले नेत्रे रक्तान्ते कृष्णतारके} %||3-44-17||

\twolineshloka
{विशालं जघनं पीनमूरू करिकरोपमौ}
{एतावुपचितौ वृत्तौ सहितौ सम्प्रगल्भितौ} %||3-44-18||

\twolineshloka
{पीनोन्नतमुखौ कान्तौ स्निग्धतालफलोपमौ}
{मणिप्रवेकाभरणौ रुचिरौ ते पयोधरौ} %||3-44-19||

\twolineshloka
{चारुस्मिते चारुदति चारुनेत्रे विलासिनि}
{मनो हरसि मे रामे नदीकूलमिवाम्भसा} %||3-44-20||

\twolineshloka
{करान्तमितमध्यासि सुकेशी संहतस्तनी}
{नैव देवी न गन्धर्वी न यक्षी न च किंनरी} %||3-44-21||

\twolineshloka
{नैवंरूपा मया नारी दृष्टपूर्वा महीतले}
{इह वासश्च कान्तारे चित्तमुन्माथयन्ति मे} %||3-44-22||

\twolineshloka
{सा प्रतिक्राम भद्रं ते न त्वं वस्तुमिहार्हसि}
{राक्षसानामयं वासो घोराणां कामरूपिणाम्} %||3-44-23||

\twolineshloka
{प्रासादाग्र्याणि रम्याणि नगरोपवनानि च}
{सम्पन्नानि सुगन्धीनि युक्तान्याचरितुं त्वया} %||3-44-24||

\twolineshloka
{वरं माल्यं वरं पानं वरं वस्त्रं च शोभने}
{भर्तारं च वरं मन्ये त्वद्युक्तमसितेक्षणे} %||3-44-25||

\twolineshloka
{का त्वं भवसि रुद्राणां मरुतां वा शुचिस्मिते}
{वसूनां वा वरारोहे देवता प्रतिभासि मे} %||3-44-26||

\twolineshloka
{नेह गच्छन्ती गन्धर्वा न देवा न च किंनराः}
{राक्षसानामयं वासः कथं नु त्वमिहागता} %||3-44-27||

\twolineshloka
{इह शाखामृगाः सिंहा द्वीपिव्याघ्रमृगास्तथा}
{ऋक्षास्तरक्षवः कङ्काः कथं तेभ्यो न बिभ्यसे} %||3-44-28||

\twolineshloka
{मदान्वितानां घोराणां कुञ्जराणां तरस्विनाम्}
{कथमेका महारण्ये न बिभेषि वनानने} %||3-44-29||

\twolineshloka
{कासि कस्य कुतश्च त्वं किंनिमित्तं च दण्डकान्}
{एका चरसि कल्याणि घोरान्राक्षससेवितान्} %||3-44-30||

\threelineshloka
{इति प्रशस्ता वैदेही रावणेन दुरात्मना}
{द्विजातिवेषेण हि तं दृष्ट्वा रावणमागतम्}
{सर्वैरतिथिसत्कारैः पूजयामास मैथिली} %||3-44-31||

\twolineshloka
{उपानीयासनं पूर्वं पाद्येनाभिनिमन्त्र्य च}
{अब्रवीत्सिद्धमित्येव तदा तं सौम्यदर्शनम्} %||3-44-32||

\fourlineindentedshloka
{द्विजातिवेषेण समीक्ष्य मैथिली}
{ तमागतं पात्रकुसुम्भधारिणम्}
{अशक्यमुद्द्वेष्टुमुपायदर्शना}
{न्न्यमन्त्रयद्ब्राह्मणवद्यथागतम्} %||3-44-33||

\fourlineindentedshloka
{इयं बृसी ब्राह्मण काममास्यताम्}
{ इदं च पाद्यं प्रतिगृह्यतामिति}
{इदं च सिद्धं वनजातमुत्तमम्}
{ त्वदर्थमव्यग्रमिहोपभुज्यताम्} %||3-44-34||

\fourlineindentedshloka
{निमन्त्र्यमाणः प्रतिपूर्णभाषिणीम्}
{ नरेन्द्रपत्नीं प्रसमीक्ष्य मैथिलीम्}
{प्रहस्य तस्या हरणे धृतं मनः}
{ समर्पयामास वधाय रावणः} %||3-44-35||

\fourlineindentedshloka
{ततः सुवेषं मृगया गतं पतिम्}
{ प्रतीक्षमाणा सहलक्ष्मणं तदा}
{निरीक्षमाणा हरितं ददर्श त}
{न्महद्वनं नैव तु रामलक्ष्मणौ} %||3-45-36||


{॥इत्यार्षे श्रीमद्रामायणे वाल्मीकीये आदिकाव्ये अरण्यकाण्डे चतुश्चत्वारिंशः सर्गः॥}

\sect{पञ्चचत्वारिंशः सर्गः}

\twolineshloka
{रावणेन तु वैदेही तदा पृष्टा जिहीर्षुणा}
{परिव्राजकरूपेण शशंसात्मानमात्मना} %||3-45-1||

\twolineshloka
{ब्राह्मणश्चातिथिश्चैष अनुक्तो हि शपेत माम्}
{इति ध्यात्वा मुहूर्तं तु सीता वचनमब्रवीत्} %||3-45-2||

\twolineshloka
{दुहिता जनकस्याहं मैथिलस्य महात्मनः}
{सीता नाम्नास्मि भद्रं ते रामभार्या द्विजोत्तम} %||3-45-3||

\twolineshloka
{संवत्सरं चाध्युषिता राघवस्य निवेशने}
{भुञ्जाना मानुषान्भोगान्सर्वकामसमृद्धिनी} %||3-45-4||

\twolineshloka
{ततः संवत्सरादूर्ध्वं सममन्यत मे पतिम्}
{अभिषेचयितुं रामं समेतो राजमन्त्रिभिः} %||3-45-5||

\twolineshloka
{तस्मिन्सम्भ्रियमाणे तु राघवस्याभिषेचने}
{कैकेयी नाम भर्तारं ममार्या याचते वरम्} %||3-45-6||

\threelineshloka
{प्रतिगृह्य तु कैकेयी श्वशुरं सुकृतेन मे}
{मम प्रव्राजनं भर्तुर्भरतस्याभिषेचनम्}
{द्वावयाचत भर्तारं सत्यसन्धं नृपोत्तमम्} %||3-45-7||

\twolineshloka
{नाद्य भोक्ष्ये न च स्वप्स्ये न पास्येऽहं कदाचन}
{एष मे जीवितस्यान्तो रामो यद्यभिषिच्यते} %||3-45-8||

\twolineshloka
{इति ब्रुवाणां कैकेयीं श्वशुरो मे स मानदः}
{अयाचतार्थैरन्वर्थैर्न च याच्ञां चकार सा} %||3-45-9||

\threelineshloka
{मम भर्ता महातेजा वयसा पञ्चविंशकः}
{रामेति प्रथितो लोके गुणवान्सत्यवाक्शुचिः}
{विशालाक्षो महाबाहुः सर्वभूतहिते रतः} %||3-45-10||

\twolineshloka
{अभिषेकाय तु पितुः समीपं राममागतम्}
{कैकेयी मम भर्तारमित्युवाच द्रुतं वचः} %||3-45-11||

\twolineshloka
{तव पित्रा समाज्ञप्तं ममेदं शृणु राघव}
{भरताय प्रदातव्यमिदं राज्यमकण्टकम्} %||3-45-12||

\twolineshloka
{त्वया तु खलु वस्तव्यं नव वर्षाणि पञ्च च}
{वने प्रव्रज काकुत्स्थ पितरं मोचयानृतात्} %||3-45-13||

\twolineshloka
{तथेत्युवाच तां रामः कैकेयीमकुतोभयः}
{चकार तद्वचस्तस्या मम भर्ता दृढव्रतः} %||3-45-14||

\twolineshloka
{दद्यान्न प्रतिगृह्णीयात्सत्यब्रूयान्न चानृतम्}
{एतद्ब्राह्मण रामस्य व्रतं ध्रुवमनुत्तमम्} %||3-45-15||

\twolineshloka
{तस्य भ्राता तु वैमात्रो लक्ष्मणो नाम वीर्यवान्}
{रामस्य पुरुषव्याघ्रः सहायः समरेऽरिहा} %||3-45-16||

\twolineshloka
{स भ्राता लक्ष्मणो नाम धर्मचारी दृढव्रतः}
{अन्वगच्छद्धनुष्पाणिः प्रव्रजन्तं मया सह} %||3-45-17||

\twolineshloka
{ते वयं प्रच्युता राज्यात्कैलेय्यास्तु कृते त्रयः}
{विचराम द्विजश्रेष्ठ वनं गम्भीरमोजसा} %||3-45-18||

\twolineshloka
{समाश्वस मुहूर्तं तु शक्यं वस्तुमिह त्वया}
{आगमिष्यति मे भर्ता वन्यमादाय पुष्कलम्} %||3-45-19||

\twolineshloka
{स त्वं नाम च गोत्रं च कुलमाचक्ष्व तत्त्वतः}
{एकश्च दण्डकारण्ये किमर्थं चरसि द्विज} %||3-45-20||

\twolineshloka
{एवं ब्रुवत्यां सीतायां रामपत्न्यां महाबलः}
{प्रत्युवाचोत्तरं तीव्रं रावणो राक्षसाधिपः} %||3-45-21||

\twolineshloka
{येन वित्रासिता लोकाः सदेवासुरपन्नगाः}
{अहं स रावणो नाम सीते रक्षोगणेश्वरः} %||3-45-22||

\twolineshloka
{त्वां तु काञ्चनवर्णाभां दृष्ट्वा कौशेयवासिनीम्}
{रतिं स्वकेषु दारेषु नाधिगच्छाम्यनिन्दिते} %||3-45-23||

\twolineshloka
{बह्वीनामुत्तमस्त्रीणामाहृतानामितस्ततः}
{सर्वासामेव भद्रं ते ममाग्रमहिषी भव} %||3-45-24||

\twolineshloka
{लङ्का नाम समुद्रस्य मध्ये मम महापुरी}
{सागरेण परिक्षिप्ता निविष्टा गिरिमूर्धनि} %||3-45-25||

\twolineshloka
{तत्र सीते मया सार्धं वनेषु विचरिष्यसि}
{न चास्यारण्यवासस्य स्पृहयिष्यसि भामिनि} %||3-45-26||

\twolineshloka
{पञ्चदास्यः सहस्राणि सर्वाभरणभूषिताः}
{सीते परिचरिष्यन्ति भार्या भवसि मे यदि} %||3-45-27||

\twolineshloka
{रावणेनैवमुक्ता तु कुपिता जनकात्मजा}
{प्रत्युवाचानवद्याङ्गी तमनादृत्य राक्षसम्} %||3-45-28||

\twolineshloka
{महागिरिमिवाकम्प्यं महेन्द्रसदृशं पतिम्}
{महोदधिमिवाक्षोभ्यमहं राममनुव्रता} %||3-45-29||

\twolineshloka
{महाबाहुं महोरस्कं सिंहविक्रान्तगामिनम्}
{नृसिंहं सिंहसङ्काशमहं राममनुव्रता} %||3-45-30||

\twolineshloka
{पूर्णचन्द्राननं वीरं राजवत्सं जितेन्द्रियम्}
{पृथुकीर्तिं महाबाहुमहं राममनुव्रता} %||3-45-31||

\twolineshloka
{त्वं पुनर्जम्बुकः सिंहीं मामिहेच्छसि दुर्लभाम्}
{नाहं शक्या त्वया स्प्रष्टुमादित्यस्य प्रभा यथा} %||3-45-32||

\twolineshloka
{पादपान्काञ्चनान्नूनं बहून्पश्यसि मन्दभाक्}
{राघवस्य प्रियां भार्यां यस्त्वमिच्छसि रावण} %||3-45-33||

\twolineshloka
{क्षुधितस्य च सिंहस्य मृगशत्रोस्तरस्विनः}
{आशीविषस्य वदनाद्दंष्ट्रामादातुमिच्छसि} %||3-45-34||

\twolineshloka
{मन्दरं पर्वतश्रेष्ठं पाणिना हर्तुमिच्छसि}
{कालकूटं विषं पीत्वा स्वस्तिमान्गन्तुमिच्छसि} %||3-45-35||

\twolineshloka
{अक्षिसूच्या प्रमृजसि जिह्वया लेढि च क्षुरम्}
{राघवस्य प्रियां भार्यामधिगन्तुं त्वमिच्छसि} %||3-45-36||

\threelineshloka
{अवसज्य शिलां कण्ठे समुद्रं तर्तुमिच्छसि}
{सूर्या चन्द्रमसौ चोभौ प्राणिभ्यां हर्तुमिच्छसि}
{यो रामस्य प्रियां भार्यां प्रधर्षयितुमिच्छसि} %||3-45-37||

\twolineshloka
{अग्निं प्रज्वलितं दृष्ट्वा वस्त्रेणाहर्तुमिच्छसि}
{कल्याण वृत्तां रामस्य यो भार्यां हर्तुमिच्छसि} %||3-45-38||

\twolineshloka
{अयोमुखानां शूलानामग्रे चरितुमिच्छसि}
{रामस्य सदृशीं भार्यां योऽधिगन्तुं त्वमिच्छसि} %||3-45-39||

\fourlineindentedshloka
{यदन्तरं सिंहशृगालयोर्वने}
{ यदन्तरं स्यन्दनिकासमुद्रयोः}
{सुराग्र्यसौवीरकयोर्यदन्तरम्}
{ तदन्तरं दाशरथेस्तवैव च} %||3-45-40||

\fourlineindentedshloka
{यदन्तरं काञ्चनसीसलोहयोर्-}
{ यदन्तरं चन्दनवारिपङ्कयोः}
{यदन्तरं हस्तिबिडालयोर्वने}
{ तदन्तरं दशरथेस्तवैव च} %||3-45-41||

\fourlineindentedshloka
{यदन्तरं वायसवैनतेययोर्-}
{ यदन्तरं मद्गुमयूरयोरपि}
{यदन्तरं सारसगृध्रयोर्वने}
{ तदन्तरं दाशरथेस्तवैव च} %||3-45-42||

\fourlineindentedshloka
{तस्मिन्सहस्राक्षसमप्रभावे}
{ रामे स्थिते कार्मुकबाणपाणौ}
{हृतापि तेऽहं न जरां गमिष्ये}
{ वज्रं यथा मक्षिकयावगीर्णम्} %||3-45-43||

\fourlineindentedshloka
{इतीव तद्वाक्यमदुष्टभावा}
{ सुदृष्टमुक्त्वा रजनीचरं तम्}
{गात्रप्रकम्पाद्व्यथिता बभूव}
{ वातोद्धता सा कदलीव तन्वी} %||3-45-44||

\fourlineindentedshloka
{तां वेपमानामुपलक्ष्य सीताम्}
{ स रावणो मृत्युसमप्रभावः}
{कुलं बलं नाम च कर्म चात्मनः}
{ समाचचक्षे भयकारणार्थम्} %||3-46-45||


{॥इत्यार्षे श्रीमद्रामायणे वाल्मीकीये आदिकाव्ये अरण्यकाण्डे पञ्चचत्वारिंशः सर्गः॥}

\sect{षट्चत्वारिंशः सर्गः}

\twolineshloka
{एवं ब्रुवत्यां सीतायां संरब्धः परुषाक्षरम्}
{ललाटे भ्रुकुटीं कृत्वा रावणः प्रत्युवाच ह} %||3-46-1||

\twolineshloka
{भ्राता वैश्रवणस्याहं सापत्न्यो वरवर्णिनि}
{रावणो नाम भद्रं ते दशग्रीवः प्रतापवान्} %||3-46-2||

\twolineshloka
{यस्य देवाः सगन्धर्वाः पिशाचपतगोरगाः}
{विद्रवन्ति भयाद्भीता मृत्योरिव सदा प्रजाः} %||3-46-3||

\twolineshloka
{येन वैश्रवणो भ्राता वैमात्रः कारणान्तरे}
{द्वन्द्वमासादितः क्रोधाद्रणे विक्रम्य निर्जितः} %||3-46-4||

\twolineshloka
{मद्भयार्तः परित्यज्य स्वमधिष्ठानमृद्धिमत्}
{कैलासं पर्वतश्रेष्ठमध्यास्ते नरवाहनः} %||3-46-5||

\twolineshloka
{यस्य तत्पुष्पकं नाम विमानं कामगं शुभम्}
{वीर्यादावर्जितं भद्रे येन यामि विहायसम्} %||3-46-6||

\twolineshloka
{मम संजातरोषस्य मुखं दृष्ट्वैव मैथिलि}
{विद्रवन्ति परित्रस्ताः सुराः शक्रपुरोगमाः} %||3-46-7||

\twolineshloka
{यत्र तिष्ठाम्यहं तत्र मारुतो वाति शङ्कितः}
{तीव्रांशुः शिशिरांशुश्च भयात्सम्पद्यते रविः} %||3-46-8||

\twolineshloka
{निष्कम्पपत्रास्तरवो नद्यश्च स्तिमितोदकाः}
{भवन्ति यत्र यत्राहं तिष्ठामि च चरामि च} %||3-46-9||

\twolineshloka
{मम पारे समुद्रस्य लङ्का नाम पुरी शुभा}
{सम्पूर्णा राक्षसैर्घोरैर्यथेन्द्रस्यामरावती} %||3-46-10||

\twolineshloka
{प्राकारेण परिक्षिप्ता पाण्डुरेण विराजिता}
{हेमकक्ष्या पुरी रम्या वैदूर्यमय तोरणा} %||3-46-11||

\twolineshloka
{हस्त्यश्वरथसम्भाधा तूर्यनादविनादिता}
{सर्वकामफलैर्वृक्षैः सङ्कुलोद्यानशोभिता} %||3-46-12||

\twolineshloka
{तत्र त्वं वसती सीते राजपुत्रि मया सह}
{न स्रमिष्यसि नारीणां मानुषीणां मनस्विनि} %||3-46-13||

\twolineshloka
{भुञ्जाना मानुषान्भोगान्दिव्यांश्च वरवर्णिनि}
{न स्मरिष्यसि रामस्य मानुषस्य गतायुषः} %||3-46-14||

\twolineshloka
{स्थापयित्वा प्रियं पुत्रं राज्ञा दशरथेन यः}
{मन्दवीर्यः सुतो ज्येष्ठस्ततः प्रस्थापितो वनम्} %||3-46-15||

\twolineshloka
{तेन किं भ्रष्टराज्येन रामेण गतचेतसा}
{करिष्यसि विशालाक्षि तापसेन तपस्विना} %||3-46-16||

\twolineshloka
{सर्वराक्षसभर्तारं कामात्स्वयमिहागतम्}
{न मन्मथशराविष्टं प्रत्याख्यातुं त्वमर्हसि} %||3-46-17||

\twolineshloka
{प्रत्याख्याय हि मां भीरु परितापं गमिष्यसि}
{चरणेनाभिहत्येव पुरूरवसमुर्वशी} %||3-46-18||

\twolineshloka
{एवमुक्ता तु वैदेही क्रुद्धा संरक्तलोचना}
{अब्रवीत्परुषं वाक्यं रहिते राक्षसाधिपम्} %||3-46-19||

\twolineshloka
{कथं वैश्रवणं देवं सर्वभूतनमस्कृतम्}
{भ्रातरं व्यपदिश्य त्वमशुभं कर्तुमिच्छसि} %||3-46-20||

\twolineshloka
{अवश्यं विनशिष्यन्ति सर्वे रावण राक्षसाः}
{येषां त्वं कर्कशो राजा दुर्बुद्धिरजितेन्द्रियः} %||3-46-21||

\twolineshloka
{अपहृत्य शचीं भार्यां शक्यमिन्द्रस्य जीवितुम्}
{न तु रामस्य भार्यां मामपनीयास्ति जीवितम्} %||3-46-22||

\fourlineindentedshloka
{जीवेच्चिरं वज्रधरस्य हस्ता}
{च्छचीं प्रधृष्याप्रतिरूपरूपाम्}
{न मादृशीं राक्षसधर्षयित्वा}
{ पीतामृतस्यापि तवास्ति मोक्षः} %||3-47-23||


{॥इत्यार्षे श्रीमद्रामायणे वाल्मीकीये आदिकाव्ये अरण्यकाण्डे षट्चत्वारिंशः सर्गः॥}

\sect{सप्तचत्वारिंशः सर्गः}

\twolineshloka
{सीताया वचनं श्रुत्वा दशग्रीवः प्रतापवान्}
{हस्ते हस्तं समाहत्य चकार सुमहद्वपुः} %||3-47-1||

\twolineshloka
{स मैथिलीं पुनर्वाक्यं बभाषे च ततो भृशम्}
{नोन्मत्तया श्रुतौ मन्ये मम वीर्यपराक्रमौ} %||3-47-2||

\twolineshloka
{उद्वहेयं भुजाभ्यां तु मेदिनीमम्बरे स्थितः}
{आपिबेयं समुद्रं च मृत्युं हन्यां रणे स्थितः} %||3-47-3||

\twolineshloka
{अर्कं रुन्ध्यां शरैस्तीक्ष्णैर्विभिन्द्यां हि महीतलम्}
{कामरूपिणमुन्मत्ते पश्य मां कामदं पतिम्} %||3-47-4||

\twolineshloka
{एवमुक्तवतस्तस्य रावणस्य शिखिप्रभे}
{क्रुद्धस्य हरिपर्यन्ते रक्ते नेत्रे बभूवतुः} %||3-47-5||

\twolineshloka
{सद्यः सौम्यं परित्यज्य भिक्षुरूपं स रावणः}
{स्वं रूपं कालरूपाभं भेजे वैश्रवणानुजः} %||3-47-6||

\twolineshloka
{संरक्तनयनः श्रीमांस्तप्तकाञ्चनकुण्डलः}
{दशास्यः कार्मुकी बाणी बभूव क्षणदाचरः} %||3-47-7||

\twolineshloka
{स परिव्राजकच्छद्म महाकायो विहाय तत्}
{प्रतिपेदे स्वकं रूपं रावणो राक्षसाधिपः} %||3-47-8||

\twolineshloka
{संरक्तनयनः क्रोधाज्जीमूतनिचयप्रभः}
{रक्ताम्बरधरस्तस्थौ स्त्रीरत्नं प्रेक्ष्य मैथिलीम्} %||3-47-9||

\twolineshloka
{स तामसितकेशान्तां भास्करस्य प्रभामिव}
{वसनाभरणोपेतां मैथिलीं रावणोऽब्रवीत्} %||3-47-10||

\twolineshloka
{त्रिषु लोकेषु विख्यातं यदि भर्तारमिच्छसि}
{मामाश्रय वरारोहे तवाहं सदृशः पतिः} %||3-47-11||

\threelineshloka
{मां भजस्व चिराय त्वमहं श्लाघ्यस्तव प्रियः}
{नैव चाहं क्वचिद्भद्रे करिष्ये तव विप्रियम्}
{त्यज्यतां मानुषो भावो मयि भावः प्रणीयताम्} %||3-47-12||

\twolineshloka
{राज्याच्च्युतमसिद्धार्थं रामं परिमितायुषम्}
{कैर्गुणैरनुरक्तासि मूढे पण्डितमानिनि} %||3-47-13||

\twolineshloka
{यः स्त्रिया वचनाद्राज्यं विहाय ससुहृज्जनम्}
{अस्मिन्व्यालानुचरिते वने वसति दुर्मतिः} %||3-47-14||

\twolineshloka
{इत्युक्त्वा मैथिलीं वाक्यं प्रियार्हां प्रियवादिनीम्}
{जग्राह रावणः सीतां बुधः खे रोहिणीमिव} %||3-47-15||

\twolineshloka
{वामेन सीतां पद्माक्षीं मूर्धजेषु करेण सः}
{ऊर्वोस्तु दक्षिणेनैव परिजग्राह पाणिना} %||3-47-16||

\twolineshloka
{तं दृष्ट्वा गिरिशृङ्गाभं तीक्ष्णदंष्ट्रं महाभुजम्}
{प्राद्रवन्मृत्युसङ्काशं भयार्ता वनदेवताः} %||3-47-17||

\twolineshloka
{स च मायामयो दिव्यः खरयुक्तः खरस्वनः}
{प्रत्यदृश्यत हेमाङ्गो रावणस्य महारथः} %||3-47-18||

\twolineshloka
{ततस्तां परुषैर्वाक्यैरभितर्ज्य महास्वनः}
{अङ्केनादाय वैदेहीं रथमारोपयत्तदा} %||3-47-19||

\twolineshloka
{सा गृहीतातिचुक्रोश रावणेन यशस्विनी}
{रामेति सीता दुःखार्ता रामं दूरगतं वने} %||3-47-20||

\twolineshloka
{तामकामां स कामार्तः पन्नगेन्द्रवधूमिव}
{विवेष्टमानामादाय उत्पपाथाथ रावणः} %||3-47-21||

\twolineshloka
{ततः सा राक्षसेन्द्रेण ह्रियमाणा विहायसा}
{भृशं चुक्रोश मत्तेव भ्रान्तचित्ता यथातुरा} %||3-47-22||

\twolineshloka
{हा लक्ष्मण महाबाहो गुरुचित्तप्रसादक}
{ह्रियमाणां न जानीषे रक्षसा कामरूपिणा} %||3-47-23||

\twolineshloka
{जीवितं सुखमर्थांश्च धर्महेतोः परित्यजन्}
{ह्रियमाणामधर्मेण मां राघव न पश्यसि} %||3-47-24||

\twolineshloka
{ननु नामाविनीतानां विनेतासि परन्तप}
{कथमेवंविधं पापं न त्वं शाधि हि रावणम्} %||3-47-25||

\twolineshloka
{ननु सद्योऽविनीतस्य दृश्यते कर्मणः फलम्}
{कालोऽप्यङ्गी भवत्यत्र सस्यानामिव पक्तये} %||3-47-26||

\twolineshloka
{स कर्म कृतवानेतत्कालोपहतचेतनः}
{जीवितान्तकरं घोरं रामाद्व्यसनमाप्नुहि} %||3-47-27||

\twolineshloka
{हन्तेदानीं सकामा तु कैकेयी बान्धवैः सह}
{ह्रियेयं धर्मकामस्य धर्मपत्नी यशस्विनः} %||3-47-28||

\twolineshloka
{आमन्त्रये जनस्थानं कर्णिकारांश्च पुष्पितान्}
{क्षिप्रं रामाय शंसध्वं सीतां हरति रावणः} %||3-47-29||

\twolineshloka
{माल्यवन्तं शिखरिणं वन्दे प्रस्रवणं गिरिम्}
{क्षिप्रं रामाय शंसध्वं सीतां हरति रावणः} %||3-47-30||

\twolineshloka
{हंससारससङ्घुष्टां वन्दे गोदावरीं नदीम्}
{क्षिप्रं रामाय शंसध्वं सीतां हरति रावणः} %||3-47-31||

\twolineshloka
{दैवतानि च यान्त्यस्मिन्वने विविधपादपे}
{नमस्करोम्यहं तेभ्यो भर्तुः शंसत मां हृताम्} %||3-47-32||

\twolineshloka
{यानि कानिचिदप्यत्र सत्त्वानि निवसन्त्युत}
{सर्वाणि शरणं यामि मृगपक्षिगणानपि} %||3-47-33||

\twolineshloka
{ह्रियमाणां प्रियां भर्तुः प्राणेभ्योऽपि गरीयसीम्}
{विवशापहृता सीता रावणेनेति शंसत} %||3-47-34||

\twolineshloka
{विदित्वा मां महाबाहुरमुत्रापि महाबलः}
{आनेष्यति पराक्रम्य वैवस्वतहृतामपि} %||3-47-35||

\twolineshloka
{रामाय तु यथातत्त्वं जटायो हरणं मम}
{लक्ष्मणाय च तत्सर्वमाख्यातव्यमशेषतः} %||3-48-36||


{॥इत्यार्षे श्रीमद्रामायणे वाल्मीकीये आदिकाव्ये अरण्यकाण्डे सप्तचत्वारिंशः सर्गः॥}

\sect{अष्टचत्वारिंशः सर्गः}

\twolineshloka
{तं शब्दमवसुप्तस्य जटायुरथ शुश्रुवे}
{निरैक्षद्रावणं क्षिप्रं वैदेहीं च ददर्श सः} %||3-48-1||

\twolineshloka
{ततः पर्वतकूटाभस्तीक्ष्णतुण्डः खगोत्तमः}
{वनस्पतिगतः श्रीमान्व्याजहार शुभां गिरम्} %||3-48-2||

\twolineshloka
{दशग्रीवस्थितो धर्मे पुराणे सत्यसंश्रयः}
{जटायुर्नाम नाम्नाहं गृध्रराजो महाबलः} %||3-48-3||

\twolineshloka
{राजा सर्वस्य लोकस्य महेन्द्रवरुणोपमः}
{लोकानां च हिते युक्तो रामो दशरथात्मजः} %||3-48-4||

\twolineshloka
{तस्यैषा लोकनाथस्य धर्मपत्नी यशस्विनी}
{सीता नाम वरारोहा यां त्वं हर्तुमिहेच्छसि} %||3-48-5||

\threelineshloka
{कथं राजा स्थितो धर्मे परदारान्परामृशेत्}
{रक्षणीया विशेषेण राजदारा महाबलः}
{निवर्तय मतिं नीचां परदाराभिमर्शनम्} %||3-48-6||

\twolineshloka
{न तत्समाचरेद्धीरो यत्परोऽस्य विगर्हयेत्}
{यथात्मनस्तथान्येषां दारा रक्ष्या विमर्शनात्} %||3-48-7||

\twolineshloka
{अर्थं वा यदि वा कामं शिष्टाः शास्त्रेष्वनागतम्}
{व्यवस्यन्त्यनु राजानं धर्मं पौरस्त्यनन्दन} %||3-48-8||

\twolineshloka
{राजा धर्मश्च कामश्च द्रव्याणां चोत्तमो निधिः}
{धर्मः शुभं वा पापं वा राजमूलं प्रवर्तते} %||3-48-9||

\twolineshloka
{पापस्वभावश्चपलः कथं त्वं रक्षसां वर}
{ऐश्वर्यमभिसम्प्राप्तो विमानमिव दुष्कृती} %||3-48-10||

\twolineshloka
{कामस्वभावो यो यस्य न स शक्यः प्रमार्जितुम्}
{न हि दुष्टात्मनामार्य मा वसत्यालये चिरम्} %||3-48-11||

\twolineshloka
{विषये वा पुरे वा ते यदा रामो महाबलः}
{नापराध्यति धर्मात्मा कथं तस्यापराध्यसि} %||3-48-12||

\twolineshloka
{यदि शूर्पणखाहेतोर्जनस्थानगतः खरः}
{अतिवृत्तो हतः पूर्वं रामेणाक्लिष्टकर्मणा} %||3-48-13||

\twolineshloka
{अत्र ब्रूहि यथासत्यं को रामस्य व्यतिक्रमः}
{यस्य त्वं लोकनाथस्य हृत्वा भार्यां गमिष्यसि} %||3-48-14||

\twolineshloka
{क्षिप्रं विसृज वैदेहीं मा त्वा घोरेण चक्षुषा}
{दहेद्दहन भूतेन वृत्रमिन्द्राशनिर्यथा} %||3-48-15||

\twolineshloka
{सर्पमाशीविषं बद्ध्वा वस्त्रान्ते नावबुध्यसे}
{ग्रीवायां प्रतिमुक्तं च कालपाशं न पश्यसि} %||3-48-16||

\twolineshloka
{स भारः सौम्य भर्तव्यो यो नरं नावसादयेत्}
{तदन्नमुपभोक्तव्यं जीर्यते यदनामयम्} %||3-48-17||

\twolineshloka
{यत्कृत्वा न भवेद्धर्मो न कीर्तिर्न यशो भुवि}
{शरीरस्य भवेत्खेदः कस्तत्कर्म समाचरेत्} %||3-48-18||

\twolineshloka
{षष्टिवर्षसहस्राणि मम जातस्य रावण}
{पितृपैतामहं राज्यं यथावदनुतिष्ठतः} %||3-48-19||

\twolineshloka
{वृद्धोऽहं त्वं युवा धन्वी सरथः कवची शरी}
{तथाप्यादाय वैदेहीं कुशली न गमिष्यसि} %||3-48-20||

\twolineshloka
{न शक्तस्त्वं बलाद्धर्तुं वैदेहीं मम पश्यतः}
{हेतुभिर्न्यायसंयुक्तैर्ध्रुवां वेदश्रुतीमिव} %||3-48-21||

\twolineshloka
{युध्यस्व यदि शूरोऽसि मुहूर्तं तिष्ठ रावण}
{शयिष्यसे हतो भूमौ यथापूर्वं खरस्तथा} %||3-48-22||

\twolineshloka
{असकृत्संयुगे येन निहता दैत्यदानवाः}
{नचिराच्चीरवासास्त्वां रामो युधि वधिष्यति} %||3-48-23||

\twolineshloka
{किं नु शक्यं मया कर्तुं गतौ दूरं नृपात्मजौ}
{क्षिप्रं त्वं नश्यसे नीच तयोर्भीतो न संशयः} %||3-48-24||

\twolineshloka
{न हि मे जीवमानस्य नयिष्यसि शुभामिमाम्}
{सीतां कमलपत्राक्षीं रामस्य महषीं प्रियाम्} %||3-48-25||

\twolineshloka
{अवश्यं तु मया कार्यं प्रियं तस्य महात्मनः}
{जीवितेनापि रामस्य तथा दशरथस्य च} %||3-48-26||

\threelineshloka
{तिष्ठ तिष्ठ दशग्रीव मुहूर्तं पश्य रावण}
{युद्धातिथ्यं प्रदास्यामि यथाप्राणं निशाचर}
{वृन्तादिव फलं त्वां तु पातयेयं रथोत्तमात्} %||3-49-27||


{॥इत्यार्षे श्रीमद्रामायणे वाल्मीकीये आदिकाव्ये अरण्यकाण्डे अष्टचत्वारिंशः सर्गः॥}

\sect{एकोनपञ्चाशत्तमः सर्गः}

\twolineshloka
{इत्युक्तस्य यथान्यायं रावणस्य जटायुषा}
{क्रुद्धस्याग्निनिभाः सर्वा रेजुर्विंशतिदृष्टयः} %||3-49-1||

\twolineshloka
{संरक्तनयनः कोपात्तप्तकाञ्चनकुण्डलः}
{राक्षसेन्द्रोऽभिदुद्राव पतगेन्द्रममर्षणः} %||3-49-2||

\twolineshloka
{स सम्प्रहारस्तुमुलस्तयोस्तस्मिन्महावने}
{बभूव वातोद्धतयोर्मेघयोर्गगने यथा} %||3-49-3||

\twolineshloka
{तद्बभूवाद्भुतं युद्धं गृध्रराक्षसयोस्तदा}
{सपक्षयोर्माल्यवतोर्महापर्वतयोरिव} %||3-49-4||

\twolineshloka
{ततो नालीकनाराचैस्तीक्ष्णाग्रैश्च विकर्णिभिः}
{अभ्यवर्षन्महाघोरैर्गृध्रराजं महाबलः} %||3-49-5||

\twolineshloka
{स तानि शरजालानि गृध्रः पत्ररथेश्वरः}
{जटायुः प्रतिजग्राह रावणास्त्राणि संयुगे} %||3-49-6||

\twolineshloka
{तस्य तीक्ष्णनखाभ्यां तु चरणाभ्यां महाबलः}
{चकार बहुधा गात्रे व्रणान्पतगसत्तमः} %||3-49-7||

\twolineshloka
{अथ क्रोधाद्दशग्रीवो जग्राह दशमार्गणान्}
{मृत्युदण्डनिभान्घोराञ्शत्रुमर्दनकाङ्क्षया} %||3-49-8||

\twolineshloka
{स तैर्बाणैर्महावीर्यः पूर्णमुक्तैरजिह्मगैः}
{बिभेद निशितैस्तीक्ष्णैर्गृध्रं घोरैः शिलीमुखैः} %||3-49-9||

\twolineshloka
{स राक्षसरथे पश्यञ्जानकीं बाष्पलोचनाम्}
{अचिन्तयित्वा बाणांस्तान्राक्षसं समभिद्रवत्} %||3-49-10||

\twolineshloka
{ततोऽस्य सशरं चापं मुक्तामणिविभूषितम्}
{चरणाभ्यां महातेजा बभञ्ज पतगेश्वरः} %||3-49-11||

\twolineshloka
{तच्चाग्निसदृशं दीप्तं रावणस्य शरावरम्}
{पक्षाभ्यां च महातेजा व्यधुनोत्पतगेश्वरः} %||3-49-12||

\twolineshloka
{काञ्चनोरश्छदान्दिव्यान्पिशाचवदनान्खरान्}
{तांश्चास्य जवसम्पन्नाञ्जघान समरे बली} %||3-49-13||

\threelineshloka
{वरं त्रिवेणुसम्पन्नं कामगं पावकार्चिषम्}
{मणिहेमविचित्राङ्गं बभञ्ज च महारथम्}
{पूर्णचन्द्रप्रतीकाशं छत्रं च व्यजनैः सह} %||3-49-14||

\twolineshloka
{स भग्नधन्वा विरथो हताश्वो हतसारथिः}
{अङ्केनादाय वैदेहीं पपात भुवि रावणः} %||3-49-15||

\twolineshloka
{दृष्ट्वा निपतितं भूमौ रावणं भग्नवाहनम्}
{साधु साध्विति भूतानि गृध्रराजमपूजयन्} %||3-49-16||

\twolineshloka
{परिश्रान्तं तु तं दृष्ट्वा जरया पक्षियूथपम्}
{उत्पपात पुनर्हृष्टो मैथिलीं गृह्य रावणः} %||3-49-17||

\twolineshloka
{तं प्रहृष्टं निधायाङ्के गच्छन्तं जनकात्मजाम्}
{गृध्रराजः समुत्पत्य जटायुरिदमब्रवीत्} %||3-49-18||

\twolineshloka
{वज्रसंस्पर्शबाणस्य भार्यां रामस्य रावण}
{अल्पबुद्धे हरस्येनां वधाय खलु रक्षसाम्} %||3-49-19||

\twolineshloka
{समित्रबन्धुः सामात्यः सबलः सपरिच्छदः}
{विषपानं पिबस्येतत्पिपासित इवोदकम्} %||3-49-20||

\twolineshloka
{अनुबन्धमजानन्तः कर्मणामविचक्षणाः}
{शीघ्रमेव विनश्यन्ति यथा त्वं विनशिष्यसि} %||3-49-21||

\twolineshloka
{बद्धस्त्वं कालपाशेन क्व गतस्तस्य मोक्ष्यसे}
{वधाय बडिशं गृह्य सामिषं जलजो यथा} %||3-49-22||

\twolineshloka
{न हि जातु दुराधर्षौ काकुत्स्थौ तव रावण}
{धर्षणं चाश्रमस्यास्य क्षमिष्येते तु राघवौ} %||3-49-23||

\twolineshloka
{यथा त्वया कृतं कर्म भीरुणा लोकगर्हितम्}
{तस्कराचरितो मार्गो नैष वीरनिषेवितः} %||3-49-24||

\twolineshloka
{युध्यस्व यदि शूरोऽसि मुहूर्तं तिष्ठ रावण}
{शयिष्यसे हतो भूमौ यथा भ्राता खरस्तथा} %||3-49-25||

\twolineshloka
{परेतकाले पुरुषो यत्कर्म प्रतिपद्यते}
{विनाशायात्मनोऽधर्म्यं प्रतिपन्नोऽसि कर्म तत्} %||3-49-26||

\twolineshloka
{पापानुबन्धो वै यस्य कर्मणः को नु तत्पुमान्}
{कुर्वीत लोकाधिपतिः स्वयम्भूर्भगवानपि} %||3-49-27||

\twolineshloka
{एवमुक्त्वा शुभं वाक्यं जटायुस्तस्य रक्षसः}
{निपपात भृशं पृष्ठे दशग्रीवस्य वीर्यवान्} %||3-49-28||

\twolineshloka
{तं गृहीत्वा नखैस्तीक्ष्णैर्विरराद समन्ततः}
{अधिरूढो गजारोहि यथा स्याद्दुष्टवारणम्} %||3-49-29||

\twolineshloka
{विरराद नखैरस्य तुण्डं पृष्ठे समर्पयन्}
{केशांश्चोत्पाटयामास नखपक्षमुखायुधः} %||3-49-30||

\twolineshloka
{स तथा गृध्रराजेन क्लिश्यमानो मुहुर्मुहुः}
{अमर्षस्फुरितौष्ठः सन्प्राकम्पत स राक्षसः} %||3-49-31||

\twolineshloka
{सम्परिष्वज्य वैदेहीं वामेनाङ्केन रावणः}
{तलेनाभिजघानार्तो जटायुं क्रोधमूर्छितः} %||3-49-32||

\twolineshloka
{जटायुस्तमतिक्रम्य तुण्डेनास्य खराधिपः}
{वामबाहून्दश तदा व्यपाहरदरिन्दमः} %||3-49-33||

\twolineshloka
{ततः क्रुद्धो दशक्रीवः सीतामुत्सृज्य वीर्यवान्}
{मुष्टिभ्यां चरणाभ्यां च गृध्रराजमपोथयत्} %||3-49-34||

\twolineshloka
{ततो मुहूर्तं सङ्ग्रामो बभूवातुलवीर्ययोः}
{राक्षसानां च मुख्यस्य पक्षिणां प्रवरस्य च} %||3-49-35||

\twolineshloka
{तस्य व्यायच्छमानस्य रामस्यार्थेऽथ रावणः}
{पक्षौ पादौ च पार्श्वौ च खड्गमुद्धृत्य सोऽच्छिनत्} %||3-49-36||

\twolineshloka
{स छिन्नपक्षः सहसा रक्षसा रौद्रकर्मणा}
{निपपात हतो गृध्रो धरण्यामल्पजीवितः} %||3-49-37||

\twolineshloka
{तं दृष्ट्वा पतितं भूमौ क्षतजार्द्रं जटायुषम्}
{अभ्यधावत वैदेही स्वबन्धुमिव दुःखिता} %||3-49-38||

\fourlineindentedshloka
{तं नीलजीमूतनिकाशकल्पम्}
{ सुपाण्डुरोरस्कमुदारवीर्यम्}
{ददर्श लङ्काधिपतिः पृथिव्याम्}
{ जटायुषं शान्तमिवाग्निदावम्} %||3-49-39||

\fourlineindentedshloka
{ततस्तु तं पत्ररथं महीतले}
{ निपातितं रावणवेगमर्दितम्}
{पुनः परिष्वज्य शशिप्रभानना}
{ रुरोद सीता जनकात्मजा तदा} %||3-50-40||


{॥इत्यार्षे श्रीमद्रामायणे वाल्मीकीये आदिकाव्ये अरण्यकाण्डे एकोनपञ्चाशत्तमः सर्गः॥}

\sect{पञ्चाशत्तमः सर्गः}

\twolineshloka
{तमल्पजीवितं भूमौ स्फुरन्तं राक्षसाधिपः}
{ददर्श गृध्रं पतितं समीपे राघवाश्रमात्} %||3-50-1||

\twolineshloka
{सा तु ताराधिपमुखी रावणेन समीक्ष्य तम्}
{गृध्रराजं विनिहतं विललाप सुदुःखिता} %||3-50-2||

\twolineshloka
{निमित्तं लक्षणज्ञानं शकुनिस्वरदर्शनम्}
{अवश्यं सुखदुःखेषु नराणां प्रतिदृश्यते} %||3-50-3||

\twolineshloka
{न नूनं राम जानासि महद्व्यसनमात्मजः}
{धावन्ति नूनं काकुत्स्थ मदर्थं मृगपक्षिणः} %||3-50-4||

\twolineshloka
{त्राहि मामद्य काकुत्स्थ लक्ष्मणेति वराङ्गना}
{सुसन्त्रस्ता समाक्रन्दच्छृण्वतां तु यथान्तिके} %||3-50-5||

\twolineshloka
{तां क्लिष्टमाल्याभरणां विलपन्तीमनाथवत्}
{अभ्यधावत वैदेहीं रावणो राक्षसाधिपः} %||3-50-6||

\twolineshloka
{तां लतामिव वेष्टन्तीमालिङ्गन्तीं महाद्रुमान्}
{मुञ्च मुञ्चेति बहुशः प्रवदन्राक्षसाधिपः} %||3-50-7||

\twolineshloka
{क्रोशन्तीं राम रामेति रामेण रहितां वने}
{जीवितान्ताय केशेषु जग्राहान्तकसंनिभः} %||3-50-8||

\twolineshloka
{प्रधर्षितायां वैदेह्यां बभूव सचराचरम्}
{जगत्सर्वममर्यादं तमसान्धेन संवृतम्} %||3-50-9||

\twolineshloka
{दृष्ट्वा सीतां परामृष्टां दीनां दिव्येन चक्षुषा}
{कृतं कार्यमिति श्रीमान्व्याजहार पितामहः} %||3-50-10||

\twolineshloka
{प्रहृष्टा व्यथिताश्चासन्सर्वे ते परमर्षयः}
{दृष्ट्वा सीतां परामृष्टां दण्डकारण्यवासिनः} %||3-50-11||

\twolineshloka
{स तु तां राम रामेति रुदन्तीं लक्ष्मणेति च}
{जगामाकाशमादाय रावणो राक्षसाधिपः} %||3-50-12||

\twolineshloka
{तप्ताभरणसर्वाङ्गी पीतकौशेयवासनी}
{रराज राजपुत्री तु विद्युत्सौदामनी यथा} %||3-50-13||

\twolineshloka
{उद्धूतेन च वस्त्रेण तस्याः पीतेन रावणः}
{अधिकं परिबभ्राज गिरिर्दीप इवाग्निना} %||3-50-14||

\twolineshloka
{तस्याः परमकल्याण्यास्ताम्राणि सुरभीणि च}
{पद्मपत्राणि वैदेह्या अभ्यकीर्यन्त रावणम्} %||3-50-15||

\twolineshloka
{तस्याः कौशेयमुद्धूतमाकाशे कनकप्रभम्}
{बभौ चादित्यरागेण ताम्रमभ्रमिवातपे} %||3-50-16||

\twolineshloka
{तस्यास्तद्विमलं वक्त्रमाकाशे रावणाङ्कगम्}
{न रराज विना रामं विनालमिव पङ्कजम्} %||3-50-17||

\threelineshloka
{बभूव जलदं नीलं भित्त्वा चन्द्र इवोदितः}
{सुललाटं सुकेशान्तं पद्मगर्भाभमव्रणम्}
{शुक्लैः सुविमलैर्दन्तैः प्रभावद्भिरलङ्कृतम्} %||3-50-18||

\twolineshloka
{रुदितं व्यपमृष्टास्त्रं चन्द्रवत्प्रियदर्शनम्}
{सुनासं चारुताम्रौष्ठमाकाषे हाटकप्रभम्} %||3-50-19||

\twolineshloka
{राक्षसेन्द्रसमाधूतं तस्यास्तद्वचनं शुभम्}
{शुशुभे न विना रामं दिवा चन्द्र इवोदितः} %||3-50-20||

\twolineshloka
{सा हेमवर्णा नीलाङ्गं मैथिली राक्षसाधिपम्}
{शुशुभे काञ्चनी काञ्ची नीलं मणिमिवाश्रिता} %||3-50-21||

\twolineshloka
{सा पद्मगौरी हेमाभा रावणं जनकात्मजा}
{विद्युद्घनमिवाविश्य शुशुभे तप्तभूषणा} %||3-50-22||

\twolineshloka
{तस्या भूषणघोषेण वैदेह्या राक्षसाधिपः}
{बभूव विमलो नीलः सघोष इव तोयदः} %||3-50-23||

\twolineshloka
{उत्तमाङ्गच्युता तस्याः पुष्पवृष्टिः समन्ततः}
{सीताया ह्रियमाणायाः पपात धरणीतले} %||3-50-24||

\twolineshloka
{सा तु रावणवेगेन पुष्पवृष्टिः समन्ततः}
{समाधूता दशग्रीवं पुनरेवाभ्यवर्तत} %||3-50-25||

\twolineshloka
{अभ्यवर्तत पुष्पाणां धारा वैश्रवणानुजम्}
{नक्षत्रमालाविमला मेरुं नगमिवोत्तमम्} %||3-50-26||

\twolineshloka
{चरणान्नूपुरं भ्रष्टं वैदेह्या रत्नभूषितम्}
{विद्युन्मण्डलसङ्काशं पपात मधुरस्वनम्} %||3-50-27||

\twolineshloka
{तरुप्रवालरक्ता सा नीलाङ्गं राक्षसेश्वरम्}
{प्राशोभयत वैदेही गजं कष्येव काञ्चनी} %||3-50-28||

\twolineshloka
{तां महोल्कामिवाकाशे दीप्यमानां स्वतेजसा}
{जहाराकाशमाविश्य सीतां वैश्रवणानुजः} %||3-50-29||

\twolineshloka
{तस्यास्तान्यग्निवर्णानि भूषणानि महीतले}
{सघोषाण्यवकीर्यन्त क्षीणास्तारा इवाम्बरात्} %||3-50-30||

\twolineshloka
{तस्याः स्तनान्तराद्भ्रष्टो हारस्ताराधिपद्युतिः}
{वैदेह्या निपतन्भाति गङ्गेव गगनाच्च्युता} %||3-50-31||

\twolineshloka
{उत्पात वाताभिहता नानाद्विज गणायुताः}
{मा भैरिति विधूताग्रा व्याजह्रुरिव पादपाः} %||3-50-32||

\twolineshloka
{नलिन्यो ध्वस्तकमलास्त्रस्तमीनजले चराः}
{सखीमिव गतोत्साहां शोचन्तीव स्म मैथिलीम्} %||3-50-33||

\twolineshloka
{समन्तादभिसम्पत्य सिंहव्याघ्रमृगद्विजाः}
{अन्वधावंस्तदा रोषात्सीताच्छायानुगामिनः} %||3-50-34||

\twolineshloka
{जलप्रपातास्रमुखाः शृङ्गैरुच्छ्रितबाहवः}
{सीतायां ह्रियमाणायां विक्रोशन्तीव पर्वताः} %||3-50-35||

\twolineshloka
{ह्रियमाणां तु वैदेहीं दृष्ट्वा दीनो दिवाकरः}
{प्रविध्वस्तप्रभः श्रीमानासीत्पाण्डुरमण्डलः} %||3-50-36||

\twolineshloka
{नास्ति धर्मः कुतः सत्यं नार्जवं नानृशंसता}
{यत्र रामस्य वैदेहीं भार्यां हरति रावणः} %||3-50-37||

\twolineshloka
{इति सर्वाणि भूतानि गणशः पर्यदेवयन्}
{वित्रस्तका दीनमुखा रुरुदुर्मृगपोतकाः} %||3-50-38||

\twolineshloka
{उद्वीक्ष्योद्वीक्ष्य नयनैरास्रपाताविलेक्षणाः}
{सुप्रवेपितगात्राश्च बभूवुर्वनदेवताः} %||3-50-39||

\twolineshloka
{विक्रोशन्तीं दृढं सीतां दृष्ट्वा दुःखं तथा गताम्}
{तां तु लक्ष्मण रामेति क्रोशन्तीं मधुरस्वराम्} %||3-50-40||

\threelineshloka
{अवेक्षमाणां बहुषो वैदेहीं धरणीतलम्}
{स तामाकुलकेशान्तां विप्रमृष्टविशेषकाम्}
{जहारात्मविनाशाय दशग्रीवो मनस्विनाम्} %||3-50-41||

\fourlineindentedshloka
{ततस्तु सा चारुदती शुचिस्मिता}
{ विनाकृता बन्धुजनेन मैथिली}
{अपश्यती राघवलक्ष्मणावुभौ}
{ विवर्णवक्त्रा भयभारपीडिता} %||3-51-42||


{॥इत्यार्षे श्रीमद्रामायणे वाल्मीकीये आदिकाव्ये अरण्यकाण्डे पञ्चाशत्तमः सर्गः॥}

\sect{एकपञ्चाशत्तमः सर्गः}

\twolineshloka
{खमुत्पतन्तं तं दृष्ट्वा मैथिली जनकात्मजा}
{दुःखिता परमोद्विग्ना भये महति वर्तिनी} %||3-51-1||

\twolineshloka
{रोषरोदनताम्राक्षी भीमाक्षं राक्षसाधिपम्}
{रुदती करुणं सीता ह्रियमाणेदमब्रवीत्} %||3-51-2||

\twolineshloka
{न व्यपत्रपसे नीच कर्मणानेन रावण}
{ज्ञात्वा विरहितां यो मां चोरयित्वा पलायसे} %||3-51-3||

\threelineshloka
{त्वयैव नूनं दुष्टात्मन्भीरुणा हर्तुमिच्छता}
{ममापवाहितो भर्ता मृगरूपेण मायया}
{यो हि मामुद्यतस्त्रातुं सोऽप्ययं विनिपातितः} %||3-51-4||

\twolineshloka
{परमं खलु ते वीर्यं दृश्यते राक्षसाधम}
{विश्राव्य नामधेयं हि युद्धे नास्ति जिता त्वया} %||3-51-5||

\twolineshloka
{ईदृशं गर्हितं कर्म कथं कृत्वा न लज्जसे}
{स्त्रियाश्च हरणं नीच रहिते च परस्य च} %||3-51-6||

\twolineshloka
{कथयिष्यन्ति लोकेषु पुरुषाः कर्म कुत्सितम्}
{सुनृशंसमधर्मिष्ठं तव शौण्डीर्यमानिनः} %||3-51-7||

\twolineshloka
{धिक्ते शौर्यं च सत्त्वं च यत्त्वया कथितं तदा}
{कुलाक्रोशकरं लोके धिक्ते चारित्रमीदृशम्} %||3-51-8||

\twolineshloka
{किं शक्यं कर्तुमेवं हि यज्जवेनैव धावसि}
{मुहूर्तमपि तिष्ठस्व न जीवन्प्रतियास्यसि} %||3-51-9||

\twolineshloka
{न हि चक्षुःपथं प्राप्य तयोः पार्थिवपुत्रयोः}
{ससैन्योऽपि समर्तःस्त्वं मुहूर्तमपि जीवितुम्} %||3-51-10||

\twolineshloka
{न त्वं तयोः शरस्पर्शं सोढुं शक्तः कथञ्चन}
{वने प्रज्वलितस्येव स्पर्शमग्नेर्विहङ्गमः} %||3-51-11||

\threelineshloka
{साधु कृत्वात्मनः पथ्यं साधु मां मुञ्च रावण}
{मत्प्रधर्षणरुष्टो हि भ्रात्रा सह पतिर्मम}
{विधास्यति विनाशाय त्वं मां यदि न मुञ्चसि} %||3-51-12||

\twolineshloka
{येन त्वं व्यवसायेन बलान्मां हर्तुमिच्छसि}
{व्यवसायः स ते नीच भविष्यति निरर्थकः} %||3-51-13||

\twolineshloka
{न ह्यहं तमपश्यन्ती भर्तारं विबुधोपमम्}
{उत्सहे शत्रुवशगा प्राणान्धारयितुं चिरम्} %||3-51-14||

\twolineshloka
{न नूनं चात्मनः श्रेयः पथ्यं वा समवेक्षसे}
{मृत्युकाले यथा मर्त्यो विपरीतानि सेवते} %||3-51-15||

\twolineshloka
{मुमूर्षूणां हि सर्वेषां यत्पथ्यं तन्न रोचते}
{पश्यामीव हि कण्ठे त्वां कालपाशावपाशितम्} %||3-51-16||

\twolineshloka
{यथा चास्मिन्भयस्थाने न बिभेषे दशानन}
{व्यक्तं हिरण्मयान्हि त्वं सम्पश्यसि महीरुहान्} %||3-51-17||

\twolineshloka
{नदीं वैरतणीं घोरां रुधिरौघनिवाहिनीम्}
{खड्गपत्रवनं चैव भीमं पश्यसि रावण} %||3-51-18||

\twolineshloka
{तप्तकाञ्चनपुष्पां च वैदूर्यप्रवरच्छदाम्}
{द्रक्ष्यसे शाल्मलीं तीक्ष्णामायसैः कण्टकैश्चिताम्} %||3-51-19||

\twolineshloka
{न हि त्वमीदृशं कृत्वा तस्यालीकं महात्मनः}
{धारितुं शक्ष्यसि चिरं विषं पीत्वेव निर्घृणः} %||3-51-20||

\twolineshloka
{बद्धस्त्वं कालपाशेन दुर्निवारेण रावण}
{क्व गतो लप्स्यसे शर्म भर्तुर्मम महात्मनः} %||3-51-21||

\twolineshloka
{निमेषान्तरमात्रेण विना भ्रातरमाहवे}
{राक्षसा निहता येन सहस्राणि चतुर्दश} %||3-51-22||

\twolineshloka
{स कथं राघवो वीरः सर्वास्त्रकुशलो बली}
{न त्वां हन्याच्छरैस्तीक्ष्णैरिष्टभार्यापहारिणम्} %||3-51-23||

\twolineshloka
{एतच्चान्यच्च परुषं वैदेही रावणाङ्कगा}
{भयशोकसमाविष्टा करुणं विललाप ह} %||3-51-24||

\fourlineindentedshloka
{तथा भृशार्तां बहु चैव भाषिणीम्}
{ विललाप पूर्वं करुणं च भामिनीम्}
{जहार पापस्तरुणीं विवेष्टतीम्}
{ नृपात्मजामागतगात्रवेपथुम्} %||3-52-25||


{॥इत्यार्षे श्रीमद्रामायणे वाल्मीकीये आदिकाव्ये अरण्यकाण्डे एकपञ्चाशत्तमः सर्गः॥}

\sect{द्विपञ्चाशत्तमः सर्गः}

\twolineshloka
{ह्रियमाणा तु वैदेही कञ्चिन्नाथमपश्यती}
{ददर्श गिरिशृङ्गस्थान्पञ्चवानरपुङ्गवान्} %||3-52-1||

\threelineshloka
{तेषां मध्ये विशालाक्षी कौशेयं कनकप्रभम्}
{उत्तरीयं वरारोहा शुभान्याभरणानि च}
{मुमोच यदि रामाय शंसेयुरिति मैथिली} %||3-52-2||

\twolineshloka
{वस्त्रमुत्सृज्य तन्मध्ये विनिक्षिप्तं सभूषणम्}
{सम्भ्रमात्तु दशग्रीवस्तत्कर्म न च बुद्धिवान्} %||3-52-3||

\twolineshloka
{पिङ्गाक्षास्तां विशालाक्षीं नेत्रैरनिमिषैरिव}
{विक्रोशन्तीं तदा सीतां ददृशुर्वानरर्षभाः} %||3-52-4||

\twolineshloka
{स च पम्पामतिक्रम्य लङ्कामभिमुखः पुरीम्}
{जगाम रुदतीं गृह्य मैथिलीं राक्षसेश्वरः} %||3-52-5||

\twolineshloka
{तां जहार सुसंहृष्टो रावणो मृत्युमात्मनः}
{उत्सङ्गेनैव भुजगीं तीक्ष्णदंष्ट्रां महाविषाम्} %||3-52-6||

\twolineshloka
{वनानि सरितः शैलान्सरांसि च विहायसा}
{स क्षिप्रं समतीयाय शरश्चापादिव च्युतः} %||3-52-7||

\twolineshloka
{तिमिनक्रनिकेतं तु वरुणालयमक्षयम्}
{सरितां शरणं गत्वा समतीयाय सागरम्} %||3-52-8||

\twolineshloka
{सम्भ्रमात्परिवृत्तोर्मी रुद्धमीनमहोरगः}
{वैदेह्यां ह्रियमाणायां बभूव वरुणालयः} %||3-52-9||

\twolineshloka
{अन्तरिक्षगता वाचः ससृजुश्चारणास्तदा}
{एतदन्तो दशग्रीव इति सिद्धास्तदाब्रुवन्} %||3-52-10||

\twolineshloka
{स तु सीतां विवेष्टन्तीमङ्केनादाय रावणः}
{प्रविवेश पुरीं लङ्कां रूपिणीं मृत्युमात्मनः} %||3-52-11||

\twolineshloka
{सोऽभिगम्य पुरीं लङ्कां सुविभक्तमहापथाम्}
{संरूढकक्ष्या बहुलं स्वमन्तःपुरमाविशत्} %||3-52-12||

\twolineshloka
{तत्र तामसितापाङ्गीं शोकमोहपरायणाम्}
{निदधे रावणः सीतां मयो मायामिवासुरीम्} %||3-52-13||

\twolineshloka
{अब्रवीच्च दशग्रीवः पिशाचीर्घोरदर्शनाः}
{यथा नैनां पुमान्स्त्री वा सीतां पश्यत्यसम्मतः} %||3-52-14||

\twolineshloka
{मुक्तामणिसुवर्णानि वस्त्राण्याभरणानि च}
{यद्यदिच्छेत्तदेवास्या देयं मच्छन्दतो यथा} %||3-52-15||

\twolineshloka
{या च वक्ष्यति वैदेहीं वचनं किञ्चिदप्रियम्}
{अज्ञानाद्यदि वा ज्ञानान्न तस्या जीवितं प्रियम्} %||3-52-16||

\threelineshloka
{तथोक्त्वा राक्षसीस्तास्तु राक्षसेन्द्रः प्रतापवान्}
{निष्क्रम्यान्तःपुरात्तस्मात्किं कृत्यमिति चिन्तयन्}
{ददर्शाष्टौ महावीर्यान्राक्षसान्पिशिताशनान्} %||3-52-17||

\twolineshloka
{स तान्दृष्ट्वा महावीर्यो वरदानेन मोहितः}
{उवाचैतानिदं वाक्यं प्रशस्य बलवीर्यतः} %||3-52-18||

\twolineshloka
{नानाप्रहरणाः क्षिप्रमितो गच्छत सत्वराः}
{जनस्थानं हतस्थानं भूतपूर्वं खरालयम्} %||3-52-19||

\twolineshloka
{तत्रोष्यतां जनस्थाने शून्ये निहतराक्षसे}
{पौरुषं बलमाश्रित्य त्रासमुत्सृज्य दूरतः} %||3-52-20||

\twolineshloka
{बलं हि सुमहद्यन्मे जनस्थाने निवेशितम्}
{सदूषणखरं युद्धे हतं तद्रामसायकैः} %||3-52-21||

\twolineshloka
{ततः क्रोधो ममापूर्वो धैर्यस्योपरि वर्धते}
{वैरं च सुमहज्जातं रामं प्रति सुदारुणम्} %||3-52-22||

\twolineshloka
{निर्यातयितुमिच्छामि तच्च वैरमहं रिपोः}
{न हि लप्स्याम्यहं निद्रामहत्वा संयुगे रिपुम्} %||3-52-23||

\twolineshloka
{तं त्विदानीमहं हत्वा खरदूषणघातिनम्}
{रामं शर्मोपलप्स्यामि धनं लब्ध्वेव निर्धनः} %||3-52-24||

\twolineshloka
{जनस्थाने वसद्भिस्तु भवद्भी राममाश्रिता}
{प्रवृत्तिरुपनेतव्या किं करोतीति तत्त्वतः} %||3-52-25||

\twolineshloka
{अप्रमादाच्च गन्तव्यं सर्वैरेव निशाचरैः}
{कर्तव्यश्च सदा यत्नो राघवस्य वधं प्रति} %||3-52-26||

\twolineshloka
{युष्माकं हि बलज्ञोऽहं बहुशो रणमूर्धनि}
{अतश्चास्मिञ्जनस्थाने मया यूयं नियोजिताः} %||3-52-27||

\fourlineindentedshloka
{ततः प्रियं वाक्यमुपेत्य राक्षसा}
{ महार्थमष्टावभिवाद्य रावणम्}
{विहाय लङ्कां सहिताः प्रतस्थिरे}
{ यतो जनस्थानमलक्ष्यदर्शनाः} %||3-52-28||

\fourlineindentedshloka
{ततस्तु सीतामुपलभ्य रावणः}
{ सुसम्प्रहृष्टः परिगृह्य मैथिलीम्}
{प्रसज्य रामेण च वैरमुत्तमम्}
{ बभूव मोहान्मुदितः स राक्षसः} %||3-53-29||


{॥इत्यार्षे श्रीमद्रामायणे वाल्मीकीये आदिकाव्ये अरण्यकाण्डे द्विपञ्चाशत्तमः सर्गः॥}

\sect{त्रिपञ्चाशत्तमः सर्गः}

\twolineshloka
{सन्दिश्य राक्षसान्घोरान्रावणोऽष्टौ महाबलान्}
{आत्मानं बुद्धिवैक्लव्यात्कृतकृत्यममन्यत} %||3-53-1||

\twolineshloka
{स चिन्तयानो वैदेहीं कामबाणसमर्पितः}
{प्रविवेश गृहं रम्यं सीतां द्रष्टुमभित्वरन्} %||3-53-2||

\twolineshloka
{स प्रविश्य तु तद्वेश्म रावणो राक्षसाधिपः}
{अपश्यद्राक्षसीमध्ये सीतां शोकपरायणम्} %||3-53-3||

\twolineshloka
{अश्रुपूर्णमुखीं दीनां शोकभारावपीडिताम्}
{वायुवेगैरिवाक्रान्तां मज्जन्तीं नावमर्णवे} %||3-53-4||

\twolineshloka
{मृगयूथपरिभ्रष्टां मृगीं श्वभिरिवावृताम्}
{अधोमुखमुखीं दीनामभ्येत्य च निशाचरः} %||3-53-5||

\twolineshloka
{तां तु शोकवशां दीनामवशां राक्षसाधिपः}
{स बलाद्दर्शयामास गृहं देवगृहोपमम्} %||3-53-6||

\twolineshloka
{हर्म्यप्रासादसम्बधं स्त्रीसहस्रनिषेवितम्}
{नानापक्षिगणैर्जुष्टं नानारत्नसमन्वितम्} %||3-53-7||

\twolineshloka
{काञ्चनैस्तापनीयैश्च स्फाटिकै राजतैस्तथा}
{वज्रवैदूर्यचित्रैश्च स्तम्भैर्दृष्टिमनोहरैः} %||3-53-8||

\twolineshloka
{दिव्यदुन्दुभिनिर्ह्रादं तप्तकाञ्चनतोरणम्}
{सोपानं काञ्चनं चित्रमारुरोह तया सह} %||3-53-9||

\twolineshloka
{दान्तका राजताश्चैव गवाक्षाः प्रियदर्शनाः}
{हेमजालावृताश्चासंस्तत्र प्रासादपङ्क्तयः} %||3-53-10||

\twolineshloka
{सुधामणिविचित्राणि भूमिभागानि सर्वशः}
{दशग्रीवः स्वभवने प्रादर्शयत मैथिलीम्} %||3-53-11||

\twolineshloka
{दीर्घिकाः पुष्करिण्यश्च नानापुष्पसमावृताः}
{रावणो दर्शयामास सीतां शोकपरायणाम्} %||3-53-12||

\twolineshloka
{दर्शयित्वा तु वैदेहीं कृत्स्नं तद्भवनोत्तमम्}
{उवाच वाक्यं पापात्मा रावणो जनकात्मजाम्} %||3-53-13||

\twolineshloka
{दशराक्षसकोट्यश्च द्वाविंशतिरथापराः}
{वर्जयित्वा जरा वृद्धान्बालांश्च रजनीचरान्} %||3-53-14||

\twolineshloka
{तेषां प्रभुरहं सीते सर्वेषां भीमकर्मणाम्}
{सहस्रमेकमेकस्य मम कार्यपुरःसरम्} %||3-53-15||

\twolineshloka
{यदिदं राज्यतन्त्रं मे त्वयि सर्वं प्रतिष्ठितम्}
{जीवितं च विशालाक्षि त्वं मे प्राणैर्गरीयसी} %||3-53-16||

\twolineshloka
{बहूनां स्त्रीसहस्राणां मम योऽसौ परिग्रहः}
{तासां त्वमीश्वरी सीते मम भार्या भव प्रिये} %||3-53-17||

\twolineshloka
{साधु किं तेऽन्यया बुद्ध्या रोचयस्व वचो मम}
{भजस्व माभितप्तस्य प्रसादं कर्तुमर्हसि} %||3-53-18||

\twolineshloka
{परिक्षिप्ता समुद्रेण लङ्केयं शतयोजना}
{नेयं धर्षयितुं शक्या सेन्द्रैरपि सुरासुरैः} %||3-53-19||

\twolineshloka
{न देवेषु न यक्षेषु न गन्धर्वेषु नर्षिषु}
{अहं पश्यामि लोकेषु यो मे वीर्यसमो भवेत्} %||3-53-20||

\twolineshloka
{राज्यभ्रष्टेन दीनेन तापसेन गतायुषा}
{किं करिष्यसि रामेण मानुषेणाल्पतेजसा} %||3-53-21||

\twolineshloka
{भजस्व सीते मामेव भर्ताहं सदृशस्तव}
{यौवनं ह्यध्रुवं भीरु रमस्वेह मया सह} %||3-53-22||

\twolineshloka
{दर्शने मा कृथा बुद्धिं राघवस्य वरानने}
{कास्य शक्तिरिहागन्तुमपि सीते मनोरथैः} %||3-53-23||

\twolineshloka
{न शक्यो वायुराकाशे पाशैर्बद्धं महाजवः}
{दीप्यमानस्य वाप्यग्नेर्ग्रहीतुं विमलां शिखाम्} %||3-53-24||

\twolineshloka
{त्रयाणामपि लोकानां न तं पश्यामि शोभने}
{विक्रमेण नयेद्यस्त्वां मद्बाहुपरिपालिताम्} %||3-53-25||

\twolineshloka
{लङ्कायां सुमहद्राज्यमिदं त्वमनुपालय}
{अभिषेकोदकक्लिन्ना तुष्टा च रमयस्व माम्} %||3-53-26||

\twolineshloka
{दुष्कृतं यत्पुरा कर्म वनवासेन तद्गतम्}
{यश्च ते सुकृतो धर्मस्तस्येह फलमाप्नुहि} %||3-53-27||

\twolineshloka
{इह सर्वाणि माल्यानि दिव्यगन्धानि मैथिलि}
{भूषणानि च मुख्यानि तानि सेव मया सह} %||3-53-28||

\twolineshloka
{पुष्पकं नाम सुश्रोणि भ्रातुर्वैश्रवणस्य मे}
{विमानं रमणीयं च तद्विमानं मनोजवम्} %||3-53-29||

\twolineshloka
{तत्र सीते मया सार्धं विहरस्व यथासुखम्}
{वदनं पद्मसङ्काशं विमलं चारुदर्शनम्} %||3-53-30||

\twolineshloka
{शोकार्तं तु वरारोहे न भ्राजति वरानने}
{अलं व्रीडेन वैदेहि धर्मलोप कृतेन ते} %||3-53-31||

\twolineshloka
{आर्षोऽयं दैवनिष्यन्दो यस्त्वामभिगमिष्यति}
{एतौ पादौ मया स्निग्धौ शिरोभिः परिपीडितौ} %||3-53-32||

\twolineshloka
{प्रसादं कुरु मे क्षिप्रं वश्यो दासोऽहमस्मि ते}
{नेमाः शून्या मया वाचः शुष्यमाणेन भाषिताः} %||3-53-33||

\twolineshloka
{न चापि रावणः काञ्चिन्मूर्ध्ना स्त्रीं प्रणमेत ह}
{एवमुक्त्वा दशग्रीवो मैथिलीं जनकात्मजाम्} %||3-53-34||

{कृतान्तवशमापन्नो ममेयमिति मन्यते॥\devanumber{35}॥} %||3-54-35|| (Check)

\addtocounter{shlokacount}{1}


{॥इत्यार्षे श्रीमद्रामायणे वाल्मीकीये आदिकाव्ये अरण्यकाण्डे त्रिपञ्चाशत्तमः सर्गः॥}

\sect{चतुःपञ्चाशत्तमः सर्गः}

\twolineshloka
{सा तथोक्ता तु वैदेही निर्भया शोककर्षिता}
{तृणमन्तरतः कृत्वा रावणं प्रत्यभाषत} %||3-54-1||

\twolineshloka
{राजा दशरथो नाम धर्मसेतुरिवाचलः}
{सत्यसन्धः परिज्ञातो यस्य पुत्रः स राघवः} %||3-54-2||

\twolineshloka
{रामो नाम स धर्मात्मा त्रिषु लोकेषु विश्रुतः}
{दीर्घबाहुर्विशालाक्षो दैवतं स पतिर्मम} %||3-54-3||

\twolineshloka
{इक्ष्वाकूणां कुले जातः सिंहस्कन्धो महाद्युतिः}
{लक्ष्मणेन सह भ्रात्रा यस्ते प्राणां हरिष्यति} %||3-54-4||

\twolineshloka
{प्रत्यक्षं यद्यहं तस्य त्वया स्यां धर्षिता बलात्}
{शयिता त्वं हतः सङ्ख्ये जनस्थाने यथा खरः} %||3-54-5||

\twolineshloka
{य एते राक्षसाः प्रोक्ता घोररूपा महाबलाः}
{राघवे निर्विषाः सर्वे सुपर्णे पन्नगा यथा} %||3-54-6||

\twolineshloka
{तस्य ज्याविप्रमुक्तास्ते शराः काञ्चनभूषणाः}
{शरीरं विधमिष्यन्ति गङ्गाकूलमिवोर्मयः} %||3-54-7||

\twolineshloka
{असुरैर्वा सुरैर्वा त्वं यद्यवधोऽसि रावण}
{उत्पाद्य सुमहद्वैरं जीवंस्तस्य न मोक्ष्यसे} %||3-54-8||

\twolineshloka
{स ते जीवितशेषस्य राघवोऽन्तकरो बली}
{पशोर्यूपगतस्येव जीवितं तव दुर्लभम्} %||3-54-9||

\twolineshloka
{यदि पश्येत्स रामस्त्वां रोषदीप्तेन चक्षुषा}
{रक्षस्त्वमद्य निर्दग्धो गच्छेः सद्यः पराभवम्} %||3-54-10||

\twolineshloka
{यश्चन्द्रं नभसो भूमौ पातयेन्नाशयेत वा}
{सागरं शोषयेद्वापि स सीतां मोचयेदिह} %||3-54-11||

\twolineshloka
{गतायुस्त्वं गतश्रीको गतसत्त्वो गतेन्द्रियः}
{लङ्का वैधव्यसंयुक्ता त्वत्कृतेन भविष्यति} %||3-54-12||

\twolineshloka
{न ते पापमिदं कर्म सुखोदर्कं भविष्यति}
{याहं नीता विना भावं पतिपार्श्वात्त्वया वनात्} %||3-54-13||

\twolineshloka
{स हि दैवतसंयुक्तो मम भर्ता महाद्युतिः}
{निर्भयो वीर्यमाश्रित्य शून्ये वसति दण्डके} %||3-54-14||

\twolineshloka
{स ते दर्पं बलं वीर्यमुत्सेकं च तथाविधम्}
{अपनेष्यति गात्रेभ्यः शरवर्षेण संयुगे} %||3-54-15||

\twolineshloka
{यदा विनाशो भूतानां दृश्यते कालचोदितः}
{तदा कार्ये प्रमाद्यन्ति नराः कालवशं गताः} %||3-54-16||

\twolineshloka
{मां प्रधृष्य स ते कालः प्राप्तोऽयं रक्षसाधम}
{आत्मनो राक्षसानां च वधायान्तःपुरस्य च} %||3-54-17||

\twolineshloka
{न शक्या यज्ञमध्यस्था वेदिः स्रुग्भाण्ड मण्डिता}
{द्विजातिमन्त्रसम्पूता चण्डालेनावमर्दितुम्} %||3-54-18||

\threelineshloka
{इदं शरीरं निःसंज्ञं बन्ध वा घातयस्व वा}
{नेदं शरीरं रक्ष्यं मे जीवितं वापि राक्षस}
{न हि शक्ष्याम्युपक्रोशं पृथिव्यां दातुमात्मनः} %||3-54-19||

\twolineshloka
{एवमुक्त्वा तु वैदेही क्रोद्धात्सुपरुषं वचः}
{रावणं मैथिली तत्र पुनर्नोवाच किञ्चन} %||3-54-20||

\twolineshloka
{सीताया वचनं श्रुत्वा परुषं रोमहर्षणम्}
{प्रत्युवाच ततः सीतां भयसन्दर्शनं वचः} %||3-54-21||

\threelineshloka
{शृणु मैथिलि मद्वाक्यं मासान्द्वादश भामिनि}
{कालेनानेन नाभ्येषि यदि मां चारुहासिनि}
{ततस्त्वां प्रातराशार्थं सूदाश्छेत्स्यन्ति लेशशः} %||3-54-22||

\twolineshloka
{इत्युक्त्वा परुषं वाक्यं रावणः शत्रुरावणः}
{राक्षसीश्च ततः क्रुद्ध इदं वचनमब्रवीत्} %||3-54-23||

\twolineshloka
{शीघ्रमेवं हि राक्षस्यो विकृता घोरदर्शनाः}
{दर्पमस्या विनेष्यन्तु मांसशोणितभोजनाः} %||3-54-24||

\twolineshloka
{वचनादेव तास्तस्य विकृता घोरदर्शनाः}
{कृतप्राञ्जलयो भूत्वा मैथिलीं पर्यवारयन्} %||3-54-25||

\twolineshloka
{स ताः प्रोवाच राजा तु रावणो घोरदर्शनः}
{प्रचाल्य चरणोत्कर्षैर्दारयन्निव मेदिनीम्} %||3-54-26||

\twolineshloka
{अशोकवनिकामध्ये मैथिली नीयतामिति}
{तत्रेयं रक्ष्यतां गूढमुष्माभिः परिवारिता} %||3-54-27||

\twolineshloka
{तत्रैनां तर्जनैर्घोरैः पुनः सान्त्वैश्च मैथिलीम्}
{आनयध्वं वशं सर्वा वन्यां गजवधूमिव} %||3-54-28||

\twolineshloka
{इति प्रतिसमादिष्टा राक्षस्यो रावणेन ताः}
{अशोकवनिकां जग्मुर्मैथिलीं परिगृह्य ताम्} %||3-54-29||

\twolineshloka
{सर्वकामफलैर्वृक्षैर्नानापुष्पफलैर्वृताम्}
{सर्वकालमदैश्चापि द्विजैः समुपसेविताम्} %||3-54-30||

\twolineshloka
{सा तु शोकपरीताङ्गी मैथिली जनकात्मजा}
{राक्षसी वशमापन्ना व्याघ्रीणां हरिणी यथा} %||3-54-31||

\fourlineindentedshloka
{न विन्दते तत्र तु शर्म मैथिली}
{ विरूपनेत्राभिरतीव तर्जिता}
{पतिं स्मरन्ती दयितं च देवरम्}
{ विचेतनाभूद्भयशोकपीडिता} %||3-55-32||


{॥इत्यार्षे श्रीमद्रामायणे वाल्मीकीये आदिकाव्ये अरण्यकाण्डे चतुःपञ्चाशत्तमः सर्गः॥}

\sect{पञ्चपञ्चाशत्तमः सर्गः}

\twolineshloka
{राक्षसं मृगरूपेण चरन्तं कामरूपिणम्}
{निहत्य रामो मारीचं तूर्णं पथि न्यवर्तत} %||3-55-1||

\twolineshloka
{तस्य सन्त्वरमाणस्य द्रष्टुकामस्य मैथिलीम्}
{क्रूरस्वरोऽथ गोमायुर्विननादास्य पृष्ठतः} %||3-55-2||

\twolineshloka
{स तस्य स्वरमाज्ञाय दारुणं रोमहर्षणम्}
{चिन्तयामास गोमायोः स्वरेण परिशङ्कितः} %||3-55-3||

\twolineshloka
{अशुभं बत मन्येऽहं गोमायुर्वाश्यते यथा}
{स्वस्ति स्यादपि वैदेह्या राक्षसैर्भक्षणं विना} %||3-55-4||

\twolineshloka
{मारीचेन तु विज्ञाय स्वरमालक्ष्य मामकम्}
{विक्रुष्टं मृगरूपेण लक्ष्मणः शृणुयाद्यदि} %||3-55-5||

\twolineshloka
{स सौमित्रिः स्वरं श्रुत्वा तां च हित्वाथ मैथिलीम्}
{तयैव प्रहितः क्षिप्रं मत्सकाशमिहैष्यति} %||3-55-6||

\twolineshloka
{राक्षसैः सहितैर्नूनं सीताया ईप्सितो वधः}
{काञ्चनश्च मृगो भूत्वा व्यपनीयाश्रमात्तु माम्} %||3-55-7||

\twolineshloka
{दूरं नीत्वा तु मारीचो राक्षसोऽभूच्छराहतः}
{हा लक्ष्मण हतोऽस्मीति यद्वाक्यं व्यजहार ह} %||3-55-8||

\threelineshloka
{अपि स्वस्ति भवेद्द्वाभ्यां रहिताभ्यां मया वने}
{जनस्थाननिमित्तं हि कृतवैरोऽस्मि राक्षसैः}
{निमित्तानि च घोराणि दृश्यन्तेऽद्य बहूनि च} %||3-55-9||

\threelineshloka
{इत्येवं चिन्तयन्रामः श्रुत्वा गोमायुनिःस्वनम्}
{आत्मनश्चापनयनं मृगरूपेण रक्षसा}
{आजगाम जनस्थानं राघवः परिशङ्कितः} %||3-55-10||

\twolineshloka
{तं दीनमानसं दीनमासेदुर्मृगपक्षिणः}
{सव्यं कृत्वा महात्मानं घोरांश्च ससृजुः स्वरान्} %||3-55-11||

\twolineshloka
{तानि दृष्ट्वा निमित्तानि महाघोराणि राघवः}
{ततो लक्षणमायान्तं ददर्श विगतप्रभम्} %||3-55-12||

\twolineshloka
{ततोऽविदूरे रामेण समीयाय स लक्ष्मणः}
{विषण्णः स विषण्णेन दुःखितो दुःखभागिना} %||3-55-13||

\twolineshloka
{संजगर्हेऽथ तं भ्राता जेष्ठो लक्ष्मणमागतम्}
{विहाय सीतां विजने वने राक्षससेविते} %||3-55-14||

\twolineshloka
{गृहीत्वा च करं सव्यं लक्ष्मणं रघुनन्दनः}
{उवाच मधुरोदर्कमिदं परुषमार्तवत्} %||3-55-15||

\twolineshloka
{अहो लक्ष्मण गर्ह्यं ते कृतं यत्त्वं विहाय ताम्}
{सीतामिहागतः सौम्य कच्चित्स्वस्ति भवेदिति} %||3-55-16||

\twolineshloka
{न मेऽस्ति संशयो वीर सर्वथा जनकात्मजा}
{विनष्टा भक्षिता वाप राक्षसैर्वनचारिभिः} %||3-55-17||

\twolineshloka
{अशुभान्येव भूयिष्ठं यथा प्रादुर्भवन्ति मे}
{अपि लक्ष्मण सीतायाः सामग्र्यं प्राप्नुयावहे} %||3-55-18||

\fourlineindentedshloka
{इदं हि रक्षोमृगसंनिकाशम्}
{ प्रलोभ्य मां दूरमनुप्रयातम्}
{हतं कथञ्चिन्महता श्रमेण}
{ स राक्षसोऽभून्म्रियमाण एव} %||3-55-19||

\fourlineindentedshloka
{मनश्च मे दीनमिहाप्रहृष्टम्}
{ चक्षुश्च सव्यं कुरुते विकारम्}
{असंशयं लक्ष्मण नास्ति सीता}
{ हृता मृता वा पथि वर्तते वा} %||3-56-20||


{॥इत्यार्षे श्रीमद्रामायणे वाल्मीकीये आदिकाव्ये अरण्यकाण्डे पञ्चपञ्चाशत्तमः सर्गः॥}

\sect{षट्पञ्चाशत्तमः सर्गः}

\twolineshloka
{स दृष्ट्वा लक्ष्मणं दीनं शून्ये दशरथात्मजः}
{पर्यपृच्छत धर्मात्मा वैदेहीमागतं विना} %||3-56-1||

\twolineshloka
{प्रस्थितं दण्डकारण्यं या मामनुजगाम ह}
{क्व सा लक्ष्मण वैदेही यां हित्वा त्वमिहागतः} %||3-56-2||

\twolineshloka
{राज्यभ्रष्टस्य दीनस्य दण्डकान्परिधावतः}
{क्व सा दुःखसहाया मे वैदेही तनुमध्यमा} %||3-56-3||

\twolineshloka
{यां विना नोत्सहे वीर मुहूर्तमपि जीवितुम्}
{क्व सा प्राणसहाया मे सीता सुरसुतोपमा} %||3-56-4||

\twolineshloka
{पतित्वममराणां वा पृथिव्याश्चापि लक्ष्मण}
{विना तां तपनीयाभां नेच्छेयं जनकात्मजाम्} %||3-56-5||

\twolineshloka
{कच्चिज्जीवति वैदेही प्राणैः प्रियतरा मम}
{कच्चित्प्रव्राजनं सौम्य न मे मिथ्या भविष्यति} %||3-56-6||

\twolineshloka
{सीतानिमित्तं सौमित्रे मृते मयि गते त्वयि}
{कच्चित्सकामा सुखिता कैकेयी सा भविष्यति} %||3-56-7||

\twolineshloka
{सपुत्रराज्यां सिद्धार्थां मृतपुत्रा तपस्विनी}
{उपस्थास्यति कौसल्या कच्चिन्सौम्य न कैकयीम्} %||3-56-8||

\twolineshloka
{यदि जीवति वैदेही गमिष्याम्याश्रमं पुनः}
{सुवृत्ता यदि वृत्ता सा प्राणांस्त्यक्ष्यामि लक्ष्मण} %||3-56-9||

\twolineshloka
{यदि मामाश्रमगतं वैदेही नाभिभाषते}
{पुनः प्रहसिता सीता विनशिष्यामि लक्ष्मण} %||3-56-10||

\twolineshloka
{ब्रूहि लक्ष्मण वैदेही यदि जीवति वा न वा}
{त्वयि प्रमत्ते रक्षोभिर्भक्षिता वा तपस्विनी} %||3-56-11||

\twolineshloka
{सुकुमारी च बाला च नित्यं चादुःखदर्शिनी}
{मद्वियोगेन वैदेही व्यक्तं शोचति दुर्मनाः} %||3-56-12||

\twolineshloka
{सर्वथा रक्षसा तेन जिह्मेन सुदुरात्मना}
{वदता लक्ष्मणेत्युच्चैस्तवापि जनितं भयम्} %||3-56-13||

\twolineshloka
{श्रुतश्च शङ्के वैदेह्या स स्वरः सदृशो मम}
{त्रस्तया प्रेषितस्त्वं च द्रष्टुं मां शीघ्रमागतः} %||3-56-14||

\twolineshloka
{सर्वथा तु कृतं कष्टं सीतामुत्सृजता वने}
{प्रतिकर्तुं नृशंसानां रक्षसां दत्तमन्तरम्} %||3-56-15||

\twolineshloka
{दुःखिताः खरघातेन राक्षसाः पिशिताशनाः}
{तैः सीता निहता घोरैर्भविष्यति न संशयः} %||3-56-16||

\twolineshloka
{अहोऽस्मि व्यसने मग्नः सर्वथा रिपुनाशन}
{किं त्विदानीं करिष्यामि शङ्के प्राप्तव्यमीदृशम्} %||3-56-17||

\twolineshloka
{इति सीतां वरारोहां चिन्तयन्नेव राघवः}
{आजगाम जनस्थानं त्वरया सहलक्ष्मणः} %||3-56-18||

\fourlineindentedshloka
{विगर्हमाणोऽनुजमार्तरूपम्}
{ क्षुधा श्रमाच्चैव पिपासया च}
{विनिःश्वसञ्शुष्कमुखो विषण्णः}
{ प्रतिश्रयं प्राप्य समीक्ष्य शून्यम्} %||3-56-19||

\fourlineindentedshloka
{स्वमाश्रमं सम्प्रविगाह्य वीरो}
{ विहारदेशाननुसृत्य कांश्चित्}
{एतत्तदित्येव निवासभूमौ}
{ प्रहृष्टरोमा व्यथितो बभूव} %||3-57-20||


{॥इत्यार्षे श्रीमद्रामायणे वाल्मीकीये आदिकाव्ये अरण्यकाण्डे षट्पञ्चाशत्तमः सर्गः॥}

\sect{सप्तपञ्चाशत्तमः सर्गः}

\twolineshloka
{अथाश्रमादुपावृत्तमन्तरा रघुनन्दनः}
{परिपप्रच्छ सौमित्रिं रामो दुःखार्दितः पुनः} %||3-57-1||

\twolineshloka
{तमुवाच किमर्थं त्वमागतोऽपास्य मैथिलीम्}
{यदा सा तव विश्वासाद्वने विहरिता मया} %||3-57-2||

\twolineshloka
{दृष्ट्वैवाभ्यागतं त्वां मे मैथिलीं त्यज्य लक्ष्मण}
{शङ्कमानं महत्पापं यत्सत्यं व्यथितं मनः} %||3-57-3||

\twolineshloka
{स्फुरते नयनं सव्यं बाहुश्च हृदयं च मे}
{दृष्ट्वा लक्ष्मण दूरे त्वां सीताविरहितं पथि} %||3-57-4||

\twolineshloka
{एवमुक्तस्तु सौमित्रिर्लक्ष्मणः शुभलक्षणः}
{भूयो दुःखसमाविष्टो दुःखितं राममब्रवीत्} %||3-57-5||

\twolineshloka
{न स्वयं कामकारेण तां त्यक्त्वाहमिहागतः}
{प्रचोदितस्तयैवोग्रैस्त्वत्सकाशमिहागतः} %||3-57-6||

\twolineshloka
{आर्येणेव परिक्रुष्टं हा सीते लक्ष्मणेति च}
{परित्राहीति यद्वाक्यं मैथिल्यास्तच्छ्रुतिं गतम्} %||3-57-7||

\twolineshloka
{सा तमार्तस्वरं श्रुत्वा तव स्नेहेन मैथिली}
{गच्छ गच्छेति मामाह रुदन्ती भयविह्वला} %||3-57-8||

\twolineshloka
{प्रचोद्यमानेन मया गच्छेति बहुशस्तया}
{प्रत्युक्ता मैथिली वाक्यमिदं त्वत्प्रत्ययान्वितम्} %||3-57-9||

\twolineshloka
{न तत्पश्याम्यहं रक्षो यदस्य भयमावहेत्}
{निर्वृता भव नास्त्येतत्केनाप्येवमुदाहृतम्} %||3-57-10||

\twolineshloka
{विगर्हितं च नीचं च कथमार्योऽभिधास्यति}
{त्राहीति वचनं सीते यस्त्रायेत्त्रिदशानपि} %||3-57-11||

\threelineshloka
{किंनिमित्तं तु केनापि भ्रातुरालम्ब्य मे स्वरम्}
{विस्वरं व्याहृतं वाक्यं लक्ष्मण त्राहि मामिति}
{न भवत्या व्यथा कार्या कुनारीजनसेविता} %||3-57-12||

\threelineshloka
{अलं वैक्लव्यमालम्ब्य स्वस्था भव निरुत्सुका}
{न चास्ति त्रिषु लोकेषु पुमान्यो राघवं रणे}
{जातो वा जायमानो वा संयुगे यः पराजयेत्} %||3-57-13||

\twolineshloka
{एवमुक्ता तु वैदेही परिमोहितचेतना}
{उवाचाश्रूणि मुञ्चन्ती दारुणं मामिदं वचः} %||3-57-14||

\twolineshloka
{भावो मयि तवात्यर्थं पाप एव निवेशितः}
{विनष्टे भ्रातरि प्राप्ते न च त्वं मामवाप्स्यसि} %||3-57-15||

\twolineshloka
{सङ्केताद्भरतेन त्वं रामं समनुगच्छसि}
{क्रोशन्तं हि यथात्यर्थं नैनमभ्यवपद्यसे} %||3-57-16||

\twolineshloka
{रिपुः प्रच्छन्नचारी त्वं मदर्थमनुगच्छसि}
{राघवस्यान्तरप्रेप्सुस्तथैनं नाभिपद्यसे} %||3-57-17||

\twolineshloka
{एवमुक्तो हि वैदेह्या संरब्धो रक्तलोचनः}
{क्रोधात्प्रस्फुरमाणौष्ठ आश्रमादभिनिर्गतः} %||3-57-18||

\twolineshloka
{एवं ब्रुवाणं सौमित्रिं रामः सन्तापमोहितः}
{अब्रवीद्दुष्कृतं सौम्य तां विना यत्त्वमागतः} %||3-57-19||

\twolineshloka
{जानन्नपि समर्थं मां रक्षसां विनिवारणे}
{अनेन क्रोधवाक्येन मैथिल्या निःसृतो भवान्} %||3-57-20||

\twolineshloka
{न हि ते परितुष्यामि त्यक्त्वा यद्यासि मैथिलीम्}
{क्रुद्धायाः परुषं श्रुत्वा स्त्रिया यत्त्वमिहागतः} %||3-57-21||

\twolineshloka
{सर्वथा त्वपनीतं ते सीतया यत्प्रचोदितः}
{क्रोधस्य वशमागम्य नाकरोः शासनं मम} %||3-57-22||

\twolineshloka
{असौ हि राक्षसः शेते शरेणाभिहतो मया}
{मृगरूपेण येनाहमाश्रमादपवादितः} %||3-57-23||

\fourlineindentedshloka
{विकृष्य चापं परिधाय सायकम्}
{ सलील बाणेन च ताडितो मया}
{मार्गीं तनुं त्यज्य च विक्लवस्वरो}
{ बभूव केयूरधरः स राक्षसः} %||3-57-24||

\fourlineindentedshloka
{शराहतेनैव तदार्तया गिरा}
{ स्वरं ममालम्ब्य सुदूरसंश्रवम्}
{उदाहृतं तद्वचनं सुदारुणम्}
{ त्वमागतो येन विहाय मैथिलीम्} %||3-58-25||


{॥इत्यार्षे श्रीमद्रामायणे वाल्मीकीये आदिकाव्ये अरण्यकाण्डे सप्तपञ्चाशत्तमः सर्गः॥}

\sect{अष्टपञ्चाशत्तमः सर्गः}

\twolineshloka
{भृशमाव्रजमानस्य तस्याधोवामलोचनम्}
{प्रास्फुरच्चास्खलद्रामो वेपथुश्चास्य जायते} %||3-58-1||

\twolineshloka
{उपालक्ष्य निमित्तानि सोऽशुभानि मुहुर्मुहुः}
{अपि क्षेमं तु सीताया इति वै व्याजहार ह} %||3-58-2||

\twolineshloka
{त्वरमाणो जगामाथ सीतादर्शनलालसः}
{शून्यमावसथं दृष्ट्वा बभूवोद्विग्नमानसः} %||3-58-3||

\twolineshloka
{उद्भ्रमन्निव वेगेन विक्षिपन्रघुनन्दनः}
{तत्र तत्रोटजस्थानमभिवीक्ष्य समन्ततः} %||3-58-4||

\twolineshloka
{ददर्श पर्णशालां च रहितां सीतया तदा}
{श्रिया विरहितां ध्वस्तां हेमन्ते पद्मिनीमिव} %||3-58-5||

\twolineshloka
{रुदन्तमिव वृक्षैश्च म्लानपुष्पमृगद्विजम्}
{श्रिया विहीनं विध्वस्तं सन्त्यक्तवनदैवतम्} %||3-58-6||

\twolineshloka
{विप्रकीर्णाजिनकुशं विप्रविद्धबृसीकटम्}
{दृष्ट्वा शून्योटजस्थानं विललाप पुनः पुनः} %||3-58-7||

\twolineshloka
{हृता मृता वा नष्टा वा भक्षिता वा भविष्यति}
{निलीनाप्यथ वा भीरुरथ वा वनमाश्रिता} %||3-58-8||

\twolineshloka
{गता विचेतुं पुष्पाणि फलान्यपि च वा पुनः}
{अथ वा पद्मिनीं याता जलार्थं वा नदीं गता} %||3-58-9||

\twolineshloka
{यत्नान्मृगयमाणस्तु नाससाद वने प्रियाम्}
{शोकरक्तेक्षणः शोकादुन्मत्त इव लक्ष्यते} %||3-58-10||

\twolineshloka
{वृक्षाद्वृक्षं प्रधावन्स गिरींश्चापि नदीन्नदीम्}
{बभूव विलपन्रामः शोकपङ्कार्णवप्लुतः} %||3-58-11||

\twolineshloka
{अस्ति कच्चित्त्वया दृष्टा सा कदम्बप्रिया प्रिया}
{कदम्ब यदि जानीषे शंस सीतां शुभाननाम्} %||3-58-12||

\twolineshloka
{स्निग्धपल्लवसङ्काशां पीतकौशेयवासिनीम्}
{शंसस्व यदि वा दृष्टा बिल्व बिल्वोपमस्तनी} %||3-58-13||

\twolineshloka
{अथ वार्जुन शंस त्वं प्रियां तामर्जुनप्रियाम्}
{जनकस्य सुता भीरुर्यदि जीवति वा न वा} %||3-58-14||

\twolineshloka
{ककुभः ककुभोरुं तां व्यक्तं जानाति मैथिलीम्}
{लतापल्लवपुष्पाढ्यो भाति ह्येष वनस्पतिः} %||3-58-15||

\twolineshloka
{भ्रमरैरुपगीतश्च यथा द्रुमवरो ह्ययम्}
{एष व्यक्तं विजानाति तिलकस्तिलकप्रियाम्} %||3-58-16||

\twolineshloka
{अशोकशोकापनुद शोकोपहतचेतसम्}
{त्वन्नामानं कुरु क्षिप्रं प्रियासन्दर्शनेन माम्} %||3-58-17||

\twolineshloka
{यदि ताल त्वया दृष्टा पक्वतालफलस्तनी}
{कथयस्व वरारोहां कारुष्यं यदि ते मयि} %||3-58-18||

\twolineshloka
{यदि दृष्टा त्वया सीता जम्बुजाम्बूनदप्रभा}
{प्रियां यदि विजानीषे निःशङ्कं कथयस्व मे} %||3-58-19||

\twolineshloka
{अथ वा मृगशावाक्षीं मृग जानासि मैथिलीम्}
{मृगविप्रेक्षणी कान्ता मृगीभिः सहिता भवेत्} %||3-58-20||

\twolineshloka
{गज सा गजनासोरुर्यदि दृष्टा त्वया भवेत्}
{तां मन्ये विदितां तुभ्यमाख्याहि वरवारण} %||3-58-21||

\twolineshloka
{शार्दूल यदि सा दृष्टा प्रिया चन्द्रनिभानना}
{मैथिली मम विस्रब्धः कथयस्व न ते भयम्} %||3-58-22||

\twolineshloka
{किं धावसि प्रिये नूनं दृष्टासि कमलेक्षणे}
{वृक्षेणाच्छाद्य चात्मानं किं मां न प्रतिभाषसे} %||3-58-23||

\twolineshloka
{तिष्ठ तिष्ठ वरारोहे न तेऽस्ति करुणा मयि}
{नात्यर्थं हास्यशीलासि किमर्थं मामुपेक्षसे} %||3-58-24||

\twolineshloka
{पीतकौशेयकेनासि सूचिता वरवर्णिनि}
{धावन्त्यपि मया दृष्टा तिष्ठ यद्यस्ति सौहृदम्} %||3-58-25||

\twolineshloka
{नैव सा नूनमथ वा हिंसिता चारुहासिनी}
{कृच्छ्रं प्राप्तं हि मां नूनं यथोपेक्षितुमर्हति} %||3-58-26||

\twolineshloka
{व्यक्तं सा भक्षिता बाला राक्षसैः पिशिताशनैः}
{विभज्याङ्गानि सर्वाणि मया विरहिता प्रिया} %||3-58-27||

\twolineshloka
{नूनं तच्छुभदन्तौष्ठं मुखं निष्प्रभतां गतम्}
{सा हि चम्पकवर्णाभा ग्रीवा ग्रैवेय शोभिता} %||3-58-28||

\twolineshloka
{कोमला विलपन्त्यास्तु कान्ताया भक्षिता शुभा}
{नूनं विक्षिप्यमाणौ तौ बाहू पल्लवकोमलौ} %||3-58-29||

\twolineshloka
{भक्षितौ वेपमानाग्रौ सहस्ताभरणाङ्गदौ}
{मया विरहिता बाला रक्षसां भक्षणाय वै} %||3-58-30||

\twolineshloka
{सार्थेनेव परित्यक्ता भक्षिता बहुबान्धवा}
{हा लक्ष्मण महाबाहो पश्यसि त्वं प्रियां क्वचित्} %||3-58-31||

\twolineshloka
{हा प्रिये क्व गता भद्रे हा सीतेति पुनः पुनः}
{इत्येवं विलपन्रामः परिधावन्वनाद्वनम्} %||3-58-32||

\twolineshloka
{क्वचिदुद्भ्रमते वेगात्क्वचिद्विभ्रमते बलात्}
{क्वचिन्मत्त इवाभाति कान्तान्वेषणतत्परः} %||3-58-33||

\twolineshloka
{स वनानि नदीः शैलान्गिरिप्रस्रवणानि च}
{काननानि च वेगेन भ्रमत्यपरिसंस्थितः} %||3-58-34||

\fourlineindentedshloka
{तथा स गत्वा विपुलं महद्वनम्}
{ परीत्य सर्वं त्वथ मैथिलीं प्रति}
{अनिष्ठिताशः स चकार मार्गणे}
{ पुनः प्रियायाः परमं परिश्रमम्} %||3-59-35||


{॥इत्यार्षे श्रीमद्रामायणे वाल्मीकीये आदिकाव्ये अरण्यकाण्डे अष्टपञ्चाशत्तमः सर्गः॥}

\sect{एकोनषष्टितमः सर्गः}

\twolineshloka
{दृष्टाश्रमपदं शून्यं रामो दशरथात्मजः}
{रहितां पर्णशालां च विध्वस्तान्यासनानि च} %||3-59-1||

\twolineshloka
{अदृष्ट्वा तत्र वैदेहीं संनिरीक्ष्य च सर्वशः}
{उवाच रामः प्राक्रुश्य प्रगृह्य रुचिरौ भुजौ} %||3-59-2||

\twolineshloka
{क्व नु लक्ष्मण वैदेही कं वा देशमितो गता}
{केनाहृता वा सौमित्रे भक्षिता केन वा प्रिया} %||3-59-3||

\twolineshloka
{वृष्केणावार्य यदि मां सीते हसितुमिच्छसि}
{अलं ते हसितेनाद्य मां भजस्व सुदुःखितम्} %||3-59-4||

\twolineshloka
{यैः सह क्रीडसे सीते विश्वस्तैर्मृगपोतकैः}
{एते हीनास्त्वया सौम्ये ध्यायन्त्यस्राविलेक्षणाः} %||3-59-5||

\twolineshloka
{मृतं शोकेन महता सीताहरणजेन माम्}
{परलोके महाराजो नूनं द्रक्ष्यति मे पिता} %||3-59-6||

\twolineshloka
{कथं प्रतिज्ञां संश्रुत्य मया त्वमभियोजितः}
{अपूरयित्वा तं कालं मत्सकाशमिहागतः} %||3-59-7||

\twolineshloka
{कामवृत्तमनार्यं मां मृषावादिनमेव च}
{धिक्त्वामिति परे लोके व्यक्तं वक्ष्यति मे पिता} %||3-59-8||

\twolineshloka
{विवशं शोकसन्तप्तं दीनं भग्नमनोरथम्}
{मामिहोत्सृज्य करुणं कीर्तिर्नरमिवानृजुम्} %||3-59-9||

\twolineshloka
{क्व गच्छसि वरारोहे मामुत्सृज्य सुमध्यमे}
{त्वया विरहितश्चाहं मोक्ष्ये जीवितमात्मनः} %||3-59-10||

\twolineshloka
{इतीव विलपन्रामः सीतादर्शनलालसः}
{न ददर्श सुदुःखार्तो राघवो जनकात्मजाम्} %||3-59-11||

\threelineshloka
{अनासादयमानं तं सीतां दशरथात्मजम्}
{पङ्कमासाद्य विपुलं सीदन्तमिव कुञ्जरम्}
{लक्ष्मणो राममत्यर्थमुवाच हितकाम्यया} %||3-59-12||

\twolineshloka
{मा विषादं महाबाहो कुरु यत्नं मया सह}
{इदं च हि वनं शूर बहुकन्दरशोभितम्} %||3-59-13||

\twolineshloka
{प्रियकाननसञ्चारा वनोन्मत्ता च मैथिली}
{सा वनं वा प्रविष्टा स्यान्नलिनीं वा सुपुष्पिताम्} %||3-59-14||

\threelineshloka
{सरितं वापि सम्प्राप्ता मीनवञ्जुरसेविताम्}
{वित्रासयितुकामा वा लीना स्यात्कानने क्वचित्}
{जिज्ञासमाना वैदेही त्वां मां च पुरुषर्षभ} %||3-59-15||

\threelineshloka
{तस्या ह्यन्वेषणे श्रीमन्क्षिप्रमेव यतावहे}
{वनं सर्वं विचिनुवो यत्र सा जनकात्मजा}
{मन्यसे यदि काकुत्स्थ मा स्म शोके मनः कृथाः} %||3-59-16||

\threelineshloka
{एवमुक्तस्तु सौहार्दाल्लक्ष्मणेन समाहितः}
{सह सौमित्रिणा रामो विचेतुमुपचक्रमे}
{तौ वनानि गिरींश्चैव सरितश्च सरांसि च} %||3-59-17||

\twolineshloka
{निखिलेन विचिन्वन्तौ सीतां दशरथात्मजौ}
{तस्य शैलस्य सानूनि गुहाश्च शिखराणि च} %||3-59-18||

\twolineshloka
{निखिलेन विचिन्वन्तौ नैव तामभिजग्मतुः}
{विचित्य सर्वतः शैलं रामो लक्ष्मणमब्रवीत्} %||3-59-19||

\twolineshloka
{नेह पश्यामि सौमित्रे वैदेहीं पर्वते शुभे}
{ततो दुःखाभिसन्तप्तो लक्ष्मणो वाक्यमब्रवीत्} %||3-59-20||

\twolineshloka
{विचरन्दण्डकारण्यं भ्रातरं दीप्ततेजसम्}
{प्राप्स्यसि त्वं महाप्राज्ञ मैथिलीं जनकात्मजाम्} %||3-59-21||

\twolineshloka
{यथा विष्णुर्महाबाहुर्बलिं बद्ध्वा महीमिमाम्}
{एवमुक्तस्तु वीरेण लक्ष्मणेन स राघवः} %||3-59-22||

\twolineshloka
{उवाच दीनया वाचा दुःखाभिहतचेतनः}
{वनं सर्वं सुविचितं पद्मिन्यः फुल्लपङ्कजाः} %||3-59-23||

\twolineshloka
{गिरिश्चायं महाप्राज्ञ बहुकन्दरनिर्झरः}
{न हि पश्यामि वैदेहीं प्राणेभ्योऽपि गरीयसीम्} %||3-59-24||

\twolineshloka
{एवं स विलपन्रामः सीताहरणकर्शितः}
{दीनः शोकसमाविष्टो मुहूर्तं विह्वलोऽभवत्} %||3-59-25||

\twolineshloka
{स विह्वलितसर्वाङ्गो गतबुद्धिर्विचेतनः}
{विषसादातुरो दीनो निःश्वस्याशीतमायतम्} %||3-59-26||

\twolineshloka
{बहुशः स तु निःश्वस्य रामो राजीवलोचनः}
{हा प्रियेति विचुक्रोश बहुशो बाष्पगद्गदः} %||3-59-27||

\twolineshloka
{तं सान्त्वयामास ततो लक्ष्मणः प्रियबान्धवः}
{बहुप्रकारं धर्मज्ञः प्रश्रितः प्रश्रिताञ्जलिः} %||3-59-28||

\twolineshloka
{अनादृत्य तु तद्वाक्यं लक्ष्मणौष्ठपुटच्युतम्}
{अपश्यंस्तां प्रियां सीतां प्राक्रोशत्स पुनः पुनः} %||3-60-29||


{॥इत्यार्षे श्रीमद्रामायणे वाल्मीकीये आदिकाव्ये अरण्यकाण्डे एकोनषष्टितमः सर्गः॥}

\sect{षष्टितमः सर्गः}

\threelineshloka
{स दीनो दीनया वाचा लक्ष्मणं वाक्यमब्रवीत्}
{शीघ्रं लक्ष्मण जानीहि गत्वा गोदावरीं नदीम्}
{अपि गोदावरीं सीता पद्मान्यानयितुं गता} %||3-60-1||

\twolineshloka
{एवमुक्तस्तु रामेण लक्ष्मणः पुनरेव हि}
{नदीं गोदावरीं रम्यां जगाम लघुविक्रमः} %||3-60-2||

\twolineshloka
{तां लक्ष्मणस्तीर्थवतीं विचित्वा राममब्रवीत्}
{नैनां पश्यामि तीर्थेषु क्रोशतो न शृणोति मे} %||3-60-3||

\twolineshloka
{कं नु सा देशमापन्ना वैदेही क्लेशनाशिनी}
{न हि तं वेद्मि वै राम यत्र सा तनुमध्यमा} %||3-60-4||

\twolineshloka
{लक्ष्मणस्य वचः श्रुत्वा दीनः सन्ताप मोहितः}
{रामः समभिचक्राम स्वयं गोदावरीं नदीम्} %||3-60-5||

{स तामुपस्थितो रामः क्व सीतेत्येवमब्रवीत्॥\devanumber{6}॥} %||3-60-6|| (Check)

\addtocounter{shlokacount}{1}

\twolineshloka
{भूतानि राक्षसेन्द्रेण वधार्हेण हृतामपि}
{न तां शशंसू रामाय तथा गोदावरी नदी} %||3-60-7||

\twolineshloka
{ततः प्रचोदिता भूतैः शंसास्मै तां प्रियामिति}
{न च साभ्यवदत्सीतां पृष्टा रामेण शोचिता} %||3-60-8||

\twolineshloka
{रावणस्य च तद्रूपं कर्माणि च दुरात्मनः}
{ध्यात्वा भयात्तु वैदेहीं सा नदी न शशंस ताम्} %||3-60-9||

\twolineshloka
{निराशस्तु तया नद्या सीताया दर्शने कृतः}
{उवाच रामः सौमित्रिं सीतादर्शनकर्शितः} %||3-60-10||

\twolineshloka
{किं नु लक्ष्मण वक्ष्यामि समेत्य जनकं वचः}
{मातरं चैव वैदेह्या विना तामहमप्रियम्} %||3-60-11||

\twolineshloka
{या मे राज्यविहीनस्य वने वन्येन जीवतः}
{सर्वं व्यपनयच्छोकं वैदेही क्व नु सा गता} %||3-60-12||

\twolineshloka
{ज्ञातिपक्षविहीनस्य राजपुत्रीमपश्यतः}
{मन्ये दीर्घा भविष्यन्ति रात्रयो मम जाग्रतः} %||3-60-13||

\twolineshloka
{गोदावरीं जनस्थानमिमं प्रस्रवणं गिरिम्}
{सर्वाण्यनुचरिष्यामि यदि सीता हि दृश्यते} %||3-60-14||

\twolineshloka
{एवं सम्भाषमाणौ तावन्योन्यं भ्रातरावुभौ}
{वसुन्धरायां पतितं पुष्पमार्गमपश्यताम्} %||3-60-15||

\twolineshloka
{तां पुष्पवृष्टिं पतितां दृष्ट्वा रामो महीतले}
{उवाच लक्ष्मणं वीरो दुःखितो दुःखितं वचः} %||3-60-16||

\twolineshloka
{अभिजानामि पुष्पाणि तानीमामीह लक्ष्मण}
{अपिनद्धानि वैदेह्या मया दत्तानि कानने} %||3-60-17||

\twolineshloka
{एवमुक्त्वा महाबाहुर्लक्ष्मणं पुरुषर्षभम्}
{क्रुद्धोऽब्रवीद्गिरिं तत्र सिंहः क्षुद्रमृगं यथा} %||3-60-18||

\twolineshloka
{तां हेमवर्णां हेमाभां सीतां दर्शय पर्वत}
{यावत्सानूनि सर्वाणि न ते विध्वंसयाम्यहम्} %||3-60-19||

\twolineshloka
{मम बाणाग्निनिर्दग्धो भस्मीभूतो भविष्यसि}
{असेव्यः सततं चैव निस्तृणद्रुमपल्लवः} %||3-60-20||

\twolineshloka
{इमां वा सरितं चाद्य शोषयिष्यामि लक्ष्मण}
{यदि नाख्याति मे सीतामद्य चन्द्रनिभाननाम्} %||3-60-21||

\twolineshloka
{एवं स रुषितो रामो दिधक्षन्निव चक्षुषा}
{ददर्श भूमौ निष्क्रान्तं राक्षसस्य पदं महत्} %||3-60-22||

\twolineshloka
{स समीक्ष्य परिक्रान्तं सीताया राक्षसस्य च}
{सम्भ्रान्तहृदयो रामः शशंस भ्रातरं प्रियम्} %||3-60-23||

\twolineshloka
{पश्य लक्ष्मण वैदेह्याः शीर्णाः कनकबिन्दवः}
{भूषणानां हि सौमित्रे माल्यानि विविधानि च} %||3-60-24||

\twolineshloka
{तप्तबिन्दुनिकाशैश्च चित्रैः क्षतजबिन्दुभिः}
{आवृतं पश्य सौमित्रे सर्वतो धरणीतलम्} %||3-60-25||

\twolineshloka
{मन्ये लक्ष्मण वैदेही राक्षसैः कामरूपिभिः}
{भित्त्वा भित्त्वा विभक्ता वा भक्षिता वा भविष्यति} %||3-60-26||

\twolineshloka
{तस्य निमित्तं वैदेह्या द्वयोर्विवदमानयोः}
{बभूव युद्धं सौमित्रे घोरं राक्षसयोरिह} %||3-60-27||

\twolineshloka
{मुक्तामणिचितं चेदं तपनीयविभूषितम्}
{धरण्यां पतितं सौम्य कस्य भग्नं महद्धनुः} %||3-60-28||

\twolineshloka
{तरुणादित्यसङ्काशं वैदूर्यगुलिकाचितम्}
{विशीर्णं पतितं भूमौ कवचं कस्य काञ्चनम्} %||3-60-29||

\twolineshloka
{छत्रं शतशलाकं च दिव्यमाल्योपशोभितम्}
{भग्नदण्डमिदं कस्य भूमौ सौम्य निपातितम्} %||3-60-30||

\twolineshloka
{काञ्चनोरश्छदाश्चेमे पिशाचवदनाः खराः}
{भीमरूपा महाकायाः कस्य वा निहता रणे} %||3-60-31||

\twolineshloka
{दीप्तपावकसङ्काशो द्युतिमान्समरध्वजः}
{अपविद्धश्च भग्नश्च कस्य साङ्ग्रामिको रथः} %||3-60-32||

\twolineshloka
{रथाक्षमात्रा विशिखास्तपनीयविभूषणाः}
{कस्येमेऽभिहता बाणाः प्रकीर्णा घोरकर्मणः} %||3-60-33||

\twolineshloka
{वैरं शतगुणं पश्य ममेदं जीवितान्तकम्}
{सुघोरहृदयैः सौम्य राक्षसैः कामरूपिभिः} %||3-60-34||

\twolineshloka
{हृता मृता वा सीता हि भक्षिता वा तपस्विनी}
{न धर्मस्त्रायते सीतां ह्रियमाणां महावने} %||3-60-35||

\twolineshloka
{भक्षितायां हि वैदेह्यां हृतायामपि लक्ष्मण}
{के हि लोके प्रियं कर्तुं शक्ताः सौम्य ममेश्वराः} %||3-60-36||

\twolineshloka
{कर्तारमपि लोकानां शूरं करुणवेदिनम्}
{अज्ञानादवमन्येरन्सर्वभूतानि लक्ष्मण} %||3-60-37||

\twolineshloka
{मृदुं लोकहिते युक्तं दान्तं करुणवेदिनम्}
{निर्वीर्य इति मन्यन्ते नूनं मां त्रिदशेश्वराः} %||3-60-38||

\threelineshloka
{मां प्राप्य हि गुणो दोषः संवृत्तः पश्य लक्ष्मण}
{अद्यैव सर्वभूतानां रक्षसामभवाय च}
{संहृत्यैव शशिज्योत्स्नां महान्सूर्य इवोदितः} %||3-60-39||

\twolineshloka
{नैव यक्षा न गन्धर्वा न पिशाचा न राक्षसाः}
{किंनरा वा मनुष्या वा सुखं प्राप्स्यन्ति लक्ष्मण} %||3-60-40||

\twolineshloka
{ममास्त्रबाणसम्पूर्णमाकाशं पश्य लक्ष्मण}
{निःसम्पातं करिष्यामि ह्यद्य त्रैलोक्यचारिणाम्} %||3-60-41||

\twolineshloka
{संनिरुद्धग्रहगणमावारितनिशाकरम्}
{विप्रनष्टानलमरुद्भास्करद्युतिसंवृतम्} %||3-60-42||

\twolineshloka
{विनिर्मथितशैलाग्रं शुष्यमाणजलाशयम्}
{ध्वस्तद्रुमलतागुल्मं विप्रणाशितसागरम्} %||3-60-43||

\twolineshloka
{न तां कुशलिनीं सीतां प्रदास्यन्ति ममेश्वराः}
{अस्मिन्मुहूर्ते सौमित्रे मम द्रक्ष्यन्ति विक्रमम्} %||3-60-44||

\twolineshloka
{नाकाशमुत्पतिष्यन्ति सर्वभूतानि लक्ष्मण}
{मम चापगुणान्मुक्तैर्बाणजालैर्निरन्तरम्} %||3-60-45||

\twolineshloka
{अर्दितं मम नाराचैर्ध्वस्तभ्रान्तमृगद्विजम्}
{समाकुलममर्यादं जगत्पश्याद्य लक्ष्मण} %||3-60-46||

\twolineshloka
{आकर्णपूर्णैरिषुभिर्जीवलोकं दुरावरैः}
{करिष्ये मैथिलीहेतोरपिशाचमराक्षसम्} %||3-60-47||

\twolineshloka
{मम रोषप्रयुक्तानां सायकानां बलं सुराः}
{द्रक्ष्यन्त्यद्य विमुक्तानाममर्षाद्दूरगामिनाम्} %||3-60-48||

\twolineshloka
{नैव देवा न दैतेया न पिशाचा न राक्षसाः}
{भविष्यन्ति मम क्रोधात्त्रैलोक्ये विप्रणाशिते} %||3-60-49||

\threelineshloka
{देवदानवयक्षाणां लोका ये रक्षसामपि}
{बहुधा निपतिष्यन्ति बाणौघैः शकुलीकृताः}
{निर्मर्यादानिमाँल्लोकान्करिष्याम्यद्य सायकैः} %||3-60-50||

\threelineshloka
{यथा जरा यथा मृत्युर्यथाकालो यथाविधिः}
{नित्यं न प्रतिहन्यन्ते सर्वभूतेषु लक्ष्मण}
{तथाहं क्रोधसंयुक्तो न निवार्योऽस्म्यसंशयम्} %||3-60-51||

\fourlineindentedshloka
{पुरेव मे चारुदतीमनिन्दिताम्}
{ दिशन्ति सीतां यदि नाद्य मैथिलीम्}
{सदेवगन्धर्वमनुष्य पन्नगम्}
{ जगत्सशैलं परिवर्तयाम्यहम्} %||3-61-52||


{॥इत्यार्षे श्रीमद्रामायणे वाल्मीकीये आदिकाव्ये अरण्यकाण्डे षष्टितमः सर्गः॥}

\sect{एकषष्टितमः सर्गः}

\twolineshloka
{तप्यमानं तथा रामं सीताहरणकर्शितम्}
{लोकानामभवे युक्तं साम्वर्तकमिवानलम्} %||3-61-1||

\twolineshloka
{वीक्षमाणं धनुः सज्यं निःश्वसन्तं मुहुर्मुहुः}
{हन्तुकामं पशुं रुद्रं क्रुद्धं दक्षक्रतौ यथा} %||3-61-2||

\twolineshloka
{अदृष्टपूर्वं सङ्क्रुद्धं दृष्ट्वा रामं स लक्ष्मणः}
{अब्रवीत्प्राञ्जलिर्वाक्यं मुखेन परिशुष्यता} %||3-61-3||

\twolineshloka
{पुरा भूत्वा मृदुर्दान्तः सर्वभूतहिते रतः}
{न क्रोधवशमापन्नः प्रकृतिं हातुमर्हसि} %||3-61-4||

\twolineshloka
{चन्द्रे लक्ष्णीः प्रभा सूर्ये गतिर्वायौ भुवि क्षमा}
{एतच्च नियतं सर्वं त्वयि चानुत्तमं यशः} %||3-61-5||

\twolineshloka
{न तु जानामि कस्यायं भग्नः साङ्ग्रामिको रथः}
{केन वा कस्य वा हेतोः सायुधः सपरिच्छदः} %||3-61-6||

\twolineshloka
{खुरनेमिक्षतश्चायं सिक्तो रुधिरबिन्दुभिः}
{देशो निवृत्तसङ्ग्रामः सुघोरः पार्थिवात्मज} %||3-61-7||

\twolineshloka
{एकस्य तु विमर्दोऽयं न द्वयोर्वदतां वर}
{न हि वृत्तं हि पश्यामि बलस्य महतः पदम्} %||3-61-8||

\twolineshloka
{नैकस्य तु कृते लोकान्विनाशयितुमर्हसि}
{युक्तदण्डा हि मृदवः प्रशान्ता वसुधाधिपाः} %||3-61-9||

\twolineshloka
{सदा त्वं सर्वभूतानां शरण्यः परमा गतिः}
{को नु दारप्रणाशं ते साधु मन्येत राघव} %||3-61-10||

\twolineshloka
{सरितः सागराः शैला देवगन्धर्वदानवाः}
{नालं ते विप्रियं कर्तुं दीक्षितस्येव साधवः} %||3-61-11||

\twolineshloka
{येन राजन्हृता सीता तमन्वेषितुमर्हसि}
{मद्द्वितीयो धनुष्पाणिः सहायैः परमर्षिभिः} %||3-61-12||

\twolineshloka
{समुद्रं च विचेष्यामः पर्वतांश्च वनानि च}
{गुहाश्च विविधा घोरा नलिनीः पार्वतीश्च ह} %||3-61-13||

\twolineshloka
{देवगन्धर्वलोकांश्च विचेष्यामः समाहिताः}
{यावन्नाधिगमिष्यामस्तव भार्यापहारिणम्} %||3-61-14||

\twolineshloka
{न चेत्साम्ना प्रदास्यन्ति पत्नीं ते त्रिदशेश्वराः}
{कोसलेन्द्र ततः पश्चात्प्राप्तकालं करिष्यसि} %||3-61-15||

\fourlineindentedshloka
{शीलेन साम्ना विनयेन सीताम्}
{ नयेन न प्राप्स्यसि चेन्नरेन्द्र}
{ततः समुत्सादय हेमपुङ्खैर्-}
{ महेन्द्रवज्रप्रतिमैः शरौघैः} %||3-62-16||


{॥इत्यार्षे श्रीमद्रामायणे वाल्मीकीये आदिकाव्ये अरण्यकाण्डे एकषष्टितमः सर्गः॥}

\sect{द्विषष्टितमः सर्गः}

\twolineshloka
{तं तथा शोकसन्तप्तं विलपन्तमनाथवत्}
{मोहेन महताविष्टं परिद्यूनमचेतनम्} %||3-62-1||

\twolineshloka
{ततः सौमित्रिराश्वास्य मुहूर्तादिव लक्ष्मणः}
{रामं सम्बोधयामास चरणौ चाभिपीडयन्} %||3-62-2||

\twolineshloka
{महता तपसा राम महता चापि कर्मणा}
{राज्ञा दशरथेनासील्लब्धोऽमृतमिवामरैः} %||3-62-3||

\twolineshloka
{तव चैव गुणैर्बद्धस्त्वद्वियोगान्महीपतिः}
{राजा देवत्वमापन्नो भरतस्य यथा श्रुतम्} %||3-62-4||

\twolineshloka
{यदि दुःखमिदं प्राप्तं काकुत्स्थ न सहिष्यसे}
{प्राकृतश्चाल्पसत्त्वश्च इतरः कः सहिष्यति} %||3-62-5||

\twolineshloka
{दुःखितो हि भवाँल्लोकांस्तेजसा यदि धक्ष्यते}
{आर्ताः प्रजा नरव्याघ्र क्व नु यास्यन्ति निर्वृतिम्} %||3-62-6||

\twolineshloka
{लोकस्वभाव एवैष ययातिर्नहुषात्मजः}
{गतः शक्रेण सालोक्यमनयस्तं समस्पृशत्} %||3-62-7||

\twolineshloka
{महर्षयो वसिष्ठस्तु यः पितुर्नः पुरोहितः}
{अह्ना पुत्रशतं जज्ञे तथैवास्य पुनर्हतम्} %||3-62-8||

\twolineshloka
{या चेयं जगतो माता देवी लोकनमस्कृता}
{अस्याश्च चलनं भूमेर्दृश्यते सत्यसंश्रव} %||3-62-9||

\twolineshloka
{यौ चेमौ जगतां नेत्रे यत्र सर्वं प्रतिष्ठितम्}
{आदित्यचन्द्रौ ग्रहणमभ्युपेतौ महाबलौ} %||3-62-10||

\twolineshloka
{सुमहान्त्यपि भूतानि देवाश्च पुरुषर्षभ}
{न दैवस्य प्रमुञ्चन्ति सर्वभूतानि देहिनः} %||3-62-11||

\twolineshloka
{शक्रादिष्वपि देवेषु वर्तमानौ नयानयौ}
{श्रूयेते नरशार्दूल न त्वं व्यथितुमर्हसि} %||3-62-12||

\twolineshloka
{नष्टायामपि वैदेह्यां हृतायामपि चानघ}
{शोचितुं नार्हसे वीर यथान्यः प्राकृतस्तथा} %||3-62-13||

\twolineshloka
{त्वद्विधा हि न शोचन्ति सततं सत्यदर्शिनः}
{सुमहत्स्वपि कृच्छ्रेषु रामानिर्विण्णदर्शणाः} %||3-62-14||

\twolineshloka
{तत्त्वतो हि नरश्रेष्ठ बुद्ध्या समनुचिन्तय}
{बुद्ध्या युक्ता महाप्राज्ञा विजानन्ति शुभाशुभे} %||3-62-15||

\twolineshloka
{अदृष्टगुणदोषाणामधृतानां च कर्मणाम्}
{नान्तरेण क्रियां तेषां फलमिष्टं प्रवर्तते} %||3-62-16||

\twolineshloka
{मामेव हि पुरा वीर त्वमेव बहुषोऽन्वशाः}
{अनुशिष्याद्धि को नु त्वामपि साक्षाद्बृहस्पतिः} %||3-62-17||

\twolineshloka
{बुद्धिश्च ते महाप्राज्ञ देवैरपि दुरन्वया}
{शोकेनाभिप्रसुप्तं ते ज्ञानं सम्बोधयाम्यहम्} %||3-62-18||

\twolineshloka
{दिव्यं च मानुषं चैवमात्मनश्च पराक्रमम्}
{इक्ष्वाकुवृषभावेक्ष्य यतस्व द्विषतां बधे} %||3-62-19||

\twolineshloka
{किं ते सर्वविनाशेन कृतेन पुरुषर्षभ}
{तमेव तु रिपुं पापं विज्ञायोद्धर्तुमर्हसि} %||3-63-20||


{॥इत्यार्षे श्रीमद्रामायणे वाल्मीकीये आदिकाव्ये अरण्यकाण्डे द्विषष्टितमः सर्गः॥}

\sect{त्रिषष्टितमः सर्गः}

\twolineshloka
{पूर्वजोऽप्युक्तमात्रस्तु लक्ष्मणेन सुभाषितम्}
{सारग्राही महासारं प्रतिजग्राह राघवः} %||3-63-1||

\twolineshloka
{संनिगृह्य महाबाहुः प्रवृद्धं कोपमात्मनः}
{अवष्टभ्य धनुश्चित्रं रामो लक्ष्मणमब्रवीत्} %||3-63-2||

\twolineshloka
{किं करिष्यावहे वत्स क्व वा गच्छाव लक्ष्मण}
{केनोपायेन पश्येयं सीतामिति विचिन्तय} %||3-63-3||

\twolineshloka
{तं तथा परितापार्तं लक्ष्मणो राममब्रवीत्}
{इदमेव जनस्थानं त्वमन्वेषितुमर्हसि} %||3-63-4||

\twolineshloka
{राक्षसैर्बहुभिः कीर्णं नानाद्रुमलतायुतम्}
{सन्तीह गिरिदुर्गाणि निर्दराः कन्दराणि च} %||3-63-5||

\twolineshloka
{गुहाश्च विविधा घोरा नानामृगगणाकुलाः}
{आवासाः किंनराणां च गन्धर्वभवनानि च} %||3-63-6||

\twolineshloka
{तानि युक्तो मया सार्धं त्वमन्वेषितुमर्हसि}
{त्वद्विधो बुद्धिसम्पन्ना माहात्मानो नरर्षभ} %||3-63-7||

\twolineshloka
{आपत्सु न प्रकम्पन्ते वायुवेगैरिवाचलाः}
{इत्युक्तस्तद्वनं सर्वं विचचार सलक्ष्मणः} %||3-63-8||

\twolineshloka
{क्रुद्धो रामः शरं घोरं सन्धाय धनुषि क्षुरम्}
{ततः पर्वतकूटाभं महाभागं द्विजोत्तमम्} %||3-63-9||

\threelineshloka
{ददर्श पतितं भूमौ क्षतजार्द्रं जटायुषम्}
{तं दृष्ट्वा गिरिशृङ्गाभं रामो लक्ष्मणमब्रवीत्}
{अनेन सीता वैदेही भक्षिता नात्र संशयः} %||3-63-10||

\threelineshloka
{गृध्ररूपमिदं व्यक्तं रक्षो भ्रमति काननम्}
{भक्षयित्वा विशालाक्षीमास्ते सीतां यथासुखम्}
{एनं वधिष्ये दीप्ताग्रैर्घोरैर्बाणैरजिह्मगैः} %||3-63-11||

\twolineshloka
{इत्युक्त्वाभ्यपतद्गृध्रं सन्धाय धनुषि क्षुरम्}
{क्रुद्धो रामः समुद्रान्तां चालयन्निव मेदिनीम्} %||3-63-12||

\twolineshloka
{तं दीनदीनया वाचा सफेनं रुधिरं वमन्}
{अभ्यभाषत पक्षी तु रामं दशरथात्मजम्} %||3-63-13||

\twolineshloka
{यामोषधिमिवायुष्मन्नन्वेषसि महावने}
{सा देवी मम च प्राणा रावणेनोभयं हृतम्} %||3-63-14||

\twolineshloka
{त्वया विरहिता देवी लक्ष्मणेन च राघव}
{ह्रियमाणा मया दृष्टा रावणेन बलीयसा} %||3-63-15||

\twolineshloka
{सीतामभ्यवपन्नोऽहं रावणश्च रणे मया}
{विध्वंसितरथच्छत्रः पातितो धरणीतले} %||3-63-16||

\twolineshloka
{एतदस्य धनुर्भग्नमेतदस्य शरावरम्}
{अयमस्य रणे राम भग्नः साङ्ग्रामिको रथः} %||3-63-17||

\threelineshloka
{परिश्रान्तस्य मे पक्षौ छित्त्वा खड्गेन रावणः}
{सीतामादाय वैदेहीमुत्पपात विहायसम्}
{रक्षसा निहतं पूर्व्म न मां हन्तुं त्वमर्हसि} %||3-63-18||

\twolineshloka
{रामस्तस्य तु विज्ञाय सीतासक्तां प्रियां कथाम्}
{गृध्रराजं परिष्वज्य रुरोद सहलक्ष्मणः} %||3-63-19||

\twolineshloka
{एकमेकायने दुर्गे निःश्वसन्तं कथञ्चन}
{समीक्ष्य दुःखितो रामः सौमित्रिमिदमब्रवीत्} %||3-63-20||

\twolineshloka
{राज्याद्भ्रंशो वने वासः सीता नष्टा हतो द्विजः}
{ईदृशीयं ममालक्ष्मीर्निर्दहेदपि पावकम्} %||3-63-21||

\twolineshloka
{सम्पूर्णमपि चेदद्य प्रतरेयं महोदधिम्}
{सोऽपि नूनं ममालक्ष्म्या विशुष्येत्सरितां पतिः} %||3-63-22||

\twolineshloka
{नास्त्यभाग्यतरो लोके मत्तोऽस्मिन्सचराचरे}
{येनेयं महती प्राप्ता मया व्यसनवागुरा} %||3-63-23||

\twolineshloka
{अयं पितृवयस्यो मे गृध्रराजो जरान्वितः}
{शेते विनिहतो भूमौ मम भाग्यविपर्ययात्} %||3-63-24||

\twolineshloka
{इत्येवमुक्त्वा बहुशो राघवः सहलक्ष्मणः}
{जटायुषं च पस्पर्श पितृस्नेहं निदर्शयन्} %||3-63-25||

\fourlineindentedshloka
{निकृत्तपक्षं रुधिरावसिक्तम्}
{ तं गृध्रराजं परिरभ्य रामः}
{क्व मैथिलि प्राणसमा ममेति}
{ विमुच्य वाचं निपपात भूमौ} %||3-64-26||


{॥इत्यार्षे श्रीमद्रामायणे वाल्मीकीये आदिकाव्ये अरण्यकाण्डे त्रिषष्टितमः सर्गः॥}

\sect{चतुःषष्टितमः सर्गः}

\twolineshloka
{रामः प्रेक्ष्य तु तं गृध्रं भुवि रौद्रेण पातितम्}
{सौमित्रिं मित्रसम्पन्नमिदं वचनमब्रवीत्} %||3-64-1||

\twolineshloka
{ममायं नूनमर्थेषु यतमानो विहङ्गमः}
{राक्षसेन हतः सङ्ख्ये प्राणांस्त्यजति दुस्त्यजान्} %||3-64-2||

\twolineshloka
{अयमस्य शरीरेऽस्मिन्प्राणो लक्ष्मण विद्यते}
{तथा स्वरविहीनोऽयं विक्लवं समुदीक्षते} %||3-64-3||

\twolineshloka
{जटायो यदि शक्नोषि वाक्यं व्याहरितुं पुनः}
{सीतामाख्याहि भद्रं ते वधमाख्याहि चात्मनः} %||3-64-4||

\twolineshloka
{किंनिमित्तोऽहरत्सीतां रावणस्तस्य किं मया}
{अपराद्धं तु यं दृष्ट्वा रावणेन हृता प्रिया} %||3-64-5||

\twolineshloka
{कथं तच्चन्द्रसङ्काशं मुखमासीन्मनोहरम्}
{सीतया कानि चोक्तानि तस्मिन्काले द्विजोत्तम} %||3-64-6||

\twolineshloka
{कथंवीर्यः कथंरूपः किङ्कर्मा स च राक्षसः}
{क्व चास्य भवनं तात ब्रूहि मे परिपृच्छतः} %||3-64-7||

\twolineshloka
{तमुद्वीक्ष्याथ दीनात्मा विलपन्तमनन्तरम्}
{वाचातिसन्नया रामं जटायुरिदमब्रवीत्} %||3-64-8||

\twolineshloka
{सा हृता राक्षसेन्द्रेण रावणेन विहायसा}
{मायामास्थाय विपुलां वातदुर्दिनसङ्कुलाम्} %||3-64-9||

\twolineshloka
{परिश्रान्तस्य मे तात पक्षौ छित्त्वा निशाचरः}
{सीतामादाय वैदेहीं प्रयातो दक्षिणा मुखः} %||3-64-10||

\twolineshloka
{उपरुध्यन्ति मे प्राणा दृष्टिर्भ्रमति राघव}
{पश्यामि वृक्षान्सौवर्णानुशीरकृतमूर्धजान्} %||3-64-11||

\twolineshloka
{येन याति मुहूर्तेन सीतामादाय रावणः}
{विप्रनष्टं धनं क्षिप्रं तत्स्वामिप्रतिपद्यते} %||3-64-12||

\twolineshloka
{विन्दो नाम मुहूर्तोऽसौ स च काकुत्स्थ नाबुधत्}
{झषवद्बडिशं गृह्य क्षिप्रमेव विनश्यति} %||3-64-13||

\twolineshloka
{न च त्वया व्यथा कार्या जनकस्य सुतां प्रति}
{वैदेह्या रंस्यसे क्षिप्रं हत्वा तं राक्षसं रणे} %||3-64-14||

\twolineshloka
{असम्मूढस्य गृध्रस्य रामं प्रत्यनुभाषतः}
{आस्यात्सुस्राव रुधिरं म्रियमाणस्य सामिषम्} %||3-64-15||

\twolineshloka
{पुत्रो विश्रवसः साक्षाद्भ्राता वैश्रवणस्य च}
{इत्युक्त्वा दुर्लभान्प्राणान्मुमोच पतगेश्वरः} %||3-64-16||

\twolineshloka
{ब्रूहि ब्रूहीति रामस्य ब्रुवाणस्य कृताञ्जलेः}
{त्यक्त्वा शरीरं गृध्रस्य जग्मुः प्राणा विहायसम्} %||3-64-17||

\twolineshloka
{स निक्षिप्य शिरो भूमौ प्रसार्य चरणौ तदा}
{विक्षिप्य च शरीरं स्वं पपात धरणीतले} %||3-64-18||

\twolineshloka
{तं गृध्रं प्रेक्ष्य ताम्राक्षं गतासुमचलोपमम्}
{रामः सुबहुभिर्दुःखैर्दीनः सौमित्रिमब्रवीत्} %||3-64-19||

\twolineshloka
{बहूनि रक्षसां वासे वर्षाणि वसता सुखम्}
{अनेन दण्डकारण्ये विचीर्णमिह पक्षिणा} %||3-64-20||

\twolineshloka
{अनेकवार्षिको यस्तु चिरकालं समुत्थितः}
{सोऽयमद्य हतः शेते कालो हि दुरतिक्रमः} %||3-64-21||

\twolineshloka
{पश्य लक्ष्मण गृध्रोऽयमुपकारी हतश्च मे}
{सीतामभ्यवपन्नो वै रावणेन बलीयसा} %||3-64-22||

\twolineshloka
{गृध्रराज्यं परित्यज्य पितृपैतामहं महत्}
{मम हेतोरयं प्राणान्मुमोच पतगेश्वरः} %||3-64-23||

\twolineshloka
{सर्वत्र खलु दृश्यन्ते साधवो धर्मचारिणः}
{शूराः शरण्याः सौमित्रे तिर्यग्योनिगतेष्वपि} %||3-64-24||

\twolineshloka
{सीताहरणजं दुःखं न मे सौम्य तथागतम्}
{यथा विनाशो गृध्रस्य मत्कृते च परन्तप} %||3-64-25||

\twolineshloka
{राजा दशरथः श्रीमान्यथा मम मया यशाः}
{पूजनीयश्च मान्यश्च तथायं पतगेश्वरः} %||3-64-26||

\twolineshloka
{सौमित्रे हर काष्ठानि निर्मथिष्यामि पावकम्}
{गृध्रराजं दिधक्षामि मत्कृते निधनं गतम्} %||3-64-27||

\twolineshloka
{नाथं पतगलोकस्य चितामारोपयाम्यहम्}
{इमं धक्ष्यामि सौमित्रे हतं रौद्रेण रक्षसा} %||3-64-28||

\twolineshloka
{या गतिर्यज्ञशीलानामाहिताग्नेश्च या गतिः}
{अपरावर्तिनां या च या च भूमिप्रदायिनाम्} %||3-64-29||

\twolineshloka
{मया त्वं समनुज्ञातो गच्छ लोकाननुत्तमान्}
{गृध्रराज महासत्त्व संस्कृतश्च मया व्रज} %||3-64-30||

\twolineshloka
{एवमुक्त्वा चितां दीप्तामारोप्य पतगेश्वरम्}
{ददाह रामो धर्मात्मा स्वबन्धुमिव दुःखितः} %||3-64-31||

\twolineshloka
{रामोऽथ सहसौमित्रिर्वनं यात्वा स वीर्यवान्}
{स्थूलान्हत्वा महारोहीननु तस्तार तं द्विजम्} %||3-64-32||

\twolineshloka
{रोहिमांसानि चोद्धृत्य पेशीकृत्वा महायशाः}
{शकुनाय ददौ रामो रम्ये हरितशाद्वले} %||3-64-33||

\twolineshloka
{यत्तत्प्रेतस्य मर्त्यस्य कथयन्ति द्विजातयः}
{तत्स्वर्गगमनं तस्य क्षिप्रं रामो जजाप ह} %||3-64-34||

\twolineshloka
{ततो गोदावरीं गत्वा नदीं नरवरात्मजौ}
{उदकं चक्रतुस्तस्मै गृध्रराजाय तावुभौ} %||3-64-35||

\fourlineindentedshloka
{स गृध्रराजः कृतवान्यशस्करम्}
{ सुदुष्करं कर्म रणे निपातितः}
{महर्षिकल्पेन च संस्कृतस्तदा}
{ जगाम पुण्यां गतिमात्मनः शुभाम्} %||3-65-36||


{॥इत्यार्षे श्रीमद्रामायणे वाल्मीकीये आदिकाव्ये अरण्यकाण्डे चतुःषष्टितमः सर्गः॥}

\sect{पञ्चषष्टितमः सर्गः}

\twolineshloka
{कृत्वैवमुदकं तस्मै प्रस्थितौ राघवौ तदा}
{अवेक्षन्तौ वने सीतां पश्चिमां जग्मतुर्दिशम्} %||3-65-1||

\twolineshloka
{तां दिशं दक्षिणां गत्वा शरचापासिधारिणौ}
{अविप्रहतमैक्ष्वाकौ पन्थानं प्रतिपेदतुः} %||3-65-2||

\twolineshloka
{गुल्मैर्वृक्षैश्च बहुभिर्लताभिश्च प्रवेष्टितम्}
{आवृतं सर्वतो दुर्गं गहनं घोरदर्शनम्} %||3-65-3||

\twolineshloka
{व्यतिक्रम्य तु वेगेन गृहीत्वा दक्षिणां दिशम्}
{सुभीमं तन्महारण्यं व्यतियातौ महाबलौ} %||3-65-4||

\twolineshloka
{ततः परं जनस्थानात्त्रिक्रोशं गम्य राघवौ}
{क्रौञ्चारण्यं विविशतुर्गहनं तौ महौजसौ} %||3-65-5||

\twolineshloka
{नानामेघघनप्रख्यं प्रहृष्टमिव सर्वतः}
{नानावर्णैः शुभैः पुष्पैर्मृगपक्षिगणैर्युतम्} %||3-65-6||

\twolineshloka
{दिदृक्षमाणौ वैदेहीं तद्वनं तौ विचिक्यतुः}
{तत्र तत्रावतिष्ठन्तौ सीताहरणकर्शितौ} %||3-65-7||

\twolineshloka
{लक्ष्मणस्तु महातेजाः सत्त्ववाञ्शीलवाञ्शुचिः}
{अब्रवीत्प्राञ्जलिर्वाक्यं भ्रातरं दीप्ततेजसम्} %||3-65-8||

\twolineshloka
{स्पन्दते मे दृढं बाहुरुद्विग्नमिव मे मनः}
{प्रायशश्चाप्यनिष्टानि निमित्तान्युपलक्षये} %||3-65-9||

\twolineshloka
{तस्मात्सज्जीभवार्य त्वं कुरुष्व वचनं हितम्}
{ममैव हि निमित्तानि सद्यः शंसन्ति सम्भ्रमम्} %||3-65-10||

\twolineshloka
{एष वञ्चुलको नाम पक्षी परमदारुणः}
{आवयोर्विजयं युद्धे शंसन्निव विनर्दति} %||3-65-11||

\twolineshloka
{तयोरन्वेषतोरेवं सर्वं तद्वनमोजसा}
{संजज्ञे विपुलः शब्दः प्रभञ्जन्निव तद्वनम्} %||3-65-12||

\twolineshloka
{संवेष्टितमिवात्यर्थं गहनं मातरिश्वना}
{वनस्य तस्य शब्दोऽभूद्दिवमापूरयन्निव} %||3-65-13||

\twolineshloka
{तं शब्दं काङ्क्षमाणस्तु रामः कक्षे सहानुजः}
{ददर्श सुमहाकायं राक्षसं विपुलोरसम्} %||3-65-14||

\twolineshloka
{आसेदतुस्ततस्तत्र तावुभौ प्रमुखे स्थितम्}
{विवृद्धमशिरोग्रीवं कबन्धमुदरे मुखम्} %||3-65-15||

\twolineshloka
{रोमभिर्निचितैस्तीक्ष्णैर्महागिरिमिवोच्छ्रितम्}
{नीलमेघनिभं रौद्रं मेघस्तनितनिःस्वनम्} %||3-65-16||

\twolineshloka
{महापक्ष्मेण पिङ्गेन विपुलेनायतेन च}
{एकेनोरसि घोरेण नयनेनाशुदर्शिना} %||3-65-17||

\twolineshloka
{महादंष्ट्रोपपन्नं तं लेलिहानं महामुखम्}
{भक्षयन्तं महाघोरानृक्षसिंहमृगद्विपान्} %||3-65-18||

\twolineshloka
{घोरौ भुजौ विकुर्वाणमुभौ योजनमायतौ}
{कराभ्यां विविधान्गृह्य ऋष्कान्पक्षिगणान्मृगान्} %||3-65-19||

\twolineshloka
{आकर्षन्तं विकर्षन्तमनेकान्मृगयूथपान्}
{स्थितमावृत्य पन्थानं तयोर्भ्रात्रोः प्रपन्नयोः} %||3-65-20||

\twolineshloka
{अथ तौ समतिक्रम्य क्रोशमात्रे ददर्शतुः}
{महान्तं दारुणं भीमं कबन्धं भुजसंवृतम्} %||3-65-21||

\twolineshloka
{स महाबाहुरत्यर्थं प्रसार्य विपुलौ भुजौ}
{जग्राह सहितावेव राघवौ पीडयन्बलात्} %||3-65-22||

\twolineshloka
{खड्गिनौ दृढधन्वानौ तिग्मतेजौ महाभुजौ}
{भ्रातरौ विवशं प्राप्तौ कृष्यमाणौ महाबलौ} %||3-65-23||

\twolineshloka
{तावुवाच महाबाहुः कबन्धो दानवोत्तमः}
{कौ युवां वृषभस्कन्धौ महाखड्गधनुर्धरौ} %||3-65-24||

\twolineshloka
{घोरं देशमिमं प्राप्तौ मम भक्षावुपस्थितौ}
{वदतं कार्यमिह वां किमर्थं चागतौ युवाम्} %||3-65-25||

\threelineshloka
{इमं देशमनुप्राप्तौ क्षुधार्तस्येह तिष्ठतः}
{सबाणचापखड्गौ च तीक्ष्णशृङ्गाविवर्षभौ}
{ममास्यमनुसम्प्राप्तौ दुर्लभं जीवितं पुनः} %||3-65-26||

\twolineshloka
{तस्य तद्वचनं श्रुत्वा कबन्धस्य दुरात्मनः}
{उवाच लक्ष्मणं रामो मुखेन परिशुष्यता} %||3-65-27||

\twolineshloka
{कृच्छ्रात्कृच्छ्रतरं प्राप्य दारुणं सत्यविक्रम}
{व्यसनं जीवितान्ताय प्राप्तमप्राप्य तां प्रियाम्} %||3-65-28||

\threelineshloka
{कालस्य सुमहद्वीर्यं सर्वभूतेषु लक्ष्मण}
{त्वां च मां च नरव्याघ्र व्यसनैः पश्य मोहितौ}
{नातिभारोऽस्ति दैवस्य सर्वभुतेषु लक्ष्मण} %||3-65-29||

\twolineshloka
{शूराश्च बलवन्तश्च कृतास्त्राश्च रणाजिरे}
{कालाभिपन्नाः सीदन्ति यथा वालुकसेतवः} %||3-65-30||

\fourlineindentedshloka
{इति ब्रुवाणो दृढसत्यविक्रमो}
{ महायशा दाशरथिः प्रतापवान्}
{अवेक्ष्य सौमित्रिमुदग्रविक्रमम्}
{ स्थिरां तदा स्वां मतिमात्मनाकरोत्} %||3-66-31||


{॥इत्यार्षे श्रीमद्रामायणे वाल्मीकीये आदिकाव्ये अरण्यकाण्डे पञ्चषष्टितमः सर्गः॥}

\sect{षट्षष्टितमः सर्गः}

\twolineshloka
{तौ तु तत्र स्थितौ दृष्ट्वा भ्रातरौ रामलक्ष्मणौ}
{बाहुपाशपरिक्षिप्तौ कबन्धो वाक्यमब्रवीत्} %||3-66-1||

\twolineshloka
{तिष्ठतः किं नु मां दृष्ट्वा क्षुधार्तं क्षत्रियर्षभौ}
{आहारार्थं तु सन्दिष्टौ दैवेन गतचेतसौ} %||3-66-2||

\twolineshloka
{तच्छ्रुत्वा लक्ष्मणो वाक्यं प्राप्तकालं हितं तदा}
{उवाचार्तिसमापन्नो विक्रमे कृतनिश्चयः} %||3-66-3||

\twolineshloka
{त्वां च मां च पुरा तूर्णमादत्ते राक्षसाधमः}
{तस्मादसिभ्यामस्याशु बाहू छिन्दावहे गुरू} %||3-66-4||

\twolineshloka
{ततस्तौ देशकालज्ञौ खड्गाभ्यामेव राघवौ}
{अच्छिन्दतां सुसंहृष्टौ बाहू तस्यांसदेशयोः} %||3-66-5||

\twolineshloka
{दक्षिणो दक्षिणं बाहुमसक्तमसिना ततः}
{चिच्छेद रामो वेगेन सव्यं वीरस्तु लक्ष्मणः} %||3-66-6||

\twolineshloka
{स पपात महाबाहुश्छिन्नबाहुर्महास्वनः}
{खं च गां च दिशश्चैव नादयञ्जलदो यथा} %||3-66-7||

\twolineshloka
{स निकृत्तौ भुजौ दृष्ट्वा शोणितौघपरिप्लुतः}
{दीनः पप्रच्छ तौ वीरौ कौ युवामिति दानवः} %||3-66-8||

\twolineshloka
{इति तस्य ब्रुवाणस्य लक्ष्मणः शुभलक्षणः}
{शशंस तस्य काकुत्स्थं कबन्धस्य महाबलः} %||3-66-9||

\twolineshloka
{अयमिक्ष्वाकुदायादो रामो नाम जनैः श्रुतः}
{अस्यैवावरजं विद्धि भ्रातरं मां च लक्ष्मणम्} %||3-66-10||

\twolineshloka
{अस्य देवप्रभावस्य वसतो विजने वने}
{रक्षसापहृता भार्या यामिच्छन्ताविहागतौ} %||3-66-11||

\twolineshloka
{त्वं तु को वा किमर्थं वा कबन्ध सदृशो वने}
{आस्येनोरसि दीप्तेन भग्नजङ्घो विचेष्टसे} %||3-66-12||

\twolineshloka
{एवमुक्तः कबन्धस्तु लक्ष्मणेनोत्तरं वचः}
{उवाच परमप्रीतस्तदिन्द्रवचनं स्मरन्} %||3-66-13||

\twolineshloka
{स्वागतं वां नरव्याघ्रौ दिष्ट्या पश्यामि चाप्यहम्}
{दिष्ट्या चेमौ निकृत्तौ मे युवाभ्यां बाहुबन्धनौ} %||3-66-14||

\twolineshloka
{विरूपं यच्च मे रूपं प्राप्तं ह्यविनयाद्यथा}
{तन्मे शृणु नरव्याघ्र तत्त्वतः शंसतस्तव} %||3-67-15||


{॥इत्यार्षे श्रीमद्रामायणे वाल्मीकीये आदिकाव्ये अरण्यकाण्डे षट्षष्टितमः सर्गः॥}

\sect{सप्तषष्टितमः सर्गः}

\threelineshloka
{पुरा राम महाबाहो महाबलपराक्रम}
{रूपमासीन्ममाचिन्त्यं त्रिषु लोकेषु विश्रुतम्}
{यथा सोमस्य शक्रस्य सूर्यस्य च यथा वपुः} %||3-67-1||

\twolineshloka
{सोऽहं रूपमिदं कृत्वा लोकवित्रासनं महत्}
{ऋषीन्वनगतान्राम त्रासयामि ततस्ततः} %||3-67-2||

\twolineshloka
{ततः स्थूलशिरा नाम महर्षिः कोपितो मया}
{सञ्चिन्वन्विविधं वन्यं रूपेणानेन धर्षितः} %||3-67-3||

\twolineshloka
{तेनाहमुक्तः प्रेक्ष्यैवं घोरशापाभिधायिना}
{एतदेव नृशंसं ते रूपमस्तु विगर्हितम्} %||3-67-4||

\twolineshloka
{स मया याचितः क्रुद्धः शापस्यान्तो भवेदिति}
{अभिशापकृतस्येति तेनेदं भाषितं वचः} %||3-67-5||

\twolineshloka
{यदा छित्त्वा भुजौ रामस्त्वां दहेद्विजने वने}
{तदा त्वं प्राप्स्यसे रूपं स्वमेव विपुलं शुभम्} %||3-67-6||

\twolineshloka
{श्रिया विराजितं पुत्रं दनोस्त्वं विद्धि लक्ष्मण}
{इन्द्रकोपादिदं रूपं प्राप्तमेवं रणाजिरे} %||3-67-7||

\twolineshloka
{अहं हि तपसोग्रेण पितामहमतोषयम्}
{दीर्घमायुः स मे प्रादात्ततो मां विभ्रमोऽस्पृशत्} %||3-67-8||

\twolineshloka
{दीर्घमायुर्मया प्राप्तं किं मे शक्रः करिष्यति}
{इत्येवं बुद्धिमास्थाय रणे शक्रमधर्षयम्} %||3-67-9||

\twolineshloka
{तस्य बाहुप्रमुक्तेन वज्रेण शतपर्वणा}
{सक्थिनी च शिरश्चैव शरीरे सम्प्रवेशितम्} %||3-67-10||

\twolineshloka
{स मया याच्यमानः सन्नानयद्यमसादनम्}
{पितामहवचः सत्यं तदस्त्विति ममाब्रवीत्} %||3-67-11||

\twolineshloka
{अनाहारः कथं शक्तो भग्नसक्थिशिरोमुखः}
{वज्रेणाभिहतः कालं सुदीर्घमपि जीवितुम्} %||3-67-12||

\twolineshloka
{एवमुक्तस्तु मे शक्रो बाहू योजनमायतौ}
{प्रादादास्यं च मे कुक्षौ तीक्ष्णदंष्ट्रमकल्पयत्} %||3-67-13||

\twolineshloka
{सोऽहं भुजाभ्यां दीर्घाभ्यां समाकृष्य वनेचरान्}
{सिंहद्विपमृगव्याघ्रान्भक्षयामि समन्ततः} %||3-67-14||

\twolineshloka
{स तु मामब्रवीदिन्द्रो यदा रामः सलक्ष्मणः}
{छेत्स्यते समरे बाहू तदा स्वर्गं गमिष्यसि} %||3-67-15||

\twolineshloka
{स त्वं रामोऽसि भद्रं ते नाहमन्येन राघव}
{शक्यो हन्तुं यथातत्त्वमेवमुक्तं महर्षिणा} %||3-67-16||

\twolineshloka
{अहं हि मतिसाचिव्यं करिष्यामि नरर्षभ}
{मित्रं चैवोपदेक्ष्यामि युवाभ्यां संस्कृतोऽग्निना} %||3-67-17||

\twolineshloka
{एवमुक्तस्तु धर्मात्मा दनुना तेन राघवः}
{इदं जगाद वचनं लक्ष्मणस्योपशृण्वतः} %||3-67-18||

\twolineshloka
{रावणेन हृता सीता मम भार्या यशस्विनी}
{निष्क्रान्तस्य जनस्थानात्सह भ्रात्रा यथासुखम्} %||3-67-19||

\twolineshloka
{नाममात्रं तु जानामि न रूपं तस्य रक्षसः}
{निवासं वा प्रभावं वा वयं तस्य न विद्महे} %||3-67-20||

\twolineshloka
{शोकार्तानामनाथानामेवं विपरिधावताम्}
{कारुण्यं सदृशं कर्तुमुपकारे च वर्तताम्} %||3-67-21||

\twolineshloka
{काष्ठान्यानीय शुष्काणि काले भग्नानि कुञ्जरैः}
{भक्ष्यामस्त्वां वयं वीर श्वभ्रे महति कल्पिते} %||3-67-22||

\twolineshloka
{स त्वं सीतां समाचक्ष्व येन वा यत्र वा हृता}
{कुरु कल्याणमत्यर्थं यदि जानासि तत्त्वतः} %||3-67-23||

\twolineshloka
{एवमुक्तस्तु रामेण वाक्यं दनुरनुत्तमम्}
{प्रोवाच कुशलो वक्तुं वक्तारमपि राघवम्} %||3-67-24||

\twolineshloka
{दिव्यमस्ति न मे ज्ञानं नाभिजानामि मैथिलीम्}
{यस्तां ज्ञास्यति तं वक्ष्ये दग्धः स्वं रूपमास्थितः} %||3-67-25||

\twolineshloka
{अदग्धस्य हि विज्ञातुं शक्तिरस्ति न मे प्रभो}
{राक्षसं तं महावीर्यं सीता येन हृता तव} %||3-67-26||

\twolineshloka
{विज्ञानं हि महद्भ्रष्टं शापदोषेण राघव}
{स्वकृतेन मया प्राप्तं रूपं लोकविगर्हितम्} %||3-67-27||

\twolineshloka
{किं तु यावन्न यात्यस्तं सविता श्रान्तवाहनः}
{तावन्मामवटे क्षिप्त्वा दह राम यथाविधि} %||3-67-28||

\twolineshloka
{दग्धस्त्वयाहमवटे न्यायेन रघुनन्दन}
{वक्ष्यामि तमहं वीर यस्तं ज्ञास्यति राक्षसम्} %||3-67-29||

\twolineshloka
{तेन सख्यं च कर्तव्यं न्याय्यवृत्तेन राघव}
{कल्पयिष्यति ते प्रीतः साहाय्यं लघुविक्रमः} %||3-67-30||

\twolineshloka
{न हि तस्यास्त्यविज्ञातं त्रिषु लोकेषु राघव}
{सर्वान्परिसृतो लोकान्पुरा वै कारणान्तरे} %||3-68-31||


{॥इत्यार्षे श्रीमद्रामायणे वाल्मीकीये आदिकाव्ये अरण्यकाण्डे सप्तषष्टितमः सर्गः॥}

\sect{अष्टषष्टितमः सर्गः}

\twolineshloka
{एवमुक्तौ तु तौ वीरौ कबन्धेन नरेश्वरौ}
{गिरिप्रदरमासाद्य पावकं विससर्जतुः} %||3-68-1||

\twolineshloka
{लक्ष्मणस्तु महोल्काभिर्ज्वलिताभिः समन्ततः}
{चितामादीपयामास सा प्रजज्वाल सर्वतः} %||3-68-2||

\twolineshloka
{तच्छरीरं कबन्धस्य घृतपिण्डोपमं महत्}
{मेदसा पच्यमानस्य मन्दं दहति पावक} %||3-68-3||

\twolineshloka
{स विधूय चितामाशु विधूमोऽग्निरिवोत्थितः}
{अरजे वाससी विभ्रन्मालां दिव्यां महाबलः} %||3-68-4||

\twolineshloka
{ततश्चिताया वेगेन भास्वरो विरजाम्बरः}
{उत्पपाताशु संहृष्टः सर्वप्रत्यङ्गभूषणः} %||3-68-5||

\twolineshloka
{विमाने भास्वरे तिष्ठन्हंसयुक्ते यशस्करे}
{प्रभया च महातेजा दिशो दश विराजयन्} %||3-68-6||

\twolineshloka
{सोऽन्तरिक्षगतो रामं कबन्धो वाक्यमब्रवीत्}
{शृणु राघव तत्त्वेन यथा सीमामवाप्स्यसि} %||3-68-7||

\twolineshloka
{राम षड्युक्तयो लोके याभिः सर्वं विमृश्यते}
{परिमृष्टो दशान्तेन दशाभागेन सेव्यते} %||3-68-8||

\twolineshloka
{दशाभागगतो हीनस्त्वं राम सहलक्ष्मणः}
{यत्कृते व्यसनं प्राप्तं त्वया दारप्रधर्षणम्} %||3-68-9||

\twolineshloka
{तदवश्यं त्वया कार्यः स सुहृत्सुहृदां वर}
{अकृत्वा न हि ते सिद्धिमहं पश्यामि चिन्तयन्} %||3-68-10||

\twolineshloka
{श्रूयतां राम वक्ष्यामि सुग्रीवो नाम वानरः}
{भ्रात्रा निरस्तः क्रुद्धेन वालिना शक्रसूनुना} %||3-68-11||

\twolineshloka
{ऋष्यमूके गिरिवरे पम्पापर्यन्तशोभिते}
{निवसत्यात्मवान्वीरश्चतुर्भिः सह वानरैः} %||3-68-12||

\twolineshloka
{वयस्यं तं कुरु क्षिप्रमितो गत्वाद्य राघव}
{अद्रोहाय समागम्य दीप्यमाने विभावसौ} %||3-68-13||

\twolineshloka
{न च ते सोऽवमन्तव्यः सुग्रीवो वानराधिपः}
{कृतज्ञः कामरूपी च सहायार्थी च वीर्यवान्} %||3-68-14||

\twolineshloka
{शक्तौ ह्यद्य युवां कर्तुं कार्यं तस्य चिकीर्षितम्}
{कृतार्थो वाकृतार्थो वा कृत्यं तव करिष्यति} %||3-68-15||

\twolineshloka
{स ऋक्षरजसः पुत्रः पम्पामटति शङ्कितः}
{भास्करस्यौरसः पुत्रो वालिना कृतकिल्बिषः} %||3-68-16||

\twolineshloka
{संनिधायायुधं क्षिप्रमृष्यमूकालयं कपिम्}
{कुरु राघव सत्येन वयस्यं वनचारिणम्} %||3-68-17||

\twolineshloka
{स हि स्थानानि सर्वाणि कार्त्स्न्येन कपिकुञ्जरः}
{नरमांसाशिनां लोके नैपुण्यादधिगच्छति} %||3-68-18||

\twolineshloka
{न तस्याविदितं लोके किञ्चिदस्ति हि राघव}
{यावत्सूर्यः प्रतपति सहस्रांशुररिन्दम} %||3-68-19||

\twolineshloka
{स नदीर्विपुलाञ्शैलान्गिरिदुर्गाणि कन्दरान्}
{अन्विष्य वानरैः सार्धं पत्नीं तेऽधिगमिष्यति} %||3-68-20||

\twolineshloka
{वानरांश्च महाकायान्प्रेषयिष्यति राघव}
{दिशो विचेतुं तां सीतां त्वद्वियोगेन शोचतीम्} %||3-68-21||

\fourlineindentedshloka
{स मेरुशृङ्गाग्रगतामनिन्दिताम्}
{ प्रविश्य पातालतलेऽपि वाश्रिताम्}
{प्लवङ्गमानां प्रवरस्तव प्रियाम्}
{ निहत्य रक्षांसि पुनः प्रदास्यति} %||3-69-22||


{॥इत्यार्षे श्रीमद्रामायणे वाल्मीकीये आदिकाव्ये अरण्यकाण्डे अष्टषष्टितमः सर्गः॥}

\sect{एकोनसप्ततितमः सर्गः}

\twolineshloka
{निदर्शयित्वा रामाय सीतायाः प्रतिपादने}
{वाक्यमन्वर्थमर्थज्ञः कबन्धः पुनरब्रवीत्} %||3-69-1||

\twolineshloka
{एष राम शिवः पन्था यत्रैते पुष्पिता द्रुमाः}
{प्रतीचीं दिशमाश्रित्य प्रकाशन्ते मनोरमाः} %||3-69-2||

\twolineshloka
{जम्बूप्रियालपनसाः प्लक्षन्यग्रोधतिन्दुकाः}
{अश्वत्थाः कर्णिकाराश्च चूताश्चान्ये च पादपाः} %||3-69-3||

\twolineshloka
{तानारुह्याथवा भूमौ पातयित्वा च तान्बलात्}
{फलान्यमृतकल्पानि भक्षयन्तौ गमिष्यथः} %||3-69-4||

\twolineshloka
{चङ्क्रमन्तौ वरान्देशाञ्शैलाच्छैलं वनाद्वनम्}
{ततः पुष्करिणीं वीरौ पम्पां नाम गमिष्यथः} %||3-69-5||

\twolineshloka
{अशर्करामविभ्रंशां समतीर्थमशैवलाम्}
{राम संजातवालूकां कमलोत्पलशोभिताम्} %||3-69-6||

\twolineshloka
{तत्र हंसाः प्लवाः क्रौञ्चाः कुरराश्चैव राघव}
{वल्गुस्वरा निकूजन्ति पम्पासलिलगोचराः} %||3-69-7||

\twolineshloka
{नोद्विजन्ते नरान्दृष्ट्वा वधस्याकोविदाः शुभाः}
{घृतपिण्डोपमान्स्थूलांस्तान्द्विजान्भक्षयिष्यथः} %||3-69-8||

\twolineshloka
{रोहितान्वक्रतुण्डांश्च नलमीनांश्च राघव}
{पम्पायामिषुभिर्मत्स्यांस्तत्र राम वरान्हतान्} %||3-69-9||

\twolineshloka
{निस्त्वक्पक्षानयस्तप्तानकृशानेककण्टकान्}
{तव भक्त्या समायुक्तो लक्ष्मणः सम्प्रदास्यति} %||3-69-10||

\twolineshloka
{भृशं ते खादतो मत्स्यान्पम्पायाः पुष्पसञ्चये}
{पद्मगन्धि शिवं वारि सुखशीतमनामयम्} %||3-69-11||

\twolineshloka
{उद्धृत्य स तदाक्लिष्टं रूप्यस्फटिकसंनिभम्}
{अथ पुष्करपर्णेन लक्ष्मणः पाययिष्यति} %||3-69-12||

\threelineshloka
{स्थूलान्गिरिगुहाशय्यान्वराहान्वनचारिणः}
{अपां लोभादुपावृत्तान्वृषभानिव नर्दतः}
{रूपान्वितांश्च पम्पायां द्रक्ष्यसि त्वं नरोत्तम} %||3-69-13||

\twolineshloka
{सायाह्ने विचरन्राम विटपी माल्यधारिणः}
{शीतोदकं च पम्पायां दृष्ट्वा शोकं विहास्यसि} %||3-69-14||

\twolineshloka
{सुमनोभिश्चितांस्तत्र तिलकान्नक्तमालकान्}
{उत्पलानि च फुल्लानि पङ्कजानि च राघव} %||3-69-15||

\twolineshloka
{न तानि कश्चिन्माल्यानि तत्रारोपयिता नरः}
{मतङ्गशिष्यास्तत्रासन्नृषयः सुसमाहितः} %||3-69-16||

\twolineshloka
{तेषां भाराभितप्तानां वन्यमाहरतां गुरोः}
{ये प्रपेतुर्महीं तूर्णं शरीरात्स्वेदबिन्दवः} %||3-69-17||

\twolineshloka
{तानि माल्यानि जातानि मुनीनां तपसा तदा}
{स्वेदबिन्दुसमुत्थानि न विनश्यन्ति राघव} %||3-69-18||

\twolineshloka
{तेषामद्यापि तत्रैव दृश्यते परिचारिणी}
{श्रमणी शबरी नाम काकुत्स्थ चिरजीविनी} %||3-69-19||

\twolineshloka
{त्वां तु धर्मे स्थिता नित्यं सर्वभूतनमस्कृतम्}
{दृष्ट्वा देवोपमं राम स्वर्गलोकं गमिष्यति} %||3-69-20||

\twolineshloka
{ततस्तद्राम पम्पायास्तीरमाश्रित्य पश्चिमम्}
{आश्रमस्थानमतुलं गुह्यं काकुत्स्थ पश्यसि} %||3-69-21||

\twolineshloka
{न तत्राक्रमितुं नागाः शक्नुवन्ति तमाश्रमम्}
{ऋषेस्तस्य मतङ्गस्य विधानात्तच्च काननम्} %||3-69-22||

\twolineshloka
{तस्मिन्नन्दनसङ्काशे देवारण्योपमे वने}
{नानाविहगसङ्कीर्णे रंस्यसे राम निर्वृतः} %||3-69-23||

\threelineshloka
{ऋष्यमूकस्तु पम्पायाः पुरस्तात्पुष्पितद्रुमः}
{सुदुःखारोहणो नाम शिशुनागाभिरक्षितः}
{उदारो ब्रह्मणा चैव पूर्वकाले विनिर्मितः} %||3-69-24||

\twolineshloka
{शयानः पुरुषो राम तस्य शैलस्य मूर्धनि}
{यत्स्वप्ने लभते वित्तं तत्प्रबुद्धोऽधिगच्छति} %||3-69-25||

\twolineshloka
{न त्वेनं विषमाचारः पापकर्माधिरोहति}
{तत्रैव प्रहरन्त्येनं सुप्तमादाय राक्षसाः} %||3-69-26||

\twolineshloka
{ततोऽपि शिशुनागानामाक्रन्दः श्रूयते महान्}
{क्रीडतां राम पम्पायां मतङ्गारण्यवासिनाम्} %||3-69-27||

\twolineshloka
{सिक्ता रुधिरधाराभिः संहत्य परमद्विपाः}
{प्रचरन्ति पृथक्कीर्णा मेघवर्णास्तरस्विनः} %||3-69-28||

\twolineshloka
{ते तत्र पीत्वा पानीयं विमलं शीतमव्ययम्}
{निवृत्ताः संविगाहन्ते वनानि वनगोचराः} %||3-69-29||

\twolineshloka
{राम तस्य तु शैलस्य महती शोभते गुहा}
{शिलापिधाना काकुत्स्थ दुःखं चास्याः प्रवेशनम्} %||3-69-30||

\twolineshloka
{तस्या गुहायाः प्राग्द्वारे महाञ्शीतोदको ह्रदः}
{बहुमूलफलो रम्यो नानानगसमावृतः} %||3-69-31||

\twolineshloka
{तस्यां वसति सुग्रीवश्चतुर्भिः सह वानरैः}
{कदाचिच्छिखरे तस्य पर्वतस्यावतिष्ठते} %||3-69-32||

\twolineshloka
{कबन्धस्त्वनुशास्यैवं तावुभौ रामलक्ष्मणौ}
{स्रग्वी भास्करवर्णाभः खे व्यरोचत वीर्यवान्} %||3-69-33||

\twolineshloka
{तं तु खस्थं महाभागं कबन्धं रामलक्ष्मणौ}
{प्रस्थितौ त्वं व्रजस्वेति वाक्यमूचतुरन्तिकात्} %||3-69-34||

\twolineshloka
{गम्यतां कार्यसिद्ध्यर्थमिति तावब्रवीच्च सः}
{सुप्रीतौ तावनुज्ञाप्य कबन्धः प्रस्थितस्तदा} %||3-69-35||

\fourlineindentedshloka
{स तत्कबन्धः प्रतिपद्य रूपम्}
{ वृतः श्रिया भास्करतुल्यदेहः}
{निदर्शयन्राममवेक्ष्य खस्थः}
{ सख्यं कुरुष्वेति तदाभ्युवाच} %||3-70-36||


{॥इत्यार्षे श्रीमद्रामायणे वाल्मीकीये आदिकाव्ये अरण्यकाण्डे एकोनसप्ततितमः सर्गः॥}

\sect{सप्ततितमः सर्गः}

\twolineshloka
{तौ कबन्धेन तं मार्गं पम्पाया दर्शितं वने}
{आतस्थतुर्दिशं गृह्य प्रतीचीं नृवरात्मजौ} %||3-70-1||

\twolineshloka
{तौ शैलेष्वाचितानेकान्क्षौद्रकल्पफलद्रुमान्}
{वीक्षन्तौ जग्मतुर्द्रष्टुं सुग्रीवं रामलक्ष्मणौ} %||3-70-2||

\twolineshloka
{कृत्वा च शैलपृष्ठे तु तौ वासं रघुनन्दनौ}
{पम्पायाः पश्चिमं तीरं राघवावुपतस्थतुः} %||3-70-3||

\twolineshloka
{तौ पुष्करिण्याः पम्पायास्तीरमासाद्य पश्चिमम्}
{अपश्यतां ततस्तत्र शबर्या रम्यमाश्रमम्} %||3-70-4||

\twolineshloka
{तौ तमाश्रममासाद्य द्रुमैर्बहुभिरावृतम्}
{सुरम्यमभिवीक्षन्तौ शबरीमभ्युपेयतुः} %||3-70-5||

\twolineshloka
{तौ तु दृष्ट्वा तदा सिद्धा समुत्थाय कृताञ्जलिः}
{पादौ जग्राह रामस्य लक्ष्मणस्य च धीमतः} %||3-70-6||

\twolineshloka
{तामुवाच ततो रामः श्रमणीं संशितव्रताम्}
{कच्चित्ते निर्जिता विघ्नाः कच्चित्ते वर्धते तपः} %||3-70-7||

\threelineshloka
{कच्चित्ते नियतः कोप आहारश्च तपोधने}
{कच्चित्ते नियमाः प्राप्ताः कच्चित्ते मनसः सुखम्}
{कच्चित्ते गुरुशुश्रूषा सफला चारुभाषिणि} %||3-70-8||

\twolineshloka
{रामेण तापसी पृष्ठा सा सिद्धा सिद्धसम्मता}
{शशंस शबरी वृद्धा रामाय प्रत्युपस्थिता} %||3-70-9||

\twolineshloka
{चित्रकूटं त्वयि प्राप्ते विमानैरतुलप्रभैः}
{इतस्ते दिवमारूढा यानहं पर्यचारिषम्} %||3-70-10||

\twolineshloka
{तैश्चाहमुक्ता धर्मज्ञैर्महाभागैर्महर्षिभिः}
{आगमिष्यति ते रामः सुपुण्यमिममाश्रमम्} %||3-70-11||

\twolineshloka
{स ते प्रतिग्रहीतव्यः सौमित्रिसहितोऽतिथिः}
{तं च दृष्ट्वा वराँल्लोकानक्षयांस्त्वं गमिष्यसि} %||3-70-12||

\twolineshloka
{मया तु विविधं वन्यं सञ्चितं पुरुषर्षभ}
{तवार्थे पुरुषव्याघ्र पम्पायास्तीरसम्भवम्} %||3-70-13||

\twolineshloka
{एवमुक्तः स धर्मात्मा शबर्या शबरीमिदम्}
{राघवः प्राह विज्ञाने तां नित्यमबहिष्कृताम्} %||3-70-14||

\twolineshloka
{दनोः सकाशात्तत्त्वेन प्रभावं ते महात्मनः}
{श्रुतं प्रत्यक्षमिच्छामि सन्द्रष्टुं यदि मन्यसे} %||3-70-15||

\twolineshloka
{एतत्तु वचनं श्रुत्वा रामवक्त्राद्विनिःसृतम्}
{शबरी दर्शयामास तावुभौ तद्वनं महत्} %||3-70-16||

\twolineshloka
{पश्य मेघघनप्रख्यं मृगपक्षिसमाकुलम्}
{मतङ्गवनमित्येव विश्रुतं रघुनन्दन} %||3-70-17||

\twolineshloka
{इह ते भावितात्मानो गुरवो मे महाद्युते}
{जुहवांश्चक्रिरे तीर्थं मन्त्रवन्मन्त्रपूजितम्} %||3-70-18||

\twolineshloka
{इयं प्रत्यक्स्थली वेदी यत्र ते मे सुसत्कृताः}
{पुष्पोपहारं कुर्वन्ति श्रमादुद्वेपिभिः करैः} %||3-70-19||

\twolineshloka
{तेषां तपः प्रभावेन पश्याद्यापि रघूत्तम}
{द्योतयन्ति दिशः सर्वाः श्रिया वेद्योऽतुलप्रभाः} %||3-70-20||

\twolineshloka
{अशक्नुवद्भिस्तैर्गन्तुमुपवासश्रमालसैः}
{चिन्तितेऽभ्यागतान्पश्य समेतान्सप्त सागरान्} %||3-70-21||

\twolineshloka
{कृताभिषेकैस्तैर्न्यस्ता वल्कलाः पादपेष्विह}
{अद्यापि न विशुष्यन्ति प्रदेशे रघुनन्दन} %||3-70-22||

\twolineshloka
{कृत्स्नं वनमिदं दृष्टं श्रोतव्यं च श्रुतं त्वया}
{तदिच्छाम्यभ्यनुज्ञाता त्यक्तुमेतत्कलेवरम्} %||3-70-23||

\twolineshloka
{तेषामिच्छाम्यहं गन्तुं समीपं भावितात्मनाम्}
{मुनीनामाश्रम्मो येषामहं च परिचारिणी} %||3-70-24||

\twolineshloka
{धर्मिष्ठं तु वचः श्रुत्वा राघवः सहलक्ष्मणः}
{अनुजानामि गच्छेति प्रहृष्टवदनोऽब्रवीत्} %||3-70-25||

\twolineshloka
{अनुज्ञाता तु रामेण हुत्वात्मानं हुताशने}
{ज्वलत्पावकसङ्काशा स्वर्गमेव जगाम सा} %||3-70-26||

\twolineshloka
{यत्र ते सुकृतात्मानो विहरन्ति महर्षयः}
{तत्पुण्यं शबरीस्थानं जगामात्मसमाधिना} %||3-71-27||


{॥इत्यार्षे श्रीमद्रामायणे वाल्मीकीये आदिकाव्ये अरण्यकाण्डे सप्ततितमः सर्गः॥}

\sect{एकसप्ततितमः सर्गः}

\twolineshloka
{दिवं तु तस्यां यातायां शबर्यां स्वेन कर्मणा}
{लक्ष्मणेन सह भ्रात्रा चिन्तयामास राघवः} %||3-71-1||

\twolineshloka
{चिन्तयित्वा तु धर्मात्मा प्रभावं तं महात्मनाम्}
{हितकारिणमेकाग्रं लक्ष्मणं राघवोऽब्रवीत्} %||3-71-2||

\twolineshloka
{दृष्टोऽयमाश्रमः सौम्य बह्वाश्चर्यः कृतात्मनाम्}
{विश्वस्तमृगशार्दूलो नानाविहगसेवितः} %||3-71-3||

\twolineshloka
{सप्तानां च समुद्राणामेषु तीर्थेषु लक्ष्मण}
{उपस्पृष्टं च विधिवत्पितरश्चापि तर्पिताः} %||3-71-4||

\twolineshloka
{प्रनष्टमशुभं यत्तत्कल्याणं समुपस्थितम्}
{तेन त्वेतत्प्रहृष्टं मे मनो लक्ष्मण सम्प्रति} %||3-71-5||

\twolineshloka
{हृदये हि नरव्याघ्र शुभमाविर्भविष्यति}
{तदागच्छ गमिष्यावः पम्पां तां प्रियदर्शनाम्} %||3-71-6||

\threelineshloka
{ऋश्यमूको गिरिर्यत्र नातिदूरे प्रकाशते}
{यस्मिन्वसति धर्मात्मा सुग्रीवोंऽशुमतः सुतः}
{नित्यं वालिभयात्त्रस्तश्चतुर्भिः सह वानरैः} %||3-71-7||

\twolineshloka
{अभित्वरे च तं द्रष्टुं सुग्रीवं वानरर्षभम्}
{तदधीनं हि मे सौम्य सीतायाः परिमार्गणम्} %||3-71-8||

\twolineshloka
{इति ब्रुवाणं तं रामं सौमित्रिरिदमब्रवीत्}
{गच्छावस्त्वरितं तत्र ममापि त्वरते मनः} %||3-71-9||

\twolineshloka
{आश्रमात्तु ततस्तस्मान्निष्क्रम्य स विशां पतिः}
{आजगाम ततः पम्पां लक्ष्मणेन सहाभिभूः} %||3-71-10||

\threelineshloka
{समीक्षमाणः पुष्पाढ्यं सर्वतो विपुलद्रुमम्}
{कोयष्टिभिश्चार्जुनकैः शतपत्रैश्च कीचकैः}
{एतैश्चान्यैश्च विविधैर्नादितं तद्वनं महत्} %||3-71-11||

\twolineshloka
{स रामो विधिवान्वृक्षान्सरांसि विविधानि च}
{पश्यन्कामाभिसन्तप्तो जगाम परमं ह्रदम्} %||3-71-12||

\twolineshloka
{स तामासाद्य वै रामो दूरादुदकवाहिनीम्}
{मतङ्गसरसं नाम ह्रदं समवगाहत} %||3-71-13||

\twolineshloka
{स तु शोकसमाविष्टो रामो दशरथात्मजः}
{विवेश नलिनीं पम्पां पङ्कजैश्च समावृताम्} %||3-71-14||

\twolineshloka
{तिलकाशोकपुंनागबकुलोद्दाल काशिनीम्}
{रम्योपवनसम्बाधां पद्मसम्पीडितोदकाम्} %||3-71-15||

\twolineshloka
{स्फटिकोपमतोयाढ्यां श्लक्ष्णवालुकसन्तताम्}
{मत्स्यकच्छपसम्बाधां तीरस्थद्रुमशोभिताम्} %||3-71-16||

\threelineshloka
{सखीभिरिव युक्ताभिर्लताभिरनुवेष्टिताम्}
{किंनरोरगगन्धर्वयक्षराक्षससेविताम्}
{नानाद्रुमलताकीर्णां शीतवारिनिधिं शुभाम्} %||3-71-17||

\twolineshloka
{पद्मैः सौगन्धिकैस्ताम्रां शुक्लां कुमुदमण्डलैः}
{नीलां कुवलयोद्धातैर्बहुवर्णां कुथामिव} %||3-71-18||

\twolineshloka
{अरविन्दोत्पलवतीं पद्मसौगन्धिकायुताम्}
{पुष्पिताम्रवणोपेतां बर्हिणोद्घुष्टनादिताम्} %||3-71-19||

\twolineshloka
{स तां दृष्ट्वा ततः पम्पां रामः सौमित्रिणा सह}
{विललाप च तेजस्वी कामाद्दशरथात्मजः} %||3-71-20||

\twolineshloka
{तिलकैर्बीजपूरैश्च वटैः शुक्लद्रुमैस्तथा}
{पुष्पितैः करवीरैश्च पुंनागैश्च सुपुष्पितैः} %||3-71-21||

\threelineshloka
{मालतीकुन्दगुल्मैश्च भण्डीरैर्निचुलैस्तथा}
{अशोकैः सप्तपर्णैश्च केतकैरतिमुक्तकैः}
{अन्यैश्च विविधैर्वृक्षैः प्रमदेवोपशोभिताम्} %||3-71-22||

\twolineshloka
{अस्यास्तीरे तु पूर्वोक्तः पर्वतो धातुमण्डितः}
{ऋश्यमूक इति ख्यातश्चित्रपुष्पितकाननः} %||3-71-23||

\twolineshloka
{हरिरृक्षरजो नाम्नः पुत्रस्तस्य महात्मनः}
{अध्यास्ते तं महावीर्यः सुग्रीव इति विश्रुतः} %||3-71-24||

\twolineshloka
{सुग्रीवमभिगच्छ त्वं वानरेन्द्रं नरर्षभ}
{इत्युवाच पुनर्वाक्यं लक्ष्मणं सत्यविक्रमम्} %||3-71-25||

\fourlineindentedshloka
{ततो महद्वर्त्म च दूरसङ्क्रमम्}
{ क्रमेण गत्वा प्रविलोकयन्वनम्}
{ददर्श पम्पां शुभदर्श काननाम्}
{ अनेकनानाविधपक्षिसङ्कुलाम्} %||3-72-26||


{॥इत्यार्षे श्रीमद्रामायणे वाल्मीकीये आदिकाव्ये अरण्यकाण्डे एकसप्ततितमः सर्गः॥}

